\documentclass[11pt]{article}
\usepackage[english]{babel}
\usepackage[a4paper,margin=0.9in]{geometry}
\usepackage[T1]{fontenc}
\usepackage[utf8]{inputenc}
\usepackage{lmodern}
\usepackage{amsmath,amssymb,amsthm,mathtools,mathrsfs}
\usepackage{microtype}
\usepackage{fancyhdr}
\usepackage{array,longtable,enumitem,multicol}
\usepackage{tikz}

\providecommand{\dd}{\,d}
\providecommand{\tht}{\vartheta}
\providecommand{\vark}{\varkappa}

\providecommand{\sn}{\operatorname{sn}}
\providecommand{\cn}{\operatorname{cn}}
\providecommand{\dn}{\operatorname{dn}}
\providecommand{\am}{\operatorname{am}}
\providecommand{\ctg}{\operatorname{cotg}}
\providecommand{\eps}{\varepsilon}
\providecommand{\Jac}[2]{\left(\frac{#1}{#2}\right)}
\providecommand{\ii}{\mathrm{i}}
\providecommand{\cotg}{\operatorname{cotg}}

\setlength{\parindent}{1.5em}
\setlength{\parskip}{0.18em}
\setlength{\headheight}{14pt}
\emergencystretch=6em
\pagestyle{fancy}
\fancyhf{}
\cfoot{\thepage}
\lhead{Heinrich Weber, Lehrbuch der Algebra, Vol. II}
\rhead{English Translation}

\newcommand{\Om}{\Omega}
\newcommand{\Omm}{\Omega_m}
\newcommand{\Hm}{H_m}
\newcommand{\Lm}{\Omega_m}
\newcommand{\Norm}{\mathrm{N}}
\newcommand{\Normone}{\mathrm{N}_1}
\newcommand{\frp}{\mathfrak p}
\newcommand{\frP}{\mathfrak P}
\newcommand{\frG}{\mathfrak G}
\newcommand{\frA}{\mathfrak A}
\newcommand{\frB}{\mathfrak B}
\newcommand{\Gg}{\mathfrak G}
\newcommand{\Aa}{\mathfrak A}
\newcommand{\Bb}{\mathfrak B}
\newcommand{\pp}{\mathfrak p}
\providecommand{\eps}{\varepsilon}
\newcommand{\G}{\mathfrak G}
\newcommand{\B}{\mathfrak B}
\newcommand{\A}{\mathfrak A}
\newcommand{\frakA}{\mathfrak a}
\newcommand{\Na}{N_a}
\newcommand{\abs}[1]{\left|#1\right|}
\newcommand{\tg}{\operatorname{tg}}
\newcommand{\idxitem}[1]{\par\hangindent=1.4em\hangafter=1 #1}

\providecommand{\jac}[2]{\left(\frac{#1}{#2}\right)}
\providecommand{\Lzero}{\mathfrak L_0}
\providecommand{\Kgrp}{\mathfrak K}

% Batch62 macros
\providecommand{\Om}{\Omega}
\providecommand{\grpA}{\mathfrak A}
\providecommand{\grpB}{\mathfrak B}
\providecommand{\grpN}{\mathfrak N}
\providecommand{\grpP}{\mathfrak P}
\providecommand{\grpQ}{\mathfrak Q}
\providecommand{\grpX}{\mathfrak X}

\newcommand{\generalCayleyTable}{%
\[
\begin{array}{c|ccccc}
 & 1 & \alpha & \beta & \gamma & \cdots\\\hline
1&1&\alpha&\beta&\gamma&\cdots\\
\alpha&\alpha&\alpha^2&\alpha\beta&\alpha\gamma&\cdots\\
\beta&\beta&\beta\alpha&\beta^2&\beta\gamma&\cdots\\
\gamma&\gamma&\gamma\alpha&\gamma\beta&\gamma^2&\cdots\\
\vdots&\vdots&\vdots&\vdots&\vdots&
\end{array}
\]
}
\newcommand{\orderFourTables}{%
\[
\begin{array}{c|cccc}
 &1&\alpha&\beta&\gamma\\\hline
1&1&\alpha&\beta&\gamma\\
\alpha&\alpha&\beta&\gamma&1\\
\beta&\beta&\gamma&1&\alpha\\
\gamma&\gamma&1&\alpha&\beta
\end{array}
\qquad
\begin{array}{c|cccc}
 &1&\alpha&\beta&\gamma\\\hline
1&1&\alpha&\beta&\gamma\\
\alpha&\alpha&1&\gamma&\beta\\
\beta&\beta&\gamma&1&\alpha\\
\gamma&\gamma&\beta&\alpha&1
\end{array}
\]
}
\newcommand{\orderSixTable}{%
\[
\begin{array}{c|cccccc}
&1&\alpha&\beta&\gamma&\delta&\varepsilon\\\hline
1&1&\alpha&\beta&\gamma&\delta&\varepsilon\\
\alpha&\alpha&\beta&1&\varepsilon&\gamma&\delta\\
\beta&\beta&1&\alpha&\delta&\varepsilon&\gamma\\
\gamma&\gamma&\delta&\varepsilon&1&\alpha&\beta\\
\delta&\delta&\varepsilon&\gamma&\beta&1&\alpha\\
\varepsilon&\varepsilon&\gamma&\delta&\alpha&\beta&1
\end{array}
\]
}

\providecommand{\mat}[4]{\begin{pmatrix}#1&#2\\#3&#4\end{pmatrix}}
\providecommand{\mmat}[9]{\begin{pmatrix}#1&#2&#3\\#4&#5&#6\\#7&#8&#9\end{pmatrix}}
\providecommand{\ord}{\operatorname{ord}}

\providecommand{\mat}[4]{\begin{pmatrix}#1&#2\\#3&#4\end{pmatrix}}
\providecommand{\mmat}[9]{\begin{pmatrix}#1&#2&#3\\#4&#5&#6\\#7&#8&#9\end{pmatrix}}
\providecommand{\Cfield}{\mathfrak C}
\providecommand{\Lp}{L_p}
\providecommand{\Ap}{A_p}
\providecommand{\Fp}{F_p}

\providecommand{\trip}[3]{\begin{array}{ccc}#1&#2&#3\end{array}}

% Batch74 macros
\providecommand{\Cov}{\operatorname{Cov}}
\providecommand{\trip}[3]{\begin{array}{ccc}#1&#2&#3\end{array}}
\begin{document}
\begin{center}{\large Second Volume.}\\[0.4em]{\Large First Book. Groups.}\\[0.4em]\end{center}

\section*{First Section. General Group Theory.}

\section*{\S\,1. Definition of Groups.}

In the first volume, in connection with permutations, we became acquainted with the notion of a group and made important algebraic applications of it. Our next task must now be to formulate this concept, which is so exceedingly important in all more recent mathematics, in a more general way, and to become acquainted with the laws which govern it. We place the following definition at the beginning.

A system $P$ of things, or elements, of any kind becomes a group if the following assumptions are satisfied:

\begin{enumerate}[label=\arabic*.]
\item A rule is given by which, from a first and a second element of the system, a quite definite third element of the same system is derived.
\end{enumerate}
If $a$ is the first element, $b$ the second, and $c$ the third, one writes symbolically
\[
ab=c,\qquad c=ab,
\]
and says that $c$ is composed from $a$ and $b$, and that $a$ and $b$ are the components of $c$.

In this composition the commutative law is in general not assumed; that is, $ab$ may be different from $ba$. On the other hand,
\begin{enumerate}[resume,label=\arabic*.]
\item the associative law is assumed,
\end{enumerate}
that is, if $a,b,c$ are any three elements of $P$, then
\[
(ab)c=a(bc).
\]
From this it follows, by complete induction, that one always obtains the same result if, in any finite row of elements of $P$, $a,b,c,d\ldots$, one first composes two adjacent elements, then again two adjacent elements, and so on until the whole row is reduced to one element. This element is denoted by $abcd\ldots$. Thus, for example,
\[
\begin{aligned}
abcd&=(ab)cd=[(ab)c]d=(ab)(cd)\\
&=a(bc)d=[a(bc)]d=a[(bc)d]\\
&=ab(cd)=(ab)(cd)=a[b(cd)].
\end{aligned}
\]
\begin{enumerate}[resume,label=\arabic*.]
\item It is assumed that, if $ab=ab'$ or $ab=a'b$, then necessarily in the first case $b=b'$, and in the second case $a=a'$.
\end{enumerate}

If $P$ comprises a finite number of elements, then the group is called finite and the number of its elements its degree.

For finite groups 1., 2., 3. imply the following consequence:
\begin{enumerate}[resume,label=\arabic*.]
\item If two of the three elements $a,b,c$ of $P$ are given arbitrarily, then the third can always be determined, and in only one way, so that $ab=c$.
\end{enumerate}
Indeed, if $a,b$ are the given elements, the assertion 4. coincides with 1. If, however, $a$ and $c$ are given, let the second component $b$ in the composite $ab$ run through the whole system $P$, whose degree is $n$. Then by 1. and 3. the elements $ab$ so obtained are all distinct elements of $P$, and since their number is $n$, all elements of $P$, and hence also $c$, must occur among them. The same argument applies if $b$ and $c$ are given, by letting $a$ run through the whole system $P$.

For infinite groups this conclusion is no longer valid. Thus for infinite groups property 4. is also included as a requirement in the definition.

In the permutation groups considered in the first volume the marks of this general concept of a group are easily recognized.

We now draw from this definition some quite general consequences. By 4., for every given $b$ there is an element $e$ in $P$ which satisfies
\begin{equation}
eb=b,
\tag{1}
\end{equation}
and this $e$ is independent of $b$; for from (1) it follows, for every $c$, that
\[
ebc=bc,
\]
and by 4. the element $bc$ may mean any element of $P$. Similarly there is an element $e'$ satisfying, for every $b$,
\begin{equation}
be'=b.
\tag{2}
\end{equation}
But this element $e'$ is not distinct from $e$; for, putting $b=e'$ in (1) and $b=e$ in (2), one obtains
\[
ee'=e',\qquad ee'=e,
\]
and hence $e=e'$.

The element $e$ changes nothing when it is composed with arbitrary elements of $P$ and is called the identity of the group. In many cases it may, without misunderstanding, simply be denoted by ``1''.

For each element $a$ there is by 4. a definite element $a^{-1}$ satisfying
\begin{equation}
a^{-1}a=e.
\tag{3}
\end{equation}
From (3), (1), and (2) it follows that
\[
a^{-1}aa^{-1}=ea^{-1}=a^{-1}=a^{-1}e,
\]
and consequently, by 3.,
\begin{equation}
aa^{-1}=e.
\tag{4}
\end{equation}
The two elements $a,a^{-1}$ are called opposite or reciprocal to one another.

In special cases the commutative law may also hold for the composition of the elements of a group $P$; that is, for every two elements $a,b$ of the group one may have $ab=ba$.

\begin{enumerate}[resume,label=\arabic*.]
\item Groups which have this property are called commutative groups, or also Abelian groups.
\end{enumerate}

If the elements of two groups
\[
a,b,c,d\ldots \qquad\hbox{and}\qquad a',b',c',d'\ldots
\]
correspond to one another in a one-to-one manner in such a way that whenever $ab=c$ also $a'b'=c'$, then the groups are called isomorphic. The evident theorem holds that two groups isomorphic with a third group are isomorphic with one another. One may therefore collect all mutually isomorphic groups into one class of groups, which is itself again a group, whose elements are the generic concepts obtained by collecting the corresponding elements of the individual isomorphic groups into a common concept. The individual mutually isomorphic groups are then to be regarded as different representatives of one generic concept.

The properties of groups which may come into consideration are of different kinds. They may attach to particular groups and be derived from the nature of the elements of which the group consists, or from the nature of the law of composition. Or they may attach to groups as such and then have to be derived solely from the definition of the group concept. The latter properties are common to all isomorphic groups and may be called invariant properties of the group. If a comparison is allowed, one may recall the difference between metric and projective properties in geometry.

To the first kind of properties, derived from the special nature of the elements, belong for example, in permutation groups, the properties of transitivity and intransitivity, of primitivity and imprimitivity. To the invariant properties belong commutativity or noncommutativity, the degree, the divisors and their index, and the normal divisors. We shall first deal with these invariant properties, for which it is immaterial from which representative of a class of isomorphic groups they are derived.

\section*{\S\,2. The Divisors of Finite Groups.}

As we have already seen in the first volume, it is a particularly important question whether a part of the elements of a group is itself again a group. For this it is necessary and sufficient that every two elements of this part, when composed, again give an element belonging to the same part. This smaller group is then called a divisor of the whole group. The identity element ``1'' by itself forms a group in every group, and is therefore a divisor in every group.

We now restrict ourselves, as almost always in what follows, to finite groups, and suppose that
\[
P=1,a_1,a_2,\ldots,a_{n-1}
\]
is a group of degree $n$, with elements $a_0=1,a_1,a_2,\ldots,a_{n-1}$, and that
\[
Q=1,a_1,a_2,\ldots,a_{v-1}
\]
is a divisor of $P$ of degree $v$. Every such divisor $Q$ must have the fundamental properties of a group and therefore certainly contains the identity element ``1'', and with each of its elements $a_i$ also the opposite element $a_i^{-1}$. If $v<n$, we call $Q$ a proper divisor of $P$. Thus if $b$ is an element of $P$ not contained in $Q$, then the elements
\[
Q_1=b,a_1b,a_2b,\ldots,a_{v-1}b
\]
are all distinct from one another and are all not contained in $Q$; for if, say, $a_i b$ were contained in $Q$, then, since $Q$ is a group, $a_i^{-1}a_i b=b$ would also have to be contained in $Q$, contrary to the assumption.

If the whole group $P$ is not yet exhausted by $Q$ and $Q_1$, we take one of the remaining elements, which we may denote by $c$, and form
\[
Q_2=c,a_1c,a_2c,\ldots,a_{v-1}c.
\]
Then the elements of $Q_2$ are all distinct not only from one another, but also from the elements of $Q$ and $Q_1$. For if, for example, $a_i c=a_j b$, it would follow that $c=a_i^{-1}a_jb$; but $a_i^{-1}a_j$ is contained in $Q$ and is therefore identical with one of the elements $1,a_1,a_2,\ldots,a_{v-1}$, so that, contrary to the assumption, $c$ would be contained in $Q_1$. We continue in this way, forming the systems $Q,Q_1,Q_2,\ldots,Q_{\mu-1}$, the last of which is
\[
Q_{\mu-1}=g,a_1g,a_2g,\ldots,a_{v-1}g.
\]
Then $P$ must be exhausted by the systems $Q,Q_1,Q_2,\ldots,Q_{\mu-1}$, and since these systems contain only distinct elements, it follows that
\[
n=v\mu,
\]
or the theorem:
\begin{enumerate}[label=\arabic*.]
\item The degree of a group is divisible by the degree of each of its divisors. The quotient $\mu=n:v$ is called the index of the divisor $Q$.
\end{enumerate}
The systems $Q_1,Q_2,\ldots,Q_{\mu-1}$ we call, as in the special case of permutation groups, the side groups belonging to $Q$, and denote them by
\[
Q_1=Qb,
\quad Q_2=Qc,
\quad\ldots\quad
Q_{\mu-1}=Qg,
\]
so that we may also write symbolically
\begin{equation}
P=Q+Qb+Qc+\cdots+Qg.
\tag{1}
\end{equation}

At times we shall also call the whole system $Q,Q_1,\ldots,Q_{\mu-1}$ a system of side groups, and thus call the representation (1) the decomposition of $P$ into a system of side groups.

From the manner in which the side groups are formed it follows further that, if $b,c$ are any two elements of $P$, the two systems $Qb$ and $Qc$ are either wholly identical or have no element in common. For if $a_1b=a_2c$, where $a_1,a_2$ are two elements of $Q$, then for every other element $a$ of $Q$ one also has
\[
aa_1^{-1}b=aa_2^{-1}c,
\]
and when $a$ runs through the whole group $Q$, each of $aa_1^{-1}$ and $aa_2^{-1}$ also runs through the whole group $Q$. Hence $Qb$ and $Qc$ are identical.

As in Vol. I, \S\,154, 2., we may also decompose $P$ into side groups in the following way:
\[
P=Q+bQ+cQ+\cdots+gQ.
\]

The side groups do not have the marks of a group. For two elements of $Q_1$ to give, on composition, again an element of $Q_1$, one would have to have, say,
\[
a_1ba_2b=a_3b,
\]
and hence $b=a_2^{-1}a_1^{-1}a_3$. Thus, contrary to the assumption, $b$ would be contained in $Q$. The name side group is therefore to be understood only improperly.

Two divisors $Q,Q'$ of $P$ always have the element 1 in common. They may, however, also have other elements in common, and these common elements form a group, which we shall denote by $R$. For if the elements $a$ and $b$ belong both to $Q$ and to $Q'$, then, since both $Q$ and $Q'$ are groups, the same is true of the composite $ab$; hence $R$ is a group. This common group we call the greatest common divisor of $Q$ and $Q'$, or, with a geometrical resonance, the intersection of $Q$ and $Q'$. In the same way it follows that the elements common to any number of divisors of $P$ form a group, which we likewise call the greatest common divisor or the intersection of all these groups.

In every group we can form divisors by the following procedure. The identity element by itself is always a divisor.

Denote the repeated composition of an element with itself by powers, and denote the identity element by $a^0$. Then the row of elements
\[
1,a,a^2,a^3,\ldots
\]
all of which belong to $P$, cannot consist entirely of distinct elements. If the first element which occurs for the second time is
\[
a^n=a^{n+\alpha},
\]
then it follows that $a^\alpha=1$, and that $\alpha$ is the least positive number satisfying this condition. Then, if $m=q\alpha$ is any multiple of $\alpha$, one has $a^m=1$; conversely, whenever $a^m=1$, $m$ must be divisible by $\alpha$. For otherwise one could write $m=q'\alpha+\alpha'$, with $\alpha'$ positive and less than $\alpha$, and then $a^{\alpha'}=1$, contrary to the assumption that $\alpha$ is the least positive exponent for which $a^\alpha=1$.

One has
\[
a^\mu=a^{\mu'}
\]
always and only when $\mu\equiv\mu'\pmod{\alpha}$; here the exponents may also be negative, if $a^{-\nu}$ denotes the $\nu$-th power of the element $a^{-1}$, or the element opposite to $a^\nu$.

The row
\[
A=1,a,a^2,\ldots,a^{\alpha-1}
\]
then contains $\alpha$ mutually distinct elements of $P$, in which all powers of $a$ are represented. The elements $A$ plainly form a group of degree $\alpha$, since the exponents simply add under composition. This group is a divisor of $P$; hence $\alpha$ is a divisor of $n$. Thus for every element $a^n=1$.

The group $A$ is called the period of the element $a$, and $\alpha$ is also called the degree of the element $a$.

\section*{\S\,3. Normal Divisors of a Group.}

Let $P$ be a group as above, let $Q$ be a divisor of $P$, and let $b$ be an element of $P$ not contained in $Q$. Then the side group $Qb$ is not a group. On the other hand the system $b^{-1}Qb$, or written out,
\[
1,\;b^{-1}a_1b,\;b^{-1}a_2b,\ldots,\;b^{-1}a_{v-1}b
\]
certainly forms a group, since
\[
b^{-1}a_1b\cdot b^{-1}a_2b=b^{-1}a_1a_2b,
\]
and this group is isomorphic with $Q$. The group $b^{-1}Qb$ is called the group transformed by $b$. If $b$ itself belongs to $Q$, then $b^{-1}Qb$ is identical with $Q$. If $b$ is taken to be the different elements of $P$, one obtains a whole family of such groups, which we call the divisors of $P$ conjugate to $Q$, or briefly conjugate groups.

It may happen that all conjugate divisors are identical with one another. We have already seen, in the case of permutation groups, that this case is of special importance, and therefore introduce the following general definition.

If $Q$ is a divisor of $P$ which is identical with all its conjugate divisors, then $Q$ is called a normal divisor of $P$.

The group formed from the single element 1 is a normal divisor of every group. We obtain another normal divisor in the greatest common divisor $R$ of all divisors conjugate, in $P$, to any divisor $Q$.

Indeed, it was proved above that $R$, as the intersection of several divisors of $P$, is a group. If
\begin{equation}
Q,Q',Q'',\ldots
\tag{1}
\end{equation}
are the divisors of $P$ conjugate to $Q$, and if $b$ is any element of $P$, then the system of groups
\begin{equation}
b^{-1}Qb,
\quad b^{-1}Q'b,
\quad b^{-1}Q''b,
\ldots
\tag{2}
\end{equation}
is not different from the system (1). If $R$ is the intersection of the groups (1), then $b^{-1}Rb$ is the intersection of (2), and consequently $R$ is identical with $b^{-1}Rb$; that is, $R$ is a normal divisor of $P$.

If $N$ is a normal divisor of any group $P$ and $b$ is an arbitrary element of $P$, then
\begin{equation}
b^{-1}Nb=N\quad\hbox{or}\quad Nb=bN.
\tag{3}
\end{equation}
If $P$ has degree $n$, if $N$ has degree $v$ and index $\mu$, so that $n=\mu v$, then the elements $1,b_1,b_2,\ldots,b_{\mu-1}$ may be chosen so that the decomposition of $P$ into side groups gives
\[
P=N+Nb_1+Nb_2+\cdots+Nb_{\mu-1}=N+b_1N+b_2N+\cdots+b_{\mu-1}N.
\]
Thus
\begin{equation}
N_0=N,
\quad N_1=Nb_1=b_1N,
\quad N_2=Nb_2=b_2N,
\ldots,
\quad N_{\mu-1}=Nb_{\mu-1}=b_{\mu-1}N
\tag{4}
\end{equation}
is the system of side groups.

If $a$ is any element in $N$, then $Na$ is identical with $N$; hence $Nab$ is also identical with $Nb$. Since every element $c$ of $P$ occurs in $N$ or in one of its side groups, and so is contained in the form $ab_k$, it follows that $Nc=cN$ must coincide with one of the systems (4).

We also mention the following theorem, whose correctness follows immediately from the isomorphism of transformed groups. If $Q$ is a divisor of $P$ and $R$ is a divisor of $Q$, then, if $Q'$ and $R'$ are obtained from $Q$ and $R$ by transformation with the same element of $P$, $R'$ is also a divisor of $Q'$; and if $R$ is a normal divisor of $Q$, then $R'$ is also a normal divisor of $Q'$.

\section*{\S\,4. Composition of Parts.}

If $A$ and $B$ are any two rows of elements of a group $P$--whether groups or not--we shall understand by the symbolic product $AB$ the totality of all elements obtained by composing an element $a$ of $A$ with an element $b$ of $B$, according to the rule prevailing in $P$, to form $ab$. We call this kind of composition of $A$ and $B$ into $AB$ the composition of parts of $P$. We distinguish here between a part and a divisor of $P$, so that a divisor is always to be a group, whereas this is not necessary for a part.

In the same way one may compose three or more parts $A,B,C,\ldots$ of $P$, and the associative law holds for the composition of parts, as follows immediately from the fact that this law holds in $P$. Thus the meaning of a composite $ABC\ldots$ is uniquely determined.

In the composition $AB$ one of the parts $A,B$ may also consist of a single element, and then the notation $AB$ agrees with the notation used above for side groups.

For the composition of parts the following theorems hold:

\begin{enumerate}[label=\arabic*.]
\item The necessary and sufficient condition that a part $A$ of $P$ be a group is the equation
\begin{equation}
AA=A.
\tag{1}
\end{equation}
\end{enumerate}
For this equation says nothing other than that, if $a$ and $a'$ are contained in $A$, then $aa'$ also belongs to $A$, and hence that $A$ is a group.

\begin{enumerate}[resume,label=\arabic*.]
\item The two equations
\begin{equation}
AA=A,
\tag{2}
\end{equation}
\begin{equation}
AB=BA,
\tag{3}
\end{equation}
where $A$ is fixed and $B$ may be any arbitrary part of $P$, express that $A$ is a normal divisor of $P$.
\end{enumerate}
Indeed, by (2) $A$ is a divisor of $P$, and by (3), for each element $b$ of $P$, one has $b^{-1}Ab=A$; that is, $A$ is a normal divisor.

\begin{enumerate}[resume,label=\arabic*.]
\item If $A$ and $B$ are two divisors of $P$ (groups), then either of the two equations
\begin{equation}
AB=A,
\qquad BA=A,
\tag{4}
\end{equation}
is the necessary and sufficient condition that $B$ be a divisor of $A$.
\end{enumerate}
For if $B$ is a group and a divisor of the group $A$, and if $a$ is contained in $A$ and $b$ in $B$, hence also in $A$, then $ab$ is also contained in $A$. Thus $A$ contains every element of $AB$. Since secondly $B$, as a group, contains also the element 1, $AB$ contains every element of $A$; hence $AB$ is identical with $A$. In the same way one sees that $BA$ is identical with $A$. Conversely, from either of the equations (4), since $A$ as a group contains the identity element, it follows that every element of $B$ is contained in $A$, and hence that $B$ is a divisor of $A$.

\begin{enumerate}[resume,label=\arabic*.]
\item Under the composition of parts, the system of side groups belonging to a normal divisor of $P$ itself forms a group, in which the normal divisor $N$ is the identity, and $Nb_k,Nb_k^{-1}$ are opposite elements.
\end{enumerate}
For let $N$ be a normal divisor of $P$ and let
\begin{equation}
N,
\quad N_1=Nb_1,
\quad N_2=Nb_2,
\ldots
\tag{5}
\end{equation}
be the system of side groups. By the property of normal divisors $N$ is commutable with every $b_k$, hence $b_hb_kN=Nb_hb_k$, and consequently
\[
N_hN_k=NNb_hb_k.
\]
Since $b_hb_k$ is contained in $P$, $Nb_hb_k$ is identical with one of the side groups (5); and since $NN=N$ may be put, it follows, if we put $Nb_hb_k=N_i$, that
\begin{equation}
N_hN_k=N_i.
\tag{6}
\end{equation}
This proves the property \S\,1, 1. of groups for the composition of side groups. We have already emphasized that \S\,1, 2., the associative law, also holds. It remains to prove \S\,1, 3.; that is, if
\begin{equation}
N_hN_i=N_kN_l
\tag{7}
\end{equation}
is put, then from $N_h=N_k$ it follows that also $N_i=N_l$, and from $N_i=N_l$ it follows that also $N_h=N_k$.

According to the representation (5), however, we may write (7) in the two ways
\begin{equation}
Nb_hb_i=Nb_kb_l,
\qquad b_hb_iN=b_kb_lN.
\tag{8}
\end{equation}
If now $N_i=N_l$, then we may also assume $b_i=b_l$, and from (8), by composition with $b_i^{-1}$, we obtain
\[
Nb_h=Nb_k.
\]
If $Nb_h=Nb_k$, then we assume $b_h=b_k$, and obtain from (8), by composition with $b_h^{-1}$,
\[
b_iN=b_lN,
\]
and hence also $N_i=N_l$.

It is thereby proved that the system of side groups becomes a group under the composition of parts. That in this group the normal divisor $N$ is the identity element follows immediately from 1., since then $NN_k=NNb_k=N_k$. And since $Nb_kNb_k^{-1}=Nb_kb_k^{-1}=N$, the elements $Nb_k$ and $Nb_k^{-1}$ are opposite elements.

The group of the quantities $N_k$, whose existence has thereby been proved, is called the group of side groups, or the group complementary to $N$ with respect to $P$. Following Holder, we denote it by $P/N$. The degree of the complementary group is equal to the index of the group $N$.

At the end of these considerations we mention one more theorem on normal divisors.
\begin{enumerate}[resume,label=\arabic*.]
\item If $N$ is a divisor of $P$, $N'$ a divisor of $N$, and at the same time a normal divisor of $P$, then $N'$ is also a normal divisor of $N$.
\end{enumerate}
This is self-evident, since for every element $b$ of $P$ one has $b^{-1}N'b=N'$, and since every element of $N$ is also an element of $P$.

But one must not conversely conclude that, if $N$ is a normal divisor of $P$ and $N'$ is a normal divisor of $N$, then $N'$ must also be a normal divisor of $P$.

\section*{\S\,5. Multi-stage isomorphism.}

As above, let $N$ be a normal divisor of $P$ of degree $v$ and index $\mu$, so that $n=\mu v$. Then the group complementary to $P$ is $P/N$, of degree $\mu$. We shall denote this group, or a group isomorphic with it, by $Q$, and its elements, corresponding to the side-groups $N,Nb_1,Nb_2\ldots$, by $A,B,C\ldots$. To each of these elements there correspond $v$ elements of the group $P$, say
\[
\begin{array}{ll}
\hbox{to the element } A & \hbox{the elements } a,a_1\ldots a_{v-1},\\
\hbox{to the element } B & \hbox{the elements } b,b_1\ldots b_{v-1},
\end{array}
\]
\[
\cdots\cdots\cdots\cdots\cdots\cdots.
\]
This correspondence is then such that every composite element $ab$ corresponds to the composite element $AB$. For the elements $A$ and $B$ correspond to the side-groups $Na$ and $Nb$, and, since $N$ is a normal divisor,
\[
NaNb=Nab.
\]
Thus the same law holds here for the composition of the $A,B\ldots$ on the one hand and of the $a,b\ldots$ on the other as in isomorphic groups (\S\,1), except that to an element $A$ there corresponds not merely one element $a$, but several. This fact gives occasion to extend the concept of isomorphism, as is done in the following definition.

A group $P$ with elements $a,b,c\ldots$ is called (multi-stage) isomorphic with a group $Q$ with elements $A,B,C\ldots$ if the two groups can be put into relation so that to each of the elements $A,B,C\ldots$ there correspond one or more of the elements $a,b,c\ldots$, in such a way that every element of $Q$ corresponds to one and only one element of $P$, and that, whenever $a$ and $A$, $b$ and $B$ correspond to one another, $ab$ corresponds to the element $AB$.

It is first proved from this that to each of the elements $A,B,C\ldots$ there corresponds the same number of elements $a,b,c\ldots$, and hence that the degree of $Q$ is a divisor of the degree of $P$. For let $A$ be the identity element in $Q$, and suppose that the elements $a,a_1\ldots a_{v-1}$ of $P$ correspond to it. These latter elements must then form a group of degree $v$, which we shall denote by $N$, because, by the definition of isomorphism, the element $aa_1$ must correspond to the element $AA=A$. If $b$ is an element of $P$ corresponding to the element $B$, then again all the elements $Nb$ must correspond to the same element $B$ of $Q$. For every $ab$ must correspond to the element $AB=B$. If $b_1,b$ are two elements corresponding to $B$, then $b_1b^{-1}=a$ corresponds to the element $A$, and is therefore contained in $N$; hence $b_1=ab$ is contained in $Nb$. Thus all elements corresponding to $B$ are exhausted by $Nb$, and to each element of $Q$ there correspond $v$ elements of $P$. The isomorphism is called $v$-stage. But every element $b^{-1}ab$ likewise corresponds to the identity element $A$, and therefore $b^{-1}Nb=N$; that is, $N$ is a normal divisor of $P$, and $Q$ is one-stage isomorphic with $P/N$.

We thus see that there is no multi-stage isomorphism other than that between the complementary group of a normal divisor and the whole group. One could, however, extend the concept of isomorphism further by saying that two groups which are $\mu$- and $v$-stage isomorphic to a third group are $\mu v$-stage isomorphic to one another. By far the most important case is the one-stage isomorphism; this we therefore call simply isomorphism, while a multi-stage isomorphism will always be marked by an added qualification.

\section*{\S\,6. The composition series and C. Jordan's theorem.}

A group $P$ is called simple if it has, apart from itself and the identity, no normal divisor; it is called composite if further normal divisors are present.

A normal divisor which has no other normal divisor above it, that is, which is not part of another proper normal divisor of $P$, is called a greatest normal divisor of $P$. If $P_1$ is a greatest normal divisor of $P$, $P_2$ a greatest normal divisor of $P_1$, $P_3$ a greatest normal divisor of $P_2$, and so on, then the sequence of groups
\begin{equation}
P,\;P_1,\;P_2,\;P_3\ldots,\;1,
\tag{1}
\end{equation}
which necessarily ends with the group $1$ formed from the identity element alone, is called a composition series of the group $P$.

We shall denote the degrees of the groups (1) by $n,n_1,n_2,n_3\ldots,1$, and the quotients $n:n_1$, $n_1:n_2$, $n_2:n_3\ldots$, that is, the indices of $P_1$ with respect to $P$, of $P_2$ with respect to $P_1$, and so on, by $\nu_1,\nu_2,\nu_3\ldots$. The sequence of integers
\begin{equation}
\nu_1,\;\nu_2,\;\nu_3\ldots
\tag{2}
\end{equation}
is called the index series of the group $P$.

For a given group $P$ the greatest normal divisors $P_1,P_2,P_3\ldots$ can in general be chosen in several different ways, and accordingly the degrees $n_1,n_2,n_3\ldots$ and the indices $\nu_1,\nu_2,\nu_3\ldots$ may also come out differently. Nevertheless the following beautiful and important theorem of C. Jordan holds:

\begin{enumerate}[label=\Roman*.]
\item However a composition series of a group $P$ may be chosen, the index series is always the same, apart from the order of its terms.
\end{enumerate}

The proof rests on the complementary groups introduced in the preceding paragraph.

Let $Q$ and $Q_1$ be normal divisors of $P$, and let $Q_1$ at the same time be a divisor of $Q$; hence, by \S\,4, 5., it is also a normal divisor of $Q$. Then:

\begin{enumerate}[label=\arabic*.]
\item The group $Q/Q_1$ is a normal divisor of the group $P/Q_1$, and the index of $Q/Q_1$ with respect to $P/Q_1$ is equal to the index of $Q$ with respect to $P$.
\end{enumerate}

To see this, one need only examine more closely the way in which the individual groups are formed. Let $n,q,q_1$ be the degrees of $P,Q,Q_1$, and let
\begin{equation}
n=\nu q,
\qquad q=\nu_1q_1,
\qquad n=\mu q_1,
\qquad \mu=\nu\nu_1.
\tag{3}
\end{equation}
Decompose $Q$ into side-groups with respect to $Q_1$, that is, if $b_1,b_2\ldots b_{\nu_1-1}$ are suitably chosen elements in $Q$, put
\begin{equation}
Q=Q_1+Q_1b_1+\cdots+Q_1b_{\nu_1-1}.
\tag{4}
\end{equation}
The elements of the group $Q/Q_1$ are
\begin{equation}
Q_1,\;Q_1b_1,\ldots,Q_1b_{\nu_1-1}.
\tag{5}
\end{equation}
When we decompose $P$ into side-groups with respect to $Q_1$, the elements (5) certainly occur among them, and hence $Q/Q_1$ is a divisor of $P/Q_1$.

It remains only to prove that this divisor is normal. If $a$ denotes any element of $P$, then all elements of $P/Q_1$ are contained in the form $Q_1a$, and we therefore have only to show that, if $b$ is any element of $Q$, every element
\begin{equation}
Q_1a^{-1}Q_1bQ_1a
\tag{6}
\end{equation}
is contained in $Q/Q_1$, that is, is of the form $Q_1b$. This follows immediately from
\begin{equation}
Q_1a^{-1}Q_1bQ_1a=Q_1a^{-1}ba,
\tag{7}
\end{equation}
since, together with $b$, the element $a^{-1}ba$ also occurs in $Q$, because $Q$ is a normal divisor of $P$. The correctness of the statement about the indices is shown by counting: the degrees of $P/Q_1$ and $Q/Q_1$ are $\nu\nu_1$ and $\nu_1$, so the index is $\nu$.

Alongside this there is the following theorem:
\begin{enumerate}[resume,label=\arabic*.]
\item If $Q_1$ is a normal divisor of $P$ of degree $q_1$, and if the group $P/Q_1$ itself has a divisor $M$ of degree $m$, then $P$ has a divisor $Q$ of degree $q=mq_1$. Moreover $Q_1$ is a normal divisor of $Q$, and $M=Q/Q_1$. If $M$ is a normal divisor of $P/Q_1$, then $Q$ is also a normal divisor of $P$.
\end{enumerate}

Let $\mu$ be the index of $Q_1$ with respect to $P$, and decompose $P$ into side-groups:
\begin{equation}
P=Q_1+Q_1e_1+Q_1e_2+\cdots+Q_1e_{\mu-1},
\tag{8}
\end{equation}
where $e_1,e_2\ldots e_{\mu-1}$ are suitably chosen elements in $P$, and
\begin{equation}
Q_1,\;Q_1e_1,\;Q_1e_2\ldots Q_1e_{\mu-1}
\tag{9}
\end{equation}
are the elements of $P/Q_1$. If a divisor $M$ is contained in the group $P/Q_1$, then it must consist of a part of the elements (9), say of
\begin{equation}
Q_1,\;Q_1b_1,\;Q_1b_2\ldots Q_1b_{\nu_1-1},
\tag{10}
\end{equation}
and if this system is to form a group, then for any two elements $b_i,b_k$ the product
\begin{equation}
Q_1b_iQ_1b_k=Q_1b_ib_k
\tag{11}
\end{equation}
must again be contained in (10), say equal to $Q_1b_h$. If the system $1,b_1,b_2\ldots b_{\nu_1-1}$ is denoted by $B$, then, by the method of composing parts, the relation (11) may be written
\begin{equation}
Q_1BQ_1B=Q_1B,
\tag{12}
\end{equation}
and by \S\,4, 1. this says that the elements of $P$ contained in
\begin{equation}
Q=Q_1B
\tag{13}
\end{equation}
form a group which contains $Q_1$ as a divisor and is itself a divisor of $P$. That $Q_1$ is a normal divisor of $Q$ follows from \S\,4, 5., since $Q_1$ is assumed to be a normal divisor of $P$; and from the representation (10) of the group $M$ it follows that $M=Q/Q_1$.

It remains to show that, when $M$ is a normal divisor of $P/Q_1$, $Q$ is also a normal divisor of $P$. This follows as follows. If $e$ is any element of $P$ and $Q_1b_k$ an element of $M$, the assumption that $M$ is a normal divisor of $P/Q_1$ gives
\[
Q_1e^{-1}Q_1b_kQ_1e=Q_1b_h,
\]
which, by the composition of parts, we may express by the formula
\begin{equation}
Q_1e^{-1}Q_1BQ_1e=Q_1B.
\tag{14}
\end{equation}
Since $Q_1$ is a normal divisor of $P$, $Q_1$ can be interchanged with each of the other factors, and so (14) gives
\[
e^{-1}Q_1Be=Q_1B,
\]
or
\begin{equation}
e^{-1}Qe=Q,
\tag{15}
\end{equation}
which expresses that $Q$ is a normal divisor of $P$.

From theorems 1 and 2 one obtains the corollary:
\begin{enumerate}[resume,label=\arabic*.]
\item A normal divisor $Q$ of $P$ is a greatest normal divisor if and only if the complementary group $P/Q$ is simple.
\end{enumerate}

If $Q$ and $Q'$ are two normal divisors of $P$, then the symbolic product $QQ'$ is also a normal divisor. For, since $Q$ is a normal divisor,
\[
QQ'=Q'Q,
\]
and therefore
\[
QQ'QQ'=QQQ'Q'=QQ'.
\]
By \S\,4, 1. this is the criterion that $QQ'$ is a group. That it is a normal divisor of $P$ then follows likewise from \S\,4, 2. from
\[
QQ'B=QBQ'=BQQ',
\]
where $B$ is any part of $P$.

Now if $Q$ is a greatest normal divisor of $P$, and $Q'$ is different from $Q$, then $QQ'$ is a normal divisor of $P$ which contains both $Q$ and $Q'$, and is therefore different from $Q$. Consequently, since there is no normal divisor of $P$ above $Q$ except $P$ itself,
\begin{equation}
QQ'=P,
\tag{16}
\end{equation}
and likewise one must have
\begin{equation}
Q'Q=P.
\tag{17}
\end{equation}
These theorems hold if only one of the two groups $Q,Q'$ is a greatest normal divisor of $P$. We now assume that both have this property, and denote by $R$ the greatest common divisor of $Q$ and $Q'$. This group $R$ is then a normal divisor of $P$, and also of $Q$ and of $Q'$. For if $a$ is any element of $P$, and $c$ is contained both in $Q$ and in $Q'$, then $a^{-1}ca$ is also in $Q$ and in $Q'$, hence also in $R$. Therefore $a^{-1}Ra=R$, and $R$ is a normal divisor of $P$ and consequently a normal divisor of $Q$ and of $Q'$ (\S\,4, 5.). To decompose $Q$ into side-groups of $R$, choose a suitable system of elements $1,b_1,b_2,b_3\ldots$ in $Q$ so that $R,Rb_1,Rb_2,Rb_3\ldots$ are all distinct, and put
\[
Q=R+Rb_1+Rb_2+Rb_3+\cdots.
\]
If we set
\begin{equation}
B=1,b_1,b_2,b_3\ldots,
\tag{18}
\end{equation}
we may write this briefly as
\begin{equation}
Q=RB.
\tag{19}
\end{equation}
Now by (13) $Q'Q=P$, so also
\begin{equation}
P=Q'RB=Q'B
\tag{20}
\end{equation}
or
\[
P=Q'+Q'b_1+Q'b_2+Q'b_3\cdots.
\]
The side-groups
\begin{equation}
Q',\;Q'b_1,\;Q'b_2,\;Q'b_3\ldots
\tag{21}
\end{equation}
are all distinct. For if, for example,
\[
Q'b_2=Q'b_1,
\]
then, since $Q'$ contains the identity, it would follow that there is an element $e$ in $Q'$ such that
\[
b_2=eb_1.
\]
But this element $e$ is equal to $b_2b_1^{-1}$, and since the $b$'s are contained in $Q$, it is likewise contained in $Q$ and hence also in $R$. Then one would also have $Rb_2=Rb_1$, contrary to the assumption.

We may now form the elements of the group complementary to $R$ with respect to $Q$, and of the group complementary to $Q'$ with respect to $P$. We obtain these two groups:
\begin{equation}
\begin{array}{ll}
Q/R &= R,\;Rb_1,\;Rb_2,\;Rb_3\ldots,\\
P/Q' &= Q',\;Q'b_1,\;Q'b_2,\;Q'b_3\ldots.
\end{array}
\tag{22}
\end{equation}
From this one sees readily that these groups are not only of the same degree, but are isomorphic. For at the same time
\[
Rb_1Rb_2=Rb_1b_2
\]
and
\[
Q'b_1Q'b_2=Q'b_1b_2.
\]
In the same way one may also conclude that the two groups
\[
Q'/R,
\qquad P/Q
\]
are isomorphic.

But since $Q,Q'$ are greatest normal divisors, the groups $P/Q$ and $P/Q'$ are simple by 3), and therefore the isomorphic groups $Q'/R$ and $Q/R$ are simple. Thus we have:
\begin{enumerate}[resume,label=\arabic*.]
\item If $Q$ and $Q'$ are two distinct greatest normal divisors of $P$, and if $R$ is the intersection of $Q$ and $Q'$, then $R$ is a greatest normal divisor both of $Q$ and of $Q'$, and the index of $R$ with respect to $Q$ is equal to the index of $Q'$ with respect to $P$.
\end{enumerate}

We have thereby obtained the basis for proving the theorem on the constancy of the index series very simply.

For better oversight we write the different composition series to be compared in such a way that the index of each member with respect to the preceding one is placed under that member. Thus let any two given composition series of $P$, with their associated indices, be
\begin{equation}
\begin{array}{cccccc}
P,&Q,&Q_1,&Q_2,&Q_3&\ldots\\
&\nu,&\nu_1,&\nu_2,&\nu_3&\ldots
\end{array}
\tag{23}
\end{equation}
\begin{equation}
\begin{array}{cccccc}
P,&Q',&Q'_1,&Q'_2,&Q'_3&\ldots\\
&\nu',&\nu'_1,&\nu'_2,&\nu'_3&\ldots
\end{array}
\tag{24}
\end{equation}
Let $R$ be the intersection of $Q$ and $Q'$, and let $\mu$ and $\mu'$ be the indices of $R$ with respect to $Q$ and $Q'$. Because of the isomorphism of the groups $P/Q$ and $Q'/R$, one has $\mu'=\nu$, and for the corresponding reason $\mu=\nu'$. Form a composition series of $R$ with its index series
\begin{equation}
\begin{array}{ccccc}
R,&R_1,&R_2,&R_3&\ldots,\\
&\mu_1,&\mu_2,&\mu_3&\ldots.
\end{array}
\tag{25}
\end{equation}
Since $R$ is a greatest normal divisor of $Q$ and of $Q'$, we can form from this two new composition series of $P$, namely
\begin{equation}
\begin{array}{ccccccc}
P,&Q,&R,&R_1,&R_2,&R_3&\ldots\\
&\nu,&\nu',&\mu_1,&\mu_2,&\mu_3&\ldots
\end{array}
\tag{26}
\end{equation}
\begin{equation}
\begin{array}{ccccccc}
P,&Q',&R,&R_1,&R_2,&R_3&\ldots\\
&\nu',&\nu,&\mu_1,&\mu_2,&\mu_3&\ldots.
\end{array}
\tag{27}
\end{equation}
The two composition series (26) and (27) of $P$ therefore have the same index series.

Now suppose that the theorem to be proved has already been proved for groups whose degree is $\leq \frac12 n$, where $n$ is the degree of $P$. Since the degree of $Q$ is a proper divisor of $n$, and hence at most $\frac12n$, it follows that all index series of $Q$ agree with one another; that is, the series
\[
\nu',\;\mu_1,\;\mu_2,\;\mu_3\ldots,
\qquad
\nu_1,\;\nu_2,\;\nu_3,\;\nu_4\ldots,
\]
agree up to order. Similarly the index series of $Q'$
\[
\nu,\;\mu_1,\;\mu_2,\;\mu_3\ldots,
\qquad
\nu'_1,\;\nu'_2,\;\nu'_3,\;\nu'_4\ldots
\]
agree. Hence the two index series of $P$
\[
\nu,\;\nu_1,\;\nu_2,\;\nu_3,\;\nu_4\ldots,
\qquad
\nu',\;\nu'_1,\;\nu'_2,\;\nu'_3,\;\nu'_4\ldots
\]
also agree with one another, since the first agrees with $\nu,\nu',\mu_1,\mu_2,\mu_3\ldots$, and the second with $\nu',\nu,\mu_1,\mu_2,\mu_3\ldots$.

For groups of degree 2, which have only the one normal divisor 1, the theorem is evident; by complete induction it is therefore proved in general.

If in two composition series of $P$, which we shall now write with their index series as
\begin{equation}
\begin{array}{cccccc}
P,&P_1,&P_2,&P_3&\ldots&P_{\mu-1},1\\
&j_1,&j_2,&j_3&\ldots&j_{\mu-1},j_\mu
\end{array}
\tag{28}
\end{equation}
\begin{equation}
\begin{array}{cccccc}
P,&P'_1,&P'_2,&P'_3&\ldots&P'_{\mu-1},1\\
&j'_1,&j'_2,&j'_3&\ldots&j'_{\mu-1},j'_\mu
\end{array}
\tag{29}
\end{equation}
a common member $P_r=P'_s$ occurs, then the theorem on the constancy of the index series also holds for the group $P_r=P'_s$. It follows first that $s=r$, and that the index series $j_{r+1},j_{r+2}\ldots j_\mu$ and $j'_{r+1},j'_{r+2}\ldots j'_\mu$ agree up to order. Consequently the preceding parts of the index series,
\[
j_1,j_2\ldots j_r
\qquad\hbox{and}\qquad
j'_1,j'_2\ldots j'_r,
\]
must also agree.

\section*{\S\,7. Further theorems on composition series.}

Many further theorems hold for composition series; we derive here the most important of them.

\begin{enumerate}[label=\Roman*.]
\setcounter{enumi}{1}
\item If $P_r$ is any normal divisor of $P$, then there is a composition series of $P$ in which $P_r$ occurs.
\end{enumerate}

If $P_r$ is a greatest normal divisor of $P$, then we may set $P_r=P_1$ and continue the composition series arbitrarily from there. If $P_r$ is not a greatest normal divisor of $P$, we seek a greatest normal divisor above $P_r$ and denote it by $P_1$. Then $P_r$ is a normal divisor of $P_1$. If it is greatest, we set $P_r=P_2$ and continue the composition series arbitrarily from there. If $P_r$ is still not a greatest normal divisor of $P_1$, we seek a greatest normal divisor of $P_1$ above $P_r$ and proceed in this way until we finally reach $P_r$, which must occur after finitely many steps. If we then append a composition series of $P_r$, we have a composition series of $P$ in which $P_r$ occurs, as theorem II requires.

In a composition series of a group $P$, by definition, each member is a normal divisor of each preceding member. The first member $P_1$ and the last member 1 are at the same time normal divisors of $P$ itself. This will not in general be the case for the intermediate members $P_2,P_3\ldots$. Thus in the composition series there are certain distinguished members which are at the same time normal divisors of $P$ itself and which play a special role. We now derive an important theorem about them.

Assume that $Q$ is a member of a composition series of $P$ which is also a normal divisor of $P$, whereas the member $Q_1$ following $Q$ is not a normal divisor of $P$. Then among the groups conjugate to $Q_1$ with respect to $P$, that is, among the groups $a^{-1}Q_1a$, there are several distinct ones; denote them by
\begin{equation}
Q_1,\;Q'_1,\;Q''_1,\;Q'''_1\ldots.
\tag{1}
\end{equation}
The greatest common divisor $R$ of all these groups is again a normal divisor of $P$.

All the groups (1) are, moreover, greatest normal divisors of $Q$. This follows from the final theorem of \S\,3, since $Q,Q_1$ are isomorphic with $a^{-1}Qa,a^{-1}Q_1a$, while $Q$, being a normal divisor of $P$, is identical with $a^{-1}Qa$. We may therefore let any one of the groups (1) follow $Q$ in the composition series; they all have the same index $\nu$.

If $Q_2$ is the greatest common divisor of $Q_1,Q'_1$ and is therefore different from $Q_1$, then by \S\,6, 4., $Q_2$ is a greatest normal divisor of $Q_1$ and $Q'_1$, and its index is likewise $\nu$. We may therefore continue the composition series after $Q$ in the following two ways:
\begin{equation}
\begin{array}{ccc@{\qquad}ccc}
Q,&Q_1,&Q_2&Q,&Q'_1,&Q_2\\
&\nu,&\nu&&\nu,&\nu
\end{array}
\tag{2}
\end{equation}
We shall generalize the law expressed here and prove it by complete induction.

If $Q_2$ is not yet equal to $R$, add to $Q_1,Q'_1$ a third group $Q''_1$ from the sequence (1), so that the intersection $Q_3$ of $Q_1,Q'_1,Q''_1$ becomes a proper divisor of $Q_2$, and proceed in this way until the greatest common divisor $R$ of all the groups (1) is reached.

Thus let
\begin{equation}
\begin{array}{rcl}
Q_2&\hbox{be the intersection of}&Q_1,Q'_1,\\
Q_3&\hbox{ ,, ,, }&Q_1,Q'_1,Q''_1,\\
\cdots&&\cdots\\
Q_r&\hbox{ ,, ,, }&Q_1,Q'_1,Q''_1\ldots Q_1^{(r-1)},
\end{array}
\tag{3}
\end{equation}
and choose the ordering of the groups $Q_1,Q'_1\ldots Q_1^{(r-1)}$ so that, among $Q_1,Q_2\ldots Q_r$, each is a proper divisor of the preceding one. This ordering can be continued so long as $Q_r$ is not equal to $R$.

We shall prove that the composition series of $P$ can be chosen so that it contains
\begin{equation}
\begin{array}{ccccc}
Q,&Q_1,&Q_2&\ldots&Q_r\\
&\nu,&\nu&&\nu
\end{array}
\tag{4}
\end{equation}
and that all the indices in this part of the series are equal to $\nu$.

The assertion is true for $r=2$. We prove it in general by complete induction, assuming that it has already been proved for $r$.

If $Q_r$ is not equal to $R$, choose $Q_1^{(r)}$ so that $Q_{r+1}$ is the intersection of $Q_1,Q'_1\ldots Q_1^{(r-1)},Q_1^{(r)}$, and also denote by $Q'_r$ the intersection of $Q_1,Q'_1\ldots Q_1^{(r-2)},Q_1^{(r)}$. Then $Q'_r$ is different from $Q_r$, and $Q_{r+1}$ is the intersection of $Q_r$ and $Q'_r$. By the hypothesis (4) we now have two corresponding pieces of the composition series with the constant index series $\nu$:
\begin{equation}
\begin{array}{ccccc}
Q,&Q_1&\ldots&Q_{r-1},&Q_r,\\
Q,&Q_1&\ldots&Q_{r-1},&Q'_r;
\end{array}
\tag{5}
\end{equation}
for the second of these series is formed exactly as the first according to rule (3), except that $Q_1^{(r)}$ has replaced $Q_1^{(r-1)}$. By \S\,6, 2., $Q_{r+1}$ is therefore a greatest normal divisor both of $Q_r$ and of $Q'_r$, of index $\nu$; and hence it follows, as required, that the composition series can be continued as
\begin{equation}
Q,\;Q_1\ldots Q_{r-1},\;Q_r,\;Q_{r+1},
\tag{6}
\end{equation}
and that the newly added index is again $\nu$.

This can be continued until the group $R$ is reached, and we obtain a piece of the composition series
\begin{equation}
\begin{array}{cccccc}
Q,&Q_1,&Q_2&\ldots&Q_s,&R,\\
&\nu,&\nu&&\nu&\nu,
\end{array}
\tag{7}
\end{equation}
whose last member $R$ is again a normal divisor of $P$, and in which all indices agree.

From $R$ onward we continue the composition series. The index may then change; but from there on it can again be kept equal until another normal divisor of $P$ is reached.

Thus the following theorem is proved:
\begin{enumerate}[label=\Roman*.]
\setcounter{enumi}{2}
\item One can arrange the composition series of $P$ with its index series
\end{enumerate}
\begin{equation}
\begin{array}{cccccc}
P,&P_1,&P_2,&P_3&\ldots\\
&j_1,&j_2,&j_3&\ldots
\end{array}
\tag{8}
\end{equation}
in such a way that, whenever a change occurs in the index series, that is, whenever $j_{r+1}$ is different from $j_r$, the member $P_r$ is a normal divisor of $P$.

\section*{\S\,8. Metacyclic groups.}

In \S\,176 of the first volume, in connection with permutation groups, we defined metacyclic groups and learned their importance for the solution of equations. The concept, however, is independent of the special meaning of the group elements, and we therefore now define it in general.

A group whose index series consists only of prime numbers is called a metacyclic group.

For these groups there is an important theorem for applications to algebra; but it is independent of those applications, and we therefore derive it here. It is as follows:
\begin{enumerate}[label=\Roman*.]
\setcounter{enumi}{3}
\item A metacyclic group has a commutative normal divisor different from the identity group.
\end{enumerate}

Let $P$ be a metacyclic group and let
\begin{equation}
\begin{array}{cccccc}
P,&P_1,&P_2&\ldots&P_{\mu-1},&1\\
&j_1,&j_2&\ldots&j_{\mu-1},&j_\mu
\end{array}
\tag{1}
\end{equation}
be its composition series, whose index series $j_1,j_2\ldots$ consists only of prime numbers. Then $P_{\mu-1}$ is certainly a commutative group. For if $\pi$ is any element different from 1 in $P_{\mu-1}$, then the powers
\begin{equation}
1,\pi,\pi^2\ldots \pi^{j_\mu-1}
\tag{2}
\end{equation}
are, since $j_\mu$ is prime, all distinct, and hence make up the whole group $P_{\mu-1}$; the cyclic group (2) is certainly commutative, since for any two exponents $h,k$ one always has
\[
\pi^h\pi^k=\pi^k\pi^h=\pi^{h+k}.
\]
At the same time, by the definition of the composition series, $P_{\mu-1}$ is a normal divisor of $P_{\mu-2}$.

We now assume that we have a commutative normal divisor $Q_\nu$, different from the identity, of some group $P_\nu$ in the series (1), and we assume in addition that, if there are several divisors of $P_\nu$ with the properties of $Q_\nu$, one of lowest possible degree has been chosen as $Q_\nu$.

If we show how one can derive from $Q_\nu$ a commutative normal divisor $Q_{\nu-1}$ of $P_{\nu-1}$ which contains $Q_\nu$, then this process can be continued until a commutative normal divisor $Q$ of $P$ itself is obtained, as the theorem requires.

Choose an element $\gamma$ in the group $P_{\nu-1}$ which is not contained in $P_\nu$. Since $P_\nu$ is a normal divisor of $P_{\nu-1}$, one has
\begin{equation}
\gamma P_\nu=P_\nu\gamma.
\tag{3}
\end{equation}
If $h$ is the least exponent for which $\gamma^h$ is contained in $P_\nu$, then $h$ is greater than 1 and
\begin{equation}
\gamma^hP_\nu=P_\nu.
\tag{4}
\end{equation}
The elements contained in the system of side-groups
\begin{equation}
P_\nu,\;\gamma P_\nu,\;\gamma^2P_\nu\ldots \gamma^{h-1}P_\nu
\tag{5}
\end{equation}
are all distinct and form a group whose degree, if the degree of $P_\nu$ is denoted by $p_\nu$, is $hp_\nu$. The group (5) is also a divisor of $P_{\nu-1}$, and therefore $hp_\nu$ must be a divisor of $p_{\nu-1}=j_\nu p_\nu$; hence $h$ is a divisor of $j_\nu$. But $j_\nu$ is a prime number, so $h=j_\nu$, and (5) exhausts the whole group $P_{\nu-1}$:
\begin{equation}
P_{\nu-1}=P_\nu+\gamma P_\nu+\gamma^2P_\nu+\cdots+\gamma^{j_\nu-1}P_\nu.
\tag{6}
\end{equation}

We now consider the system of transformed groups
\begin{equation}
\gamma^{-1}Q_\nu\gamma,
\tag{7}
\end{equation}
where $\gamma$ runs through all elements of $P_{\nu-1}$.

All these groups are isomorphic to one another, and therefore, like $Q_\nu$, commutative. They are all contained in $P_\nu$, and since $\gamma^{-1}P_\nu\gamma=P_\nu$, they are, by the final theorem of \S\,3, normal divisors of $P_\nu$.

If all the groups (7) are identical, then $Q_\nu$ is a normal divisor of $P_{\nu-1}$, and the desired determination of $Q_{\nu-1}$ is achieved by taking $Q_\nu$ itself, or, if there should still be a commutative normal divisor of lower degree in $P_{\nu-1}$, by taking the latter for $Q_{\nu-1}$.

In the other case denote the distinct groups among (7) by
\begin{equation}
Q_\nu,\;Q'_\nu,\;Q''_\nu\ldots
\tag{8}
\end{equation}
and derive from them a group
\begin{equation}
R=(Q_\nu,Q'_\nu,Q''_\nu\ldots)
\tag{9}
\end{equation}
which is defined to be the smallest group containing each of the groups (8) as a divisor, and which may be called the least common multiple of the groups $Q_\nu,Q'_\nu,Q''_\nu\ldots$.

That there is one and only one such group $R$, and that every group divisible by $Q_\nu,Q'_\nu,Q''_\nu\ldots$ must also be divisible by $R$, is easy to see. First there are finite groups, for example $P_\nu$, which contain all the groups (8) as divisors; and second, if two groups $R,R'$ contain all these groups, the same is true of the intersection of $R$ and $R'$. Thus, if $R$ has the lowest possible degree, this intersection must be identical with $R$, and so $R$ is a divisor of $R'$. If $R'$ also has the lowest possible degree, then $R'$ is not distinct from $R$.

We now prove of this group $R$:
\begin{enumerate}[label=\alph*)]
\item that it is a normal divisor of $P_{\nu-1}$,
\item that it is a commutative group.
\end{enumerate}
When both assertions have been proved, the goal is reached; for then, if no commutative normal divisor of lower degree exists in $P_{\nu-1}$, one can take $R$ itself for $Q_{\nu-1}$.

The first assertion is immediate. For if $\gamma$ denotes any element of $P_{\nu-1}$, then
\[
\gamma^{-1}R\gamma=(\gamma^{-1}Q_\nu\gamma,\gamma^{-1}Q'_\nu\gamma,\gamma^{-1}Q''_\nu\gamma\ldots)
\]
is the least common multiple of the groups
\[
\gamma^{-1}Q_\nu\gamma,
\quad \gamma^{-1}Q'_\nu\gamma,
\quad \gamma^{-1}Q''_\nu\gamma\ldots.
\]
But by the definition (7), these groups, taken all together, are not different from the groups
\[
Q_\nu,\;Q'_\nu,\;Q''_\nu\ldots;
\]
hence
\[
\gamma^{-1}R\gamma=R,
\]
that is, $R$ is a normal divisor of $P_{\nu-1}$.

To prove that $R$ is a commutative group, we must first make two auxiliary observations.

1) It must be proved that any two of the groups $Q_\nu,Q'_\nu,Q''_\nu\ldots$ have no common divisor except the identity.

It suffices to prove this property for $Q_\nu,Q'_\nu$. For if
\[
Q'_\nu=\gamma'^{-1}Q_\nu\gamma',
\qquad
Q''_\nu=\gamma''^{-1}Q_\nu\gamma''
\]
have greatest common divisor $T$, then $Q_\nu$ and $Q'''_\nu=\gamma'Q''_\nu\gamma'^{-1}$ have greatest common divisor $\gamma'T\gamma'^{-1}$.

Assume then that $T$ is the greatest common divisor of $Q_\nu$ and $Q'_\nu$. Then $T$ is certainly a commutative group, since $Q_\nu,Q'_\nu$ are such groups. But $T$ is also a normal divisor of $P_\nu$. For $Q_\nu$ and $Q'_\nu$ are such divisors; and if $\tau$ is an element of $T$ and $\pi$ an element of $P_\nu$, then $\pi^{-1}\tau\pi$ is contained both in $Q_\nu$ and in $Q'_\nu$, and hence in $T$. Therefore
\[
\pi^{-1}T\pi=T.
\]
But we have assumed that $Q_\nu$ is a commutative normal divisor of $P_\nu$ above the identity and of the least possible degree; moreover $Q_\nu$ and $Q'_\nu$ were distinct. Therefore $T$ can only be the identity group, as was to be proved.

2) To form the group $R$, one may proceed by composing elements of $Q_\nu,Q'_\nu,Q''_\nu\ldots$ with one another in arbitrary order and repetition as long as new elements arise. For all composites so obtained form a group $C$ contained in $R$, since $R$ contains the individual elements of $Q_\nu,Q'_\nu,Q''_\nu\ldots$. Conversely the groups $Q_\nu,Q'_\nu,Q''_\nu\ldots$ are also contained in $C$, and therefore $R$ is contained in $C$ and $C$ is identical with $R$.

Thus, to prove finally that $R$ is commutative, it remains only to show that, if $\alpha,\beta$ are any two elements from $Q_\nu,Q'_\nu,Q''_\nu\ldots$, then the commutativity relation
\begin{equation}
\alpha\beta=\beta\alpha
\tag{10}
\end{equation}
holds.

If $\alpha,\beta$ both belong to the same group $Q_\nu$, then (10) follows from our assumption that $Q_\nu$ is commutative; and likewise if both occur in one of the other groups $Q'_\nu,Q''_\nu\ldots$.

If, however, $\alpha$ and $\beta$ are contained in two different groups, say in $Q_\nu$ and $Q'_\nu$, we reason as follows. Since $\beta$ is contained in $P_\nu$ and $Q_\nu$ is a normal divisor of $P_\nu$, one also has
\begin{equation}
\beta^{-1}\alpha\beta=\alpha'
\tag{11}
\end{equation}
with $\alpha'$ contained in $Q_\nu$. From this it follows that
\begin{equation}
\beta\alpha'\alpha^{-1}=\alpha\beta\alpha^{-1}=\beta',
\tag{12}
\end{equation}
and since $\alpha$ is contained in $P_\nu$ and $Q'_\nu$ is a normal divisor of $P_\nu$, $\beta'$ is also contained in $Q'_\nu$. But from (12) it follows that
\[
\alpha'\alpha^{-1}=\beta^{-1}\beta'=\delta,
\]
and hence the element $\delta$ is contained both in $Q_\nu$ and in $Q'_\nu$. By what was proved under 1), it can therefore only be the identity element; that is,
\[
\alpha'=\alpha,
\qquad
\beta'=\beta.
\]
It then follows from (11) that $\alpha\beta=\beta\alpha$, as was to be proved.

\clearpage

\begin{center}
{\large Second Section.}\\[0.35em]
{\Large Abelian Groups.}
\end{center}

\section*{\S\,9. Representation of Abelian groups by a basis.}

In \S\,1 we have already called those groups in which the commutative law holds for composition commutative or Abelian groups. In these groups, which are of special importance in the applications, much simpler laws prevail than in general groups, as we shall now see. The same rules hold for composing elements in these groups as for multiplying numbers; accordingly, as in multiplication, we call the composites of elements of such a group products and powers.

Here also we restrict ourselves to the consideration of finite groups.

From the definition of Abelian groups it follows that the commutative law also holds for the composition of the parts of such a group (\S\,4). Every divisor of an Abelian group is itself an Abelian group, and is at the same time a normal divisor, so that for these groups we may omit the addition ``normal''.

A group consisting only of repetitions of one element, such as
\[
1,\ A,\ A^2,\ldots,\ A^{a-1},
\]
was already called a cyclic group earlier (Vol. I, \S\,148), and we retain this designation here. If $a$ is the least positive number for which $A^a=1$, hence $a$ is the degree of the group, then $a$ is also called the degree of the element $A$. If $m$ is any exponent for which $A^m=1$, then $m$ is a multiple of $a$. The element $1$ has degree $1$. All cyclic groups are commutative.

For every Abelian group $S$ the following fundamental theorem holds: one can always select a system of elements $A,B,C,\ldots$ of degrees $a,b,c,\ldots$ such that, in the form
\[
\Theta=A^\alpha B^\beta C^\gamma\cdots,
\]
every element $\Theta$ of $S$, and each only once, is represented, when $\alpha$ runs through a complete residue system modulo $a$, $\beta$ a complete residue system modulo $b$, $\gamma$ a complete residue system modulo $c$, and so on.

A system of elements suitable for such a representation is called a basis of the group, and the theorem to be proved may therefore be stated as follows:

\begin{enumerate}[label=\Roman*.]
\item Every Abelian group can be represented by a basis.
\end{enumerate}

The proof of the theorem is based on the following sequence of conclusions.

\begin{enumerate}[label=\arabic*.]
\item Let $n$ be the degree of the group $S$, let $A_1,A_2,\ldots,A_n$ be its elements and $a_1,a_2,\ldots,a_n$ their degrees; let also $\alpha_1,\alpha_2,\ldots,\alpha_n$ run through complete residue systems modulo $a_1,a_2,\ldots,a_n$. Then in the form
\begin{equation}
\Theta=A_1^{\alpha_1}A_2^{\alpha_2}\cdots A_n^{\alpha_n}
\tag{1}
\end{equation}
every element of $S$ is represented, and each equally often.
\end{enumerate}

That every element is obtained in the form (1) at all is immediate; for example, one obtains $A_1$ by putting $\alpha_1=1$, $\alpha_2=\cdots=\alpha_n=0$.

It remains to prove that every element is obtained equally often. If
\begin{equation}
1=A_1^{\lambda_1}A_2^{\lambda_2}\cdots A_n^{\lambda_n}
\tag{2}
\end{equation}
is a representation of the element $1$, then $\Theta$ is not changed if $\alpha_1,\alpha_2,\ldots,\alpha_n$ in (1) are replaced by $\alpha_1+\lambda_1$, $\alpha_2+\lambda_2$, \ldots, $\alpha_n+\lambda_n$. Thus every element $\Theta$ is represented by (1) at least as many times as the element $1$ is represented by (2).

Conversely, if the two rows of exponents $\alpha_1,\alpha_2,\ldots,\alpha_n$ and $\alpha'_1,\alpha'_2,\ldots,\alpha'_n$ give two representations of the element $\Theta$, then
\begin{equation}
1=A_1^{\alpha_1-\alpha'_1}A_2^{\alpha_2-\alpha'_2}\cdots A_n^{\alpha_n-\alpha'_n},
\tag{3}
\end{equation}
and hence a representation of the element $1$. Therefore, according to (2), we may put
\[
\alpha_1=\alpha'_1+\lambda_1,
\quad \alpha_2=\alpha'_2+\lambda_2,
\quad \ldots,
\quad \alpha_n=\alpha'_n+\lambda_n.
\]
Consequently there cannot be more distinct representations of the element $\Theta$ than the number of representations of the element $1$. We denote the number of these representations by $h$. Altogether the number of possible choices of the $\alpha$ in (1) is the product $a_1a_2\cdots a_n$, and since $n$ is the number of elements of $S$, it follows that
\begin{equation}
nh=a_1a_2\cdots a_n.
\tag{4}
\end{equation}

Formula (4) gives the theorem:

\begin{enumerate}[resume,label=\arabic*.]
\item If $r$ is a prime number dividing the degree $n$ of the group, then there is an element of degree $r$ in $S$.
\end{enumerate}

For from (4) one first sees that one of the degrees $a_1,a_2,\ldots,a_n$ is divisible by $r$. If, say, $a_1=rk$, then $A_1^k$ is an element of degree $r$.

\begin{enumerate}[resume,label=\arabic*.]
\item If $A,B,\ldots$ are arbitrary elements of $S$ of degrees $a,b,\ldots$, then $AB\ldots$ is also an element of $S$, and the degree of this product is a divisor of the least common multiple $m$ of $a,b,\ldots$.
\end{enumerate}

For
\[
(AB\cdots)^m=A^mB^m\cdots=1,
\]
and hence $m$ is a multiple of the degree of $AB\cdots$.

\begin{enumerate}[resume,label=\arabic*.]
\item If the degree $n$ of the group $S$ is decomposed into two factors $n=ab$ such that $a$ and $b$ are relatively prime, then in $S$ there are exactly $a$ elements $A$ whose degrees divide $a$, and $b$ elements $B$ whose degrees divide $b$; and in the form
\begin{equation}
\Theta=AB
\tag{5}
\end{equation}
all elements of $S$, and each only once, are contained.
\end{enumerate}

To see this, make the following consideration. The totality $\mathfrak A$ of the elements $A$ whose degree is a divisor of $a$ is a group contained in $S$; for if $A,A'$ are two such elements, then by 3. the degree of $AA'$ is a divisor of $a$, and therefore $AA'$ also belongs to $\mathfrak A$.

Similarly, the totality $\mathfrak B$ of the elements $B$ whose degree is a divisor of $b$ is a group. The element $1$ occurs both in $\mathfrak A$ and in $\mathfrak B$, but otherwise the two have no element in common.

The degree $a'$ of $\mathfrak A$ is relatively prime to $b$. For if a prime $r$ divides $a'$, then by 2. there is in $\mathfrak A$ an element of degree $r$. But since the degree of every element of $\mathfrak A$ divides $a$, $r$ is also a divisor of $a$ and not of $b$. In the same way one proves that the degree $b'$ of $\mathfrak B$ is relatively prime to $a$.

Now determine two integers $x$ and $y$ (according to Vol. I, \S\,118) such that
\begin{equation}
ax+by=1
\tag{6}
\end{equation}
and take any element $\Theta$ of $S$. Then
\begin{equation}
\Theta=\Theta^{ax}\Theta^{by},
\tag{7}
\end{equation}
and since $(\Theta^{by})^a=1$, $\Theta^{by}$ belongs to $\mathfrak A$, and for the same reason $\Theta^{ax}$ belongs to $\mathfrak B$. Thus from (7) follows
\begin{equation}
\Theta=AB.
\tag{8}
\end{equation}

Accordingly every element $\Theta$ is contained in the form $AB$. Such a representation is possible in only one way. For if $A',B'$ are two elements of $\mathfrak A$ and $\mathfrak B$, and
\[
AB=A'B',
\]
then, on raising to the power $by=1-ax$, one obtains $A=A'$ and consequently also $B=B'$. The number of distinct products of the form $AB$ is $a'b'$, and hence
\[
n=ab=a'b',
\]
and since $a$ is relatively prime to $b'$ and $b$ relatively prime to $a'$,
\[
a'=a,
\qquad
b'=b.
\]
This proves statement 4 in all its parts.

If the two groups $\mathfrak A$ and $\mathfrak B$ are represented by bases, formula (8) shows that $S$ is also represented by a basis, and the basis of $S$ contains the elements of the bases of $\mathfrak A$ and $\mathfrak B$ and no others.

If $a$ and $b$ can be further decomposed into factors which are pairwise relatively prime, we may proceed in the same way with $\mathfrak A$ and $\mathfrak B$. Thus our Theorem I is proved in general as soon as it is proved for groups whose degree is a power of a prime number.

Let now the degree $n$ of $S$ be a power of a prime $p$,
\begin{equation}
n=p^\mu.
\tag{9}
\end{equation}
The degrees of all elements of $S$, being divisors of $n$, must also be powers of $p$. We must prove that such a group can be represented by a basis.

Choose in $S$ an element $A$ of the largest possible degree $a$. Then $a$ is a power of $p$, and the degrees of all other elements divide $a$, so that for every element $\Theta$ of $S$
\begin{equation}
\Theta^a=1.
\tag{10}
\end{equation}
The elements
\begin{equation}
1,A,A^2,\ldots,A^{a-1}
\tag{11}
\end{equation}
are all distinct, and their totality is a divisor of $S$. If these already exhaust $S$, that is, if every element is of the form $A^\alpha$, then $S$ is represented by a one-element basis.

If, however, $S$ is not exhausted by (11), then for each element $\Theta$ of $S$ there will still be some exponent $h$ such that $\Theta^h$ is contained in the series (11). This certainly occurs when $h$ is the degree of $\Theta$, so that $\Theta^h=1$. Among such numbers $h$ there is a least positive one, which we denote by $b$. Thus for every element $\Theta$ there is a certain least positive number $b$ such that
\begin{equation}
\Theta^b=A^\lambda
\tag{12}
\end{equation}
occurs in the series (11).

This number $b$ is a divisor of $a$, hence also a power of $p$, and at the same time a divisor of $\lambda$. For if $a=qb+b'$, where $0\le b'<b$, then from (10) and (12) it follows that
\[
1=\Theta^a=\Theta^{qb+b'}=A^{q\lambda}\Theta^{b'},
\]
or $\Theta^{b'}=A^{-q\lambda}$; if $b'$ were not zero, there would, contrary to the assumption, be a still smaller number than $b$, namely $b'$, satisfying (12).

From (12) one further obtains
\[
1=\Theta^a=A^{\lambda a/b};
\]
therefore $\lambda a:b$ must be a multiple of $a$, and consequently $\lambda$ a multiple of $b$. Thus if we put
\begin{equation}
B=\Theta A^{-\lambda/b},
\tag{13}
\end{equation}
then $B^b=1$, and at the same time $B^b$ is the lowest power of $B$ equal to a power of $A$, because if $B^{b'}$ is a power of $A$, then by (13) the same is true of $\Theta^{b'}$. We choose the element $\Theta$ so that $b$ is as large as possible. Then for every other element $\Theta_1$ of $S$, $\Theta_1^b$ is also a power of $A$ whose exponent is divisible by $b$.

Now let
\[
\alpha=0,1,2,\ldots,a-1,
\qquad
\beta=0,1,2,\ldots,b-1.
\]
Then the elements
\begin{equation}
A^\alpha B^\beta
\tag{14}
\end{equation}
are all distinct in number $ab$. They form a divisor $S'$ contained in $S$, represented by the basis $A,B$. If $S'$ is identical with $S$, we have reached the goal.

If $S$ is not yet exhausted by (14), we continue in the same way.

We shall now argue in general by complete induction. Suppose a divisor $S_{\nu-1}$ of $S$ has been found, represented by a basis in the form
\begin{equation}
S_{\nu-1}=A_1^{\alpha_1}A_2^{\alpha_2}\cdots A_{\nu-1}^{\alpha_{\nu-1}},
\tag{15}
\end{equation}
and suppose it satisfies the following assumptions:

1) The degrees of $A_1,A_2,
\ldots,A_{\nu-1}$ are $a_1,a_2,
\ldots,a_{\nu-1}$, and
\[
a_1\ge a_2\ge\cdots\ge a_{\nu-1}.
\]

2) For every element $\Theta$ contained in $S$,
\begin{equation}
\Theta^{a_{\nu-1}}=A_1^{\lambda_1}A_2^{\lambda_2}\cdots A_{\nu-2}^{\lambda_{\nu-2}},
\tag{16}
\end{equation}
that is, it belongs to $S_{\nu-2}$, and the exponents $\lambda_1,\lambda_2,
\ldots,
\lambda_{\nu-2}$ are divisible by $a_{\nu-1}$.

The group $S'$ constructed above satisfies these requirements when $\nu=3$ and $a_{\nu-1}=b$. It remains to show how, if $S_{\nu-1}$ does not yet exhaust $S$, one can derive from it a similarly constituted larger group $S_\nu$.

For every element $\Theta$ of $S$ there is a certain least positive exponent $a_\nu$ for which $\Theta^{a_\nu}$ is contained in $S_{\nu-1}$; by (16) this $a_\nu$ is equal to or less than $a_{\nu-1}$, and is moreover a power of $p$, since it must be a divisor of the degree of $\Theta$. We choose $\Theta$ so that the exponent $a_\nu$ is as large as possible. If $\Theta_1$ is any other element in $S$ and $\Theta_1^h$ is the lowest power of $\Theta_1$ contained in $S_{\nu-1}$, then $h$ is also a power of $p$ and is equal to or less than $a_\nu$. Hence $\Theta_1^{a_\nu}$ is contained in $S_{\nu-1}$.

Put
\begin{equation}
\Theta^{a_\nu}=A_1^{\lambda_1}A_2^{\lambda_2}\cdots A_{\nu-1}^{\lambda_{\nu-1}}.
\tag{17}
\end{equation}
Then all exponents $\lambda$ are divisible by $a_\nu$. For
\[
\Theta^{a_{\nu-1}}=A_1^{\lambda_1 a_{\nu-1}/a_\nu}
A_2^{\lambda_2 a_{\nu-1}/a_\nu}
\cdots
A_{\nu-1}^{\lambda_{\nu-1}a_{\nu-1}/a_\nu},
\]
and by assumption 2) the exponents of $A_1,A_2,
\ldots,A_{\nu-1}$ on the right-hand side must be integers divisible by $a_{\nu-1}$. Consequently $\lambda_1,
\lambda_2,
\ldots,
\lambda_{\nu-1}$ are divisible by $a_\nu$. Returning to the special $\Theta$ considered above, and denoting integers by $h_1,h_2,
\ldots,h_{\nu-1}$, we may therefore write
\[
\Theta^{a_\nu}=A_1^{h_1a_\nu}A_2^{h_2a_\nu}\cdots A_{\nu-1}^{h_{\nu-1}a_\nu}.
\]
Then, if we put
\[
A_\nu=\Theta A_1^{-h_1}A_2^{-h_2}\cdots A_{\nu-1}^{-h_{\nu-1}},
\]
we have
\begin{equation}
A_\nu^{a_\nu}=1,
\tag{18}
\end{equation}
and this is the lowest power of $A_\nu$ contained in $S_{\nu-1}$, since otherwise a still lower power of $\Theta$ would be contained in $S_{\nu-1}$.

Then
\begin{equation}
A_1^{\alpha_1}A_2^{\alpha_2}\cdots A_\nu^{\alpha_\nu},
\qquad
\alpha_i=0,1,2,\ldots,a_i-1,
\tag{19}
\end{equation}
is equal to $1$ only when $\alpha_1,\alpha_2,
\ldots,
\alpha_\nu$ are divisible by $a_1,a_2,
\ldots,a_\nu$, and hence are zero. Therefore the elements represented in (19) are all distinct. They form a group $S_\nu$ with basis $A_1,A_2,
\ldots,A_\nu$, satisfying requirement 1).

At the same time, as formula (17) shows, if $\Theta_1$ is any element in $S$, then $\Theta_1^{a_\nu}$ is contained in $S_{\nu-1}$, and the exponents $\lambda_1,
\lambda_2,
\ldots,
\lambda_{\nu-1}$ are divisible by $a_\nu$. Thus requirement 2) is also satisfied.

This proves Theorem I. At the same time this derivation shows that, for an arbitrary Abelian group $S$, the elements of a basis may always be chosen so that their degrees are prime powers.

From representability by a basis it also follows that every Abelian group is metacyclic. For if the elements of such a group $S$ are represented in the form
\[
A=A_1^{\alpha_1}A_2^{\alpha_2}\cdots A_\nu^{\alpha_\nu},
\]
and if $p$ is a prime dividing the degree $a_1$ of $A_1$, then the elements
\[
A'=A_1^{p\alpha_1}A_2^{\alpha_2}\cdots A_\nu^{\alpha_\nu},
\]
with $\alpha_1$ running through a complete residue system modulo $a_1:p$, form a divisor $S'$ of $S$ of index $p$, represented by the basis $A'_1,A_2,
\ldots,A_\nu$. From $S'$ one can again find, in the same manner, a divisor of prime index, and so on. Thus $S$ is metacyclic (\S\,8).

\section*{\S\,10. The invariants of Abelian groups.}

The proof given in the preceding paragraph for the possibility of representing an Abelian group by a basis also contains a way to find such a basis, and indeed one in which the degrees of the basis elements are prime powers. Nevertheless it can happen that one and the same group may be represented by bases in different ways.

For example, consider two elements $A,B$ whose degrees $a,b$ are relatively prime. As elements of a two-element basis they represent a group
\begin{equation}
A^\alpha B^\beta,
\qquad
\alpha=0,1,2,\ldots,a-1,
\qquad
\beta=0,1,2,\ldots,b-1.
\tag{1}
\end{equation}
The same group can also be represented by the one-element basis $AB$ in the form
\begin{equation}
(AB)^s,
\qquad
s=0,1,2,\ldots,ab-1.
\tag{2}
\end{equation}
For if $\alpha$ and $\beta$ are arbitrarily given, then $s$ can be chosen modulo $ab$ so that
\[
s\equiv\alpha\pmod a,
\qquad
s\equiv\beta\pmod b,
\]
which makes (2) identical with (1). In a certain sense, however, the constitution of the basis is nevertheless completely determined by the nature of the groups, according to the following theorem:

\begin{enumerate}[label=\Roman*.,start=2]
\item Let
\[
A_1,A_2,\ldots,A_\nu,
\qquad
B_1,B_2,\ldots,B_\mu
\]
with degrees
\[
a_1,a_2,\ldots,a_\nu,
\qquad
b_1,b_2,\ldots,b_\mu
\]
be two bases of an Abelian group $S$ of degree $n$. Let $p$ be a prime divisor of $n$ and
\begin{align}
a_1&=p_1a'_1,\quad a_2=p_2a'_2,
\quad\ldots\quad
 a_\nu=p_\nu a'_\nu,
\tag{1}\\
b_1&=p'_1b'_1,
\quad b_2=p'_2b'_2,
\quad\ldots\quad
 b_\mu=p'_\mu b'_\mu,
\tag{2}
\end{align}
where $p_1,p_2,
\ldots,p_\nu,p'_1,p'_2,
\ldots,p'_\mu$ are the highest powers of $p$ dividing $a_1,a_2,
\ldots,a_\nu,b_1,b_2,
\ldots,b_\mu$. Then all prime powers $p_1,p_2,
\ldots,p_\nu$ greater than $1$ also occur among $p'_1,p'_2,
\ldots,p'_\mu$, and conversely.
\end{enumerate}

To prove this theorem, arrange the elements of the two bases $A$ and $B$ so that
\begin{equation}
p_1\ge p_2\ge\cdots\ge p_\nu,
\qquad
p'_1\ge p'_2\ge\cdots\ge p'_\mu.
\tag{3}
\end{equation}
The integers $a'_1,a'_2,
\ldots,a'_\nu,b'_1,b'_2,
\ldots,b'_\mu$ occurring in (1) and (2) are, by definition, not divisible by $p$. If $m$ denotes the least common multiple of $a'_1,a'_2,
\ldots,a'_\nu$, then $m$ is also not divisible by $p$. If
\begin{equation}
\Theta=A_1^{\alpha_1}A_2^{\alpha_2}\cdots A_\nu^{\alpha_\nu}
\tag{4}
\end{equation}
is an arbitrary element of $S$, it follows that
\[
\Theta^{p_1m}=1.
\]
Putting $B_1,B_2,
\ldots,B_\mu$ in place of $\Theta$, we get that $p_1m$ is divisible by each of $b_1,b_2,
\ldots,b_\mu$. Thus $p_1$ is divisible by $p'_1$ and $m$ by the least common multiple of $b'_1,b'_2,
\ldots,b'_\mu$. In this argument we may interchange $A$ and $B$ throughout. Hence $p'_1$ must also be divisible by $p_1$, and therefore
\begin{equation}
p_1=p'_1,
\tag{5}
\end{equation}
and it also follows that $m$ is the least common multiple of $b'_1,b'_2,
\ldots,b'_\mu$ as well.

To prove the theorem in general from this, assume that for some $s$ it has been proved that
\begin{equation}
p_1=p'_1,
\quad p_2=p'_2,
\quad\ldots\quad
p_{s-1}=p'_{s-1}.
\tag{6}
\end{equation}
By (4), for every element $\Theta$,
\begin{equation}
\Theta^{p_sm}
=A_1^{\alpha_1p_sm}A_2^{\alpha_2p_sm}\cdots A_{s-1}^{\alpha_{s-1}p_sm},
\tag{7}
\end{equation}
where $m$ is, as above, the least common multiple both of $a'_1,a'_2,
\ldots,a'_\nu$ and of $b'_1,b'_2,
\ldots,b'_\mu$, and hence is not divisible by $p$.

We now determine the number of distinct elements contained in the form (7). If $\alpha_1$ runs through the numbers $0,1,2,
\ldots,p_1/p_s-1$, then the elements
\begin{equation}
1,
A_1^{p_sm},
A_1^{2p_sm},
\ldots,
A_1^{\left(\frac{p_1}{p_s}-1\right)p_sm}
\tag{8}
\end{equation}
are all distinct, while $A_1^{(p_1/p_s)p_sm}$ is again equal to $1$. Thus all other powers $A_1^{\alpha_1p_sm}$ occur again in the series (8). The same reasoning applies to the other factors of (7), giving the exact number of distinct elements contained in $\Theta^{p_sm}$ as
\begin{equation}
z=\frac{p_1}{p_s}\frac{p_2}{p_s}\cdots\frac{p_{s-1}}{p_s}.
\tag{9}
\end{equation}

Of the two numbers $p_s,p'_s$, if they are unequal, one is the greater. We shall assume that
\begin{equation}
p'_s>p_s.
\tag{10}
\end{equation}
Now express $\Theta$ in terms of the basis $B$ and put
\begin{equation}
\Theta^{p_sm}=B_1^{\beta_1p_sm}B_2^{\beta_2p_sm}\cdots B_\mu^{\beta_\mu p_sm}.
\tag{11}
\end{equation}
We count again how many distinct elements are contained in this form. The two values for $z$ must agree.

Counting as above in (7), the number of distinct elements contained in the form
\begin{equation}
B_1^{\beta_1p_sm}B_2^{\beta_2p_sm}\cdots B_{s-1}^{\beta_{s-1}p_sm}
\tag{12}
\end{equation}
is, by assumption (6), also $z$. Therefore the form
\begin{equation}
B_s^{\beta_sp_sm}\cdots B_\mu^{\beta_\mu p_sm}
\tag{13}
\end{equation}
can contain only the single element $1$, from which it follows that $p_s$ must be divisible by $p'_s$. This is compatible with (10) only if
\begin{equation}
p_s=p'_s.
\tag{14}
\end{equation}
This proves Theorem II.

The prime powers contained in the degree numbers of a basis of $S$ are therefore wholly independent of the particular choice of the basis, and we call them the invariants of the group. The product of all invariants is equal to the degree $n$ of the group.

By \S\,9 we can choose a basis for $S$ in which the degrees of the elements are all prime powers. By Theorem II these degrees are the invariants of the group and are therefore completely determined by the group.

That the invariants also determine the nature of the group completely follows from the theorem:

\begin{enumerate}[label=\Roman*.,start=3]
\item Two groups with the same invariants are isomorphic, and isomorphic groups have the same invariants.
\end{enumerate}

For if two groups $S$ and $S'$ have basis elements corresponding in number and degree, say if
\[
\Theta=A_1^{\alpha_1}A_2^{\alpha_2}\cdots A_\nu^{\alpha_\nu}
\]
are the elements of $S$, and
\[
\Theta'=B_1^{\alpha_1}B_2^{\alpha_2}\cdots B_\nu^{\alpha_\nu}
\]
are the elements of $S'$, where in both cases $\alpha_1,\alpha_2,
\ldots,
\alpha_\nu$ run through residue systems modulo $a_1,a_2,
\ldots,a_\nu$, then it suffices to let $\Theta$ and $\Theta'$ correspond when the exponents $\alpha$ are the same in both. Then the two groups are placed in isomorphic correspondence.

Thus if two groups have the same invariants, these invariants may in both groups be made the degrees of the basis elements, and we are in the case just described.

Conversely, if two groups are isomorphic, then the elements of the other group corresponding to the basis elements of the one form a basis of the latter, and therefore the invariants must also be the same in both.\footnote{On the theory of Abelian groups compare: Gauss, Demonstration de quelques theoremes concernant les periodes des classes des formes binaires du second degre. Werke, Vol. II, p. 266. Schering, Die Fundamentalclassen der zusammensetzbaren arithmetischen Formen, Göttinger Abhandlungen, Vol. 14. Frobenius and Stickelberger, Ueber Gruppen vertauschbarer Elemente. Crelle's Journal, Vol. 86, p. 217. Weber, ``Theorie der Abel'schen Zahlkörper'', Acta mathematica, Vols. 8 and 9. The concept of invariants was first introduced in the paper of Frobenius and Stickelberger, but was defined somewhat differently there than here. The author followed that older definition in the cited paper in Acta mathematica; for application to cyclotomic numbers it proves less suitable.}

\section*{\S\,11. Group characters.}

A principal problem in the theory of Abelian groups is to find all divisors of an Abelian group. The solution of this problem is considerably facilitated by introducing the concept of group characters, which is also important in many other respects.

Let $S$ be an Abelian group of degree $n$, whose elements are denoted by $A$, and let
\begin{equation}
A_1,A_2,\ldots,A_\nu
\tag{1}
\end{equation}
be the elements of a basis of $S$ of degrees $a_1,a_2,
\ldots,a_\nu$. Every element $A$ is therefore, and in only one way, representable in the form
\begin{equation}
A=A_1^{\alpha_1}A_2^{\alpha_2}\cdots A_\nu^{\alpha_\nu},
\tag{2}
\end{equation}
where
\begin{equation}
0\le \alpha_1<a_1,
\quad 0\le \alpha_2<a_2,
\quad\ldots\quad
0\le \alpha_\nu<a_\nu.
\tag{3}
\end{equation}
The exponents $\alpha_1,\alpha_2,
\ldots,
\alpha_\nu$, or any numbers congruent to them modulo $a_1,a_2,
\ldots,a_\nu$, are called the indices of the element $A$, and the following theorem holds:

\begin{enumerate}[label=\arabic*.]
\item The indices of a composite $AA'$ are obtained by adding the corresponding indices of the two factors.
\end{enumerate}

If to every element $A$ we assign in any way a numerical value, we may call this assignment a function of $A$. Such a function $\chi(A)$ shall now be called a group character if for any two elements $A,A'$ of $S$ it satisfies
\begin{equation}
\chi(AA')=\chi(A)\chi(A'),
\tag{4}
\end{equation}
and consequently also the more general relation
\[
\chi(AA'A''\cdots)=\chi(A)\chi(A')\chi(A'')\cdots.
\]
Two such functions $\chi$ and $\chi_1$ are regarded as different if there is at least one element $A$ for which $\chi(A)$ differs from $\chi_1(A)$. We now determine these characters more closely.

Putting $A'=1$, (4) gives
\begin{equation}
\chi(A)=\chi(A)\chi(1),
\tag{5}
\end{equation}
and if we exclude once and for all the case in which all $\chi(A)=0$,
\begin{equation}
\chi(1)=1,
\tag{6}
\end{equation}
which is the first theorem:

\begin{enumerate}[resume,label=\arabic*.]
\item Every character has value $1$ at the identity element.
\end{enumerate}

Putting further $A'=A$, one obtains from (4)
\[
\chi(A^2)=[\chi(A)]^2,
\]
and hence, by complete induction, for every exponent $h$,
\begin{equation}
\chi(A^h)=[\chi(A)]^h.
\tag{7}
\end{equation}
If $a$ denotes the degree of the element $A$, then, putting $h=a$ in (7) and using (6), one obtains
\begin{equation}
[\chi(A)]^a=1,
\tag{8}
\end{equation}
and hence:

\begin{enumerate}[resume,label=\arabic*.]
\item The values of a character $\chi(A)$ are roots of unity whose degree divides the degree of $A$.
\end{enumerate}

Thus all characters are easily determined. Put
\begin{equation}
\chi(A_1)=\omega_1,
\quad \chi(A_2)=\omega_2,
\quad\ldots\quad
\chi(A_\nu)=\omega_\nu.
\tag{9}
\end{equation}
Then
\begin{equation}
\omega_1^{a_1}=1,
\quad \omega_2^{a_2}=1,
\quad\ldots\quad
\omega_\nu^{a_\nu}=1,
\tag{10}
\end{equation}
and by (2) and (4)
\begin{equation}
\chi(A)=\omega_1^{\alpha_1}\omega_2^{\alpha_2}\cdots\omega_\nu^{\alpha_\nu}.
\tag{11}
\end{equation}
Conversely, if $\omega_1,\omega_2,
\ldots,
\omega_\nu$ denotes any solution of the equations (10), then (11) defines a function of $A$ which, by 1., satisfies (4) and is therefore a character of the group.

Each of the equations (10) has $a_1,a_2,
\ldots,a_\nu$ roots, and combining each with each gives
\[
a_1a_2\cdots a_\nu=n
\]
combinations. All these combinations lead to different characters $\chi(A)$. For let $\omega'_1,\omega'_2,
\ldots,
\omega'_\nu$ be a second combination of roots of (10), and suppose that for all elements $A$
\[
\omega_1^{\alpha_1}\omega_2^{\alpha_2}\cdots\omega_\nu^{\alpha_\nu}
=(\omega'_1)^{\alpha_1}(\omega'_2)^{\alpha_2}\cdots(\omega'_\nu)^{\alpha_\nu}.
\]
Putting $\alpha_1=1$, $\alpha_2=\cdots=\alpha_\nu=0$, gives $\omega_1=\omega'_1$, and similarly $\omega_2=\omega'_2$, \ldots, $\omega_\nu=\omega'_\nu$. Hence:

\begin{enumerate}[resume,label=\arabic*.]
\item There are exactly $n$ different characters of an Abelian group of degree $n$, all represented by formula (11).
\end{enumerate}

Let $\varepsilon_1,\varepsilon_2,
\ldots,\varepsilon_\nu$ be a system of primitive roots of the equations (10). Then we may put
\begin{equation}
\omega_1=\varepsilon_1^{\beta_1},
\quad \omega_2=\varepsilon_2^{\beta_2},
\quad\ldots\quad
\omega_\nu=\varepsilon_\nu^{\beta_\nu},
\tag{12}
\end{equation}
where $\beta_1,\beta_2,
\ldots,
\beta_\nu$, just as $\alpha_1,\alpha_2,
\ldots,
\alpha_\nu$, are each taken from a complete residue system modulo $a_1,a_2,
\ldots,a_\nu$. With fixed $\varepsilon_1,\varepsilon_2,
\ldots,
\varepsilon_\nu$ all $n$ group characters are obtained in the form
\begin{equation}
\chi(A)=\varepsilon_1^{\alpha_1\beta_1}
\varepsilon_2^{\alpha_2\beta_2}
\cdots
\varepsilon_\nu^{\alpha_\nu\beta_\nu}.
\tag{13}
\end{equation}

Among these characters is the one obtained by putting $\beta_1=0$, $\beta_2=\cdots=\beta_\nu=0$, which has value $1$ for all elements $A$. We shall call this the identity character.

The $n$ characters of $S$ can now themselves be combined into an Abelian group, in fact one isomorphic with $S$.

By (13), each of the $n$ characters is determined by a number system $\beta_1,\beta_2,
\ldots,
\beta_\nu$ taken modulo $a_1,a_2,
\ldots,a_\nu$. This system of $\beta$ is at the same time the system of indices of a definite element $B$ of $S$, namely
\begin{equation}
B=A_1^{\beta_1}A_2^{\beta_2}\cdots A_\nu^{\beta_\nu},
\tag{14}
\end{equation}
so that to each element $B$ of $S$ there corresponds a definite character, which we denote by $\chi_B(A)$, or, omitting $A$, by $\chi_B$. This assignment of characters $\chi$ to elements $B$ changes if another basis is used, or if the roots of unity $\varepsilon_1,\varepsilon_2,
\ldots,
\varepsilon_\nu$ are chosen differently.

If $B'$ is a second element of $S$ with indices $\beta'_1,
\beta'_2,
\ldots,
\beta'_\nu$, then in the same way we obtain the character $\chi_{B'}$. If $\chi_{BB'}$ denotes the character determined for every $A$ by
\begin{equation}
\chi_{BB'}(A)=
\varepsilon_1^{\alpha_1(\beta_1+\beta'_1)}
\varepsilon_2^{\alpha_2(\beta_2+\beta'_2)}
\cdots
\varepsilon_\nu^{\alpha_\nu(\beta_\nu+\beta'_\nu)},
\tag{15}
\end{equation}
then the characters $\chi$ are thereby combined into a group isomorphic with $S$. The identity element in this character group is the identity character.

Accordingly the group of characters can also be represented by a basis, obtained by putting
\begin{equation}
\chi_1(A)=\varepsilon_1^{\alpha_1},
\quad \chi_2(A)=\varepsilon_2^{\alpha_2},
\quad\ldots\quad
\chi_\nu(A)=\varepsilon_\nu^{\alpha_\nu}.
\tag{16}
\end{equation}
Then every character is contained in the form
\begin{equation}
\chi_B=\chi_1^{\beta_1}\chi_2^{\beta_2}\cdots\chi_\nu^{\beta_\nu},
\tag{17}
\end{equation}
and if $B,B'$ are two elements in $S$,
\begin{equation}
\chi_B\chi_{B'}=\chi_{BB'}.
\tag{18}
\end{equation}
If $\chi_1,\chi_2$ are any two characters and $\chi_1\chi_2$ is their composite, then by (15), for every element $A$,
\begin{equation}
\chi_1(A)\chi_2(A)=(\chi_1\chi_2)(A).
\tag{19}
\end{equation}

We state the result obtained as a theorem:

\begin{enumerate}[resume,label=\arabic*.]
\item The characters of a group $S$ can be combined into a group isomorphic with $S$.
\end{enumerate}

We also mention the theorem following immediately from the definition of $\chi_B$ by formula (13):
\begin{equation}
\chi_B(A)=\chi_A(B).
\tag{20}
\end{equation}

Finally, the following theorem also holds:

\begin{enumerate}[resume,label=\arabic*.]
\item If $A$ runs through all elements of the group $S$, then for a fixed character $\chi$
\begin{equation}
\sum_A\chi(A)=n\quad\hbox{or}\quad =0,
\tag{21}
\end{equation}
according as $\chi$ is the identity character or not. Likewise, if the element $A$ is fixed and $\chi$ runs through all characters,
\begin{equation}
\sum_\chi \chi(A)=n\quad\hbox{or}\quad =0,
\tag{22}
\end{equation}
according as $A$ is the identity element or not.
\end{enumerate}

When $\chi$ or $A$ is the identity element of its group, (21) and (22) are evident, since then each of the $n$ terms in the sum has value $1$. If, however, $\chi$ is not the identity character, then there is an element $B$ in $S$ such that $\chi(B)$ is not $1$. Put
\[
\sigma=\sum_A\chi(A).
\]
Multiplication by $\chi(B)$ gives
\[
\sum_A\chi(AB)=\sigma\chi(B).
\]
But $AB$ runs through the entire group $S$ together with $A$, and hence also $\sum_A\chi(AB)=\sigma$. Thus $\sigma\{1-\chi(B)\}=0$, whence $\sigma=0$, as was to be proved. Formula (22) is proved similarly; it also follows immediately from (21) by (20).

\section*{\S\,12. Divisors of an Abelian group. Reciprocal groups.}

A divisor $T$ of an Abelian group $S$ is, as already noted, always a normal divisor, and therefore by \S\,4 there is a group complementary to $T$,
\[
S/T,
\]
whose elements are the cosets of $T$, namely
\begin{equation}
T,
\ TA',
\ TA'',
\ldots,
\tag{1}
\end{equation}
where $A',A'',\ldots$ are certain elements of $S$. The number of elements (1), that is, the degree of the group $S/T$, is, if $t$ is the degree of $T$, equal to the index of the divisor $T$ in $S$, hence
\[
j=\frac nt.
\]
This group $S/T$ is itself Abelian; for
\[
TA'\,TA''=TA'A'',
\qquad
TA''\,TA'=TA''A',
\]
and these are equal.

We now determine the characters of the group $S/T$. Denote the elements of this group by
\begin{equation}
T,
\ T_1,
\ T_2,
\ldots,
T_{j-1},
\tag{2}
\end{equation}
and one of its characters by $\xi(T_i)$.

From this function $\xi(T_i)$ we may derive a function $\xi(A)$ of the elements of $S$ by putting
\begin{equation}
\xi(A_i)=\xi(T_i)
\tag{3}
\end{equation}
if $A_i$ is any element occurring in the coset $T_i$. For the elements of the group $T$ itself, then, $\xi(A)=1$.

This function $\xi(A_i)$ is contained among the characters of $S$. For if $A_i$ lies in $T_i$ and $A_k$ in $T_k$, then $A_iA_k$ lies in $T_iT_k$, and hence
\begin{equation}
\xi(A_i)\xi(A_k)=\xi(T_i)\xi(T_k)=\xi(T_iT_k)=\xi(A_iA_k).
\tag{4}
\end{equation}
But by \S\,11, (4) this is precisely the definition of a character of $S$.

Conversely, if one of the characters $\xi$ of the group $S$ has the same value on all elements of the group $T$, this value can only be $1$, since the identity element of $S$ certainly occurs in $T$ (\S\,11, 2.). If $A_i$ is any element of $T_i$, and $A$ runs through the entire group $T$, then $AA_i$ runs through the coset $T_i$. But then $\xi(A)=1$, and so
\begin{equation}
\xi(AA_i)=\xi(A_i),
\tag{5}
\end{equation}
that is, $\xi(A)$ has one and the same value for all elements of a coset $T_i$, and may therefore be regarded as a function of $T_i$ and denoted by $\xi(T_i)$.

Since, if $A_i$ lies in $T_i$ and $A_k$ in $T_k$, then $A_iA_k$ lies in $T_iT_k$, it follows that
\begin{equation}
\xi(T_i)\xi(T_k)=\xi(T_iT_k),
\tag{6}
\end{equation}
that is, $\xi(T_i)$ is among the characters of the group $S/T$. Thus there are exactly $j$ characters $\xi(A)$ defined by the property that they have value $1$ on every element of $T$. These $j$ characters $\xi$ form a group under composition of characters; for if $\xi_1(A)=1$ and $\xi_2(A)=1$, then also $\xi_1\xi_2(A)=1$ [\S\,11, (18)]. Since, by (4), the functions $\xi$ can also be regarded as the characters of the group $S/T$, the group of the $\xi$ is isomorphic with $S/T$ by \S\,11, 5.

We have thus proved the following theorem:

\begin{enumerate}[label=\arabic*.,start=7]
\item If an Abelian group $S$ has a divisor $T$ of degree $t$ and index $j$, then among the characters of $S$ there are exactly $j$, and no more, which have value $1$ for all elements of $T$, while for each element of $S$ not contained in $T$ at least one of them differs from $1$. These characters form a group isomorphic with $S/T$.
\end{enumerate}

We shall call these characters the characters belonging to the group $T$.

Since every character $\chi_B$ corresponds to a definite element $B$ of $S$, when $\chi_B$ runs through the group of characters belonging to $T$, the element $B$, by formula \S\,11, (18), also runs through a group of degree $j$, isomorphic with the group of the $\chi_B$ and hence also with $S/T$. We denote this group of the $B$ by $U$ and call it the group reciprocal to $T$. Here too, however, it must be noted that the concept of reciprocal group depends in general on the choice of the basis and of the roots of unity $\varepsilon$, and therefore does not belong to the groups $S$ and $T$ as such.

The group reciprocal to $T$ is characterized by the theorem:

\begin{enumerate}[resume,label=\arabic*.]
\item Let $A$ run through the elements of a divisor $T$ of $S$, and seek all elements $B$ of $S$ which for every $A$ satisfy
\begin{equation}
\chi_B(A)=1.
\tag{7}
\end{equation}
Then $B$ runs through the elements of the group $U$ reciprocal to $T$, whose degree is equal to the index of $T$.
\end{enumerate}

By \S\,11, (20), $T$ is the group reciprocal to $U$; the relation between the two groups is therefore mutual.

For these reciprocal groups there is also the theorem:

\begin{enumerate}[resume,label=\arabic*.]
\item If $T'$ is a divisor of $T$, then conversely the group $U$ reciprocal to $T$ is a divisor of the group $U'$ reciprocal to $T'$.
\end{enumerate}

For every element $A'$ of $T'$ is also contained in $T$, and hence, when $B$ runs through $U$, $\chi_B(A')=1$. Thus $B$ must occur in $U'$.

If from the totality of characters $\chi$ one selects any number $\xi_1,\xi_2,
\ldots,
\xi_\mu$, whether or not they form a group, then all elements $A$ satisfying the $\mu$ conditions
\begin{equation}
\xi_1(A)=1,
\quad \xi_2(A)=1,
\quad\ldots\quad
\xi_\mu(A)=1
\tag{8}
\end{equation}
form a group, since from $\xi_i(A)=1$ and $\xi_i(A')=1$ it follows that also $\xi_i(AA')=1$. If $\xi_1,\xi_2,
\ldots,
\xi_\mu$ do not form a group, then from equations (8) there follow further equations $\xi_{\mu+1}(A)=1,
\ldots$ until $\xi_1,\xi_2,
\ldots,
\xi_\mu$ are completed to a group.

\begin{enumerate}[resume,label=\arabic*.]
\item From Theorem 7 it follows that all possible divisors of $S$ can be obtained in this way.
\end{enumerate}

It may be useful to illustrate these theorems with a simple example. Let the degree of $S$ be the square of a prime, $n=p^2$. Then either $S$ has only one invariant, $p^2$, and is cyclic,
\begin{equation*}
\text{a)}\qquad S=1,A,A^2,
\ldots,A^{p^2-1},
\end{equation*}
or $S$ has two invariants, both equal to $p$; then $S$ consists of the elements
\begin{equation*}
\text{b)}\qquad S=A_1^{\alpha_1}A_2^{\alpha_2},
\qquad \alpha_1,\alpha_2=0,1,\ldots,p-1.
\end{equation*}

In case a) we obtain the $p^2$ characters
\[
\chi(A^\alpha)=\varepsilon^{\alpha\beta},
\qquad \alpha,
\beta=0,1,
\ldots,p^2-1,
\]
where $\varepsilon$ is a primitive root of $\varepsilon^{p^2}=1$. To form the divisors of $S$ according to Theorem 10, take one of the characters $\chi_B$ in which $\beta$ is not divisible by $p$. Then $\chi_B(A^\alpha)$ can be $1$ only when $\alpha$ is divisible by $p^2$, hence $=0$; that is, we obtain only the divisor of $S$ consisting of the identity element. But if $\beta=p\beta'$ is divisible by $p$ and not by $p^2$, then $\chi_{p\beta'}(A^\alpha)=1$ if and only if $\alpha=p\alpha'$ is divisible by $p$. We obtain the divisor
\[
T=1,A^p,A^{2p},
\ldots,A^{(p-1)p},
\]
and the divisor $U$ reciprocal to $T$ is here identical with $T$.

In case b), to form the characters we must take two $p$-th roots of unity $\varepsilon_1,\varepsilon_2$, and if
\[
B=A_1^{\beta_1}A_2^{\beta_2},
\]
we obtain
\[
\chi_B(A_1^{\alpha_1}A_2^{\alpha_2})=
\varepsilon^{\alpha_1\beta_1+\alpha_2\beta_2}.
\]
Putting two of these characters equal to $1$, that is,
\[
\alpha_1\beta_1+\alpha_2\beta_2\equiv0,
\qquad
\alpha_1\beta'_1+\alpha_2\beta'_2\equiv0 \pmod p,
\]
we see that, if the determinant $\beta_1\beta'_2-\beta_2\beta'_1$ is not divisible by $p$, then $\alpha_1$ and $\alpha_2$ must be divisible by $p$. If the determinant is divisible by $p$, one of the two congruences follows from the other. In the first case we therefore obtain a group consisting only of the identity element. Thus all divisors $T$ and $U$ different from $1$ and $S$ are obtained by determining $\alpha_1,
\alpha_2$, for fixed $\beta_1,
\beta_2$, in all possible ways according to the congruence
\begin{equation}
\alpha_1\beta_1+\alpha_2\beta_2\equiv0\pmod p.
\tag{9}
\end{equation}
If $\alpha_1,
\alpha_2$ is a solution of this congruence in which not both numbers are zero, then all solutions are obtained by letting the factor $h$ in $h\alpha_1,h\alpha_2$ run through $0,1,
\ldots,p-1$, and the group $T$ is therefore, with fixed $\alpha_1,
\alpha_2$ satisfying (9),
\[
(A_1^{\alpha_1}A_2^{\alpha_2})^h,
\qquad h=0,1,\ldots,p-1.
\]
The reciprocal group $U$ is obtained by seeking all values of $\beta$ satisfying (9), and therefore has the form
\[
(A_1^{\beta_1}A_2^{\beta_2})^h,
\qquad h=0,1,\ldots,p-1,
\]
where $\alpha_1,
\alpha_2,
\beta_1,
\beta_2$ are now four fixed numbers satisfying (9). Putting $\beta_1=0$, $\beta_2=1$, $\alpha_1=1$, $\alpha_2=0$, gives the special case of the two reciprocal groups
\[
A_1^h,
\qquad A_2^h,
\qquad h=0,1,\ldots,p-1.
\]

\section*{\S\,13. Genera in an Abelian group.}

An element of an Abelian group $S$ which is identical with its inverse, so that its square is the identity element, is called a two-sided element.\footnote{Compare Vol. I, p. 390, note.} The identity element therefore always belongs among the two-sided elements.

Represent the elements of $S$ by a basis $A_1,A_2,
\ldots,A_\nu$, whose elements have degrees $a_1,a_2,
\ldots,a_\nu$. Then
\begin{equation}
A=A_1^{\alpha_1}A_2^{\alpha_2}\cdots A_\nu^{\alpha_\nu}
\tag{1}
\end{equation}
is a two-sided element if and only if
\begin{equation}
2\alpha_1\equiv0\pmod{a_1},
\quad 2\alpha_2\equiv0\pmod{a_2},
\quad\ldots\quad
2\alpha_\nu\equiv0\pmod{a_\nu}.
\tag{2}
\end{equation}
If among the invariants of the group a power of $2$ occurs $\lambda$ times, then the basis of $S$ may be ordered so that $a_1,a_2,
\ldots,a_\lambda$ are even, while $a_{\lambda+1},
\ldots,a_\nu$ are odd. The solutions of (2) are then
\[
\alpha_1=\frac{\eta_1a_1}{2},
\quad
\alpha_2=\frac{\eta_2a_2}{2},
\quad\ldots\quad
\alpha_\lambda=\frac{\eta_\lambda a_\lambda}{2},
\quad
\alpha_{\lambda+1}=\cdots=\alpha_\nu=0,
\]
where $\eta_1,
\eta_2,
\ldots,
\eta_\lambda$ may be either $0$ or $1$. Thus all two-sided elements are obtained in the form
\[
A_1^{\eta_1a_1/2}A_2^{\eta_2a_2/2}\cdots A_\lambda^{\eta_\lambda a_\lambda/2},
\]
and their number is, since all combinations of $0$ and $1$ for the exponents $\eta$ are permissible,
\[
2^\lambda.
\]

If $\chi$ is any character of $S$, then, when $A$ is a two-sided element, $\chi(A)=\pm1$, since
\[
\chi(A)^2=\chi(1)=1.
\]

Likewise we call a character two-sided if it is a two-sided element in the group of characters. Since the group of characters is isomorphic with $S$, there are as many two-sided characters as two-sided elements, namely $2^\lambda$. By \S\,11, (17) they are represented as
\[
\chi_1^{\eta_1a_1/2}\chi_2^{\eta_2a_2/2}\cdots\chi_\lambda^{\eta_\lambda a_\lambda/2}.
\]

The two-sided characters have value $\pm1$ for every element.

The two-sided elements form by themselves a group $T$ of degree $2^\lambda$. Likewise the two-sided characters form a group isomorphic with it, and one obtains the group $U$ reciprocal to $T$ by seeking all elements $A$ for which all two-sided characters have value $+1$. This group, which by \S\,12, 7. is a divisor of $S$ of index $2^\lambda$, is by its last definition independent both of the choice of basis and of the roots of unity used to represent the characters, and is completely determined by the nature of the group $S$. Denoting it by $G$, the system of cosets
\begin{equation}
G,
\quad G_1,
\quad G_2,
\ldots,
G_{2^\lambda-1}
\tag{3}
\end{equation}
is characterized by the fact that for all elements of any one of these systems the two-sided characters give one and the same system of values. The systems $G,G_1,G_2,
\ldots$ are called the genera contained in $S$. The group $G$ itself is called the principal genus.

The number of genera is therefore as large as the number of two-sided elements.\footnote{Gauss first introduced these concepts into the theory of quadratic forms and used the expression ``genera''. Disqu. arithm. art. 228 ff.}

\section*{\S\,14. Indices modulo an odd prime power.}

The most important example of a commutative group is furnished by the natural numbers when they are combined by ordinary multiplication.

To derive a finite group from them, take an arbitrary positive integer $m$ as modulus and decompose $m$ into its prime factors,
\begin{equation}
m=2^\lambda q_1^{\lambda_1}q_2^{\lambda_2}\cdots,
\tag{1}
\end{equation}
where the $q_i$ are distinct odd primes, $\lambda,\lambda_1,
\lambda_2,
\ldots$ positive exponents, with $\lambda$ possibly zero, namely when $m$ is odd.

All numbers, positive as well as negative, that are congruent to one another modulo $m$ are then collected into one residue class, and each such class is represented by a representative, for example by the remainder on division, hence by one of the numbers $0,1,2,
\ldots,m-1$. If $a$ and $a'$ are two numbers of one such class, and $b$ and $b'$ two numbers of a second class, then the products $ab$ and $a'b'$ also belong to the same class. For if $a\equiv a'$ and $b\equiv b'$, then $ab\equiv a'b'\pmod m$. Thus multiplication composes not only numbers but also classes.

Nevertheless the totality of these residue classes does not yet form a group; from a congruence
\[
ab\equiv ac\pmod m
\]
it follows necessarily that $b\equiv c$ only when $a$ is relatively prime to $m$. Thus condition \S\,1, 3. is not generally satisfied.

The greatest common divisor that a number $a$ has with $m$ is the same throughout the class represented by $a$, and may be called the divisor of the class. If $d,d'$ are the divisors of two classes $a,a'$, then the greatest common divisor of $m$ and $dd'$ is the divisor of the class $aa'$. It follows that the residue classes whose numbers are relatively prime to $m$ again give such classes under composition, and for these the condition \S\,1, 3. is satisfied.

\begin{enumerate}[label=\arabic*.]
\item The residue classes whose individuals are relatively prime to the modulus therefore form, under composition by multiplication, an Abelian group.
\end{enumerate}

This group is the object of our consideration, and our chief aim is to determine a basis for it. We shall denote by $N$ every residue class containing only numbers prime to $m$, and the group of the classes $N$ whose existence has just been proved by $\mathfrak N$. By $n$ we shall denote any number prime to $m$.

The degree of this group is the number of positive integers not exceeding $m$ and relatively prime to $m$, a number determined in \S\,132 of the first volume. It is
\begin{equation}
\Pi=\varphi(m)=2^{\lambda-1}q_1^{\lambda_1-1}(q_1-1)q_2^{\lambda_2-1}(q_2-1)\cdots
\tag{2}
\end{equation}
or
\[
\Pi=\varphi(m)=q_1^{\lambda_1-1}(q_1-1)q_2^{\lambda_2-1}(q_2-1)\cdots,
\]
when $\lambda=0$.

The degree of an element $N$ of this group is the least positive exponent $e$ to which one must raise a representative $n$ of $N$ so that $n^e$ becomes congruent to the identity modulo $m$. Since every such $e$ must divide the degree $\varphi(m)$ of the group, one obtains the generalized theorem of Fermat:
\begin{equation}
n^{\varphi(m)}\equiv1\pmod m,
\tag{3}
\end{equation}
a formula valid for every $n$ prime to $m$.

Let $q$ be one of the odd prime factors of $m$, and let $q^\varkappa$ be the highest power of $q$ dividing $m$. We take a primitive root $g$ of $q$, chosen, if $\varkappa>1$, so that $g^q-g$ is not divisible by $q^2$ (Vol. I, \S\,184). For brevity, a number $g$ satisfying this last condition will be called a primitive root of $q^2$.

If $\varkappa=1$, then by the definition of $g$ the numbers
\[
1,g,g^2,
\ldots,g^{q-2}
\]
are all distinct modulo $q$, while $g^{q-1}$ is again congruent to $1$.

If $\varkappa>1$, then
\begin{equation}
g^{q-1}=1+hq,
\tag{4}
\end{equation}
and by our assumption on $g$, $h$ is not divisible by $q$. Raising (4) to the $q$-th power and applying the binomial theorem on the right gives
\[
g^{q(q-1)}=(1+hq)^q=1+hq^2+\cdots,
\]
hence
\begin{equation}
g^{q(q-1)}=1+h_1q^2,
\tag{5}
\end{equation}
where $h_1$ is not divisible by $q$. This would no longer be correct for $q=2$, which is why the prime $2$ requires different treatment.

Raising (5) once more to the $q$-th power gives
\[
g^{q^2(q-1)}=1+h_2q^3,
\]
and so continuing, for every positive exponent $\varkappa$ we obtain
\begin{equation}
g^{q^{\varkappa-1}(q-1)}=1+hq^\varkappa,
\tag{6}
\end{equation}
where $h$ is an integer not divisible by $q$. For abbreviation put
\begin{equation}
q^{\varkappa-1}(q-1)=\varphi(q^\varkappa)=c.
\tag{7}
\end{equation}
Then
\begin{equation}
g^c\equiv1\pmod{q^\varkappa},
\tag{8}
\end{equation}
and it remains to prove that $c$ is the least positive exponent for which (8) is satisfied. Suppose $e$ is this least exponent, so that
\begin{equation}
g^e\equiv1\pmod{q^\varkappa}.
\tag{9}
\end{equation}
Then $e$ must divide $c$, since otherwise (8) would also be satisfied if $c$ were replaced by the remainder after division by $e$, which is smaller than $e$.

On the other hand, $e$ must be divisible by $q-1$, since the congruence $g^e\equiv1\pmod q$ contained in (9) holds only for exponents divisible by $q-1$. Hence $e$ is of the form $q^{\lambda-1}(q-1)$ with a positive $\lambda$.

But $\lambda$ cannot be smaller than $\varkappa$, as follows from (6), which shows that $q^\lambda$ is the highest power of $q$ dividing the difference
\[
g^{q^{\lambda-1}(q-1)}-1.
\]
Thus $c$ is the least positive number for which (8) is satisfied, equivalently the numbers
\begin{equation}
1,g,g^2,
\ldots,g^{c-1}
\tag{10}
\end{equation}
are all incongruent modulo $q^\varkappa$. Thus $g$ may also be called a primitive root of $q^\varkappa$.

The number of terms in (10) is $c$. This is also the number of incongruent numbers modulo $q^\varkappa$ not divisible by $q$, and hence:

\begin{enumerate}[resume,label=\arabic*.]
\item If $q$ is an odd prime, $n$ any number not divisible by $q$, $g$ a primitive root of $q^2$, and $\varkappa$ a positive exponent, then one and only one number $\gamma$ modulo $c$ can be determined such that
\[
n\equiv g^\gamma\pmod{q^\varkappa}.
\]
\end{enumerate}

The number $c$ is always divisible by $2$, and we therefore also emphasize the special theorem
\begin{equation}
g^{c/2}\equiv-1\pmod{q^\varkappa}.
\tag{11}
\end{equation}
Indeed, in the product divisible by $q^\varkappa$,
\[
g^c-1=(g^{c/2}-1)(g^{c/2}+1),
\]
the first factor is not divisible by $q^\varkappa$, and therefore the second factor must be divisible by $q$. It then follows that the first factor cannot be divisible by $q$, since otherwise the sum of the two factors, and hence $g$ itself, would be divisible by $q$. Therefore the second factor must be divisible by $q^\varkappa$.

There is thus a complete analogy with the theorems on primitive roots and index systems that we learned in \S\,136 of the first volume, and here also $\gamma$ may be called the index of $n$ for the modulus $q^\varkappa$.

\section*{\S\,15. Indices for a power of 2 as modulus.}

A similar theorem must now be established for the powers $2^\lambda$ of the prime $2$, which, as already noted, behaves differently.

If $\lambda=1$, then every number $n$ not divisible by $2$ satisfies $n\equiv1\pmod2$. If $\lambda=2$, then $-1$ is to be regarded as a primitive root of $4$, for every odd number $n$ satisfies one of the congruences
\[
n\equiv(-1)^\alpha\pmod4,
\]
where $\alpha=0$ or $1$.

Already for modulus $8$ there is no primitive root, that is, no number whose powers represent all residue classes of odd numbers modulo $8$. For if $g$ is any odd number, then among its powers at most the two $1,g$ are distinct modulo $8$, because always $g^2\equiv1\pmod8$. If $g$ is taken neither congruent to $1$ nor to $-1\pmod8$, thus $g\equiv3$ or $5$, then every odd number is congruent modulo $8$ to one of the four numbers
\[
(-1)^\alpha g^\beta,
\qquad
\alpha=0,1,
\quad
\beta=0,1.
\]

The number of classes of odd numbers modulo $2^\lambda$ is $2^{\lambda-1}$. But for every odd number $g$, when $\lambda>2$,
\begin{equation}
g^{2^{\lambda-2}}\equiv1\pmod{2^\lambda}.
\tag{1}
\end{equation}
For assume (1) is true and put accordingly
\[
g^{2^{\lambda-2}}=1+h2^\lambda.
\]
Squaring gives
\[
g^{2^{\lambda-1}}=1+h2^{\lambda+1}+h^2 2^{2\lambda}
\equiv1\pmod{2^{\lambda+1}}.
\]
Thus if (1) is true for $\lambda$, it is also true for $\lambda+1$; since it is true for $\lambda=3$, it holds generally. It follows that among the powers of an odd number $g$ at most $2^{\lambda-2}$ can be distinct modulo $2^\lambda$.

On the other hand it follows easily for $g=5$ that $5^{2^{\lambda-3}}$ is the lowest power congruent to the identity modulo $2^\lambda$. For if $5^e$ is the lowest power congruent to the identity, then by (1) $e$ is a divisor of $2^{\lambda-2}$, hence a power of $2$; and if $e<2^{\lambda-2}$, then one would have
\[
5^{2^{\lambda-3}}\equiv1\pmod{2^\lambda}.
\]
This is impossible, for for every $\lambda>2$ the formula
\[
5^{2^{\lambda-3}}=1+\mu 2^{\lambda-1}
\]
holds with $\mu$ odd; it is proved by complete induction in the same way as (1). Thus, for $\lambda\ge3$, the $2^{\lambda-2}$ powers
\begin{equation}
1,5,5^2,
\ldots,
5^{2^{\lambda-2}-1}
\tag{2}
\end{equation}
are all distinct modulo $2^\lambda$. A relation of the form $5^\mu\equiv-5^\nu$ modulo $4$, and therefore even more so modulo any higher power of $2$, is impossible. Hence the quantities
\begin{equation}
-1,-5,-5^2,
\ldots,
-5^{2^{\lambda-2}-1}
\tag{3}
\end{equation}
are distinct modulo $2^\lambda$ both among themselves and from the quantities (2); since their total number is $2^{\lambda-1}$, every odd number is congruent modulo $2^\lambda$ to one and only one of the quantities (2), (3).

Putting
\begin{equation}
a=2,
\qquad b=2^{\lambda-2},
\tag{4}
\end{equation}
we therefore obtain the following theorem for $\lambda\ge3$:

\begin{enumerate}[label=\arabic*.,start=3]
\item For every odd number $n$ there is a pair of numbers $\alpha,\beta$, completely determined modulo the pair of moduli $a,b$, such that
\[
n\equiv(-1)^\alpha5^\beta\pmod{2^\lambda}.
\]
\end{enumerate}

The case $\lambda=2$ may be subsumed under this by putting $b=1$ and $\beta=0$. To include also the case $\lambda=1$, which is of no particular interest, one would have to put $a=1$, $b=1$.\footnote{If one wished to take the number 3 as basis in place of 5, then this formula would not yet be valid for $\lambda=3$, but it would be valid for every larger $\lambda$; hence 3 could serve as basis just as well as 5. The reason for preferring the basis 5 lies in the fact that all powers of this basis are $=1\pmod4$.}



% --- Batch62 insertion: Volume II §§16-21 ---

\section*{\S\,16. The groups of residue classes modulo a composite modulus.}

Combining the two preceding propositions gives the following result, which completely solves the problem posed at the beginning of \S\,14.

\medskip
\noindent\textbf{4.} Let
\begin{equation}
 m=2^\lambda q_1^{\chi_1}q_2^{\chi_2}\cdots
\end{equation}
be the modulus, and let $g_1,g_2,\ldots$ be primitive roots of the squares of the primes $q_1,q_2,\ldots$. Put
\[
 a=2,
 \qquad b=\frac12\varphi(2^\lambda),
 \qquad c_1=\varphi(q_1^{\chi_1}),\quad c_2=\varphi(q_2^{\chi_2}),\ldots.
\]
For every integer $n$ prime to $m$ there is a uniquely determined system of integers
\[
 \alpha,\ \beta,\ \gamma_1,\ \gamma_2,
\ldots
\]
modulo $a,b,c_1,c_2,
\ldots$, satisfying
\begin{equation}
\begin{aligned}
 n&\equiv (-1)^\alpha 5^\beta &&\pmod {2^\lambda},\\
 n&\equiv g_1^{\gamma_1} &&\pmod {q_1^{\chi_1}},\\
 n&\equiv g_2^{\gamma_2} &&\pmod {q_2^{\chi_2}},\\
 &\hspace{1.0em}\cdots.
\end{aligned}
\tag{1}
\end{equation}
The numbers $\alpha,\beta,\gamma_1,\gamma_2,
\ldots$ are called the system of indices of $n$ for the modulus $m$.

Here $\lambda\ge3$ is assumed first. The same statement remains valid for the remaining values of $\lambda$ if
\[
\begin{array}{c|cc}
\lambda& a&b\\ \hline
0,1&1&1\\
2&2&1
\end{array}
\]
is used. For $\lambda=0,1$ the indices $\alpha$ and $\beta$ are zero or omitted; for $\lambda=2$ the index $\beta$ disappears, and $\alpha$ is $0$ or $1$.

To represent the group $\grpN$ of residue classes prime to $m$ by a basis, determine the residue classes
\begin{equation}
 A,\ B,\ C_1,\ C_2,
\ldots,
\tag{2}
\end{equation}
or representatives of them, by the congruences
\[
\begin{array}{rcll@{\quad}rcll@{\quad}rcll}
A&\equiv&-1&\pmod {2^\lambda},& A&\equiv&1&\pmod {q_1^{\chi_1}},&A&\equiv&1&\pmod {q_2^{\chi_2}},\ldots,\\
B&\equiv&5&\pmod {2^\lambda},& B&\equiv&1&\pmod {q_1^{\chi_1}},&B&\equiv&1&\pmod {q_2^{\chi_2}},\ldots,\\
C_1&\equiv&1&\pmod {2^\lambda},&C_1&\equiv&g_1&\pmod {q_1^{\chi_1}},&C_1&\equiv&1&\pmod {q_2^{\chi_2}},\ldots,\\
C_2&\equiv&1&\pmod {2^\lambda},&C_2&\equiv&1&\pmod {q_1^{\chi_1}},&C_2&\equiv&g_2&\pmod {q_2^{\chi_2}},\ldots.
\end{array}
\]
Then every $n$ in the group is represented by
\begin{equation}
 n\equiv A^\alpha B^\beta C_1^{\gamma_1}C_2^{\gamma_2}\cdots\pmod m .
\tag{3}
\end{equation}
Thus $A,B,C_1,C_2,
\ldots$ form a basis of the group $\grpN$ with orders $a,b,c_1,c_2,
\ldots$.

If $m$ is odd, or divisible by only the first power of $2$, then $A$ and $B$ are absent from the basis. If $m$ is divisible by $4$ but not by $8$, $B$ is absent and the element $A$ of order $2$ remains.

\medskip
\noindent\textbf{5.} If one decomposes the numbers $a,b,c_1,c_2,
\ldots$ into powers of distinct primes, one obtains the invariants of the group $\grpN$.

Write the elements of the basis as
\begin{equation}
 C_{-1},\ C_0,\ C_1,
\ldots,C_u.
\tag{4}
\end{equation}
Here $C_{-1}$ and $C_0$ stand for $A$ and $B$, with the usual exceptional conventions for $\lambda=0,1,2$. Their orders are denoted by
\begin{equation}
 c_{-1},\ c_0,\ c_1,
\ldots,c_u,
\tag{5}
\end{equation}
and the indices of an element of $\grpN$ by
\begin{equation}
 \nu_{-1},\ \nu_0,\ \nu_1,
\ldots,\nu_u.
\tag{6}
\end{equation}
For $\lambda=0,1$ the first two indices are zero; for $\lambda=2$ only the index corresponding to $A$ remains. For $\lambda\ge3$, $\nu_{-1}$ has the two values $0,1$, while $\nu_0$ runs modulo $c_0$.

If $\varepsilon_{-1},\varepsilon_0,
\varepsilon_1,
\ldots,
\varepsilon_u$ are primitive roots of unity of degrees $c_{-1},c_0,c_1,
\ldots,c_u$, and if $\beta_{-1},\beta_0,
\ldots,
\beta_u$ are the indices of an element $b$, then the characters of $\grpN$ have the form
\begin{equation}
 \chi_b(n)=\varepsilon_{-1}^{\beta_{-1}\nu_{-1}}
 \varepsilon_0^{\beta_0\nu_0}
 \varepsilon_1^{\beta_1\nu_1}\cdots
 \varepsilon_u^{\beta_u\nu_u}.
\tag{7}
\end{equation}
For $\lambda=0,1,2$ one has $\varepsilon_{-1}=1$; otherwise $\varepsilon_{-1}=-1$. The remaining roots $\varepsilon_i$ have degrees $c_i=\varphi(q_i^{\chi_i})$.

\section*{Third Section. The group of cyclotomic fields.}

\section*{\S\,17. The resolvents of cyclotomy.}

We now apply Abelian groups to cyclotomy in the case where the degree of the root of unity is not a prime, but a higher power of a prime. Let $q$ be an odd prime and let $m=q^\chi$, with $\chi>1$. Choose a primitive congruence root $g$ modulo $m$, and define the index $\nu$ of each integer $n$ prime to $q$ by
\begin{equation}
 n\equiv g^\nu\pmod m.
\tag{1}
\end{equation}
As $n$ runs through the group $\grpN$ of residue classes prime to $q$ modulo $m$, whose order is
\begin{equation}
 c=\varphi(q^\chi)=q^{\chi-1}(q-1),
\tag{2}
\end{equation}
the index $\nu$ runs through a complete residue system modulo $c$. Let $r$ be a primitive $m$-th root of unity and $\varepsilon$ a primitive $c$-th root. Form the Lagrange resolvents
\begin{equation}
 (\varepsilon^\beta,r)=\sum_n \varepsilon^{\beta\nu}r^n.
\tag{3}
\end{equation}
The question is when such a resolvent can vanish.

For $\chi=1$ it never vanishes. In the present case let $n'$ be another element of $\grpN$, with index $\nu'$. Replacing $n$ by $nn'$ and then summing gives
\begin{equation}
 (\varepsilon^{-\beta},r)(\varepsilon^\beta,r)
 =\sum_n \varepsilon^{\beta\nu}
 \sum_{n'} r^{n'(n+1)}.
\tag{4}
\end{equation}
Thus one has first to evaluate
\begin{equation}
 \delta=\sum_{n'}r^{n'(n+1)}.
\tag{5}
\end{equation}
This sum vanishes unless $n+1$ is divisible by $q^{\chi-1}$. The exceptional values are
\begin{equation}
 n\equiv -1+tq^{\chi-1}\pmod {q^\chi},
\tag{6}
\end{equation}
with
\begin{equation}
 t=0,1,2,
\ldots,q-1.
\tag{7}
\end{equation}
For $t=0$ all terms are $1$, so
\begin{equation}
 \delta_0=\varphi(m)=q^\chi-q^{\chi-1};
\tag{8}
\end{equation}
for $t\ne0$ each primitive $q$-th root occurs equally often, giving
\begin{equation}
 \delta_t=-q^{\chi-1}.
\tag{9}
\end{equation}

For the corresponding index one has
\begin{equation}
 g^{\frac12\varphi(q^\chi)}\equiv -1\pmod q
\tag{10}
\end{equation}
and hence
\begin{equation}
 \nu\equiv \frac12\varphi(q^\chi)+\tau\varphi(q^{\chi-1})\pmod {\varphi(q^\chi)}.
\tag{11}
\end{equation}
Consequently
\begin{equation}
 (\varepsilon^{-\beta},r)(\varepsilon^\beta,r)=
 \begin{cases}
 0,&\beta\equiv0\pmod q,\\
 (-1)^\beta q^\chi,&\beta\not\equiv0\pmod q.
 \end{cases}
\tag{12}
\end{equation}
If $\beta$ is not divisible by $q$, neither factor can vanish. If $\beta$ is divisible by $q$, a direct summation shows that each factor already vanishes:
\begin{equation}
 (\varepsilon^\beta,r)=0.
\tag{13}
\end{equation}

\medskip
\noindent\textbf{1.} If the degree of $r$ is a power of an odd prime, then $(\varepsilon^\beta,r)$ vanishes if and only if $\beta$ is divisible by $q$.

For a power of $2$, write $m=2^\lambda$ with $\lambda\ge3$. Each odd $n$ has indices $\nu,\nu_1$ given by
\[
 n\equiv (-1)^{\nu_1}5^\nu\pmod m.
\]
The resolvents are
\begin{equation}
 ((-1)^{\beta_1},\Theta^\beta,r)=
 \sum_n (-1)^{\beta_1\nu_1}\Theta^{\beta\nu}r^n,
\tag{14}
\end{equation}
where $\Theta$ is a primitive root of degree $2^{\lambda-2}$. The same computation gives
\begin{equation}
 ((-1)^{-\beta_1},\Theta^{-\beta},r)
 ((-1)^{\beta_1},\Theta^\beta,r)=
 \begin{cases}
 (-1)^{\beta_1}2^\lambda,&\beta\equiv1\pmod2,\\
 0,&\beta\equiv0\pmod2.
 \end{cases}
\tag{17}
\end{equation}
For even $\beta$ one has in fact
\begin{equation}
 ((-1)^{\beta_1},\Theta^\beta,r)=0.
\tag{18}
\end{equation}

\medskip
\noindent\textbf{2.} If the degree of $r$ is a power of $2$ greater than $4$, then the resolvent $((-1)^{\beta_1},\Theta^\beta,r)$ vanishes if and only if $\beta$ is even. For $m=4$ this reduces to the two evident resolvents $i+i^3$ and $i-i^3$.

\section*{\S\,18. Cyclotomic fields.}

The theory of Abelian groups gives a deeper view of roots of unity and of the algebraic numbers arising from them. Let
\begin{equation}
 m=2^\lambda q_1^{\chi_1}q_2^{\chi_2}\cdots
\tag{1}
\end{equation}
and let $r$ be a primitive $m$-th root of unity. We may ignore $\lambda=1$, since for odd $m$ the primitive $m$-th roots and their negatives give the primitive $2m$-th roots.

The root $r$ satisfies an irreducible Abelian equation of degree
\begin{equation}
 \nu=\varphi(m).
\tag{2}
\end{equation}
The field of all rational functions of $r$ is a number field $\Omega(r)$ of degree $\nu$. More generally, any field whose elements are rational functions of roots of unity is called a cyclotomic field; the field generated by all rational functions of a primitive $m$-th root is the full cyclotomic field $\Omega_m$ of order $m$.

Every finite set of roots of unity can be regarded as powers of a single root whose degree is the least common multiple of their degrees. Hence all cyclotomic fields occur as subfields of full cyclotomic fields. The Galois group of $\Omega_m$ consists of the substitutions
\begin{equation}
 (r,r^n),
\tag{3}
\end{equation}
where $n$ ranges over the residue classes prime to $m$; it is isomorphic to the residue class group $\grpN$.

If $\grpA\subset\grpN$ and $\rho$ is fixed by the substitutions attached to $\grpA$, then $\Omega(\rho)$ is a subfield of $\Omega_m$. Conversely every subfield arises in this way.

\medskip
\noindent\textbf{1.} Subfields of $\Omega_m$ correspond to subgroups $\grpA$ of $\grpN$: a number belongs to $\Omega(\rho)$ precisely when it is fixed by all substitutions of $\grpA$.

The Galois group of the field corresponding to $\grpA$ is isomorphic to $\grpN/\grpA$, and also to the reciprocal group $\grpB$. If $|\grpA|=a$ and $|\grpB|=b$, then $ab=\nu$ and $\rho$ satisfies an irreducible Abelian equation of degree $b$.

\medskip
\noindent\textbf{2.} All cyclotomic fields are obtained by forming the subgroups $\grpA$ of the groups $\grpN$ for all moduli $m$, choosing a function belonging to $\grpA$, and adjoining it.

\medskip
\noindent\textbf{3.} If $\Omega'$ and $\Omega''$ correspond to $\grpA'$ and $\grpA''$, then $\Omega''\subset\Omega'$ if and only if $\grpA'\subset\grpA''$.

For a divisor $m_1$ of $m$, the field $\Omega_{m_1}$ corresponds to the subgroup
\begin{equation}
 a\equiv1\pmod {m_1},
\tag{4}
\end{equation}
which is denoted symbolically by
\begin{equation}
 \grpA_{m_1}=1\pmod {m_1}.
\tag{5}
\end{equation}
If $m_1,m_2$ are divisors of $m$ and $\mu=(m_1,m_2)$, then the least common multiple of the two subgroups
\begin{equation}
 \grpA_{m_1}=1\pmod {m_1},\qquad
 \grpA_{m_2}=1\pmod {m_2}
\tag{6}
\end{equation}
is the group
\begin{equation}
 \grpA=1\pmod \mu.
\tag{7}
\end{equation}
Therefore:

\medskip
\noindent\textbf{4.} $\Omega_\mu$ is the greatest common subfield of $\Omega_{m_1}$ and $\Omega_{m_2}$, and $\Omega_{m_1m_2/\mu}$ is their least common multiple.

A subfield of $\Omega_m$ is called primary if it is not contained in a full cyclotomic field of lower order. If $q_1,
\ldots,q_r$ are the distinct primes dividing $m$ and
\begin{equation}
 m=q_1m_1=q_2m_2=\cdots=q_rm_r,
\tag{8}
\end{equation}
then the non-primary subfields are obtained from the subfields of
\begin{equation}
 \Omega_{m_1},\Omega_{m_2},
\ldots,\Omega_{m_r},
\tag{9}
\end{equation}
corresponding to the groups
\begin{equation}
 \grpQ_1=1\pmod {m_1},\quad
 \grpQ_2=1\pmod {m_2},\quad
 \ldots,
 \quad\grpQ_r=1\pmod {m_r}.
\tag{10}
\end{equation}

\medskip
\noindent\textbf{5.} To obtain all primary subfields of $\Omega_m$, find all primary subgroups of $\grpN$ and form the corresponding fields.

\medskip
\noindent\textbf{6.} If one lists all primary subfields of all full cyclotomic fields, each cyclotomic field appears exactly once, and then in terms of roots of unity of lowest possible degree.

\section*{\S\,19. Primary and non-primary divisors of $\grpN$.}

Let $q$ be one of the primes dividing $m$, write $m=qm'$, and consider
\[
 \grpQ\equiv1\pmod {m'}.
\]
Using the index notation from \S\,16, let the reciprocal group of $\grpQ$ be denoted by $P$. The indices of an element of $\grpQ$ are
\[
 \gamma_{-1},\gamma_0,
\gamma_1,
\ldots,
\gamma_\mu,
\]
and those of an element of $P$ by
\[
 \delta_{-1},\delta_0,
\delta_1,
\ldots,
\delta_\mu.
\]
The defining property of $P$ is that the index corresponding to $q$ is constrained: if $q$ divides $m$ more than once, that index is divisible by $q$; if $q$ divides $m$ only once, it is divisible by $q-1$:
\begin{equation}
\begin{array}{ll}
 \beta_q\equiv0\pmod q,& q^2\mid m,\\
 \beta_q\equiv0\pmod {q-1},& q\mid m,
 q^2\nmid m.
\end{array}
\tag{1}
\end{equation}

\medskip
\noindent\textbf{7.} A divisor $\grpA$ of $\grpN$ is primary if and only if, for no prime $q$ dividing $m$, do all elements of the reciprocal group $\grpB$ satisfy the corresponding condition (1).

This is the form in which primary divisors of the residue class group are used in the construction of cyclotomic fields.

\section*{\S\,20. Cyclotomic periods.}

Let $\grpA$ be a divisor of $\grpN$ of order $e$, and let $r$ be a primitive $m$-th root of unity. Put
\begin{equation}
 \eta=\sum_{a\in\grpA} r^a.
\tag{1}
\end{equation}
This is a cyclotomic period. If $b$ belongs to the reciprocal group $\grpB$, the corresponding Lagrange resolvent is
\begin{equation}
 (\chi_b,r)=\sum_n \chi_b(n)r^n.
\tag{2}
\end{equation}
Factoring $r$ according to the prime-power factors of $m$,
\begin{equation}
 r=r_0r_1\cdots r_u,
\tag{12}
\end{equation}
one obtains a product decomposition of the resolvent. The vanishing criterion of \S\,17 implies that if a period does not belong to the group $\grpA$, then among the primes dividing $m$ to higher than the first power there is one such that every element of $\grpB$ has the corresponding index divisible by that prime:
\begin{equation}
 \beta_0\equiv0\pmod2,
 \quad\beta_1\equiv0\pmod {q_1},
 \quad\ldots,
 \quad\beta_u\equiv0\pmod {q_u}.
\tag{15}
\end{equation}

\medskip
\noindent\textbf{1.} If the period $\eta$ does not belong to $\grpA$, then $\grpA$ is not primary.

\medskip
\noindent\textbf{2.} Consequently, if $\grpA$ is primary, then the period $\eta$ belongs to $\grpA$.

\medskip
\noindent\textbf{3.} Every primary subfield of $\Omega_m$ is representable in the form $\Omega(\eta)$.

Since non-primary subfields reappear as primary subfields of lower full cyclotomic fields, we obtain:

\medskip
\noindent\textbf{4.} All cyclotomic fields are representable as $\Omega(\eta)$.

\medskip
\noindent\textbf{5.} Every rational function of roots of unity can be expressed as a rational function of a cyclotomic period.

\medskip
\noindent\textbf{6.} If $\grpA$ is any primary or non-primary divisor of $\grpN$ of index $e$, then there is a divisor $m_1$ of $m$ and a period $\eta\in\Omega_{m_1}$ of index $e$ which, regarded as a function of $r$, belongs to $\grpA$.

For
\[
 \eta_h=\sum_{a\in\grpA}r^{ha}
\]
all such periods are rationally expressible in terms of $\eta$. Products are again linear combinations of the periods, since
\[
 \eta_h\eta_k=\sum_{a,a'\in\grpA}r^{ha+ka'}
\]
can be rearranged by replacing $a'$ with $aa'$ for fixed $a$.

If $q$ occurs only once in $m$ and $\grpA$ is divisible by the group $\grpQ$, then the period may be represented by roots of unity of lower degree. The numbers of $\grpQ$ have the form $1+hm'$ except for the single value with
\begin{equation}
 1+hm'\equiv0\pmod q.
\tag{18}
\end{equation}
Consequently, with
\begin{equation}
 a=a'(1+hm'),
\tag{19}
\end{equation}
the period is, apart from sign, a period formed from $m'$-th roots.

\section*{\S\,21. Cyclotomic fields with prescribed group.}

We have shown that every period
\[
 \eta=\sum_{a\in\grpA}r^a
\]
where $\grpA$ is a primary divisor of $\grpN$, is a root of an Abelian equation whose group is isomorphic to the reciprocal group $\grpB$.

The problem is now to determine $m$ and $\grpA$ in all possible ways so that $\grpB$ has a prescribed system of invariants. Equivalently, determine all cyclotomic fields with a given Galois group. We assume $\grpA$ primary.

Let the indices of $a\in\grpA$ be
\[
 \alpha_1,
\ldots,\alpha_\mu,
\]
with index moduli $c_1,
\ldots,c_\mu$, and let the indices of $b\in\grpB$ be $\beta_1,
\ldots,
\beta_\mu$. The reciprocal relation is
\begin{equation}
 \omega_1^{\alpha_1\beta_1}\omega_2^{\alpha_2\beta_2}
 \cdots\omega_\mu^{\alpha_\mu\beta_\mu}=1.
\tag{1}
\end{equation}
Let the invariants of $\grpB$ be
\[
 i_1,i_2,
\ldots,i_v,
\]
and let
\begin{equation}
 J=i_1j_1=i_2j_2=\cdots=i_vj_v
\tag{2}
\end{equation}
be their least common multiple. If
\[
 b=g_1^{x_1}g_2^{x_2}\cdots g_v^{x_v},
\tag{3}
\]
then the $g_k$ form a basis of the prescribed group precisely when
\begin{equation}
 g_1^{x_1}g_2^{x_2}\cdots g_v^{x_v}\equiv1\pmod m
\tag{4}
\end{equation}
implies and is implied by
\begin{equation}
 x_1\equiv0\pmod {i_1},\quad
 x_2\equiv0\pmod {i_2},\quad\ldots,
 x_v\equiv0\pmod {i_v}.
\tag{5}
\end{equation}

The number $J$ is therefore the least positive integer satisfying
\begin{equation}
 J\beta_1\equiv0\pmod {c_1},\quad
 J\beta_2\equiv0\pmod {c_2},\quad\ldots,
 J\beta_\mu\equiv0\pmod {c_\mu}
\tag{6}
\end{equation}
for all elements of $\grpB$. It follows that: an odd prime not dividing $J$ can occur in $m$ only to the first power; an odd prime dividing $J$ can occur in $m$ at most one power more often than in $J$; if $J$ is odd then $m$ is odd; and if $J$ is divisible by a power of $2$, then $m$ can contain the factor $2$ at most two powers more often than $J$ does. A prime that occurs simply in $m$ must have $q-1$ divisible by one of the prime factors of $J$.

To determine $\grpB$, let
\[
 \gamma_{k,h}\qquad (k=1,
\ldots,v,
\ h=1,
\ldots,
\mu)
\]
be the $h$-th index of the basis element $g_k$. Then the indices of a general element are
\begin{equation}
 \beta_h=\gamma_{1,h}x_1+\gamma_{2,h}x_2+\cdots+\gamma_{v,h}x_v\pmod {c_h}.
\tag{8}
\end{equation}
The necessary and sufficient basis condition is that the corresponding simultaneous congruences
\begin{equation}
 \gamma_{1,h}x_1+\cdots+\gamma_{v,h}x_v\equiv0\pmod {c_h}
\quad(h=1,
\ldots,
\mu)
\tag{9}
\end{equation}
are equivalent to (5). In addition, the primarity condition requires that the entries in each column are not all divisible by the forbidden prime divisor.

Let $d_{k,h}=(i_k,c_h)$ and write
\begin{equation}
 i_k=d_{k,h}i_{k,h},\qquad c_h=d_{k,h}c_{k,h}.
\tag{11}
\end{equation}
Then $i_{k,h}$ and $c_{k,h}$ are coprime. One obtains
\begin{equation}
 \gamma_{k,h}=c_{k,h}e_{k,h},
\tag{15}
\end{equation}
and the congruences reduce to
\begin{equation}
 e_{1,h}i_{1,h}j_1x_1+e_{2,h}i_{2,h}j_2x_2+
\cdots+e_{v,h}i_{v,h}j_vx_v\equiv0\pmod J.
\tag{17}
\end{equation}
The $e_{k,h}$ must be chosen so that these congruences have only the solutions prescribed by (5), and so that the primarity divisibility condition remains satisfied.

Fix a prime $p$ occurring in the invariants, and suppose
\begin{equation}
 i_1=p^{\pi_1},\quad i_2=p^{\pi_2},\quad\ldots,
 i_\rho=p^{\pi_\rho}
\tag{18}
\end{equation}
are precisely the invariants divisible by $p$. The relevant part of (17) is a system
\begin{equation}
 e_{1,h}i_{1,h}y_1+e_{2,h}i_{2,h}y_2+
\cdots+e_{\rho,h}i_{\rho,h}y_\rho\equiv0\pmod {p^\pi}.
\tag{20}
\end{equation}
This system is treated by determinants. In matrix form it is
\begin{equation}
\begin{pmatrix}
 a_{1,1}&a_{2,1}&\cdots&a_{\rho,1}\\
 a_{1,2}&a_{2,2}&\cdots&a_{\rho,2}\\
 \vdots&\vdots&&\vdots\\
 a_{1,\mu}&a_{2,\mu}&\cdots&a_{\rho,\mu}
\end{pmatrix}.
\tag{22}
\end{equation}
The criterion is:

\medskip
\noindent\textbf{7.} The number $\rho$ of invariants divisible by $p$ must be no larger than the number $\mu$ of index moduli, and from the matrix
\begin{equation}
\begin{pmatrix}
 c_{1,1}e_{1,1}&c_{2,1}e_{2,1}&\cdots&c_{\rho,1}e_{\rho,1}\\
 c_{1,2}e_{1,2}&c_{2,2}e_{2,2}&\cdots&c_{\rho,2}e_{\rho,2}\\
 \vdots&\vdots&&\vdots\\
 c_{1,\mu}e_{1,\mu}&c_{2,\mu}e_{2,\mu}&\cdots&c_{\rho,\mu}e_{\rho,\mu}
\end{pmatrix}
\tag{24}
\end{equation}
one must be able to form a determinant of $\rho$ rows not divisible by $p$.

Thus, if a prime $p$ occurs in $\rho$ invariants $i_1,
\ldots,i_\rho$, there must be at least $\rho$ index moduli, and $\rho$ of them must be selectable so that they are respectively divisible by those invariants. This is the final condition on $m$. The remaining coefficients $e_{k,h}$ can then be chosen to satisfy the determinant condition and the primarity condition; these requirements merely prescribe divisibility by some primes and non-divisibility by others.

\subsection*{Summary of the results.}

To represent all cyclotomic fields, each exactly once and as simply as possible, choose an arbitrary system of prime-power invariants
\[
 i_1,i_2,
\ldots,i_v.
\]
Choose a modulus $m$ as follows. If $p^\pi$ is the highest power of a prime $p$ occurring among the invariants, take $p$ at most $(\pi+1)$ times in $m$, or, when $p=2$, at most $(\pi+2)$ times. A prime $q$ not occurring in the invariants may occur in $m$ only simply, and then only if $q-1$ is divisible by one of the prime factors of the invariants.

The index moduli $c_1,
\ldots,c_\mu$ must be numerous enough that for each prime $p$ and each cluster of invariants that are powers of $p$, one can select equally many distinct $c_h$ divisible by those invariants. Then determine $i_{k,h}$ and $c_{k,h}$ by (11), form the matrices (24), choose the $e_{k,h}$ so that the determinant and non-divisibility conditions are fulfilled, and put
\[
 \gamma_{k,h}=c_{k,h}e_{k,h}.
\]
The congruences (8) then give the indices of a basis of an Abelian group $\grpB$, whose reciprocal group $\grpA$ is a primary divisor of the full residue class group $\grpN$.

\section*{\S\,22. Determination of the group $\grpX$.}

The final purpose of the preceding investigation is not so much the determination of the group $\grpB$ as that of the reciprocal group $\grpX$, because from $\grpX$ the cyclotomic periods
\[
   \eta=\sum r^a
\]
can be formed directly.  This problem is solved as follows.  By the relation defining the reciprocal indices, the indices $\alpha$ of the numbers in $\grpX$ are connected with the indices $\beta$ of the numbers in $\grpB$ by
\begin{equation}
  \omega_1^{\alpha_1\beta_1}\omega_2^{\alpha_2\beta_2}\cdots
  \omega_\mu^{\alpha_\mu\beta_\mu}=1,
  \tag{1}
\end{equation}
where $\omega_1,\omega_2,\ldots,\omega_\mu$ are fixed primitive roots of unity of degrees $c_1,c_2,\ldots,c_\mu$.  Let $C$ be the least common multiple of $c_1,c_2,\ldots,c_\mu$, and write
\begin{equation}
  C=c_1c'_1=c_2c'_2=\cdots=c_\mu c'_\mu .
  \tag{2}
\end{equation}
If $\omega$ is a primitive $C$-th root of unity, (1) may be replaced by
\[
   \omega^{c'_1\alpha_1\beta_1+c'_2\alpha_2\beta_2+\cdots+c'_\mu\alpha_\mu\beta_\mu}=1.
\]
Hence the $\alpha_h$ have to satisfy the congruence
\[
   c'_1\alpha_1\beta_1+c'_2\alpha_2\beta_2+
   \cdots+c'_\mu\alpha_\mu\beta_\mu\equiv0\pmod C.
\]
Equivalently, the sums
\[
   \frac{\alpha_1\beta_1}{c_1}+\frac{\alpha_2\beta_2}{c_2}+
   \cdots+\frac{\alpha_\mu\beta_\mu}{c_\mu}
   =\sum_{h=1}^{\mu}\frac{\alpha_h\beta_h}{c_h}
\]
must be integers.  Substituting for $\beta_h$ the expressions found in \S\,21 gives
\[
   \sum_k x_k\sum_{h=1}^{\mu}\frac{\alpha_h\gamma_{k,h}}{c_h},
\]
and since the integers $x_k$ are arbitrary, each of the sums
\[
   \sum_{h=1}^{\mu}\frac{\alpha_h\gamma_{k,h}}{c_h},
   \,\qquad k=1,2,\ldots,\nu,
\]
must be an integer.  With the notation
\[
   \gamma_{k,h}=e_{k,h}c_{k,h},\qquad
   c_h=\delta_{k,h}c_{k,h},\qquad
   i_k=\delta_{k,h}i_{k,h},
\]
this becomes the system of congruences
\begin{equation}
   \sum_{h=1}^{\mu} \alpha_h i_{k,h}e_{k,h}\equiv0\pmod {i_k},
   \qquad k=1,2,\ldots,\nu .
   \tag{3}
\end{equation}
Written out, this is
\[
\begin{array}{rcl}
 e_{1,1}i_{1,1}\alpha_1+e_{1,2}i_{1,2}\alpha_2+
     \cdots+e_{1,\mu}i_{1,\mu}\alpha_\mu&\equiv&0\pmod {i_1},\\
 e_{2,1}i_{2,1}\alpha_1+e_{2,2}i_{2,2}\alpha_2+
     \cdots+e_{2,\mu}i_{2,\mu}\alpha_\mu&\equiv&0\pmod {i_2},\\
 &&\vdots\\
 e_{\nu,1}i_{\nu,1}\alpha_1+e_{\nu,2}i_{\nu,2}\alpha_2+
     \cdots+e_{\nu,\mu}i_{\nu,\mu}\alpha_\mu&\equiv&0\pmod {i_\nu}.
\end{array}
\]
These congruences determine, as they must, the numbers
$\alpha_1,\alpha_2,\ldots,\alpha_\mu$ only modulo $c_1,c_2,\ldots,c_\mu$, because $i_{k,h}c_h=i_kc_{k,h}$ is divisible by $i_k$.  Therefore, to find the indices of all numbers of the group $\grpX$, one has simply to determine all the solutions of (3).

\section*{Fourth Section. Cubic and biquadratic Abelian fields.}

\section*{\S\,23. Cubic cyclotomic fields.}

The general investigations just completed permit many applications to special cases.  It is characteristic that even the simplest cases, when attacked directly, lead almost into the same difficulties that the general theory has overcome.

We first determine all cubic cyclotomic fields, i.e. all cyclotomic periods which satisfy a rational equation of the third degree, choosing in each case the simplest expression among the equivalent ones.  Here there is only one invariant, $i_1=3$.  According to the list of \S\,21 the modulus $m$ may contain the factor $9$, but not the factor $3$, and may contain arbitrarily many primes of the form $3N+1$.  Thus, below $200$, the admissible values of $m$ are
\[
\begin{gathered}
7,9,13,19,31,37,43,61,63,67,73,79,91,97,103,109,\\
117,127,133,139,151,157,163,171,181,193,199.
\end{gathered}
\]
The smallest value of $m$ containing more than two prime factors is $7\cdot9\cdot13=819$.

Let $q_1,q_2,\ldots,q_\mu$ be chosen from $9$ and from the primes
$7,13,19,31,\ldots$.  Then
\begin{equation}
   m=q_1q_2\cdots q_\mu .
   \tag{1}
\end{equation}
If $q_1=9$, the corresponding index modulus is $c_1=6$; otherwise the index moduli are $c_h=q_h-1$.  The numbers $\delta_{1,h}$ are all equal to $3$, and hence
\[
   i_{1,h}=1,
   \qquad c_{1,h}=\frac13c_h .
\]
The matrix of \S\,21 consists here of one vertical row
\[
  e_{1,1},e_{1,2},\ldots,e_{1,\mu}.
\]
These numbers may be chosen arbitrarily modulo $3$, provided none of them is divisible by $3$; hence each may be replaced by $+1$ or $-1$.  The indices of the numbers $a$ of the group $\grpX$ are therefore obtained from the single congruence
\begin{equation}
   e_1\alpha_1+e_2\alpha_2+
   \cdots+e_\mu\alpha_\mu\equiv0\pmod3,
   \tag{2}
\end{equation}
where each $e_h=\pm1$.  Simultaneously changing all signs gives the same group; the independent possibilities are thus represented by all sign systems modulo this equivalence.

Let $r$ be a primitive $m$-th root of unity.  For each solution of (2) one obtains a period
\begin{equation}
   \eta=\sum r^n,
   \tag{3}
\end{equation}
where the exponents $n$ run through the corresponding group $\grpX$.  All cubic cyclotomic equations arise in this way.  It remains to write them explicitly.  If $\rho$ denotes an imaginary third root of unity, and if
\begin{equation}
   (\rho,\eta)=\eta+\rho\eta'+\rho^2\eta'',
   \tag{4}
\end{equation}
then, as in the first volume, these resolvents split into factors belonging to the separate prime factors of $m$:
\begin{equation}
   (\rho,\eta)=(\rho,\eta_1)(\rho,\eta_2)\cdots(\rho,\eta_\mu).
   \tag{5}
\end{equation}
The individual factors are the resolvents already calculated for prime or ninth roots of unity.  Consequently
\begin{equation}
   (\rho,
\eta)(\rho^2,\eta)=m,
   \tag{6}
\end{equation}
and
\begin{equation}
   (\rho,
\eta)^3=m(a+b\rho),
   \qquad
   (\rho^2,
\eta)^3=m(a+b\rho^2),
   \tag{7}
\end{equation}
for certain integers $a,b$ depending on the chosen sign system.  From these equations follows
\begin{equation}
   m=a^2-ab+b^2 .
   \tag{8}
\end{equation}

Let the cubic equation with roots $\eta,\eta',\eta''$ be
\begin{equation}
   \eta^3+\alpha\eta^2+\beta\eta+\gamma=0.
   \tag{9}
\end{equation}
Here $\alpha$ is the sum of all primitive $m$-th roots of unity, namely $0$ or $-(-1)^\mu$ according as $9$ occurs among the factors of $m$ or not.  Introducing
\[
   \Delta=\eta+\rho\eta'+\rho^2\eta'',
\]
so that $\Delta=(\rho,\eta)$, one obtains the reduced equation
\begin{equation}
   \Delta^3+P\Delta-Q=0,
   \tag{10}
\end{equation}
with
\begin{equation}
   P=-\frac{m}{3},
   \qquad
   27Q=m(2a-b).
   \tag{11}
\end{equation}
Moreover
\begin{equation}
   m b=(\eta-\eta')(
\eta''-\eta)(\eta''-\eta')
   \tag{12}
\end{equation}
up to the standard factor indicated by the resolvent calculation, and therefore $b$ is divisible by $3$.  Writing
\begin{equation}
   2a-b=A,
   \qquad b=3B,
\end{equation}
one obtains
\begin{equation}
   4m=A^2+27B^2,
   \tag{13}
\end{equation}
while the equation for the transformed quantity becomes
\begin{equation}
   \Delta^3-\frac{m}{3}\Delta-\frac{mA}{27}=0,
   \tag{14}
\end{equation}
with discriminant satisfying
\begin{equation}
   \sqrt D=mB.
   \tag{15}
\end{equation}

For example, using the known decompositions for $q=9,7,13$, Weber obtains the following cubic equations:
\[
\begin{array}{rcl}
 m=63:& \eta^3-21\eta-28=0,&
        \eta^3-21\eta+35=0,\\[2mm]
 m=91:& \eta^3-\eta^2-30\eta+64=0,&
        \eta^3-\eta^2-30\eta-27=0,\\[2mm]
 m=819:& \eta^3-273\eta+91=0,&
        \eta^3-273\eta+91\cdot19=0,\\
       & \eta^3-273\eta+91\cdot8=0,&
        \eta^3-273\eta+91\cdot17=0.
\end{array}
\]
In the case $m=63$ the two groups $\grpX$ can be represented by the powers of $2$ and of $11$.  The corresponding exponents are
\[
 \pm1,\\ \pm8,
\pm2,
\pm16,
\pm4,
\pm32,
\]
respectively
\[
 \pm1,
\pm8,
\pm11,
\pm25,
\pm5,
\pm23 .
\]

These equations, together with their Tschirnhausen transforms, contain all cubic cyclotomic equations.  Conversely, they are distinct and cannot be transformed into one another by Tschirnhausen transformations, because they belong either to primary divisors of different full cyclotomic fields or, within the same full field, to different groups.

Periods formed from other roots of unity may also lead to cubic equations, but then they are expressible by means of lower roots of unity and are not primary.  For example, if $r$ is a $35$-th root, a period with exponents generated by the powers of $8$ may be reduced to a period formed from seventh roots; this is exactly the phenomenon excluded by the conditions of \S\,21.

\section*{\S\,24. Biquadratic cyclotomic fields.}

For the biquadratic cyclotomic equations two cases have to be distinguished.  Either there are two invariants $i_1=i_2=2$, or there is only one invariant $i_1=4$.

First consider the case $i_1=i_2=2$.  The modulus $m$ may be odd, divisible by $4$, or divisible by $8$, but by no higher power of $2$.  No odd prime may occur to a power higher than the first, and at least two index moduli must be present; hence, except for $m=8$, the modulus must contain at least two different prime factors.  Write, according to the three cases,
\[
   m=q_1q_2\cdots q_\mu,
   \qquad
   m=4q_1q_2\cdots q_\mu,
   \qquad
   m=8q_1q_2\cdots q_\mu,
\]
where in the second formula $q_1$ may be $4$, and in the third formula $q_1=8$.  The index moduli are respectively
\[
\begin{array}{cccc}
q_1-1,& q_2-1,&\ldots,&q_\mu-1,\\
2,& q_2-1,&\ldots,&q_\mu-1,\\
2,&2,&q_2-1,&\ldots,
q_\mu-1.
\end{array}
\]
All are divisible by $2$.  Thus
\begin{equation}
   \delta_{k,h}=2,
   \qquad i_{k,h}=1,
   \qquad c_{k,h}=\frac12 c_h,
   \quad k=1,2,\\ h=1,\ldots,\mu .
   \tag{1}
\end{equation}
The entries $e_{k,h}$ must be chosen so that from the matrix
\begin{equation}
\begin{pmatrix}
 e_{1,1}&e_{1,2}&\cdots&e_{1,\mu}\\
 e_{2,1}&e_{2,2}&\cdots&e_{2,\mu}
\end{pmatrix}
\tag{2}
\end{equation}
there can be formed at least one odd determinant of two rows and two columns.  If $\grpX$ is to be a primary divisor of $\grpN$, then in no pair
\[
   (e_{1,h},e_{2,h})
\]
which corresponds to a prime factor of $m$ may both entries be even.  In the case in which $m$ is divisible by $8$ there is one index modulus not attached to any prime factor; its pair is exempt from this last restriction.

The indices $\alpha_h$ of the numbers of $\grpX$ are determined by
\begin{equation}
\begin{aligned}
 e_{1,1}\alpha_1+e_{1,2}\alpha_2+
    \cdots+e_{1,\mu}\alpha_\mu&\equiv0\pmod2,\\
 e_{2,1}\alpha_1+e_{2,2}\alpha_2+
    \cdots+e_{2,\mu}\alpha_\mu&\equiv0\pmod2.
\end{aligned}
\tag{3}
\end{equation}
By replacing the two equations by independent linear combinations one may normalize the first column to $(1,0)^t$.  The number of different groups which arise is
\begin{equation}
   \frac{3^\mu-1}{2}.
   \tag{4}
\end{equation}

For the second case, with invariant $4$, the possible moduli are determined by the conditions of \S\,21.  There is one congruence
\begin{equation}
   e_{1,1}i_{1,1}\alpha_1+
   e_{1,2}i_{1,2}\alpha_2+
   \cdots+e_{1,\mu}i_{1,\mu}\alpha_\mu\equiv0\pmod4.
   \tag{5}
\end{equation}
The coefficients belonging to the odd prime factors must not be divisible by $4$ in the first part of the list and not by $2$ in the second; again the factor corresponding only to the power of $2$ is exceptional when $m$ is divisible by $8$ or by $16$.  For the example $m=48$ one obtains two possibilities,
\[
 (e_{1,1},e_{1,2},e_{1,3})=(1,0,1),
 \qquad
 (1,1,1),
\]
which give the two groups
\begin{equation}
  \grpX_1=\{1,17,31,47\},
  \qquad
  \grpX_2=\{1,7,41,47\}.
  \tag{6}
\end{equation}

\section*{\S\,25. Cubic Abelian equations.}

The preceding results on cubic and biquadratic cyclotomic equations are completed by the following theorem:

\medskip
\noindent\textbf{Theorem.}  All cubic and biquadratic Abelian equations over the rational number field are cyclotomic equations.

\medskip
Thus all Abelian number fields of degree three and four are represented by the cyclotomic periods constructed above.  The proof in these two special cases is possible because in the fields $\mathbb Q(\rho)$ and $\mathbb Q(i)$ of third and fourth roots of unity unique factorization into prime factors holds, as was shown in the first volume.

Let first $x_0,x_1,x_2$ be the roots of an Abelian equation of third degree.  Since the group is cyclic, its cyclic functions are rational.  With $\rho$ an imaginary third root of unity, form the resolvents
\begin{equation}
  (\rho,x_0)=x_0+\rho x_1+\rho^2x_2,
  \qquad
  (\rho^2,x_0)=x_0+\rho^2x_1+\rho x_2.
  \tag{1}
\end{equation}
Their cubes lie in $\mathbb Q(\rho)$ and their product is rational.  Hence
\begin{equation}
   (\rho,x_0)^3=a+b\rho,
   \qquad
   (\rho^2,x_0)^3=a+b\rho^2,
   \qquad
   (a+b\rho)(a+b\rho^2)=c^3 .
   \tag{2}
\end{equation}
In the ring $\mathbb Z[\rho]$ the units are $\pm1,\pm\rho,\pm\rho^2$; the prime factors are the factor $\sqrt{-3}=\rho-\rho^2$, the real primes $q\equiv2\pmod3$, and the conjugate complex factors of the real primes $p\equiv1\pmod3$,
\begin{equation}
   p=\pi\pi',
   \qquad
   \pi=\psi_1(\rho),
   \quad
   \pi'=\psi_1(\rho^2).
   \tag{3}
\end{equation}
The factor $a+b\rho$ can therefore be written as a product of powers of these primes.  Because the product with its conjugate is a third power of a rational number, the exponents occur in sums divisible by $3$.  By means of the cyclotomic resolvent formulas for the factors $\pi,\pi'$, one may write
\begin{equation}
   (\rho,x_0)=(-\epsilon)^\lambda Q_1
   (\sqrt{-3})^n q_1^{\sigma_1}q_2^{\sigma_2}\cdots
   p_1^{\tau_1}p_2^{\tau_2}\cdots H_1,
   \tag{4}
\end{equation}
where $H_1$ is a cyclotomic number, and the conjugate formula gives $(\rho^2,x_0)$.  The remaining roots of unity satisfy the necessary product condition.  With
\begin{equation}
   x_0+x_1+x_2=A,
   \tag{5}
\end{equation}
we finally have
\begin{equation}
\begin{aligned}
 x_0&=\frac13\{A+(\rho,x_0)+(\rho^2,x_0)\},\\
 x_1&=\frac13\{A+\rho^2(\rho,x_0)+\rho(\rho^2,x_0)\},\\
 x_2&=\frac13\{A+\rho(\rho,x_0)+\rho^2(\rho^2,x_0)\}.
\end{aligned}
\tag{6}
\end{equation}
Thus all three roots are cyclotomic numbers.  The detailed construction gives no cubic Abelian fields beyond those already described in \S\,23.

\section*{\S\,26. Biquadratic Abelian equations.}

The proof for biquadratic Abelian equations is parallel.  Two cases occur, according as the group is non-cyclic of order four or cyclic of order four.

In the non-cyclic case the group acting on the roots $x_0,x_1,x_2,x_3$ is
\begin{equation}
  1,
  \quad (0,1)(2,3),
  \quad (0,2)(1,3),
  \quad (0,3)(1,2).
  \tag{1}
\end{equation}
Consequently the three squares
\begin{equation}
\begin{aligned}
 (x_0+x_1-x_2-x_3)^2&=a,\\
 (x_0-x_1+x_2-x_3)^2&=b,\\
 (x_0-x_1-x_2+x_3)^2&=c
\end{aligned}
\tag{2}
\end{equation}
and the product of the three linear factors, together with
\begin{equation}
   x_0+x_1+x_2+x_3=4A,
   \tag{3}
\end{equation}
are rational.  On extracting square roots one obtains
\begin{equation}
   x_0=A+B\sqrt b+C\sqrt c+D\sqrt{bc},
   \tag{4}
\end{equation}
with rational $A,B,C,D$, and the other roots follow by changing the signs of $\sqrt b$ and $\sqrt c$.  By the quadratic theory of the first volume every square root is expressible, possibly after adjoining $i$, by cyclotomic periods; hence the assertion follows for the non-cyclic case.

In the cyclic case the roots may be ordered so that the group is
\begin{equation}
 1,
 \quad (0,1,2,3),
 \quad (0,2)(1,3),
 \quad (0,3,2,1).
 \tag{5}
\end{equation}
Set
\begin{equation}
\begin{aligned}
 4A&=x_0+x_1+x_2+x_3,\\
 (i,x_0)&=x_0+i x_1-x_2-i x_3,\\
 (-1,x_0)&=x_0-x_1+x_2-x_3,\\
 (-i,x_0)&=x_0-i x_1-x_2+i x_3.
\end{aligned}
\tag{6}
\end{equation}
Then $A$ and $(-1,x_0)^2=m$ are rational, and
\begin{equation}
  (i,x_0)^4=a+bi,
  \qquad
  (-i,x_0)^4=a-bi,
  \tag{7}
\end{equation}
are numbers of the field $\mathbb Q(i)$.  Since
\begin{equation}
   (i,x_0)(-i,x_0)=c
   \tag{8}
\end{equation}
is rational, the product $(a+bi)(a-bi)$ is a fourth power of a rational number.

In $\mathbb Z[i]$ the units are $\pm1,\pm i$; the factor $1+i$ lies above $2$; real primes $q\equiv3\pmod4$ remain prime, and primes $p\equiv1\pmod4$ split as
\begin{equation}
  p=\pi\pi',
  \qquad
  \pi=\psi_1(i),
  \quad
  \pi'=\psi_1(-i).
  \tag{9}
\end{equation}
Moreover the cyclotomic resolvents satisfy
\begin{equation}
   (i,\eta)^4=\pi^3\pi',
   \qquad
   (-i,\eta)^4=\pi(\pi')^3,
   \qquad
   (i,\eta)(-i,
\eta)=p.
   \tag{10}
\end{equation}
Factoring $a+bi$ and $a-bi$ in $\mathbb Z[i]$ therefore represents them as fourth powers of cyclotomic numbers:
\begin{equation}
   a+bi=c^4H_1^4,
   \qquad
   a-bi=c^4H_2^4,
   \tag{11}
\end{equation}
up to the appropriate unit factors.  Taking fourth roots yields
\begin{equation}
   (i,x_0)=i^\lambda cH_1,
   \qquad
   (-i,x_0)=i^{-\lambda}cH_2,
   \qquad
   (-1,x_0)=\sqrt m.
   \tag{12}
\end{equation}
Thus, since $\sqrt m$, $c$, and the $H_j$ are cyclotomic numbers, each of
$x_0,x_1,x_2,x_3$ is cyclotomic.  The same consideration as in the cubic case shows that this gives no fields beyond those already obtained in \S\,24.

Consequently all Abelian fields of degrees three and four are cyclotomic fields, and all of them are represented rationally by cyclotomic periods.


\clearpage

\section*{Fifth Section. Constitution of the General Groups.}

\section*{\S\,27. Formation of Groups according to Cayley.}

The general definition of a group given in \S\,1 still leaves much about the nature of the concept in obscurity, and the particular groups met in the course of our investigations give only indications of general laws. They show, however, that the concept of a group contains no contradiction in itself. The definition contains more than appears at first sight, and the number of possible groups that can be composed from a given number of elements is very restricted. The general laws which prevail here are known only in their smallest part; therefore every new special group, especially of small order, has its own interest and calls for detailed study.

Which groups are possible among a given number of elements, that is, for a given order? This is the general question with which we are concerned, although we are still far from a complete solution. Cayley was the first to attack this problem for the lowest orders.

For every order $n$ there is always one group, namely the cyclic group obtained by putting $a^n=1$ and taking the $n$ elements
\[
1,a,a^2,\ldots,a^{n-1}
\]
to be distinct. If $n$ is a prime, this is the only group of order $n$, when isomorphic groups are regarded as identical. For the order of every element of a group divides the order of the group, and hence in a group of prime order every element except the identity has order $n$.

To represent a group of order $n$ completely, one would have to construct a square table with $n^2$ fields, arranged in $n$ rows and $n$ columns. Each row and each column is labelled by one of the elements; in the field where they meet one places the product element, the row element being the first component and the column element the second:
\generalCayleyTable
The entries of such a table cannot be chosen arbitrarily. They must obey the associative law. Thus in finding $\alpha\beta\gamma$ one must obtain the same result whether one first looks up $(\alpha\beta)$ and then $(\alpha\beta)\gamma$, or first $(\beta\gamma)$ and then $\alpha(\beta\gamma)$.

For $n=4$, if an element of order four exists, there is the cyclic group
\[
1,\alpha,\alpha^2,\alpha^3.
\]
If all elements have order one or two, there is the group
\[
1,\alpha,\beta,\alpha\beta,
\qquad \alpha\beta=\beta\alpha.
\]
Thus there are only these two groups of order four. If the elements are denoted by $1,\alpha,\beta,\gamma$, their two tables are:
\orderFourTables

The case $n=6$ can be disposed of completely as follows. Denote the six elements by $1,\alpha,\beta,\gamma,\delta,\varepsilon$, with $1$ the identity. If one of them has order $6$, the group is cyclic. Apart from this case, the orders of $\alpha,\beta,\gamma,\delta,\varepsilon$ can only be $2$ or $3$. If $\alpha$ and $\beta$ have order $2$, then $\alpha\beta$ cannot also have order $2$; otherwise
\[
\alpha=\alpha^{-1},\quad \beta=\beta^{-1},\quad \alpha\beta=\beta^{-1}\alpha^{-1}=\beta\alpha,
\]
and $1,\alpha,\beta,\alpha\beta$ would form a subgroup of order four, impossible in a group of order six.

There must therefore be at least one element of order three. Denote it by $\alpha$, and choose $\gamma$ outside $1,\alpha,\alpha^2$. Then the group is arranged as
\begin{equation}
S=1,\alpha,\alpha^2,\gamma,\gamma\alpha,\gamma\alpha^2.
\tag{1}
\end{equation}
In order that this be a group, form
\[
\gamma S=\gamma,\gamma\alpha,\gamma\alpha^2,\gamma^2,\gamma^2\alpha,\gamma^2\alpha^2;
\]
it must coincide with $S$, so $\gamma^2$ must be $1$, $\alpha$, or $\alpha^2$. If $\gamma^2=\alpha$ or $\alpha^2$, then $S$ is cyclic, namely $S=(1,\gamma,\gamma^2,\ldots,\gamma^5)$. Thus only
\begin{equation}
\gamma^2=1
\tag{2}
\end{equation}
remains. Since $\gamma$ may be replaced by $\gamma\alpha$ or by $\gamma\alpha^2$, these too must have order two, whence
\begin{equation}
\gamma\alpha=\alpha^2\gamma,
\qquad
\gamma\alpha^2=\alpha\gamma.
\tag{3}
\end{equation}
With these relations every product of powers of $\alpha$ and $\gamma$ can be reduced to exactly one element of $S$. Renaming
\[
1,\alpha,\alpha^2,\gamma,\gamma\alpha,\gamma\alpha^2
\quad\hbox{as}\quad
1,\alpha,\beta,\gamma,\delta,\varepsilon,
\]
one obtains the table
\orderSixTable
and, as required, every element occurs in every row and every column.

\section*{\S\,28. Relation of General Groups to Permutation Groups.}

The groups most used so far have been permutation groups. This kind of group gains increased importance for general group theory through the following considerations.

Let
\begin{equation}
P=a_0,a_1,a_2,\ldots,a_{n-1}
\tag{1}
\end{equation}
be any group of order $n$. If one chooses an element $b$ of $P$, then the complex $Pb$ is identical with $P$. Hence the two rows
\[
P=a_0,a_1,a_2,\ldots,a_{n-1},
\qquad
Pb=a_0b,a_1b,a_2b,\ldots,a_{n-1}b
\]
differ only in their order. The passage from $P$ to $Pb$ is thus a permutation of the $n$ symbols $0,1,2,
\ldots,n-1$. To each element $b$ of $P$ corresponds such a permutation, denoted by $\pi_b$. Distinct elements give distinct permutations, and the identity gives the identity permutation. If
\[
\pi_b=(P,Pb),\qquad \pi_c=(P,Pc),
\]
then $\pi_c$ may also be written $(Pb,Pbc)$, and so
\begin{equation}
\pi_b\pi_c=\pi_{bc}.
\tag{2}
\end{equation}
The permutations $\pi$ form a permutation group isomorphic to $P$. It is transitive, since $a_0$ can be carried to any element of $P$. Thus:
\begin{enumerate}[label=\arabic*.,leftmargin=2.2em]
\item Every group of order $n$ is isomorphic to a transitive permutation group on $n$ letters.
\end{enumerate}

There are of course other permutation groups isomorphic to a given group. Let $R$ be a divisor of $P$ of index $j$, and write a decomposition into right cosets
\begin{equation}
P=R+Rb_1+Rb_2+\cdots+Rb_{j-1}.
\tag{3}
\end{equation}
For an arbitrary element $a$ of $P$, the systems
\begin{equation}
\begin{aligned}
R&=R,Rb_1,Rb_2,\ldots,Rb_{j-1},\\
Ra&=Ra,Rb_1a,Rb_2a,\ldots,Rb_{j-1}a
\end{aligned}
\tag{4}
\end{equation}
again coincide up to order. Thus each $a$ defines a permutation $\pi_a$ of the $j$ letters,
\begin{equation}
\pi_a=(R,Ra),
\tag{5}
\end{equation}
and
\begin{equation}
\pi_a\pi_{a'}=\pi_{aa'}.
\tag{6}
\end{equation}
The elements $a_0$ satisfying
\begin{equation}
\pi_{a_0}=1
\tag{7}
\end{equation}
are precisely the intersection of all conjugates of $R$,
\[
R,
 b_1^{-1}Rb_1,
 b_2^{-1}Rb_2,
\ldots,b_{j-1}^{-1}Rb_{j-1},
\]
which is a normal divisor of $P$. If this intersection is $1$, then:
\begin{enumerate}[label=\arabic*.,leftmargin=2.2em,start=2]
\item If a group $P$ has a divisor of index $j$ which has no element in common with its conjugate divisors except the identity, then $P$ is simply isomorphic to a transitive permutation group on $j$ letters.
\end{enumerate}

Let now $R$ and $Q$ be two divisors of $P$, $R$ of index $j$ and $Q$ of order $q$. Decompose $P$ as in (3) and consider the permutations $\pi_e=(R,Re)$ induced by the elements $e$ of $Q$. These form a group $\Pi$ one- or many-stage isomorphic to $Q$.

For $a_1\in P$, the complex $Ra_1Q$ contains a certain number $h_1$ of cosets of $R$. This number is the index, in $Q$, of the intersection of $Q$ with $a_1^{-1}Ra_1$. Continuing until $P$ is exhausted gives
\begin{equation}
P=Ra_1Q+Ra_2Q+\cdots=P_1+P_2+
\cdots,
\tag{9}
\end{equation}
and also
\begin{equation}
P=Qa_1^{-1}R+Qa_2^{-1}R+
\cdots .
\tag{10}
\end{equation}
Therefore:
\begin{enumerate}[label=\arabic*.,leftmargin=2.2em,start=3]
\item If $P$ is a group of order $n$, $R$ a divisor of index $j$, and $Q$ a divisor of order $q$, then a permutation group $\Pi$ on $j$ letters is one- or many-stage isomorphic to $Q$. The $j$ letters split into orbits of lengths $h_1,h_2,h_3,
\ldots$, each transitive under $\Pi$, with
\[
j=h_1+h_2+h_3+\cdots,
\tag{11}
\]
and the numbers $h_i$ are divisors of $q$. They are the indices of divisors of $Q$ obtained as intersections of $Q$ with conjugate groups $a^{-1}Ra$.
\end{enumerate}
If those intersections have no common divisor except the identity, as in particular for a simple group $P$, the isomorphism is simple.

\section*{\S\,29. The First Theorem of Sylow.}

It is difficult to proceed far by Cayley's direct construction of groups. For further investigations one needs a theorem, due in special form to Cauchy and in general to Sylow.
\begin{enumerate}[label=\Roman*.,leftmargin=2.4em]
\item If $P$ is a group of order $n$ and $p^x$ is a prime power dividing $n$, then $P$ has a divisor of order $p^x$.
\end{enumerate}

Cauchy proved this for $x=1$. For prime $n$ it is evident, and for the small orders one reads it from \S\,27. We use complete induction and assume it known for groups of order smaller than $n$.

Collect first all elements $a$ commuting with every element $x$ of $P$,
\begin{equation}
x^{-1}ax=a \quad\hbox{or}\quad ax=xa.
\tag{1}
\end{equation}
They form an Abelian group $A$ and a normal divisor of $P$. Let its order and index be $v$ and $j$, so
\begin{equation}
n=vj.
\tag{2}
\end{equation}
If $c$ is not in $A$, the elements $b$ satisfying
\begin{equation}
b^{-1}cb=c
\tag{3}
\end{equation}
form a proper group $B$ of order $\mu$ and index $\varepsilon$,
\begin{equation}
n=\mu\varepsilon.
\tag{4}
\end{equation}
Writing
\begin{equation}
P=B+Bg_1+Bg_2+\cdots+Bg_{\varepsilon-1},
\tag{6}
\end{equation}
one obtains from $g_i^{-1}cg_i$ a system $C$ of $\varepsilon$ elements. Continuing with new elements $c',c'',\ldots$ gives
\begin{equation}
n=vj=\mu\varepsilon=\mu'\varepsilon'=\mu''\varepsilon''=\cdots,
\tag{9}
\end{equation}
\begin{equation}
n=v+\varepsilon+\varepsilon'+\varepsilon''+\cdots .
\tag{10}
\end{equation}

If $v$ is divisible by $p$, then $A$ contains an element $a$ of order $p$, and
\[
Q=1,a,a^2,\ldots,a^{p-1}
\]
is a normal divisor of $P$. The quotient $P/Q$ then gives, by induction, a divisor of $P$ of order $p^x$. If $v$ is not divisible by $p$, not all of $\varepsilon,\varepsilon',\ldots$ can be divisible by $p$ in (10). Choose one, say $\varepsilon$, not divisible by $p$. Then $\mu$ is divisible by $p^x$, and the smaller group $B$ has, by induction, a divisor of order $p^x$. This proves the theorem.

\section*{\S\,30. The Second Theorem of Sylow.}

Let $P$ be of order $n$, and let $p^x$ be the highest power of $p$ dividing $n$. Let $Q$ be a divisor of order $p^x$. The elements $c$ satisfying
\begin{equation}
c^{-1}Qc=Q
\tag{1}
\end{equation}
form a group $R$; $Q$ is a normal divisor of $R$. If $r$ is the index of $Q$ in $R$ and $j$ the index of $R$ in $P$, then
\begin{equation}
n=p^x rj,
\tag{2}
\end{equation}
where $r$ and $j$ are prime to $p$.

\begin{enumerate}[label=\Roman*.,leftmargin=2.4em,start=2]
\item The index $j$ of $R$ is of the form $ph+1$; hence $j\equiv1\pmod p$. There are exactly $j$ different divisors of $P$ of order $p^x$, all conjugate. Every divisor of $P$ whose order is a power of $p$ is contained in one of these conjugates.
\end{enumerate}

The proof is as follows. If $c\in R$ has order $y$, and $s$ is the least positive exponent with $c^s\in Q$, then $s$ divides $y$. The elements $Q,Qc,
\ldots,Qc^{s-1}$ form a group of order $p^xs$, so $s$ is not divisible by $p$. Therefore $R$ contains no element outside $Q$ whose order is a power of $p$.

Choose representatives so that
\begin{equation}
P=R+Rb_1+Rb_2+\cdots+Rb_{j-1}.
\tag{10}
\end{equation}
Then
\begin{equation}
Q,
\ b_1^{-1}Qb_1,
\ b_2^{-1}Qb_2,
\ldots,b_{j-1}^{-1}Qb_{j-1}
\tag{11}
\end{equation}
are $j$ distinct conjugate groups of order $p^x$. Applying \S\,28, theorem 3, gives
\[
j=h_1+h_2+h_3+
\cdots,
\]
with $h_1=1$ and the remaining $h_i$ powers of $p$; hence $j\equiv1\pmod p$. Reapplying the same argument to an arbitrary $p$-power divisor $Q_1$ shows that it is contained in one of the groups (11).

\section*{\S\,31. Groups of Order $p^\alpha$.}

Sylow derived a notable property of groups whose order is a power of a prime. Let $P$ be of order $p^\alpha$ and $Q$ a divisor of order $p^\beta$, $\beta<\alpha$. From \S\,28, theorem 3, applied to $Q$, several conjugates of $Q$ must coincide with $Q$. Hence there is an element $b\notin Q$ such that
\[
b^{-1}Qb=Q.
\]
If $b^r$ is the least power of $b$ lying in $Q$, then $r$ is a power of $p$, and
\[
b^{r/p}=c
\]
is outside $Q$, has its $p$-th power in $Q$, and satisfies
\[
cQc^{-1}=Q.
\]
Thus
\[
R=Q+Qc+Qc^2+
\cdots+Qc^{p-1}
\]
is a group of order $p^{\beta+1}$ and $Q$ is a normal divisor of it. Therefore:
\begin{enumerate}[label=\Roman*.,leftmargin=2.4em,start=3]
\item If $P$ is of order $p^\alpha$ and $Q$ is a divisor of order $p^\beta$, $\beta<\alpha$, then there is a divisor $R$ of $P$ of order $p^{\beta+1}$ in which $Q$ is normal.
\end{enumerate}
For $\beta=\alpha-1$ it follows that $Q$ is a normal divisor of $P$. Repeated application gives:
\begin{enumerate}[label=\Roman*.,leftmargin=2.4em,start=4]
\item Every group whose order is a prime power is metacyclic.
\end{enumerate}

\section*{\S\,32. Frobenius's Theorem.}

Frobenius stated a theorem which may be regarded as a counterpart to Sylow's fourth theorem:
\begin{enumerate}[label=\Roman*.,leftmargin=2.4em,start=5]
\item Every group whose order is a product of distinct primes is metacyclic.
\end{enumerate}
Suppose then that the order $n$ of $P$ is divisible by no square. Let $t$ be the greatest prime divisor of $n$. By Cauchy-Sylow there is a divisor $T$ of order $t$, and it is cyclic,
\[
T=1,a,a^2,\ldots,a^{t-1}.
\]
It is enough to show that $T$ is a normal divisor, for then the same argument applied to $P/T$ yields a composition series with prime indices.

The proof uses the following general assertion. If $n=\mu v$, $\mu$ is square-free, and every prime divisor of $v$ is greater than every prime divisor of $\mu$, then every group of order $\mu v$ has exactly $v$ elements whose order divides $v$. This is proved by induction over the number of prime divisors of $\mu$.

Let $p$ be the greatest prime divisor of $\mu$, and let $Q$ have order $p$. Let $R$ be the group of elements satisfying $c^{-1}Qc=Q$. Put its order $pr$ and its index in $P$ equal to $j$. Write $r=r_1r_2$, where $r_1$ is the greatest common divisor of $\mu$ and $r$, and $r_2$ that of $v$ and $r$. Then
\[
n=\mu v=r_1 p r_2 j.
\]
The induction hypothesis gives exactly $pv$ elements of $P$ whose order divides $pv$; call their set $U$. In $R$ there are exactly $pr_2$ such elements; call their set $V$.

If $v\in V$, all elements of $vQ$ again lie in $V$. Among them $p-1$ have order divisible by $p$ and one has order not divisible by $p$. For from
\[
v^{-1}av=a^s
\tag{1}
\]
one gets
\[
v^{-h}av^h=a^{s^h}
\tag{2}
\]
and hence, if $h$ is the order of $v$,
\[
s^h\equiv1\pmod p.
\tag{3}
\]
Since $p-1$ is prime to $h$, Fermat's theorem gives $s\equiv1\pmod p$, so $av=va$. Splitting $V$ into systems $v_iQ$ shows that $V$ contains exactly $(p-1)r_2$ elements whose orders are divisible by $p$.

Each element of $P$ whose order is divisible by $p$ lies in exactly one of the conjugates of $R$. There are $j$ such conjugates, and consequently $U$ contains exactly $jr_2(p-1)$ elements whose orders are divisible by $p$. It remains only to prove $jr_2=v$. Since $jr_2(p-1)<pv$ and $jr_2$ is divisible by $v$, this follows. The general assertion is proved, and with $v=t$ it shows that $T$ is normal in $P$.

\section*{\S\,33. Groups of Order $p^\alpha q$.}

Frobenius stated:
\begin{enumerate}[label=\Roman*.,leftmargin=2.4em,start=6]
\item Every group of order $p^\alpha q$, where $p$ and $q$ are distinct primes and $\alpha$ is positive, is metacyclic.
\end{enumerate}
It is enough to show that no group of such order is simple. For if $P$ has a proper normal divisor $Q$ different from the identity, then $Q$ and $P/Q$ are metacyclic by induction or by the preceding theorem, and hence so is $P$.

Assume therefore that $P$ is simple. If $P$ has a divisor of order $p^\beta q$, with $0\le\beta<\alpha$ and $\beta$ chosen maximal, then by \S\,28, theorem 2, $P$ is represented as a permutation group on $p^{\alpha-\beta}$ letters. Since it contains an element with order divisible by $q$,
\begin{equation}
p^{\alpha-\beta}>q.
\tag{1}
\end{equation}

Let $Q$ be a divisor of order $p^\alpha$ and index $q$. The group of elements normalizing $Q$ cannot be all of $P$; since the index is prime, it is $Q$ itself. Hence there are $q$ conjugates
\[
Q,Q_1,Q_2,\ldots,Q_{q-1}.
\]
Choose $Q$ and $Q_1$ so that their intersection $Q_0$ has maximal order $p^\gamma$. Then $\gamma<\alpha$. Applying \S\,28, theorem 3, to $Q$ gives
\[
q=h_1+h_2+h_3+
\cdots,
\]
with $h_1=1$ and the other summands divisible by $p^{\alpha-\gamma}$. Thus
\begin{equation}
q>p^{\alpha-\gamma},
\tag{2}
\end{equation}
and from (1) and (2)
\begin{equation}
\gamma>\beta,
\tag{3}
\end{equation}
so in particular $\gamma>0$.

By \S\,31 there are divisors $R\subset Q$ and $R_1\subset Q_1$ of order $p^{\gamma+1}$ in which $Q_0$ is normal. Let $S$ be the least group containing both. Its order cannot be a power of $p$, for then $R$ and $R_1$ would have to lie in one of the conjugates $Q,Q_1,\ldots$. Hence $|S|=p^\lambda q$. Since $S$ contains a group of order $p^{\gamma+1}$,
\[
\lambda\ge\gamma+1>\beta.
\]
By maximality of $\beta$, one must have $\lambda=\alpha$, so $S=P$. Since $Q_0$ is normalized by both $R$ and $R_1$, it is normalized by all of $P$. But $Q_0\ne1$, contradicting the simplicity of $P$.

\section*{\S\,34. Simple Groups.}

The preceding theorems permit far-reaching conclusions about groups. In low orders metacyclic and cyclic groups strongly predominate; the non-metacyclic groups, and especially the non-cyclic simple groups, are therefore of special interest.

No general law for the orders of simple groups is known here. Hölder carried out the determination up to order $200$, Cole up to order $500$. Among non-cyclic simple groups only the already known orders
\[
60,
\quad 168,
\quad 360
\]
occurred, with further known examples of orders $660$ and $504$. We restrict ourselves to the first hundred. The Sylow and Frobenius theorems leave only
\[
36,
\quad 60,
\quad 72,
\quad 84,
\quad 90,
\quad 100.
\]
The orders $36,72,84,100$ are excluded by the preceding arguments and the Sylow theorems.

It remains to consider $60$ and $90$. A simple group of order $60$ exists: the alternating group on five letters. We show it is unique. Let $P$ be simple of order $60$. It contains divisors of orders $5$ and $3$. Using the second Sylow theorem for the divisor of order $5$, the normalizer has order $10$, so $P$ can be represented as a transitive permutation group on six letters. It cannot contain a 3-cycle; otherwise, being primitive, it would contain the full alternating group on six letters. Hence we may take an element of order three as
\[
a=(1,2,3)(4,5,6).
\]
Add a cyclic permutation $b$ of order five. Omitting the letter $6$ and taking a suitable power of $b$, there are six initial possibilities:
\[
\begin{array}{lll}
(1,2,3,4,5),&(1,2,4,3,5),&(1,2,4,5,3),\\
(1,2,3,5,4),&(1,2,5,3,4),&(1,2,5,4,3).
\end{array}
\]
Up to replacing generators by powers and renaming letters, only two remain:
\[
\begin{array}{ll}
a=(1,2,3)(4,5,6),& b=(1,2,3,4,5),\\
a=(1,2,3)(4,5,6),& b=(1,2,3,5,4).
\end{array}
\]
The first is impossible since $a^2b=(1,4,6)$ would be a 3-cycle. Thus only
\[
a=(1,2,3)(4,5,6),\qquad b=(1,2,3,5,4)
\]
remains, and there is only one simple group of order $60$ up to equivalence.

Its elements, besides the identity, are: the permutations of order five
\[
\begin{array}{lll}
(2,3,6,5,4),&(1,3,5,6,4),&(1,2,5,4,6),\\
(1,2,6,3,5),&(1,2,4,6,3),&(1,2,3,5,4),
\end{array}
\]
which with their powers give $24$ elements; the permutations of order three
\[
\begin{array}{ll}
(1,2,3)(4,5,6),&(1,3,5)(2,4,6),\\
(1,2,4)(3,6,5),&(1,3,6)(2,4,5),\\
(1,2,5)(3,6,4),&(1,4,5)(2,3,6),\\
(1,2,6)(3,4,5),&(1,4,6)(2,3,5),\\
(1,3,4)(2,6,5),&(1,5,6)(2,3,4),
\end{array}
\]
which with their squares give $20$ elements; and the $15$ permutations of order two
\[
\begin{array}{lll}
(1,2)(3,4),&(1,2)(5,6),&(3,4)(5,6),\\
(1,3)(4,5),&(1,3)(2,6),&(2,6)(4,5),\\
(1,4)(2,5),&(1,4)(3,6),&(2,5)(3,6),\\
(1,5)(2,3),&(1,5)(4,6),&(2,3)(4,6),\\
(1,6)(2,4),&(1,6)(3,5),&(2,4)(3,5).
\end{array}
\]
This simple group of order $60$ is called the icosahedral group. The same consideration excludes a simple group of order $90$, since it too would have to be a transitive group on six letters and would lead to the same generators, hence to the icosahedral group, which cannot be a divisor of a group of order $90$.

\section*{\S\,35. Groups of Order $pq$.}

It remains useful to see how groups are constituted for a given order. We treat the simple case in which the order is $pq$, the product of two primes, possibly equal.

Such a group $P$ is metacyclic and has a normal divisor $Q$ of order $p$, a cyclic group
\begin{equation}
Q=1,a,a^2,\ldots,a^{p-1},\qquad a^p=1.
\tag{1}
\end{equation}
If $b\notin Q$, then $P$ has the form
\begin{equation}
P=Q+Qb+Qb^2+\cdots+Qb^{q-1}.
\tag{2}
\end{equation}
Let $g$ be a primitive root modulo $p$ and $v$ an exponent modulo $p-1$. Since $b^{-1}ab\in Q$,
\begin{equation}
b^{-1}ab=a^{g^v}.
\tag{3}
\end{equation}
Then
\begin{equation}
b^{-1}a^\alpha b=a^{\alpha g^v},
\tag{4}
\end{equation}
and
\begin{equation}
b^{-\beta}a^\alpha b^\beta=a^{\alpha g^{\beta v}}.
\tag{5}
\end{equation}
Since one may assume $b^q=1$,
\begin{equation}
g^{qv}\equiv1\pmod p,
\qquad qv\equiv0\pmod{p-1}.
\tag{6}
\end{equation}
If $p-1$ is not divisible by $q$, then $v\equiv0$, so the group is Abelian. If
\[
p\equiv1\pmod q,
\]
then $v$ may take the $q$ values
\begin{equation}
v=0,\frac{p-1}{q},2\frac{p-1}{q},\ldots,(q-1)\frac{p-1}{q}.
\tag{7}
\end{equation}
The value $v=0$ gives the Abelian group; the other values give non-commutative groups, isomorphic among themselves. They are represented by the elements
\begin{equation}
\Theta=a^\alpha b^\beta,
\tag{8}
\end{equation}
where $\alpha=0,1,\ldots,p-1$ and $\beta=0,1,\ldots,q-1$, with rules
\begin{equation}
b^{-\beta}a^\alpha=a^{\alpha g^{-\beta v}}b^{-\beta},\qquad a^p=1,
\quad b^q=1,
\tag{9}
\end{equation}
and multiplication
\begin{equation}
a^\alpha b^\beta a^{\alpha'}b^{\beta'}=a^{\alpha+\alpha'g^{-\beta v}}b^{\beta+\beta'}.
\tag{10}
\end{equation}
This formula immediately gives cancellation and associativity.

\section*{\S\,36. Bounds for the Index of a Divisor of the Symmetric Permutation Group.}

We close these considerations with a theorem on permutation groups. Let $P$ be the symmetric group on $n$ letters. Besides the alternating group of index $2$, there is the intransitive divisor of index $n$ fixing one letter.

\begin{enumerate}[label=\Roman*.,leftmargin=2.4em]
\item The index of an imprimitive divisor of $P$ is always greater than $n$; the index of an intransitive divisor is at least $n$, and is equal to $n$ only when the divisor fixes one letter and permutes the remaining $n-1$ letters arbitrarily.
\end{enumerate}
If $Q$ is imprimitive of index $j$, with $r$ systems of imprimitivity of $s$ letters, $n=rs$, then
\begin{equation}
j\ge {\Pi(n)\over [\Pi(s)]^r\Pi(r)}.
\tag{1}
\end{equation}
Writing the quotient as in the source gives
\begin{equation}
j>{n\over2}\left({r+1\over2}\right)^{r-1}\left({2r\over3}\right)^r\cdots
\left({(s-1)r\over s}\right)^r>n.
\tag{2}
\end{equation}
If $Q$ is intransitive and splits the letters into two systems of $a$ and $b$ letters, $n=a+b$, then
\begin{equation}
j\ge {\Pi(n)\over \Pi(a)\Pi(b)}={n\over1}{n-1\over2}\cdots {b+1\over a}.
\tag{3}
\end{equation}
For $b\ge a$ this is greater than $n$ except for $a=1,b=n-1$.

Next:
\begin{enumerate}[label=\Roman*.,leftmargin=2.4em,start=2]
\item Apart from the alternating group, there is no transitive primitive divisor $Q$ of $P$ whose index is $\le n$, except for $n=4$ and $n=6$.
\end{enumerate}
We use the earlier theorem that a primitive transitive divisor not containing the alternating group contains neither a transposition nor a 3-cycle. Decompose
\begin{equation}
P=Q+Q\pi_1+Q\pi_2+\cdots+Q\pi_{j-1}.
\tag{4}
\end{equation}
The transpositions
\begin{equation}
(0,1),(0,2),\ldots,(0,n-1)
\tag{5}
\end{equation}
are not in $Q$. If $j<n$, two of them lie in the same coset and their quotient gives a 3-cycle in $Q$, contradiction. Thus $j\not<n$.

If $j=n$, then
\begin{equation}
P=Q+Q(0,1)+Q(0,2)+\cdots+Q(0,n-1).
\tag{6}
\end{equation}
Considering another transposition, say $(2,3)$, one may assume it lies in $Q(0,1)$, and hence $Q$ contains
\begin{equation}
(0,1)(2,3).
\tag{7}
\end{equation}
Let $Q_0$ be the divisor fixing $0$. Since $Q$ is transitive,
\begin{equation}
Q=Q_0+Q_0\chi_1+\cdots+Q_0\chi_{n-1},
\tag{8}
\end{equation}
and
\begin{equation}
g={\Pi(n-1)\over n}.
\tag{9}
\end{equation}
Thus $j=n$ cannot occur for prime $n$. More generally, for $n\ne4$, $Q_0$ is transitive on $1,2,\ldots,n-1$.

Let $Q_{0,1}$ fix $0$ and $1$. Then
\begin{equation}
g_1={\Pi(n-1)\over n(n-1)}={\Pi(n-2)\over n}.
\tag{10}
\end{equation}
For $n\ne6$, $Q_{0,1}$ is transitive on $2,3,
\ldots,n-1$. If $n>6$, then $Q_{0,1,2}$ has order
\[
g_2={g_1\over n-2}={\Pi(n-3)\over n}
\]
and cannot fix a fourth letter. Hence there is $\chi\in Q$ carrying $0,1,2,3$ to $0,1,2,4$. From (7),
\[
\chi^{-1}(0,1)(2,3)\chi=(0,1)(2,4),
\]
and consequently
\[
(0,1)(2,3)(0,1)(2,4)=(2,3,4),
\]
a 3-cycle, contradiction. The cases $n=4$ and $n=6$ are genuine exceptions, as shown earlier: the symmetric group on four letters has a transitive divisor of index $3$, and that on six letters has one of index $6$.

\clearpage
\section*{Sixth Section. Groups of Linear Substitutions.}

\section*{\S\,37. Linear substitutions and their composition.}

One of the most effective means for forming groups, and one which at the same time leads to many applications, is supplied by linear substitutions and their composition. We have already met such linear substitutions more than once: for instance in the second section of the first volume in connection with the multiplication law for determinants, and again in the eleventh section in connection with equivalent numbers.

By a linear substitution in $n$ variables we understand a system of equations by which a system of $n$ variables $y_1,y_2,\ldots,y_n$ is expressed linearly in terms of another system $x_1,x_2,\ldots,x_n$. According to the number of variables we speak of unary, binary, ternary, quaternary substitutions, and in general we shall call the number of variables the dimension of the substitution. For the present we restrict ourselves to homogeneous substitutions, so that, if the coefficients are denoted by $a_i^{(k)}$, the substitution is represented by the system
\begin{equation}
y_k=\sum_{i=1}^{n} a_i^{(k)}x_i,\qquad k=1,2,\ldots,n.
\tag{1}
\end{equation}
Often the variables themselves are immaterial, so that such a substitution is sufficiently characterized by its coefficients $a_i^{(k)}$. One also uses a single letter, say $A$, to denote the substitution, writing
\begin{equation}
A=\begin{pmatrix}
a_1^{(1)}&a_2^{(1)}&\cdots&a_n^{(1)}\\
a_1^{(2)}&a_2^{(2)}&\cdots&a_n^{(2)}\\
\cdots&\cdots&&\cdots\\
a_1^{(n)}&a_2^{(n)}&\cdots&a_n^{(n)}
\end{pmatrix}.
\tag{2}
\end{equation}
We shall also write briefly $A=(a_i^{(k)})$, and denote the determinant of the substitution-coefficients by $|A|$. This determinant will always be assumed different from zero. We shall also represent the system (1) symbolically by
\begin{equation}
(y_1,y_2,\ldots,y_n)=A(x_1,x_2,\ldots,x_n),
\tag{3}
\end{equation}
or, still more briefly,
\begin{equation}
(y)=A(x).
\tag{4}
\end{equation}
The substitution
\begin{equation}
y_1=x_1,
\quad y_2=x_2,
\quad \ldots,
\quad y_n=x_n
\tag{5}
\end{equation}
or
\begin{equation}
J=\begin{pmatrix}
1&0&\cdots&0\\
0&1&\cdots&0\\
\cdots&\cdots&&\cdots\\
0&0&\cdots&1
\end{pmatrix}
\tag{6}
\end{equation}
is called the identical substitution. A substitution of the form
\begin{equation}
y_1=\mu_1x_1,
\quad y_2=\mu_2x_2,
\quad\ldots,
\quad y_n=\mu_nx_n
\tag{7}
\end{equation}
or
\begin{equation}
M=\begin{pmatrix}
\mu_1&0&\cdots&0\\
0&\mu_2&\cdots&0\\
\cdots&\cdots&&\cdots\\
0&0&\cdots&\mu_n
\end{pmatrix}
\tag{8}
\end{equation}
will be called a multiplicative substitution, or more briefly a multiplication; the elements $\mu_1,\mu_2,\ldots,\mu_n$ are called its multipliers.

Let a second linear substitution of dimension $n$ be given, introducing a new system of variables $z$:
\begin{equation}
z_k=\sum_{h=1}^{n}b_h^{(k)}y_h,
\qquad
(z)=B(y).
\tag{9}
\end{equation}
If we insert for $y$ the values from (1), we obtain $z$ expressed in terms of $x$ by a new linear substitution $E$, denoted by
\begin{equation}
z=E(x)=BA(x).
\tag{10}
\end{equation}
The substitution $E=BA$ is said to be composed from $B$ and $A$. One must distinguish $AB$ from $BA$. Two substitutions $A,B$ having the special property that $AB=BA$ are said to be interchangeable, or commutative. The coefficients of $E$ are obtained by substituting (1) in (9):
\begin{equation}
z_k=\sum_{i=1}^{n}e_i^{(k)}x_i,
\qquad
 e_i^{(k)}=\sum_{h=1}^{n}b_h^{(k)}a_i^{(h)}.
\tag{11}
\end{equation}
Thus the law of composition of substitutions can be stated as follows:
\begin{enumerate}[label=\arabic*.,leftmargin=2.2em]
\item To form the substitution $BA$ composed from two substitutions $B,A$, one proceeds exactly as if the two determinants $|B|,|A|$ were to be multiplied according to the determinant multiplication rule. To form the elements of a row, the elements of a row of the first component are multiplied by the corresponding elements of the columns of the second component and then added. Briefly: in the first component one sums by rows, in the second by columns.
\item The determinant of a composed substitution is the product of the determinants of the components.
\end{enumerate}

To compose three substitutions of the same dimension, $C,B,A$, one introduces the expressions (10) for $z$ into a new linear substitution
\begin{equation}
(u)=C(z)
\tag{12}
\end{equation}
and expresses $u$ in terms of $x$:
\begin{equation}
(u)=CBA(x).
\tag{13}
\end{equation}
It is plainly immaterial whether one inserts (10) in (12), or first expresses $z$ in terms of $y$ by (9) and then $y$ in terms of $x$ by (4). Hence this composition satisfies the associative law
\begin{equation}
C(BA)=(CB)A=CBA.
\tag{14}
\end{equation}
For the elements in both cases one finds the same expression,
\[
\sum_i\sum_h c_i^{(k)}b_h^{(i)}a_j^{(h)}.
\]
Thus:
\begin{enumerate}[label=\arabic*.,leftmargin=2.2em,start=3]
\item In the composition of linear substitutions the associative law holds, but the commutative law need not hold.
\end{enumerate}

A multiplication in which all multipliers are equal,
\[
N=\begin{pmatrix}
\nu&0&\cdots&0\\
0&\nu&\cdots&0\\
\cdots&\cdots&&\cdots\\
0&0&\cdots&\nu
\end{pmatrix},
\]
will be called a similarity substitution. Two substitutions which, like $A$ and $AN$, or $A$ and $NA$, can be derived from one another by composition with a similarity substitution are called similar substitutions. From the law of composition one immediately obtains:
\begin{enumerate}[label=\arabic*.,leftmargin=2.2em,start=4]
\item A similarity substitution is interchangeable with every substitution $A$ of the same dimension; two similarity substitutions composed with each other again give a similarity substitution; and two substitutions similar to a third are also similar to one another.
\item Composition with the identical substitution $J$ changes no substitution.
\end{enumerate}
The substitution $J$ is therefore regarded as the unit for composition and may be omitted wherever it occurs composed with other substitutions.

Further:
\begin{enumerate}[label=\arabic*.,leftmargin=2.2em,start=6]
\item To every substitution $A$ there is one and only one inverse substitution $A^{-1}$ satisfying
\begin{equation}
AA^{-1}=A^{-1}A=J.
\tag{15}
\end{equation}
\end{enumerate}
This follows from the fundamental formulae of determinant theory. If the elements $\alpha_h^{(k)}$ of $A^{-1}$ are determined from the linear equations
\begin{equation}
\sum_i a_i^{(h)}\alpha_k^{(i)}=0\quad(h\ne k),\qquad
\sum_i a_i^{(h)}\alpha_h^{(i)}=1,
\tag{16}
\end{equation}
then, writing $A_k^{(h)}$ for the cofactors of $|A|$, one obtains
\begin{equation}
|A|\alpha_h^{(k)}=A_k^{(h)},
\tag{17}
\end{equation}
and these equations imply the other system
\begin{equation}
\sum_i a_k^{(i)}\alpha_i^{(h)}=0\quad(h\ne k),\qquad
\sum_i a_h^{(i)}\alpha_i^{(h)}=1.
\tag{18}
\end{equation}
The inverse substitution to $A$ is nothing but the solution of the system (1) of the direct substitution $A$; thus from $(y)=A(x)$ follows
\begin{equation}
(x)=A^{-1}(y).
\tag{19}
\end{equation}
For the composition of inverse substitutions, $ABB^{-1}A^{-1}=J$ gives
\begin{equation}
(AB)^{-1}=B^{-1}A^{-1}.
\tag{20}
\end{equation}
If $A,B,C$ are three substitutions, then composition with $A^{-1}$ in either of the equations $AB=AC$ and $BA=CA$ gives $B=C$. Hence the characteristic properties of a group are fulfilled for composition of substitutions, and the totality of all substitutions of a fixed dimension is an infinite group. Any subset of it closed under composition is likewise a group, finite or infinite.

If $L$ is a fixed substitution, then from any substitution $A$ of the same dimension one obtains
\begin{equation}
A'=L^{-1}AL\qquad(LA'=AL),
\tag{21}
\end{equation}
called the transform of $A$ by $L$. If
\begin{equation}
(x)=L(x'),\qquad (y)=L(y'),
\tag{22}
\end{equation}
then (4) gives
\begin{equation}
(y')=A'(x'),
\tag{23}
\end{equation}
so that passing to the transformed substitution is equivalent to simultaneously transforming both systems of variables by $L^{-1}$. If
\[
A'=L^{-1}AL,
\qquad
B'=L^{-1}BL,
\]
then $A'B'=L^{-1}ABL$, and therefore:
\begin{enumerate}[label=\arabic*.,leftmargin=2.2em,start=7]
\item If $A$ runs through the substitutions of a group, then for fixed $L$ the transforms $L^{-1}AL$ run through the substitutions of an isomorphic group.
\end{enumerate}

Inverse substitutions must be distinguished from transposed substitutions. The transposed substitution to $A$, denoted here by $A_1$, is obtained by making the rows of $A$ into columns and conversely, i.e. by interchanging the upper and lower indices of the coefficients in (2).

From the coefficient formula (11) one obtains the composition law for transposed substitutions:
\begin{enumerate}[label=\arabic*.,leftmargin=2.2em,start=8]
\item If $A_1,B_1$ are the transposes of $A,B$, then $A_1B_1$ is the transpose of $BA$:
\begin{equation}
(BA)_1=A_1B_1.
\tag{24}
\end{equation}
\end{enumerate}

Multiplying the formulae (1) by arbitrary factors $\eta_k$ and summing over $k$ gives
\begin{equation}
\sum_k y_k\eta_k=\sum_i\sum_k x_i a_i^{(k)}\eta_k.
\tag{25}
\end{equation}
If
\begin{equation}
\sum_k a_i^{(k)}\eta_k=\xi_i,
\tag{26}
\end{equation}
then
\begin{equation}
\sum_k y_k\eta_k=\sum_i x_i\xi_i.
\tag{27}
\end{equation}
Thus (26) expresses the substitution $(\xi)=A_1(\eta)$. Hence:
\begin{enumerate}[label=\arabic*.,leftmargin=2.2em,start=9]
\item If $A,A_1$ are transposed substitutions of dimension $n$, and if $x,y,\xi,\eta$ are four systems of variables connected by
\begin{equation}
(y)=A(x),\qquad (\xi)=A_1(\eta),
\tag{28}
\end{equation}
then
\begin{equation}
y_1\eta_1+y_2\eta_2+
\cdots+y_n\eta_n=x_1\xi_1+x_2\xi_2+
\cdots+x_n\xi_n.
\tag{29}
\end{equation}
\end{enumerate}
When two rows of variables
\[
x_1,x_2,\ldots,x_n,
\qquad
\xi_1,\xi_2,\ldots,\xi_n
\]
are simultaneously transformed by (28) into
\[
y_1,y_2,\ldots,y_n,
\qquad
\eta_1,\eta_2,\ldots,\eta_n,
\]
the two rows are called contragredient, and (28) are contragredient transformations.

\section*{\S\,38. Substitution of ratios.}

Since two substitutions similar to a third are similar to one another, all mutually similar substitutions of dimension $n$ may be collected into one class, and every substitution belongs to one and only one such class. The individual substitutions in a class are its representatives; a class is completely determined by any one representative.

If $A$ is similar to $A'$ and $B$ to $B'$, then $AB$ is similar to $A'B'$. Thus if $\mathfrak A,\mathfrak B$ denote the classes containing $A,A'$ and $B,B'$, the class represented by $AB$ or $A'B'$ is independent of the representatives chosen. We call it the composition of $\mathfrak A$ and $\mathfrak B$ and denote it by $\mathfrak A\mathfrak B$. Under this composition the same rules hold as for substitutions themselves; the totality of the classes is therefore also a group.

This way of regarding substitutions is useful whenever only ratios of variables matter, as in projective geometry. We call the resulting classes substitutions of ratios. We shall denote them in the same way as ordinary substitutions, by choosing a representative of the class. The representative may be selected according to various principles. Often it is convenient to choose it so that the determinant is $1$. In dimension $n$ this determines the multiplier only up to an $n$th root of unity; but if representatives are chosen in this way, the property is preserved under composition, and these representatives themselves form a substitution group.

In the binary case the substitution
\[
(y_1,y_2)=\begin{pmatrix}a&b\\c&d\end{pmatrix}(x_1,x_2)
\]
is, as a substitution of ratios, equivalent to the fractional linear substitution
\[
\eta={a\xi+b\over c\xi+d},
\]
where $\xi=x_1:x_2$ and $\eta=y_1:y_2$. If a second substitution
\[
(z_1,z_2)=\begin{pmatrix}a'&b'\\c'&d'\end{pmatrix}(y_1,y_2),
\qquad
\zeta={a'\eta+b'\over c'\eta+d'}
\]
is added, the composed substitution $A''=A'A$ expresses $\zeta$ in terms of $\xi$ as
\[
\zeta={a''\xi+b''\over c''\xi+d''},
\]
with
\[
\begin{pmatrix}a''&b''\\c''&d''\end{pmatrix}
=\begin{pmatrix}a'&b'\\c'&d'\end{pmatrix}
\begin{pmatrix}a&b\\c&d\end{pmatrix}
=\begin{pmatrix}a'a+b'c&a'b+b'd\\c'a+d'c&c'b+d'd\end{pmatrix}.
\]
If such a substitution is multiplicative, i.e. of the form
\[
\begin{pmatrix}a&0\\0&d\end{pmatrix},
\]
we shall also call the ratio $a:d$ its multiplier.

\section*{\S\,39. Permutations as linear substitutions.}

The importance of linear substitutions, and especially of finite groups formed from them, for algebra follows from the fact that permutation groups on $n$ elements may be regarded as special cases of groups of linear substitutions.

Let $x_1,x_2,\ldots,x_n$ be $n$ variables, and let $\alpha_1,\alpha_2,\ldots,\alpha_n$ be an arrangement of the numbers $1,2,\ldots,n$. Then
\begin{equation}
x_1=x'_{\alpha_1},
\quad
x_2=x'_{\alpha_2},
\quad\ldots,\quad
x_n=x'_{\alpha_n}
\tag{1}
\end{equation}
defines a linear substitution of dimension $n$,
\begin{equation}
(x)=A(x'),
\tag{2}
\end{equation}
whose matrix has exactly one non-zero coefficient in each row and each column, that coefficient being $1$. Hence $|A|=\pm1$. After performing the substitution, if $x'_i$ is again written as $x_i$, the result is just the permutation
\[
\begin{pmatrix}1&2&\cdots&n\\\alpha_1&\alpha_2&\cdots&\alpha_n\end{pmatrix}
\]
of the indices of $x$.

Let $B$ be a second such substitution,
\begin{equation}
(x')=B(x''),
\tag{3}
\end{equation}
that is
\begin{equation}
x'_1=x''_{\beta_1},
\quad x'_2=x''_{\beta_2},
\quad\ldots,
\quad x'_n=x''_{\beta_n}.
\tag{4}
\end{equation}
Composition gives
\begin{equation}
x_1=x''_{\beta_{\alpha_1}},
\quad x_2=x''_{\beta_{\alpha_2}},
\quad\ldots,
\quad x_n=x''_{\beta_{\alpha_n}},
\tag{5}
\end{equation}
or briefly
\begin{equation}
(x)=AB(x'').
\tag{6}
\end{equation}
But by the rule for composing permutations,
\begin{equation}
\begin{pmatrix}1&2&\cdots&n\\\alpha_1&\alpha_2&\cdots&\alpha_n\end{pmatrix}
\begin{pmatrix}1&2&\cdots&n\\\beta_1&\beta_2&\cdots&\beta_n\end{pmatrix}
=
\begin{pmatrix}1&2&\cdots&n\\\beta_{\alpha_1}&\beta_{\alpha_2}&\cdots&\beta_{\alpha_n}\end{pmatrix}.
\tag{7}
\end{equation}
Thus if the substitutions $A,B$ are regarded as permutations of indices, the substitution $AB$ is precisely the composed permutation $AB$.

The permutation groups of $n$ letters are therefore special cases of finite groups of linear substitutions of dimension $n$. Even permutations correspond to substitutions with determinant $+1$, and odd permutations to substitutions with determinant $-1$. This follows because a transposition, say $(1,2)$, corresponds to the substitution
\[
\begin{pmatrix}
0&1&0&\cdots&0\\
1&0&0&\cdots&0\\
0&0&1&\cdots&0\\
\cdots&\cdots&\cdots&&\cdots\\
0&0&0&\cdots&1
\end{pmatrix},
\]
whose determinant is $-1$, and because every even permutation is a product of an even number of transpositions, every odd one a product of an odd number.

\section*{\S\,40. The invariants of finite groups of linear substitutions.}

We now restrict attention to finite groups of homogeneous linear substitutions. The permutation groups show that such groups exist; later we shall meet others. Denote such a group by $S$ and its substitutions by $A,B,C,\ldots$. If in a function of the variables $x_1,x_2,\ldots,x_n$ we replace the variables $(x)$ by $A(x)$, we say that the substitution $A$ has been applied to $(x)$.

\begin{enumerate}[label=\arabic*.,leftmargin=2.2em]
\item A form $\Phi(x_1,x_2,\ldots,x_n)$ is called an invariant, or invariant form, of the group $S$ if it remains unchanged when all substitutions $A,B,C,\ldots$ of the group are applied to $(x)$.
\end{enumerate}
Thus
\begin{equation}
\Phi[A(x)]=\Phi[B(x)]=\Phi[C(x)]=\cdots.
\tag{1}
\end{equation}
Here, as before, a form means an integral homogeneous function of the variables.

For every finite group $S$ there are invariant forms in abundance. Choose any form $\varphi(x)$ and form
\begin{equation}
\varphi[A(x)],\quad\varphi[B(x)],\quad\varphi[C(x)],\ldots
\tag{2}
\end{equation}
for all substitutions of the group; any symmetric function of these forms is an invariant, because applying a substitution of $S$ only permutes the set (2).

The notion of invariant may be generalized:
\begin{enumerate}[label=\arabic*.,leftmargin=2.2em,start=2]
\item A form $\Psi(x_1,x_2,\ldots,x_n)$ is also called an invariant of $S$ if it is multiplied by constant factors under the substitutions of $S$.
\end{enumerate}
Thus
\begin{equation}
\Psi[A(x)]=\alpha\Psi(x),
\quad
\Psi[B(x)]=\beta\Psi(x),
\quad
\Psi[C(x)]=\gamma\Psi(x),\ldots,
\tag{3}
\end{equation}
where $\alpha,\beta,\gamma,\ldots$ are independent of $x$. When distinction is needed, the forms in 1. are called absolute invariants, those in 2. relative invariants.

From (3) follows
\begin{equation}
\Psi[AB(x)]=\alpha\beta\Psi(x),
\tag{4}
\end{equation}
so the factors $\alpha,\beta,\gamma,\ldots$ form, under multiplication, a group. This abelian group is generally only multiply isomorphic to $S$, since different substitutions of $S$ may give the same factor.

If $\mu$ is the order of $S$, then the order of every element of $S$ divides $\mu$, and hence $A^\mu=B^\mu=\cdots=1$. Therefore all factors $\alpha,\beta,\gamma,\ldots$ are roots of unity. Let $e$ be the least positive integer satisfying
\begin{equation}
\alpha^e=\beta^e=\gamma^e=\cdots=1.
\tag{5}
\end{equation}
Then the factors are $e$th roots of unity; $e$ is called the index of the invariant $\Psi(x)$. If $\varepsilon$ is a primitive $e$th root, write
\[
\alpha=\varepsilon^a,
\quad
\beta=\varepsilon^b,
\quad
\gamma=\varepsilon^c,
\ldots.
\]
The exponents $a,b,c,\ldots$ have no common divisor with $e$, so integers $x,y,z,\ldots$ may be chosen such that
\[
ax+by+cz+\cdots\equiv1\pmod e.
\]
Since the factors themselves form a group, the factor
\[
\alpha^x\beta^y\gamma^z\cdots=\varepsilon
\]
also occurs. Hence all factors are precisely the powers
\[
1,\varepsilon,\varepsilon^2,\ldots,\varepsilon^{e-1}.
\]

The substitutions $A$ of $S$ for which
\begin{equation}
\Psi[A(x)]=\Psi(x)
\tag{4'}
\end{equation}
form a subgroup $T$ of $S$. If $E$ is a substitution with
\begin{equation}
\Psi[E(x)]=\varepsilon\Psi(x),
\tag{5'}
\end{equation}
then all substitutions $AE$ and $EA$ have the same property; hence
\begin{equation}
S=T+TE+TE^2+\cdots+TE^{e-1}.
\tag{6}
\end{equation}
Thus the index of $T$ is $e$. Since every $E^{-1}AE$ again satisfies (4'),
\begin{equation}
E^{-1}TE=T,
\tag{7}
\end{equation}
so $T$ is a normal subgroup of $S$. Thus:
\begin{enumerate}[label=\arabic*.,leftmargin=2.2em,start=3]
\item If $\Psi(x)$ is an invariant of $S$ of index $e$, then the substitutions of $S$ leaving $\Psi(x)$ unchanged form a normal subgroup $T$ of index $e$.
\end{enumerate}
For $T$, $\Psi(x)$ is an absolute invariant. Index $1$ invariants are precisely the absolute invariants of $S$.

As an example, for the symmetric permutation group $P$ on $x_1,x_2,\ldots,x_n$, the absolute invariants are the symmetric functions. The product of differences
\[
V_\Delta=(x_1-x_2)(x_1-x_3)\cdots(x_{n-1}-x_n)
\]
is a relative invariant of index $2$, with associated subgroup the alternating group. Since $P$ has no normal subgroup except the alternating group and is not cyclic, there are no relative invariants of higher index.

The absolute invariants here are rationally expressible through finitely many forms, the elementary symmetric functions; the invariants of index $2$ are products of $V_\Delta$ with absolute invariants. The analogous theorem will be proved in general.

We finish with a useful general observation. In the first volume covariants $G(x)$ of a form $F(x)$ were characterized by the rule that if $F(x)$ is transformed into $F'(y)$ by a linear substitution, then
\[
G'(y)=r^\lambda G(x),
\]
where $r$ is the determinant of the substitution. If $F(x)$ is an invariant of $S$, and $y$ is obtained from $x$ by a substitution in $S$, then the coefficients of $F'$ differ from those of $F$ only by a common constant factor; the same holds for $G'$ and $G$. Therefore $G$ is again an invariant of $S$:
\begin{enumerate}[label=\arabic*.,leftmargin=2.2em,start=4]
\item Any covariant formed from an invariant form of $S$ is again an invariant form of $S$.
\end{enumerate}

\section*{\S\,41. Hilbert's theorem.}

The proof that the invariant system of a linear substitution group is finite rests on a very general theorem on systems of forms, discovered by Hilbert and used with great success in investigations on finiteness of invariant systems.

\begin{enumerate}[label=\Roman*.,leftmargin=2.4em]
\item Let $\mathcal S$ be any system, finite or infinite, of forms in $x_1,x_2,\ldots,x_n$. From $\mathcal S$ one can select finitely many forms $F_1,F_2,\ldots,F_u$ such that every form $F$ of $\mathcal S$ has the representation
\begin{equation}
F=A_1F_1+A_2F_2+\cdots+A_uF_u,
\tag{1}
\end{equation}
where $A_1,A_2,\ldots,A_u$ are forms in the variables $x_1,x_2,\ldots,x_n$.
\end{enumerate}

The definition of the system $\mathcal S$ must be complete enough to decide of each individual form whether it belongs to $\mathcal S$, but otherwise it is unrestricted. If $\mathcal S$ is finite the assertion is immediate. We therefore consider only infinite systems and call $F_1,F_2,\ldots,F_u$ a basis of $\mathcal S$. The products $A_iF_i$ must all have the same degree, namely that of $F$; in general the converse, that all such combinations belong to $\mathcal S$, is not required. Systems for which this converse holds are modules.

The assertion is evident for forms in one variable. For the inductive step assume it true for systems of forms in $n$ variables $x$. Consider a system $\mathcal S_r$ of forms in $x$ and one further variable $y$, where $y$ occurs in degree at most $r$. Every form $F$ can be written uniquely as
\begin{equation}
F=y^r\varphi+\psi,
\tag{2}
\end{equation}
where $\varphi$ contains no $y$ and $\psi$ contains $y$ only up to degree $r-1$.

As $F$ runs through $\mathcal S_r$, the corresponding $\varphi$ form a system $\mathcal S_0$, and by the induction hypothesis
\begin{equation}
\varphi=a_1\varphi_1+a_2\varphi_2+\cdots+a_u\varphi_u.
\tag{3}
\end{equation}
Since each $\varphi_i$ occurs from some form in $\mathcal S_r$, choose
\begin{equation}
F_i=y^r\varphi_i+\psi_i
\quad(i=1,\ldots,u),
\tag{4}
\end{equation}
with $\psi_i$ of $y$-degree at most $r-1$. Then
\begin{equation}
F=a_1F_1+a_2F_2+\cdots+a_uF_u+
\Psi,
\tag{5}
\end{equation}
where $\Psi$ has $y$-degree at most $r-1$. These remainders form a system $\mathcal S_{r-1}$. Assuming the theorem for such systems, write
\begin{equation}
\Psi=b_1\Psi_1+b_2\Psi_2+\cdots+b_v\Psi_v.
\tag{6}
\end{equation}
Choose forms in $\mathcal S_r$
\begin{align}
F'_1&=a_{1,1}F_1+a_{2,1}F_2+\cdots+a_{u,1}F_u+\Psi_1,\notag\\
F'_2&=a_{1,2}F_1+a_{2,2}F_2+\cdots+a_{u,2}F_u+\Psi_2,\notag\\
&\hspace{3em}\vdots\notag\\
F'_v&=a_{1,v}F_1+a_{2,v}F_2+\cdots+a_{u,v}F_u+\Psi_v.
\tag{7}
\end{align}
Combining (5) and (6) yields
\begin{equation}
F=b_1F'_1+b_2F'_2+\cdots+b_vF'_v+A_1F_1+A_2F_2+\cdots+A_uF_u,
\tag{8}
\end{equation}
with
\begin{equation}
A_i=a_i-a_{i,1}b_1-a_{i,2}b_2-\cdots-a_{i,v}b_v.
\tag{9}
\end{equation}
Thus $F_1,\ldots,F_u,F'_1,\ldots,F'_v$ form a finite basis of $\mathcal S_r$.

For a general system in $n+1$ variables choose a form $F_0$ of degree $r$. Make the linear change
\begin{equation}
y,
\quad x_1+\lambda_1y,
\quad x_2+\lambda_2y,
\quad\ldots,
\quad x_n+\lambda_ny,
\tag{10}
\end{equation}
choosing the constants so that the coefficient of $y^r$ in $F_0$ is non-zero. Any form $F$ of the system can then be divided by $F_0$ as a polynomial in $y$:
\begin{equation}
F=a_0F_0+\Phi,
\tag{11}
\end{equation}
where $\Phi$ has $y$-degree at most $r-1$. The remainders form a system already covered, so
\begin{equation}
\Phi=A_1\Phi_1+\cdots+A_u\Phi_u.
\tag{12}
\end{equation}
Choose forms
\begin{equation}
F_i=a_iF_0+\Phi_i
\quad(i=1,\ldots,u)
\tag{13}
\end{equation}
from the original system, and set
\begin{equation}
A_0=a_0-a_1A_1-\cdots-a_uA_u.
\tag{14}
\end{equation}
Then
\begin{equation}
F=A_0F_0+A_1F_1+\cdots+A_uF_u,
\tag{15}
\end{equation}
which proves the theorem.

\section*{\S\,42. Finiteness of the invariant system of a finite linear substitution group.}

The preceding proof has no exception in itself. Discarding identically zero forms and constant non-zero forms, we may reformulate the result:
\begin{enumerate}[label=\Roman*.,leftmargin=2.4em,start=2]
\item All forms of positive degree in a system $\mathcal S$ can be expressed through a basis $F_1,F_2,\ldots,F_u$ of positive-degree elements by
\begin{equation}
F=\Theta_1F_1+\Theta_2F_2+\cdots+\Theta_uF_u.
\tag{1}
\end{equation}
The degrees of the coefficients $\Theta_i$ are lower than the degree of $F$.
\end{enumerate}
A form system is called finite if all its forms can be expressed as integral rational functions of finitely many among them. The following criterion holds:
\begin{enumerate}[label=\Roman*.,leftmargin=2.4em,start=3]
\item If the coefficients $\Theta_i$ in (1) can always be chosen so that, unless constant, they themselves belong to $\mathcal S$, then $\mathcal S$ is finite.
\end{enumerate}
Indeed, repeated use of (1) lowers the degree and must end with constants.

It follows, by a simple argument communicated orally by Hurwitz, that the invariant system $\mathfrak I$ of a finite group $S$ of linear substitutions is finite. Let $F_1,F_2,\ldots,F_u$ be a basis of $\mathfrak I$ as in II, and let $F$ be any other non-constant invariant. Then
\begin{equation}
F=A_1F_1+A_2F_2+\cdots+A_uF_u.
\tag{2}
\end{equation}
Apply all substitutions of $S$ to this identity. The forms $F,F_1,\ldots,F_u$ remain fixed, while $A_i$ goes into conjugates $A_i,A'_i,A''_i,\ldots$. If $m$ is the order of the group, put
\begin{equation}
m\Theta_i=A_i+A'_i+A''_i+\cdots.
\tag{3}
\end{equation}
Averaging the identities gives
\begin{equation}
F=\Theta_1F_1+\Theta_2F_2+\cdots+\Theta_uF_u.
\tag{4}
\end{equation}
The coefficients $\Theta_i$ are invariants of $S$, so criterion III applies.

The same averaging applies to relative invariants. If $F(x)$ runs through all functions satisfying
\[
F[A(x)]=\alpha F(x),
\quad F[B(x)]=\beta F(x),
\quad\ldots,
\]
for fixed root-of-unity factors $\alpha,\beta,\ldots$, then Hilbert's theorem supplies a finite system $F_1,\ldots,F_m$ with
\begin{equation}
F=A_1F_1+A_2F_2+\cdots+A_mF_m.
\tag{5}
\end{equation}
Averaging as above gives
\[
F=\Phi_1F_1+\Phi_2F_2+\cdots+\Phi_mF_m,
\]
where the $\Phi_i$ are absolute invariants.

An inhomogeneous function can be an invariant only when its homogeneous parts are invariants. For rational invariants, suppose
\[
{F(x)\over F_1(x)}
\]
is an absolute rational invariant and that numerator and denominator are coprime. Then
\[
{F(x)\over F_1(x)}={F[A(x)]\over F_1[A(x)]}
\]
forces
\[
F[A(x)]=\alpha F(x),
\qquad
F_1[A(x)]=\alpha F_1(x),
\]
so numerator and denominator are relative invariants. Conversely, if the factors are $e$th roots of unity, then
\[
{F(x)\over F_1(x)}={F(x)F_1(x)^{e-1}\over F_1(x)^e}
\]
represents the same absolute rational invariant as a quotient of absolute integral invariants. Consequently all relative invariants with a fixed factor system can be written as the product of one of them with absolute invariants, possibly rational ones.

\section*{\S\,43. The form problem.}

The preceding paragraph showed that a finite group of linear substitutions has a finite number of independent invariants. If $n$ is the dimension, these forms are homogeneous functions of $n$ variables, so no more than $n$ of them can be independent. This does not say that all invariants are rational functions of $n$ among them; but among $n+1$ invariants there is always a rational equation obtained by eliminating the $n$ variables.

We now ask the converse question: whether a linear substitution group of dimension $n$ always has $n$ independent invariants. If the variables receive fixed values, all invariants receive fixed values. The question is whether for a given system of invariant values there are finitely many corresponding systems of variable values. If so, the variables are algebraic, multivalued functions of the invariants, and the number of independent invariants cannot be smaller than the number of variables.

If $(x)$ is one system of variable values belonging to a given system of invariant values, and $A$ is a substitution of $S$, then $A(x)$ belongs to the same invariant values. Thus the problem has at least as many solutions as the order of $S$. That this is exactly the number follows as below. Coincidence of several such systems can occur only when $(x)=A(x)$ for some substitution $A$.

Choose a homogeneous linear function
\begin{equation}
\Theta=\Theta(x)=c_1x_1+c_2x_2+\cdots+c_nx_n.
\tag{1}
\end{equation}
For a substitution $A$ of $S$,
\begin{equation}
\Theta[A(x)]=c'_1x_1+c'_2x_2+\cdots+c'_nx_n,
\tag{2}
\end{equation}
where the coefficients $c'_i$ are related to the original $c_i$ by the transpose of $A$:
\[
(c')=A_1(c).
\]
If $A,B,C,\ldots$ run through all substitutions of $S$, we obtain
\begin{equation}
\Theta[A(x)],\quad\Theta[B(x)],\quad\Theta[C(x)],\ldots,
\tag{3}
\end{equation}
which we denote by
\begin{equation}
\Theta,\quad\Theta_1,\quad\Theta_2,\quad\ldots,\quad\Theta_{\mu-1},
\tag{4}
\end{equation}
where $\mu$ is the order of $S$. The coefficients $c_i$ may be chosen so that no two of these functions are identical. Then:
\begin{enumerate}[label=\arabic*.,leftmargin=2.2em]
\item Applying any substitution of $S$ simultaneously to the $\mu$ functions (4) only permutes them.
\item Every symmetric function of the quantities (4) is an absolute invariant of $S$.
\end{enumerate}
The function $\Theta$ could be replaced by another function, even non-linear, rational, or inhomogeneous, provided the $\mu$ functions obtained are distinct.

Form
\begin{align}
\Phi(t)&=(t-\Theta)(t-\Theta_1)\cdots(t-\Theta_{\mu-1})\notag\\
&=t^\mu+A_1t^{\mu-1}+A_2t^{\mu-2}+\cdots+A_\mu.
\tag{5}
\end{align}
The coefficients $A_1,A_2,\ldots,A_\mu$ are invariants of $S$, and $\Theta,\Theta_1,\ldots,\Theta_{\mu-1}$ are the roots of
\begin{equation}
\Phi(t)=0.
\tag{6}
\end{equation}

Let the rationality domain $\Omega$ be the field generated by the absolute invariants of $S$ and by the numerical constants. The coefficients of $\Phi(t)$ lie in $\Omega$.
\begin{enumerate}[label=\arabic*.,leftmargin=2.2em,start=3]
\item The function $\Phi(t)$ is irreducible over $\Omega$.
\end{enumerate}
Indeed, if $\Psi(t)\in\Omega[t]$ vanishes at $t=\Theta$, then replacing $(x)$ by $A(x)$ for each $A\in S$ shows that $\Psi$ vanishes at every conjugate $\Theta_i$. Hence $\Psi(t)$ is divisible by $\Phi(t)$.

Let $\omega$ be any rational function of the variables, and let $\omega,\omega_1,\ldots,\omega_{\mu-1}$ be the functions obtained from it by applying $S$; these need not be distinct. Then
\begin{equation}
\Phi(t)\left({\omega\over t-\Theta}+{\omega_1\over t-\Theta_1}+\cdots+{\omega_{\mu-1}\over t-\Theta_{\mu-1}}\right)=\Psi(t)
\tag{7}
\end{equation}
is a polynomial of degree $\mu-1$ in $t$ and is invariant under $S$, hence lies in $\Omega[t]$. Setting $t=\Theta$ gives
\begin{equation}
\omega={\Psi(\Theta)\over\Phi'(\Theta)}.
\tag{8}
\end{equation}
Therefore:
\begin{enumerate}[label=\arabic*.,leftmargin=2.2em,start=4]
\item Every rational function of the variables $(x)$ can be expressed rationally in terms of $\Theta$; it belongs to $\Omega(\Theta)$.
\end{enumerate}

This gives a second proof that all invariants of $S$ are rational functions of finitely many of them, namely the coefficients of $\Phi(t)$. Since the variables themselves are among the possible functions $\omega$, the original question is answered affirmatively: the variables are algebraic functions of the invariants.

The problem of representing the variables $x$ as algebraic functions of the invariants of $S$, equivalently of determining the field $\Omega(\Theta)$, is called, following F. Klein, the form problem of the group $S$. If there are $n$ variables, it is the form problem of dimension $n$.

The field $\Omega(\Theta)$ is an algebraic field over $\Omega$ determined by $S$. It is normal, since all conjugates $\Theta,\Theta_1,\ldots,\Theta_{\mu-1}$ lie in it. Thus $\Phi(t)=0$ is a normal equation and the Galois resolvent of the form problem.

\begin{enumerate}[label=\arabic*.,leftmargin=2.2em,start=5]
\item The Galois group of the form problem, i.e. of the equation $\Phi(t)=0$, is isomorphic to the group $S$.
\end{enumerate}
If $A,B\in S$ and $AB=C$, and if
\[
\Theta_1=\Theta[A(x)],
\quad \Theta_2=\Theta[B(x)],
\quad \Theta_3=\Theta[C(x)],
\]
then the substitution $(\Theta,\Theta_2)$ sends $\Theta_1$ to $\Theta_3$, so
\[
(\Theta,\Theta_2)=(\Theta_1,\Theta_3),
\]
and hence
\[
(\Theta,\Theta_1)(\Theta,\Theta_2)=(\Theta,\Theta_1)(\Theta_1,\Theta_3)=(\Theta,\Theta_3).
\]
Thus the substitutions $(\Theta,\Theta_1),(\Theta,\Theta_2),\ldots$ form a group isomorphic to $S$.

If instead of $\Theta$ one chooses a function $\eta$ which assumes the same value under the substitutions of a subgroup $S'$ of $S$ of index $j$, and under no others, then $\eta$ satisfies an equation of degree $j$, a resolvent of the form problem. Every function invariant under $S'$ is rational in $\eta$ and the invariants of $S$. The resolvent of $\eta$ is a partial or total resolvent according as $S'$ and its conjugate subgroups have a common subgroup or are relatively prime.

\section*{\S\,44. Klein's extension of the fundamental algebraic problem.}

The preceding paragraph is a direct generalization of Galois theory for a general equation of degree $n$; that theory is contained as a special case in the theory of linear substitution groups.

For by \S\,39 the symmetric functions of $n$ independent variables $x_1,x_2,\ldots,x_n$ are the invariants of the group $S$ consisting of all permutations of these variables, and the equation $\Phi(t)=0$ is precisely the Galois resolvent of the general equation of degree $n$ with roots $x_i$. Thus the general algebraic problem of solving an equation of degree $n$ may be regarded as the form problem for a linear substitution group of dimension $n$. But there are special equations which can be solved by form problems of lower dimension, especially pure equations, which are solvable by a form problem of the first dimension.

Homogeneous linear substitutions of one dimension have only the form
\begin{equation}
x'=\alpha x,
\tag{1}
\end{equation}
and if they are to form a finite group of order $\mu$, the coefficients $\alpha$ must be $\mu$th roots of unity. Conversely, if $\alpha$ runs through all $\mu$th roots of unity in (1), we obtain a group of order $\mu$. This group has the absolute invariant $x^\mu$, and when this is given, $x$ is obtained as its $\mu$th root.

\clearpage


The solution of the general equations of the second, third, and fourth degrees is therefore reducible to form-problems of only one dimension. We shall see in a later section that the general equation of the fifth degree is reducible to a binary form-problem, and one may now put, as an immediate extension of the problem of reducing the solution of an equation to pure equations, the question: what is the least number of dimensions of a form-problem through which a given equation can be solved? The problem would then have to be formulated as follows.

From the roots of a given equation one is to form as few rational functions as possible, in such a way that they pass into homogeneous linear functions of themselves when the roots are subjected to the permutations of the Galois group of the given equation.

In this way F. Klein enlarged the problem of the algebraic solution of an equation. In order to open the way to answering this question, he proved for the general equations of the sixth and seventh degrees that they can be reduced to quaternary form-problems.

As we have seen, the general equation of degree $n$ is immediately equivalent to a form-problem of $n$ dimensions. The question whether the general equation of degree $n$, when $n$ is greater than $7$, corresponds to a form-problem of fewer than $n$ dimensions has not yet been answered; it must probably be answered in the negative.

\section*{\S\,45. Influence of relative invariants.}

In the definition of the form-problem in \S\,43 we used only the absolute invariants of the group $S$. But, as we have seen, there are also cases in which relative invariants exist besides the absolute invariants, and these may be used to simplify the form-problem.

If $r$ is such a relative invariant, then a certain power of $r$, whose degree $e$ is a divisor of the degree $n$ of $S$, is an absolute invariant; hence $r$ is determined in $\Omega$ by a pure equation.

All substitutions of $S$ by which $r$ remains unchanged form a group $T_r$; and, as we saw in \S\,40, the group $T_r$ is a normal divisor of $S$.

Moreover, every function of the variables $x$ that remains unchanged under the substitutions of the group $T_r$ can be expressed rationally in terms of $r$, that is, it is contained in the field $\Omega_r$ obtained from $\Omega$ by adjoining $r$.

For, by \S\,40, there is a substitution $E$ in $S$ and a primitive $e$-th root of unity $\eps$ such that $r$ is changed by $E$ into $\eps r$, and all values which $r$ can assume,
\[
 r,\ \eps r,\ \eps^2r,\ldots,\eps^{e-1}r,
\]
are obtained by repeating the substitution $E$. If $R$ denotes a function of the $x$ which admits the substitutions of $T_r$, and if $E$ and its powers change it into
\[
 R,R_1,\ldots,R_{e-1},
\]
then all these functions likewise admit the substitutions of $T_r$, since $T_r$ is a normal divisor of $S$. The function
\begin{equation}
(\eps^{-\mu}R)=R+\eps^{-\mu}R_1+\cdots+\eps^{-\mu(e-1)}R_{e-1}\qquad(\mu=0,1,2,\ldots,e-1),
\tag{1}
\end{equation}
formed in analogy with Lagrange resolvents, acquires the factor $\eps^\mu$ when the substitution $E$ is applied. Thus, if
\begin{equation}
(\eps^{-\mu}R)=r^\mu W_\mu,
\tag{2}
\end{equation}
then $W_\mu$ is an absolute invariant and hence belongs to the field $\Omega$.

From (1) and (2) one obtains
\begin{equation}
eR=W_0+rW_1+r^2W_2+\cdots+r^{e-1}W_{e-1},
\tag{3}
\end{equation}
which proves the assertion.

Put
\begin{equation}
n=e\nu.
\tag{4}
\end{equation}
Then $\nu$ is the degree of the group $T_r$. If the function $\Theta$ used in \S\,43 for solving the form-problem is carried by the substitutions of $T_r$ into
\[
\Theta,\Theta_1,\ldots,\Theta_{\nu-1},
\]
then
\[
\Phi(t)=(t-\Theta)(t-\Theta_1)\cdots(t-\Theta_{\nu-1})
\]
is a function of $t$ in the enlarged field of rationality $\Omega_r$. The field $\Omega(\Theta)$ is identical with $\Omega_r(\Theta)$. It is therefore an algebraic field of degree $n$ over $\Omega$ and of degree $\nu$ over $\Omega_r$. Thus, by adjoining one radical, the form-problem of degree $n$ is reduced to a form-problem of degree $\nu$.

\section*{\S\,46. The enlarged notion of invariant.}

Relative invariants are, at bottom, a special case of a more general notion, which can be used for the reduction of the form-problem in a manner similar to relative invariants.

We look for systems of homogeneous forms of the same degree in the variables $x_1,x_2,\ldots,x_n$,
\begin{equation}
X_1,X_2,\ldots,X_m,
\tag{1}
\end{equation}
with the property that, under the substitutions of the group $S$, the functions $X_1,X_2,\ldots,X_m$ undergo a system $\Sigma$ of linear substitutions. The substitutions $\Sigma$ then likewise form a group, namely a group one- or many-fold isomorphic with $S$. If $m=1$, one obtains the notion of relative invariant.

We now look for all substitutions of $S$ that leave each of the functions (1) unchanged. These substitutions form a group $T$. We shall show that it is a normal divisor of $S$. Let $s$ be a substitution of $S$ under which the $X$ undergo the substitution $\sigma$; then under $s^{-1}$ they undergo $\sigma^{-1}$. If, further, $\tau$ is a substitution of $T$, then $\tau$ leaves all the $X$ unchanged. Hence $s^{-1}\tau s$ effects on the $X$ the substitution $\sigma^{-1}\sigma=1$, so that it also leaves them unchanged. Thus $s^{-1}\tau s$ belongs to $T$, and $T$ is a normal divisor of $S$.

Let $n$ and $\nu$ be the degrees of $S$ and $T$, and write
\[
n=e\nu.
\]
Then $e$ is the index of $T$ and at the same time the degree of the group $\Sigma$.

Every absolute invariant of the group $T$, that is, every function of $x$ admitting the substitutions of $T$, can be expressed rationally in $\Omega$ by means of the $X$.

This follows by the argument already used several times. Let $\rho$ be a function of the $X$ which assumes $e$ distinct values $\rho,\rho_1,
\ldots,\rho_{e-1}$, and put
\begin{equation}
\Phi(t)=(t-\rho)(t-\rho_1)\cdots(t-\rho_{e-1}).
\tag{2}
\end{equation}
Then $\Phi(t)$ is a rational function of $t$ in $\Omega$.

An absolute invariant $R$ of $T$ assumes at most $e$ distinct values corresponding to the values $\rho,\rho_1,\ldots,\rho_{e-1}$, say $R,R_1,\ldots,R_{e-1}$, with repetitions possibly occurring. The function
\[
\frac{R}{t-\rho}+\frac{R_1}{t-\rho_1}+\cdots+\frac{R_{e-1}}{t-\rho_{e-1}}
\]
then belongs to $\Omega$; and at $t=\rho$ it gives $R$ rationally in terms of $\rho$, since $\Phi'(\rho)\neq0$. Therefore $R$ is contained in the field $\Omega(X_1,X_2,\ldots,X_m)$.

The invariants of $T$ are at the same time invariants of $S$. By solving the form-problem for the group $\Sigma$, the functions $X_1,X_2,
\ldots,X_m$ are known, and hence so are the invariants of $T$. Thus the form-problem for $S$ is reduced to the successive solution of the two form-problems of the groups $\Sigma$ and $T$, each of lower degree. A form-problem that can be decomposed in this way into two form-problems of lower degree may be called imprimitive. Since such a reduction requires that $T$ be a normal divisor of $S$, different both from $S$ and from the identity, the form-problem corresponding to a simple group is always primitive.

Conversely, if $S$ possesses a normal divisor of index $e$, the form-problem is imprimitive. To obtain a system $X_1,\ldots,X_m$, one need only take an invariant $\rho$ of $T$ which assumes $e$ distinct values under $S$, and choose as the $X$ precisely the values $\rho,\rho_1,\ldots,\rho_{e-1}$. The group $\Sigma$ is then the permutation group induced by $S$ on these values.

In the sense of \S\,44, however, the essential point is always to make $m$ as small as possible; a reduction in this sense has been achieved only when $m<n$.

\section*{\S\,47. Normal forms.}

One further general consideration is needed before we pass to special applications. We have already seen in \S\,37 that from every group $S$ of linear substitutions one can derive infinitely many isomorphic groups by transforming with an arbitrary substitution $L$ of the same dimension, that is, by forming
\begin{equation}
L^{-1}SL.
\tag{1}
\end{equation}
This is equivalent to introducing other variables $y$ in place of $x$ by the substitution
\begin{equation}
(x)=L(y).
\tag{2}
\end{equation}

We may use this transformation to represent $S$ in a simple normal form. Then the invariants also assume fixed normal forms, which in some cases have very simple shapes. For the normal form let us write the resolvent of the form-problem as
\begin{equation}
\Phi(\Theta,A_1,A_2,\ldots)=0,
\tag{3}
\end{equation}
where $A_1,A_2,\ldots$ are invariants of $S$. In normal form this equation too will take a simple shape. The variables $x_i$ can be expressed rationally in terms of $\Theta$ and the $A_i$, and we put
\begin{equation}
x_i=\varphi_i(\Theta,A_1,A_2,\ldots).
\tag{4}
\end{equation}
These expressions too acquire fixed shapes for the normal form once $\Theta$ has been chosen.

It may also happen that the invariants are given not in normal form but in some other form, which we shall call the general form. Then the task is to find the substitution $L$ which transforms the general form into a normal form known in advance to be possible. This problem is no other than the form-problem itself for the normal form.

Assume that the invariants in the general form, as functions of $y$, are $B_1,B_2,\ldots$, so that the substitution (2) establishes the identities
\begin{equation}
A_1=B_1,
\qquad A_2=B_2,
\qquad \ldots
\tag{5}
\end{equation}
The problem is to determine the substitution $L$, given the forms of the functions $A_i$ and $B_i$, so that the equations (5) become identities. This problem is solved if the equation (3) is supposed solved for a suitably chosen special system of values of the $A_i$.

Taking total differentials of (3) and (4) and eliminating $d\Theta$ gives
\begin{equation}
dx_i=-\sum_s\frac{\Phi'(A_s)\varphi_i'(\Theta)-\Phi'(\Theta)\varphi_i'(A_s)}{\Phi'(\Theta)}\,dA_s.
\tag{6}
\end{equation}
In consequence of (5),
\begin{equation}
dA_s=\sum_hB_s'(y_h)\,dy_h.
\tag{7}
\end{equation}
If
\begin{align}
x_i&=\alpha_{1i}y_1+\cdots+\alpha_{ni}y_n,\tag{8}\\
dx_i&=\alpha_{1i}dy_1+\cdots+\alpha_{ni}dy_n,\tag{9}
\end{align}
comparison gives
\begin{equation}
\alpha_{hi}=-\sum_s\frac{\Phi'(A_s)\varphi_i'(\Theta)-\Phi'(\Theta)\varphi_i'(A_s)}{\Phi'(\Theta)}\,B_s'(y_h).
\tag{10}
\end{equation}
The right-hand side must reduce to a constant. We can find its value by assigning to the $y$ any special set of values subject only to the condition $\Phi'(\Theta)\neq0$. For this special system the $A_i$ are determined by (5), and $\Theta$ is known if the form-problem (3) has been solved for these values. Thus (10) determines the coefficients of $L$.

\section*{Seventh Section. Groups of binary linear substitutions.}
\section*{\S\,48. Ternary orthogonal substitutions.}

We now have to single out narrower groups from the totality of linear substitutions and ultimately arrive at finite groups. One obtains such narrower groups, which may still be infinite, by requiring that given homogeneous functions of the variables remain invariant. We shall not pursue this task in full generality, but pass at once to the most important special case. We restrict ourselves to ternary substitutions and require that a quadratic form of non-vanishing determinant be invariant. Since every such quadratic form can be transformed linearly into a sum of squares, we do not further restrict the problem by taking this quadratic form to be the sum of squares. Such substitutions are called orthogonal.

The substitution
\begin{equation}
(y_1,y_2,y_3)=A(x_1,x_2,x_3)
\tag{1}
\end{equation}
is orthogonal if the coefficients are so determined that
\begin{equation}
y_1^2+y_2^2+y_3^2=x_1^2+x_2^2+x_3^2.
\tag{2}
\end{equation}
Substitution into (2) gives six relations among the nine coefficients of $A$. If
\begin{equation}
A=\begin{pmatrix}a_1&b_1&c_1\\a_2&b_2&c_2\\a_3&b_3&c_3\end{pmatrix},
\tag{3}
\end{equation}
these relations are
\begin{align}
a_1^2+b_1^2+c_1^2&=1,& a_2a_3+b_2b_3+c_2c_3&=0,\notag\\
a_2^2+b_2^2+c_2^2&=1,& a_3a_1+b_3b_1+c_3c_1&=0,\tag{4}\\
a_3^2+b_3^2+c_3^2&=1,& a_1a_2+b_1b_2+c_1c_2&=0.\notag
\end{align}
From these relations the determinant satisfies $|A|^2=1$, hence $|A|=\pm1$. The inverse of $A$ is its transpose; this property could also be used as the definition of orthogonal substitutions.

All orthogonal substitutions form a group. Inside it lies the narrower group distinguished by
\begin{equation}
|A|=+1,
\tag{5}
\end{equation}
which we shall call the group of proper orthogonal substitutions.

The linear substitution (1) may be interpreted as the passage from one rectangular coordinate system to another with the same origin. If (5) holds, the first coordinate system can be brought into coincidence with the second by a rotation about a fixed axis through a definite angle. The fixed axis is obtained from the non-trivial solution of the linear equations expressing that a direction $\lambda,\mu,\nu$ is unchanged by $A$. When $|A|=-1$ this conclusion no longer follows.

There is also a second interpretation: $(x)$ and $(y)$ may be the coordinates of two different points in one and the same coordinate system. If $x$ ranges over a line, surface, or solid, then the corresponding point $y$ ranges over a congruent region, or over a mirror-image congruent region when $|A|=-1$.

The group of proper orthogonal substitutions is equivalent to the group formed by the possible positions of a rigid body rotating about a fixed point. A position $E$ is taken as the identity; any other position $A$ is reached by a rotation about an axis through an angle. If $B$ is a third position, the product $AB$ means the position obtained by carrying out the rotation that led to $A$, not from the identity position, but from the position $B$.

The improper substitutions are obtained from the proper ones by composition with one improper substitution, for instance the reflection $(x_1,x_2,x_3)=(y_1,y_2,-y_3)$.

Of special interest are the finite groups contained in the group of proper orthogonal substitutions. To such a group of degree $g$ corresponds a finite group of positions of a body. If one fixes the body simultaneously in these different positions and regards the whole as a new rigid body, one obtains a figure which can be brought to coincide with itself in exactly $g$ different positions. Examples are regular pyramids, double pyramids, and the regular solids. Counting rotations gives $12$ positions for the tetrahedron, $24$ for the octahedron and cube, and $60$ for the dodecahedron and icosahedron. Instead of the positions of the body, the rotations themselves may be taken as the elements of the group.

\section*{\S\,49. Linear fractional substitutions.}

The group of ternary orthogonal substitutions is isomorphic with the group of linear fractional substitutions in one variable, or with the binary substitutions of ratios.

Let
\begin{align}
(y_1,y_2,y_3)&=A(x_1,x_2,x_3),\tag{1}\\
y_1^2+y_2^2+y_3^2&=x_1^2+x_2^2+x_3^2.\tag{2}
\end{align}
Apply the fixed non-orthogonal substitution
\begin{equation}
L=\begin{pmatrix}
\ii&0&0\\
0&\frac1{\sqrt2}&\frac1{\sqrt2}\\
0&\frac{\ii}{\sqrt2}&-\frac{\ii}{\sqrt2}
\end{pmatrix},\qquad
L^{-1}=\begin{pmatrix}
-\ii&0&0\\
0&\frac1{\sqrt2}&-\frac{\ii}{\sqrt2}\\
0&\frac1{\sqrt2}&\frac{\ii}{\sqrt2}
\end{pmatrix},
\tag{3}
\end{equation}
where $\ii=\sqrt{-1}$ and the determinant is $1$. Put
\begin{align}
(y_1,y_2,y_3)&=L(y'_1,y'_2,y'_3),&(x_1,x_2,x_3)&=L(x'_1,x'_2,x'_3),\tag{4}\\
(y'_1,y'_2,y'_3)&=A'(x'_1,x'_2,x'_3),\qquad A'=L^{-1}AL.\tag{5--6}
\end{align}
Then (2) becomes
\begin{equation}
-y_1'^2+2y'_2y'_3=-x_1'^2+2x'_2x'_3.
\tag{7}
\end{equation}
Set
\begin{equation}
x'_1=\sqrt2\,\xi_1\xi_2,
\qquad x'_2=\xi_1^2,
\qquad x'_3=\xi_2^2.
\tag{8}
\end{equation}
Then $x_1'^2-2x'_2x'_3$ vanishes identically, and the transformed $y'$ must satisfy
\begin{equation}
y_1'^2-2y'_2y'_3=0.
\tag{9}
\end{equation}
Thus $y'_2$ and $y'_3$ are squares of linear functions. Writing
\begin{equation}
\eta_1=\alpha\xi_1+\beta\xi_2,
\qquad
\eta_2=\gamma\xi_1+\delta\xi_2,
\tag{10}
\end{equation}
one may put
\begin{equation}
y'_1=\sqrt2\,\eta_1\eta_2,
\qquad y'_2=\eta_1^2,
\qquad y'_3=\eta_2^2.
\tag{11}
\end{equation}
This gives
\begin{equation}
A'=\begin{pmatrix}
\alpha\delta+\beta\gamma&\sqrt2\,\alpha\gamma&\sqrt2\,\beta\delta\\
\sqrt2\,\alpha\beta&\alpha^2&\beta^2\\
\sqrt2\,\gamma\delta&\gamma^2&\delta^2
\end{pmatrix}.
\tag{12}
\end{equation}
The condition is fully satisfied when
\begin{equation}
\alpha\delta-\beta\gamma=1
\tag{13}
\end{equation}
for proper orthogonal substitutions, and
\begin{equation}
\alpha\delta-\beta\gamma=-1
\tag{14}
\end{equation}
for improper ones. Hence $A'$ is reduced to the binary substitution
\begin{equation}
(\eta_1,\eta_2)=\begin{pmatrix}\alpha&\beta\\\gamma&\delta\end{pmatrix}(\xi_1,\xi_2),
\qquad \alpha\delta-\beta\gamma=\pm1,
\tag{15}
\end{equation}
and, on passing to ratios, to
\begin{equation}
\eta=\frac{\alpha\xi+\beta}{\gamma\xi+\delta}.
\tag{16}
\end{equation}

The two matrices differing by an overall sign give the same $A'$. Thus they are combined into one concept, denoted $A''$. For the fractional substitution $A'''$, an arbitrary common factor in $\alpha,\beta,\gamma,\delta$ may be chosen so that either (13) or (14) is satisfied. Therefore each $A'''$ gives two orthogonal substitutions $A$, one proper and one improper. The triples $A,A',A''$ correspond bijectively, while $A,A',A'''$ correspond bijectively only after requiring $A$ to be proper.

If $B,B',B''$ is a second triple, then $AB,A'B',A''B''$ is again a corresponding triple. This is verified by writing
\begin{equation}
(\eta_1,\eta_2)=A''(\xi_1,\xi_2),
\qquad
(\xi_1,\xi_2)=B''(\zeta_1,\zeta_2),
\tag{23}
\end{equation}
so that
\begin{equation}
(\eta_1,\eta_2)=A''B''(\zeta_1,\zeta_2),
\tag{24}
\end{equation}
and comparing the induced ternary substitutions on the triples $\sqrt2\xi_1\xi_2,\xi_1^2,\xi_2^2$. Thus the groups of the $A$, $A'$, and $A''$ are isomorphic; for proper $A$, the group of the $A'''$ is likewise isomorphic.

\section*{\S\,50. Reality conditions.}

It remains to answer one question. So far we have nowhere taken account of the reality of the coefficients. For geometric applications the orthogonal substitution $A$ must be real. We must therefore determine what conditions the coefficients $\alpha,\beta,\gamma,\delta$ in $A''$ must satisfy in order that the coefficients of $A$ be real.

From $A=LA'L^{-1}$ one obtains
\begin{equation}
A=\begin{pmatrix}
\alpha\delta+\beta\gamma & \ii(\alpha\gamma+\beta\delta) & \alpha\gamma-\beta\delta\\
-\ii(\alpha\beta+\gamma\delta)&\frac{\alpha^2+\beta^2+\gamma^2+\delta^2}{2}&\ii\frac{-\alpha^2+\beta^2-\gamma^2+\delta^2}{2}\\
\alpha\beta-\gamma\delta&\ii\frac{\alpha^2+\beta^2-\gamma^2-\delta^2}{2}&\frac{\alpha^2-\beta^2-\gamma^2+\delta^2}{2}
\end{pmatrix},
\tag{1}
\end{equation}
with
\begin{equation}
\alpha\delta-\beta\gamma=\pm1.
\tag{2}
\end{equation}
The coefficients in (1) are real precisely when $\alpha,\pm\delta$ and $\beta,\mp\gamma$ are pairs of conjugate imaginary numbers; the upper signs correspond to proper, and the lower signs to improper, orthogonal substitutions. This includes the special cases in which two of $\alpha,\beta,\gamma,\delta$ vanish.

\section*{\S\,51. Finite groups of linear fractional substitutions. Poles of the groups.}

From what has been developed, once all finite groups of linear fractional substitutions have been found, all finite groups of proper ternary orthogonal substitutions and all finite groups of binary linear substitutions with determinant $1$ follow immediately. We therefore first determine all finite groups of linear fractional substitutions.

Let $G$ be such a group of degree $n$, with elements
\begin{equation}
x,\ \Theta_1(x),\ \Theta_2(x),\ldots,\Theta_{n-1}(x),
\tag{1}
\end{equation}
where
\begin{equation}
\Theta(x)=\frac{\alpha x+\beta}{\gamma x+\delta}
\tag{2}
\end{equation}
means a linear fractional function. Groups transformed into one another by a linear fractional transformation are regarded as of the same kind.

Since the group is finite, every non-identity element $\Theta$ has a finite order $t$, and $t$ divides $n$. For every non-identity substitution we consider the values fixed by it, namely the roots of
\begin{equation}
x=\Theta(x).
\tag{5}
\end{equation}
These values are called the poles of the substitution.

Equation (5) is quadratic:
\begin{equation}
\gamma x^2+(\delta-\alpha)x-\beta=0.
\tag{6}
\end{equation}
It has a double root only if
\begin{equation}
(\delta-\alpha)^2+4\beta\gamma=0,
\tag{7}
\end{equation}
or equivalently
\begin{equation}
(\alpha+\delta)^2=4(\alpha\delta-\beta\gamma).
\tag{8}
\end{equation}
Such a case cannot occur for a substitution of finite order, except for the identity. Thus every one of the $n-1$ non-identity substitutions has two distinct poles.

Counting a value $h$ times when it is a pole of $h$ substitutions, the group has $2n-2$ poles. Multiplicative substitutions have the poles $0$ and $\infty$. Under transformation by $L$, the poles are transformed by $L$ as well. Consequently one can choose a transform of a given group so that any chosen non-identity substitution becomes a multiplication; $L$ is then determined up to a multiplication.

\section*{\S\,52. The different possible kinds of groups.}

Let $a$ be a pole of the group $G$, and suppose that exactly $\nu-1$ non-identity elements of $G$ fix $a$:
\begin{equation}
a=\Theta_1(a)=\Theta_2(a)=\cdots=\Theta_{\nu-1}(a).
\tag{1}
\end{equation}
Then $a$ is called a $\nu$-fold pole. The elements
\begin{equation}
1,\Theta_1,\Theta_2,\ldots,\Theta_{\nu-1}
\tag{2}
\end{equation}
form a subgroup $Q$ of $G$, and
\begin{equation}
n=\nu\mu.
\tag{3}
\end{equation}
This subgroup is cyclic. After transforming $a$ to $\infty$, its substitutions have the form $x\mapsto \eps x+c$, and the coefficients $\eps$ must be the distinct $\nu$-th roots of unity.

The images
\begin{equation}
a,a_1,a_2,\ldots,a_{\mu-1}
\tag{12}
\end{equation}
of $a$ under representatives of the cosets of $Q$ form a system of conjugate $\nu$-fold poles. Every substitution of $G$ merely permutes this system. Distinct systems of conjugate poles have no pole in common.

Hence all poles of $G$ may be arranged in systems of conjugate poles, and their total number is $2n-2$ when each $\nu$-fold pole is counted $\nu-1$ times. Thus
\begin{equation}
2n-2=\mu(\nu-1)+\mu'(\nu'-1)+\mu''(\nu''-1)+\cdots,
\tag{13}
\end{equation}
or
\begin{equation}
2n-2=nh-\mu-\mu'-\mu''-\cdots,
\tag{14}
\end{equation}
where $h$ is the number of systems. It follows that $h$ is either $2$ or $3$.

For $h=2$ one obtains the cyclic group:
\[
\nu=\nu'=n,
\qquad \mu=\mu'=1.
\]
For $h=3$ one obtains, besides the dihedral group,
\[
\nu=\nu'=2,\quad \nu''=m,
\qquad n=2m,
\qquad \mu=\mu'=m,
\quad \mu''=2,
\]
only three further possibilities:
\[
\begin{array}{llll}
\text{tetrahedral group:}&\nu=2,&\nu'=3,&\nu''=3,\quad n=12,\\
\text{octahedral group:}&\nu=2,&\nu'=3,&\nu''=4,\quad n=24,\\
\text{icosahedral group:}&\nu=2,&\nu'=3,&\nu''=5,\quad n=60.
\end{array}
\]
These are all possibilities, though their actual existence and multiplicity remain to be proved. They are collectively called the polyhedral groups.

\section*{\S\,53. Transforming the substitutions of $G$ into simple forms.}

Let $a$ be a $\nu$-fold pole of $G$ and
\begin{equation}
a=\Theta(a)=\Theta^2(a)=\cdots=\Theta^{\nu-1}(a).
\tag{1}
\end{equation}
The second pole $a'$ of $\Theta$ is likewise $\nu$-fold:
\begin{equation}
a'=\Theta(a')=\Theta^2(a')=\cdots=\Theta^{\nu-1}(a').
\tag{2}
\end{equation}
If there is only one system of conjugate $\nu$-fold poles, then $a$ and $a'$ belong to the same system. There is therefore a substitution $\psi$ not among the powers of $\Theta$ with
\begin{equation}
a'=\psi(a).
\tag{3}
\end{equation}
It follows that
\begin{equation}
\psi(a')=a.
\tag{4}
\end{equation}

After transforming $a$ and $a'$ to $0$ and $\infty$, the substitution $\Theta$ has the form $x\mapsto\eps x$, where $\eps$ is a primitive $\nu$-th root of unity, and $\psi$ has the form $x\mapsto c/x$. A further transformation may give $c$ any prescribed value, for instance $1$. Thus, whenever $G$ has only one system of conjugate $\nu$-fold poles, it may be transformed so as to contain the substitutions
\[
x\mapsto\eps x,
\qquad x\mapsto c/x.
\]
This applies in the cases listed in \S\,52, except for the cyclic group and the dihedral group with $m=2$.

Finally, if $a,a'$ are the poles of a substitution $\Theta$ and $\chi$ is any substitution of $G$, then $\chi(a),\chi(a')$ are the poles of the conjugate substitution $\chi\Theta\chi^{-1}$.

\section*{\S\,54. The fundamental forms.}

To give the final construction of all finite groups $G$, it is necessary to examine more closely the invariants of the corresponding binary substitution groups.

Besides the fractional substitutions
\begin{equation}
y=\frac{\alpha x+\beta}{\gamma x+\delta},
\tag{1}
\end{equation}
we consider the binary substitutions
\begin{equation}
(y_1,y_2)=\begin{pmatrix}\alpha&\beta\\\gamma&\delta\end{pmatrix}(x_1,x_2),
\tag{2}
\end{equation}
where
\begin{equation}
\alpha\delta-\beta\gamma=1.
\tag{3}
\end{equation}
The two matrices differing by an overall sign are regarded as identical.

If
\begin{equation}
a,a_1,a_2,\ldots,a_{\mu-1}
\tag{4}
\end{equation}
is a system of conjugate poles of $G$, then
\begin{equation}
f(x)=(x-a)(x-a_1)\cdots(x-a_{\mu-1})
\tag{5}
\end{equation}
is a polynomial of degree $\mu$ whose roots are those poles. For any substitution $\psi(x)=(\alpha x+\beta)/(\gamma x+\delta)$ of $G$, the polynomial
\begin{equation}
(\gamma x+\delta)^\mu f(\psi(x))
\tag{6}
\end{equation}
has as roots the inverse images of these poles, and therefore the same roots, apart from order. Thus it differs from $f$ only by a constant factor.

On homogenizing,
\begin{equation}
f(x_1,x_2)=x_2^\mu f\left(\frac{x_1}{x_2}\right),
\tag{7}
\end{equation}
this form is invariant up to a constant factor under $G$. Thus each system of conjugate poles gives an invariant form whose degree is the number of poles in the system and whose roots are precisely these poles. These forms, denoted $f_1,f_2,\ldots$, are called the fundamental forms of the group.

If an invariant form $F(x_1,x_2)$ of $G$ has a common factor with one of the fundamental forms $f(x_1,x_2)$, then all roots of $f$ are roots of $F$, and hence $F$ is divisible by $f$.

Finally, let $F$ be an invariant form of $G$ of degree lower than the degree of the group. If $\xi$ is a root of $F(x,1)=0$, then all $n$ values $\psi(\xi)$ obtained by applying the group substitutions are also roots. Since the degree is less than $n$, these cannot all be distinct, so $\xi$ must be a pole of the group. Hence $F$ is divisible by one of the fundamental forms. Repeating the argument gives: an invariant form of degree lower than the degree of $G$ is, up to a constant factor, a product of powers of the fundamental forms. Equivalently, an invariant form of lower degree which is not such a product must vanish identically.

\clearpage

\clearpage
\section*{Eighth Section. The Polyhedral Groups.}

\section*{\S\,55. The cyclic groups and the dihedral groups.}

We now pass to the actual construction of the several polyhedral groups from the general principles. First we have to show that all the kinds of groups recognized as possible in \S\,52 do in fact exist.

For the cyclic groups of degree $n$ there are only two systems of conjugate poles, each of which contains a single pole of multiplicity $n-1$. These two poles must therefore be the common poles of all substitutions of the group, and by \S\,51 they may be taken to be $0$ and $\infty$. We then have only the two linear ground forms
\begin{equation}
 f_1=x_1,\qquad f_2=x_2 .
\tag{1}
\end{equation}
All substitutions of the group are multiplicative, and the multiplier must be an $n$-th root of unity. They must therefore have the form
\begin{equation}
 x,\ \eps x,\ \eps^2x,\ldots,\eps^{n-1}x,
\tag{2}
\end{equation}
where $\eps$ is a primitive $n$-th root of unity. Thus for every $n$ we really obtain a cyclic group, denoted by $C_n$,
\begin{equation}
 C_n=\left\{\mat{\eps^h}{0}{0}{1}\ ;\ h=0,1,\ldots,n-1\right\},
\tag{3}
\end{equation}
or, written with determinant $1$,
\[
 \left\{\mat{\eps^h}{0}{0}{\eps^{-h}}\ ;\ h=0,1,\ldots,n-1\right\}.
\]

For the dihedral group of degree $n=2m$, which we shall denote by $D_m$, the preceding paragraph gives two systems of $m$ conjugate twofold poles and one system of two conjugate $m$-fold poles. The latter must therefore be the poles of one substitution. If we make them coincide with $0$ and $\infty$, we obtain the ground form of the second degree
\begin{equation}
 f_1=x_1x_2 .
\tag{4}
\end{equation}
This can be an invariant form of the group only when all substitutions occur in one of the two forms
\[
 \mat{\alpha}{0}{0}{\delta},\qquad
 \mat{0}{\beta}{\gamma}{0};
\]
in these, $\alpha:\delta$ and $-\beta:\gamma$ must be $m$-th roots of unity. Thus, if $\eps$ is a primitive $m$-th root of unity, the substitutions of the group are
\[
x,\ \eps x,\ \eps^2x,\ldots,\eps^{m-1}x,\quad
\frac1x,\ \frac{\eps}{x},\ \frac{\eps^2}{x},\ldots,\frac{\eps^{m-1}}x .
\]
With determinant $1$ this is
\begin{equation}
D_m=\left\{
 \mat{\eps^h}{0}{0}{\eps^{-h}},
 \mat{0}{\eps^h}{\eps^{-h}}{0}
 \ ;\ h=0,1,\ldots,m-1
\right\}.
\tag{5}
\end{equation}
It is immediate that these substitutions form a group.

Besides $0$ and $\infty$ there are the poles obtained from $x^m=\pm1$. They give the two ground forms of degree $m$,
\begin{equation}
 f_2=x_1^m+x_2^m,\qquad
 f_3=x_1^m-x_2^m .
\tag{6}
\end{equation}
This shows in particular that the dihedral and cyclic groups are wholly different, since their ground forms have different degrees.

A multiplicative transformation merely changes this representation of $D_m$ into one of the same kind and changes the ground forms only by constant factors. The dihedral group $D_m$ contains the cyclic group $C_m$ as a divisor, and $C_m$ is a normal divisor of $D_m$. Apart from such a multiplicative transformation, $D_m$ is completely determined by the group $C_m$ contained in it. For, if
\[
D_m=C_m+\varphi C_m,
\]
then, since $C_m$ must be normal, one has $\varphi C_m=C_m\varphi$; hence $\varphi$ has to be of the non-multiplicative form, and it can be brought to
\[
 \mat{0}{1}{1}{0}.
\]

Among the dihedral groups the group $D_2$ is especially important. It has three systems, each consisting of two twofold poles. It consists of the four substitutions
\begin{equation}
 \mat{1}{0}{0}{1},\quad
 \mat{-1}{0}{0}{1},\quad
 \mat{0}{1}{1}{0},\quad
 \mat{0}{1}{-1}{0},
\tag{7}
\end{equation}
and is called the four-group.

\section*{\S\,56. The tetrahedral group.}

For the tetrahedral group, \S\,52, III gives a system of six conjugate twofold poles. We may therefore assume, by \S\,53, that the group contains the two substitutions
\begin{equation}
 \theta(x)=-x,\qquad \psi(x)=\frac1x,
\tag{1}
\end{equation}
and hence also
\[
 \psi_1(x)=\theta\psi(x)=-\frac1x .
\]
The substitutions $\psi$ and $\psi_1$ have the poles $\pm1$ and $\pm\ii$; since these substitutions are of the second order and only twofold or threefold poles occur in $G$, these poles must be twofold. The substitution $\theta$ gives the twofold poles $0,\infty$. Thus the equation on which the twofold poles depend is
\[
 x(x^4-1)=0,
\]
and the first ground form, of degree six, is
\begin{equation}
 f=x_1x_2(x_1^4-x_2^4).
\tag{2}
\end{equation}

To find the remaining substitutions of the group we use the theorem of \S\,53: we may assume that a substitution $\chi$ in $G$ satisfies
\[
 \chi(-1)=0,\qquad \chi(+1)=\infty.
\]
It therefore has the form
\begin{equation}
 \chi(x)=\lambda\,\frac{x+1}{x-1},
\tag{3}
\end{equation}
where $\lambda$ is constant. Substitution of the six conjugate poles
\[
0,\infty,+1,-1,+\ii,-\ii
\]
must give the same six values in another order. Hence
\[
 \lambda=\pm\ii .
\]
We take $\lambda=\ii$. With
\[
 \theta(x)=-x,\qquad \psi(x)=\frac1x,\qquad
 \chi(x)=\ii\,\frac{x+1}{x-1}
\]
we obtain a group of twelve substitutions. They can be represented in the form
\begin{equation}
 \theta^\lambda\psi^\mu\chi^\nu,\qquad
 \lambda,\mu=0,1,\quad \nu=0,1,2,
\tag{4}
\end{equation}
and their products are governed by the relations
\begin{equation}
 \theta^2=\psi^2=\chi^3=1,\qquad
 \psi\theta=\theta\psi,\qquad
 \chi\theta=\psi\chi,\qquad
 \chi\psi=\theta\chi .
\tag{5}
\end{equation}
It follows from (5) that the four-group
\begin{equation}
 1,\quad \theta,\quad \psi,\quad \theta\psi
\tag{6}
\end{equation}
is a normal divisor of the tetrahedral group.

To obtain the two ground forms of the fourth degree $\Phi_1,\Phi_2$ belonging to the threefold poles, one may proceed as follows. Since these forms must admit $\theta$ and $\psi$ and cannot be divisible by $x_1x_2(x_1^4-x_2^4)$, they can only have the form
\[
 \Phi=x_1^4+m x_1^2x_2^2+x_2^4 .
\]
Application of $\chi$ determines $m$ and gives
\begin{equation}
 \Phi_1=x_1^4+2\sqrt{-3}\,x_1^2x_2^2+x_2^4,\qquad
 \Phi_2=x_1^4-2\sqrt{-3}\,x_1^2x_2^2+x_2^4.
\tag{7}
\end{equation}
The three ground forms $f,\Phi_1,\Phi_2$ satisfy the relation
\begin{equation}
 12\sqrt{-3}\, f^2=\Phi_1^3-\Phi_2^3 .
\tag{8}
\end{equation}
If the generating substitutions are written with determinant $1$, one obtains the corresponding real orthogonal ternary group. The form $f$ is an absolute invariant, while $\Phi_1,\Phi_2$ are interchanged by substitutions.

\section*{\S\,57. The octahedral group.}

In the octahedral group there occurs a system of six conjugate fourfold poles. We therefore have a substitution of the fourth order
\[
 \theta(x)=\ii x,\qquad \theta^4=1,
\]
whose poles give one pair of these conjugate fourfold poles. By \S\,53 we may take in the group the substitutions
\begin{equation}
 \theta(x)=\ii x,\qquad \psi(x)=\frac1x .
\tag{1}
\end{equation}
Then
\begin{equation}
 \psi^2=1,\qquad \theta\psi=\psi\theta^{-1}.
\tag{2}
\end{equation}
The ground form of the sixth degree belonging to the fourfold poles must remain unchanged up to a constant factor under $\theta$ and $\psi$ and must contain the factor $x_1x_2$. Thus
\begin{equation}
 f=x_1x_2(x_1^4-x_2^4).
\tag{3}
\end{equation}

To form the whole octahedral group one adds a substitution of order three, cyclically permuting the three systems of opposite octahedral edges. It may be chosen in the form
\begin{equation}
 \chi(x)=\ii\,\frac{x+1}{x-1}.
\tag{4}
\end{equation}
The substitutions generated by $\theta,\psi,\chi$ can be brought to the normal form
\begin{equation}
 \chi^\lambda\psi^\mu\theta^\nu,\qquad
 \lambda=0,1,2,\quad \mu=0,1,\quad \nu=0,1,2,3.
\tag{5}
\end{equation}
They are all distinct; the group therefore has degree $24$ and is the octahedral group.

The tetrahedral group forms a normal divisor of degree twelve. Thus the octahedral group is obtained from the tetrahedral group by adjoining the substitution $\theta(x)=\ii x$. The roots of
\[
 f=x_1x_2(x_1^4-x_2^4)
\]
correspond to the six vertices of the octahedron, or equivalently to the six faces of the cube. The two further ground forms are
\begin{equation}
 W=(x_1^4+x_2^4)^2+14x_1^4x_2^4
   =x_1^8+14x_1^4x_2^4+x_2^8,
\tag{6}
\end{equation}
and
\begin{equation}
 K=(x_1^4+x_2^4)^3-36x_1^4x_2^4(x_1^4+x_2^4)
   =x_1^{12}-33x_1^8x_2^4-33x_1^4x_2^8+x_2^{12}.
\tag{7}
\end{equation}
Between these three forms there is the identity
\begin{equation}
 W^3-K^2=108 f^4 .
\tag{8}
\end{equation}
The octahedral group contains the tetrahedral group and is obtained from it by adjoining a substitution of the fourth order. The relative invariants of the tetrahedral group thereby become absolute or are interchanged in pairs.

\section*{\S\,58. The icosahedral group.}

Since the icosahedral group has only one system of conjugate fivefold poles, by \S\,53 we may assume that it contains the two substitutions
\begin{equation}
 \theta(x)=\eps x,\qquad \psi(x)=-\frac1x,
\tag{1}
\end{equation}
where $\eps$ denotes a primitive fifth root of unity. We choose this form of $\psi$ because it simplifies the formulae.

The ground form of degree twelve belonging to the fivefold poles must, since it admits the substitutions (1), have the form
\begin{equation}
 f=x_1x_2(x_1^{10}+m x_1^5x_2^5-x_2^{10}).
\tag{2}
\end{equation}
To determine $m$, Weber uses the first invariant of the biquadratic form in the variables $\xi_1,\xi_2$. Since this covariant is of degree sixteen, it cannot be expressed as a product of the forms of degrees $12$, $20$, and $30$ that are to occur; consequently it must vanish identically. The relevant coefficient gives
\[
 m=\pm 11.
\]
One sign is transformed into the other by interchanging $x_1$ and $x_2$; we take
\begin{equation}
 f=x_1x_2(x_1^{10}+11x_1^5x_2^5-x_2^{10}).
\tag{3}
\end{equation}

To find the substitutions themselves, one adds to $\theta,\psi$ a substitution $\chi$ that sends two prescribed poles to two other prescribed poles. With
\[
 \omega=\eps+\eps^{-1},\qquad \omega'=\eps^2+\eps^{-2},\qquad
 \omega\omega'=-1
\tag{4}
\]
one obtains a representation
\begin{equation}
 \chi=\mat{\alpha}{\beta}{\gamma}{\delta},
\qquad
 \alpha\delta-\beta\gamma\ne0,
\tag{5}
\end{equation}
in which the coefficients are determined by requiring the images of the fivefold poles to fall again in the same system. Weber thereby obtains the sixty substitutions in the forms
\begin{equation}
 \theta^h,\qquad
 \psi\theta^h,\qquad
 \theta^r\chi\theta^s,\qquad
 \theta^r\chi\psi\theta^s\qquad (h,r,s=0,1,2,3,4).
\tag{6}
\end{equation}
Among the last two families the substitutions of the second, third, and fifth orders occur in exactly the required numbers. The conditions
\[
 r+s\equiv \pm2\pmod5,\qquad r+s\equiv \pm1\pmod5
\]
select the substitutions of the second and third orders.

The group property follows from the composition laws. First
\begin{equation}
 \theta^5=\psi^2=1,\qquad \theta\psi=\psi\theta^{-1},
\tag{7}
\end{equation}
and calculation with (4) gives in particular
\begin{equation}
 \chi\theta=\theta^{-1}\chi\theta^{-1},\qquad
 \chi\theta^2=\theta^2\chi\theta^2,
\tag{8}
\end{equation}
together with the further formulae by which every product of two substitutions of the forms (6) is again brought to one of the forms (6). This proves the existence of the icosahedral group. Written with determinant $1$, and with $\eps+\eps^{-1}$ substituted for $\omega$, $\chi$ gives the usual symmetric form of the icosahedral substitution.

The detailed calculation of $\chi$ uses a quadratic equation whose coefficients are determined by the selected fivefold poles. If
\[
 \chi=\mat{\alpha}{\beta}{\gamma}{\delta},
\]
then, after multiplication by powers of $\theta$, equations of the form
\[
 \alpha\gamma x^2+(\omega\delta+\beta\gamma)x+\beta\delta=0
\]
must contain the relevant poles. This yields ratios among the coefficients; after substituting (4), the remaining constants are fixed. The conditions for substitutions of the second and third orders,
\[
 \alpha+\delta=0,\qquad
 \alpha^2+\beta\gamma+\alpha\delta+\delta^2=0,
\]
give precisely the congruences for $r+s$ stated above.

\section*{\S\,59. The divisors of the icosahedral group.}

The divisors of the icosahedral group must be sought among the lower polyhedral groups. First there are the cyclic groups $C_2,C_3,C_5$, each generated by the period of an icosahedral substitution. The numbers of these groups are
\[
 15,\qquad 10,\qquad 6.
\]
The number of dihedral groups is obtained from the fact that a dihedral group is completely determined by the cyclic group contained in it. In the four-group one must notice, however, that the same group is obtained from any one of its three cyclic groups of order two; thus the number first obtained must be divided by $3$. Hence
\[
 \#D_2=5,\qquad \#D_3=10,\qquad \#D_5=6.
\]

Especially important are the tetrahedral groups contained in the icosahedral group. Such a group $Q$ must contain a four-group $D_2$; after such a group is chosen, the adjunction of a suitable substitution of order three determines the whole tetrahedral group. Weber obtains, in ordered form,
\begin{equation}
 Q=\{\,\varphi^\lambda\psi^\mu\chi^\nu;\ \lambda,\mu=0,1,\ \nu=0,1,2\,\}.
\tag{1}
\end{equation}
The composition formulae following from \S\,58 show that this is indeed a tetrahedral group. Since the four-group $D_2$ determines the tetrahedral group in which it lies, there are altogether five tetrahedral groups. Thus, apart from the whole icosahedral group and the identity group, only the groups listed above occur.

\section*{\S\,60. The ground forms of the icosahedron.}

The icosahedral group is isomorphic with the alternating group $A_5$. This is seen by considering the cosets of a tetrahedral group $Q$ in the icosahedral group. If an element of the icosahedral group is multiplied with the five cosets
\[
 Q,\ Q\theta,\ Q\theta^2,\ Q\theta^3,\ Q\theta^4,
\]
it effects a permutation of these five cosets. For the generating substitutions the corresponding permutations are even; hence the permutation group obtained is $A_5$.

The ground form of degree twelve is
\begin{equation}
 f=x_1x_2(x_1^{10}+11x_1^5x_2^5-x_2^{10}).
\tag{1}
\end{equation}
The ground form of degree twenty is obtained as the Hessian covariant of $f$. Weber writes it as
\begin{equation}
 H=-(x_1^{20}+x_2^{20})
    +228(x_1^{15}x_2^5-x_1^5x_2^{15})
    -494x_1^{10}x_2^{10}.
\tag{2}
\end{equation}
The ground form of degree thirty is defined as the functional determinant of $H$ and $f$:
\begin{equation}
 T=\frac1{20}\left(
 \frac{\partial f}{\partial x_1}\frac{\partial H}{\partial x_2}
-\frac{\partial f}{\partial x_2}\frac{\partial H}{\partial x_1}\right).
\tag{3}
\end{equation}
One obtains
\begin{equation}
 T=(x_1^{30}+x_2^{30})
 +522(x_1^{25}x_2^5-x_1^5x_2^{25})
 -10005(x_1^{20}x_2^{10}+x_1^{10}x_2^{20}).
\tag{4}
\end{equation}
Putting
\[
 \lambda=x_1^{10}+x_2^{10},\qquad \mu=x_1^5x_2^5,
\]
one obtains
\[
 H=-\lambda^2+228\lambda\mu-496\mu^2,
\]
and the numerical calculation yields the identity
\begin{equation}
 T^2+H^3=1728\,f^5.
\tag{5}
\end{equation}

\section*{\S\,61. The invariants of the icosahedron.}

From the relation (5) of the preceding paragraph it follows first that the quotients
\[
 \frac{H^3}{f^5},\qquad \frac{T^2}{f^5}
\]
remain absolutely unchanged under the icosahedral substitutions. Since under homogeneous icosahedral substitutions of determinant $1$ the forms $f,H,T$ could only acquire powers of a common factor, and since the quotients remain fixed, $f,H,T$ themselves must be absolute invariants.

The functions $f,H,T$ are the ground forms of the icosahedral group. It remains to prove that every invariant form of the icosahedral group can be represented as an integral rational function of these three forms. Weber gives the proof for the icosahedral group; the corresponding arguments for the tetrahedral and octahedral groups are analogous.

First, no two of the forms $f,H,T$ have a common divisor. For if such a divisor existed, the corresponding equations would have a common root; then one system of conjugate poles would have to belong to two different pole systems at once, which is impossible.

Next let
\[
 J(x_1,x_2)
\]
be any invariant form of the icosahedral group, and let $\xi$ be a root of $J(x,1)=0$. Then all quantities $\psi(\xi)$ obtained from $\xi$ by the fractional icosahedral substitutions are again roots of the same equation. Consequently, if $J$ has a common divisor with one of $f,H,T$, then $J$ is divisible by that form.

After the highest possible powers of $f,H,T$ have been removed from $J$, let $\xi$ be a root of the remaining factor. A constant $\eta$ can then be chosen so that $\xi$ is a simple root of a form
\[
 \Theta=H^3-\eta f^5
\]
or of an equivalent combination of the ground forms. The whole orbit of $\xi$ is then also a set of zeros of $J$, and therefore $J$ is divisible by $\Theta$. Repetition of the argument gives
\[
 J=\sum c_{\alpha\beta\gamma} f^\alpha H^\beta T^\gamma .
\]
Only the relation
\[
 T^2+H^3=1728f^5
\]
has to be observed among $f,H,T$.



\section*{\S\,62. Polyhedral groups of the second kind. Crystallographic groups.}

We have already, in \S\,51, drawn attention to groups of linear ternary substitutions which contain, besides proper, also improper orthogonal substitutions, and which we call groups of the second kind. We shall now consider the finite ones among them as polyhedral groups of the second kind, or also as enlarged polyhedral groups, and shall call the substitutions contained in them with determinant $-1$ substitutions of the second kind.

By the results of \S\,50 these groups are isomorphic with the groups of binary linear substitutions of determinant $+1$ and $-1$, whereas in the fractional substitutions which led us to the polyhedral groups the two kinds are not distinguishable.

The finite groups of the second kind are important, apart from their general group-theoretic interest, because of their application in crystallography. Their complete determination now offers no essential difficulty. By substitutions without further qualification we understand binary linear substitutions with determinant $\pm1$, and recall that two such substitutions whose coefficients differ only by the common sign,
\[
 \begin{pmatrix}\alpha&\beta\\ \gamma&\delta\end{pmatrix},\qquad
 \begin{pmatrix}-\alpha&-\beta\\ -\gamma&-\delta\end{pmatrix},
\]
are not regarded as different (\S\,49).

Since a group of linear substitutions contains the identical substitution, which has determinant $+1$, every group $Q$ of the second kind must contain proper substitutions besides improper ones, and these proper substitutions form by themselves a group $P$. If $\varphi$ is any substitution of the second kind occurring in $Q$, then every composition of $\varphi$ with a substitution of the second kind from $Q$ occurs in $P$, and therefore we can always decompose $Q$ as
\begin{equation}
 Q=P+P\varphi . \tag{1}
\end{equation}
The group $P$ must be one of the previously considered polyhedral groups. It must satisfy $\varphi^{-1}P\varphi=P$, and thus $P$ is a normal divisor of $Q$. Conversely, if $P$ is a polyhedral group and $\varphi$ a substitution of the second kind satisfying $\varphi^{-1}P\varphi=P$, then $Q=P+P\varphi$ is a group of the second kind.

Here also we regard two groups which arise from one another by transformation as not essentially different. In order to form all $Q$, we must go through the several cases.

I. If $P$ is the unit group, which we denote as the cyclic group of first degree by $C_1$, then $P$ consists only of the identity. Hence $\varphi^2=1$, and this is also sufficient. By transformation we may bring $\varphi$ into the form
\[
 \begin{pmatrix}\alpha&0\\0&-\alpha^{-1}\end{pmatrix}
\]
and obtain as condition $\alpha^2=\pm1$. Thus there are two forms of the substitution $\varphi$:
\[
 \varphi'=\begin{pmatrix}i&0\\0&i\end{pmatrix},\qquad
 \varphi''=\begin{pmatrix}1&0\\0&-1\end{pmatrix}.
\]
The first is called inversion, the second reflection. Transferred by \S\,50, (1), to a rectangular coordinate transformation, $\varphi'$ means the simultaneous change of sign of all three coordinates, and $\varphi''$ the interchange of $x$ with $-x$. Hence the two groups of the second kind are
\begin{equation}
 C_1'=\left(\begin{pmatrix}i^h&0\\0&i^h\end{pmatrix}\right),\qquad
 C_1''=\left(\begin{pmatrix}1&0\\0&(-1)^h\end{pmatrix}\right),
 \quad h=0,1. \tag{2}
\end{equation}

II. Let $P$ be a cyclic group $C_n$, $n>1$. Put
\begin{equation}
 c=\begin{pmatrix}\varepsilon&0\\0&\varepsilon^{-1}\end{pmatrix},
 \qquad \varepsilon=e^{\pi i/n}. \tag{3}
\end{equation}
Then, by \S\,55,
\begin{equation}
 C_n=1,c,c^2,\ldots,c^{n-1},\qquad
 c^s=\begin{pmatrix}\varepsilon^s&0\\0&\varepsilon^{-s}\end{pmatrix}. \tag{4}
\end{equation}
If now
\begin{equation}
 \varphi=\begin{pmatrix}\alpha&\beta\\\gamma&\delta\end{pmatrix},
 \qquad \alpha\delta-\beta\gamma=-1, \tag{5}
\end{equation}
then first
\begin{equation}
 \varphi^{-1}c\varphi=
 \begin{pmatrix}
 \alpha\delta\varepsilon-\beta\gamma\varepsilon^{-1}&
 \beta\delta(\varepsilon-\varepsilon^{-1})\\
 -\alpha\gamma(\varepsilon-\varepsilon^{-1})&
 \alpha\delta\varepsilon^{-1}-\beta\gamma\varepsilon
 \end{pmatrix}, \tag{6}
\end{equation}
and
\begin{equation}
 \varphi^2=
 \begin{pmatrix}
 \alpha^2+\beta\gamma&(\alpha+\delta)\beta\\
 (\alpha+\delta)\gamma&\beta\gamma+\delta^2
 \end{pmatrix}. \tag{7}
\end{equation}
Both substitutions must occur in the series (4). This gives the three forms
\[
\begin{aligned}
 C_n'&=\left(
 \begin{pmatrix}e^{\pi i h/(2n)}&0\\
 0&(-1)^h e^{-\pi i h/(2n)}\end{pmatrix}
 \right), &&h=0,1,\ldots,2n-1,\\[0.4em]
 C_n''&=\left(
 \begin{pmatrix}e^{\pi i h/n}&0\\
 0&\pm e^{-\pi i h/n}\end{pmatrix}
 \right), &&h=0,1,\ldots,n-1,\\[0.4em]
 C_n'''&=\left(
 \begin{pmatrix}e^{\pi i h/n}&0\\0&-e^{-\pi i h/n}\end{pmatrix},
 \begin{pmatrix}0&e^{\pi i h/n}\\ e^{-\pi i h/n}&0\end{pmatrix}
 \right), &&h=0,1,\ldots,n-1.
\end{aligned}
\]
If $n$ is even, $C_n'$ and $C_n''$ both occur under $C_n''$; if $n$ is odd, $C_1'$ occurs in $C_n'$ and $C_1''$ in $C_n''$.

III. Let $P$ be the dihedral group $D_n$, consisting of the substitutions
\begin{equation}
 D_n=1,c,c^2,\ldots,c^{n-1},\quad
 d,cd,c^2d,\ldots,c^{n-1}d,\qquad
 d=\begin{pmatrix}0&i\\ i&0\end{pmatrix}=C_n+C_nd.
 \tag{8}
\end{equation}
The substitution $c$ is of the $n$th degree, while the $c^\lambda d$ are all of the second degree. But $\varphi^{-1}c\varphi$ is of the $n$th degree and, if $n>2$, must therefore be contained among the powers of $c$. Thus for the dihedral groups of the second kind there are the forms
\begin{equation}
 C_n'+C_n'd,\qquad C_n''+C_n''d,\qquad C_n'''+C_n'''d, \tag{9}
\end{equation}
of which the first two are
\[
\begin{aligned}
D_n'&=\left(
\begin{pmatrix}e^{\pi i h/(2n)}&0\\0&(-1)^h e^{-\pi i h/(2n)}\end{pmatrix},
\begin{pmatrix}0&-e^{\pi i h/(2n)}\\ (-1)^{h+n}e^{-\pi i h/(2n)}&0\end{pmatrix}
\right),\\[-0.2em]
&\hfill h=0,1,\ldots,2n-1,\\[0.4em]
D_n''&=\left(
\begin{pmatrix}e^{\pi i h/n}&0\\0&\pm e^{-\pi i h/n}\end{pmatrix},
\begin{pmatrix}0&e^{\pi i(2h+n)/(2n)}\\
 \pm e^{-\pi i(2h+n)/(2n)}&0\end{pmatrix}
\right),\\[-0.2em]
&\hfill h=0,1,\ldots,n-1.
\end{aligned}
\]
The third group gives
\[
\begin{gathered}
 \begin{pmatrix}e^{\pi i h/n}&0\\0&e^{-\pi i h/n}\end{pmatrix},\quad
 \begin{pmatrix}0&e^{\pi i h/n}\\e^{-\pi i h/n}&0\end{pmatrix},\\
 \begin{pmatrix}0&e^{\pi i(2h+n)/(2n)}\\-e^{-\pi i(2h+n)/(2n)}&0\end{pmatrix},\quad
 \begin{pmatrix}e^{\pi i(2h+n)/(2n)}&0\\0&-e^{-\pi i(2h+n)/(2n)}\end{pmatrix},
 \quad h=0,1,\ldots,n-1,
\end{gathered}
\]
which, when $n$ is even, agrees with $D_n''$, and when $n$ is odd with $D_n'$.

For $n=2$, where
\[
 c=\begin{pmatrix}i&0\\0&-i\end{pmatrix},
\]
there is only an apparent exception, because in that case
\begin{equation}
 \varphi^{-1}c\varphi=d\quad \hbox{or}\quad
 \varphi^{-1}c\varphi=cd \tag{10}
\end{equation}
may occur. Following these assumptions by a simple computation gives two more groups:
\begin{equation}
\begin{array}{cccccccc}
1,&c,&d,&cd,&\varphi_1,&\varphi_1c,&\varphi_1d,&\varphi_1cd,\\
1,&c,&d,&cd,&\varphi_2,&\varphi_2c,&\varphi_2d,&\varphi_2cd,
\end{array} \tag{11}
\end{equation}
where
\[
 \varphi_1=\begin{pmatrix}1/\sqrt2&1/\sqrt2\\1/\sqrt2&-1/\sqrt2\end{pmatrix},\qquad
 \varphi_2=\begin{pmatrix}1/\sqrt2&-i/\sqrt2\\i/\sqrt2&-1/\sqrt2\end{pmatrix}.
\]
But
\[
\begin{gathered}
 \varphi_2^{-1}\varphi_1\varphi_2=
 \begin{pmatrix}0&e^{\pi i/4}\\ e^{-\pi i/4}&0\end{pmatrix},\qquad
 \varphi_1^{-1}\varphi_2\varphi_1=
 \begin{pmatrix}0&e^{\pi i/4}\\ e^{-\pi i/4}&0\end{pmatrix},\\
 c\varphi_1=\varphi_1d,\quad d\varphi_1=\varphi_1c,\quad
 cd\varphi_1=\varphi_1dc,\\
 c\varphi_2=\varphi_2dc,\quad d\varphi_2=\varphi_2d,\quad
 cd\varphi_2=\varphi_2c,
\end{gathered}
\]
and these formulae show that the groups (11) are transformed into $D_2''$ by transformation with $\varphi_1$ and $\varphi_2$.

For the remaining polyhedral groups of the second kind the following general observations are useful. The inversion
\[
 j=\begin{pmatrix}i&0\\0&i\end{pmatrix}
\]
is a similarity substitution (\S\,37) and therefore commutes with every other substitution. Hence from every polyhedral group of the first kind we can derive at least one group of the second kind:
\begin{equation}
 P+Pj. \tag{12}
\end{equation}
To find the other groups of this kind which may still exist, one recalls from Sylow's theorem that every group $G$ must contain a group whose degree is the highest power of $2$ dividing the degree of $G$. The enlarged tetrahedral, octahedral and icosahedral groups have degrees $24,48,120$, and therefore must contain divisors of degrees $8,16,8$.

Among the enlarged dihedral groups,
\[
D_2'=\left(
\begin{array}{cccc}
\begin{pmatrix}1&0\\0&1\end{pmatrix},&
\begin{pmatrix}\sqrt i&0\\0&i\sqrt i\end{pmatrix},&
\begin{pmatrix}i&0\\0&-i\end{pmatrix},&
\begin{pmatrix}i\sqrt i&0\\0&\sqrt i\end{pmatrix},\\[0.4em]
\begin{pmatrix}0&-1\\1&0\end{pmatrix},&
\begin{pmatrix}0&-i\sqrt i\\\sqrt i&0\end{pmatrix},&
\begin{pmatrix}0&i\\i&0\end{pmatrix},&
\begin{pmatrix}0&-\sqrt i\\i\sqrt i&0\end{pmatrix}
\end{array}
\right)
\]
arises from $D_2$ by inversion, while composition of $D_2$ with
\[
 j_1=\begin{pmatrix}0&-\sqrt i\\ i\sqrt i&0\end{pmatrix}
\]
gives
\[
D_2''=D_2+D_2j_1 .
\]
In general,
\[
 j_1^{-1}\begin{pmatrix}\alpha&\beta\\ \gamma&\delta\end{pmatrix}j_1
 =
 \begin{pmatrix}\delta&\gamma i\\ -\beta i&\alpha\end{pmatrix}.
\]
It follows from \S\,56, (4), that $j_1^{-1}Tj_1=T$. Thus from the tetrahedral group two enlarged groups may be derived:
\[
T'=T+Tj,\qquad T''=T+Tj_1,
\]
and only two. Since the substitution
\[
\begin{pmatrix}0&e^{\pi i/8}\\e^{-\pi i/8}&0\end{pmatrix}
\]
does not commute with the octahedral group of \S\,57, (8), no further group of the second kind can be derived from the octahedral group except by inversion; the same holds for the icosahedral group.

Thus we obtain the following further groups of the second kind:
\[
\begin{array}{lll}
\mathrm{IV}.&1)&T'=T+Tj,\qquad 2)\ T''=T+Tj_1,\\
\mathrm{V}.&&O'=O+Oj,\\
\mathrm{IV}.&&J'=J+Jj.
\end{array}
\]

The thirty-two symmetry systems of the crystallographers are contained here. They are the groups
\[
\begin{array}{llll}
C_1,&C_1',&C_1'',\\
C_2,&C_2',&C_2'',&C_2''',\\
C_3,&C_3',&C_3'',&C_3''',\\
C_4,&C_4',&C_4'',\\
C_6,&C_6',&C_6'',\\[0.3em]
D_2,&D_2',&D_2'',\\
D_3,&D_3',&D_3'',\\
D_4,&D_4',\\
D_6,&D_6'',\\
T,&T',&T'',\\
O,&O'.
\end{array}
\]
The remaining polyhedral groups are excluded from crystallography by the crystallographic law of rational indices.\footnote{Cf. Schönfliess, \emph{Krystallsysteme und Krystallstructur}. Leipzig 1891.}


\section*{Ninth Section. Congruence Groups.}

\section*{\S\,63. Function congruences.}

From the linear substitutions whose construction and composition we have learned in the preceding sections, one can derive a quite different kind of finite groups, the congruence groups. These congruence groups have manifold interest. They arise first in the theory of elliptic functions, where they occur as Galois groups of modular equations. They are also important for general group theory, since they give a means of constructing whole series of simple groups. This is all the more remarkable because, as we saw in the fourth section, the simple groups are rare, at least among the lower degrees.

The theory of these congruence groups is not only extraordinarily generalized, but its governing laws emerge much more sharply if one is based on an enlargement of the congruence concept due to Gauss, which has since been developed and applied several times, especially by Galois, to problems of group theory.\footnote{Galois, ``Sur la théorie des nombres'', \emph{Bulletin des sciences mathématiques de Férussac}, 1830; cf. the note to \S\,151 of the first volume. Also Schönemann and Dedekind on higher congruences.}

To obtain a basis for the theory, consider an integral function of a variable $t$,
\begin{equation}
 f(t)=a_0t^n+a_1t^{\,n-1}+\cdots+a_{n-1}t+a_n, \tag{1}
\end{equation}
whose coefficients are integers. We choose, moreover, a prime number $p$ as modulus, and assume that $a_0$ is not divisible by $p$.

Two integral functions of $t$ in which corresponding coefficients are congruent modulo $p$ shall themselves be called congruent modulo $p$. The function $f(t)$ is called reducible modulo $p$ if there are two integral functions with integer coefficients, $\varphi(t),\psi(t)$, in each of which at least one term depending on $t$ has a coefficient not divisible by $p$, such that
\begin{equation}
 f(t)\equiv\varphi(t)\psi(t)\pmod p. \tag{2}
\end{equation}
In $\varphi(t)$ and $\psi(t)$ all terms whose coefficients are divisible by $p$ may be omitted; then the degree of $\varphi(t)\psi(t)$ must agree with the degree of $f(t)$. Thus the degrees of $\varphi(t)$ and $\psi(t)$ must be lower than $n$. If there are no such functions, $f(t)$ is called irreducible modulo $p$.

A function irreducible modulo $p$ is certainly absolutely irreducible in the field of rational numbers. The converse is not necessary. For example, the quadratic function $t^2-R$ is reducible modulo $p$ when $R$ is a quadratic residue of $p$; for, if $a^2\equiv R\pmod p$, then
\[
 t^2-R\equiv(t-a)(t+a)\pmod p.
\]
On the other hand $t^2-N$, when $N$ is a non-residue of $p$, is irreducible, because otherwise $t^2-N$ would be divisible by $p$ for a rational $t$. The functions
\[
 t^2+t+1,\qquad t^3+t+1
\]
give another example; they are irreducible for the modulus $2$.

Let
\begin{equation}
 P=t^n+p_1t^{\,n-1}+p_2t^{\,n-2}+\cdots+p_{n-1}t+p_n \tag{3}
\end{equation}
be a function of degree $n$ irreducible modulo $p$, with leading coefficient equal to $1$. From every other integral function $F(t)$ with integer coefficients one obtains a remainder $\Phi(t)$ of degree lower than $n$ by writing
\begin{equation}
 F(t)=QP+\Phi(t). \tag{4}
\end{equation}
This relation of $\Phi(t)$ to $F(t)$ is called a congruence modulo $P$. Two functions $F(t),F_1(t)$ with the same remainder are congruent modulo $P$:
\[
 F(t)\equiv F_1(t)\pmod P.
\]
If their remainders are not equal but congruent modulo $p$, then $F,F_1$ are congruent modulo the double modulus $P,p$:
\begin{equation}
 F(t)\equiv F_1(t)\pmod {P,p}. \tag{5}
\end{equation}

For this congruence the following theorem holds:

\emph{1. The product of two functions $F(t),F_1(t)$ cannot be congruent to zero unless one factor is congruent to zero.}

The proof follows from the algorithm of the greatest common divisor. If $P_1$ has degree lower than $P$ and is not congruent to zero modulo $p$, one can determine functions $P_2,P_3,\ldots$ of decreasing degree and quotients $Q_1,Q_2,\ldots$ so that
\begin{equation}
 P=Q_1P_1+P_2,\quad P_1=Q_2P_2+P_3,\quad
 \ldots,\quad P_{r-2}=Q_{r-1}P_{r-1}+P_r\pmod p. \tag{6}
\end{equation}
This chain ends when $P_r$ is a constant. That constant cannot be $0\pmod p$, since otherwise $P$ would be congruent to a product $TP_{r-1}$ and hence would not be irreducible modulo $p$. Multiplying the equations by a function $Q$ satisfying $QP_1\equiv0\pmod {P,p}$ gives successively $QP_2\equiv0,\ldots,QP_r\equiv0$, and therefore $Q\equiv0$.

The remainders of all functions $F(t)$ modulo $P$ have the form
\begin{equation}
 X=X(t)=x_0+x_1t+\cdots+x_{n-1}t^{\,n-1}. \tag{7}
\end{equation}
Taking only incongruent remainders modulo $p$, it is enough to let $x_0,x_1,\ldots,x_{n-1}$ run through $0,1,\ldots,p-1$. There are therefore exactly $p^n$ functions incongruent modulo $P,p$.

\section*{\S\,64. Congruence fields.}

The considerations of the preceding paragraph are expressed in another form when, with Galois, one introduces an imaginary number $\varepsilon$ satisfying
\begin{equation}
 P(\varepsilon)\equiv0\pmod p. \tag{1}
\end{equation}
No ordinary rational integer of this kind exists when $n>1$; nevertheless such a symbol $\varepsilon$ may be used in calculation. Addition, subtraction and multiplication are performed by the ordinary rules of symbolic calculation, while the equation $P(\varepsilon)=0$ may be used at will.

All numbers obtained in this way can be brought to the form
\begin{equation}
 \alpha=a_0+a_1\varepsilon+\cdots+a_{n-1}\varepsilon^{\,n-1}. \tag{2}
\end{equation}
Since the modulus $p$ is always retained, two numbers congruent modulo $p$ are simply regarded as equal. These numbers $\alpha$ are called Galois imaginaries. The system belonging to a fixed $\varepsilon$ will be denoted by $\mathfrak C$.

First,

\emph{2. There are $p^n$ and no more different numbers in $\mathfrak C$.}

From theorem 1 one further obtains:

\emph{3. The product of two or more numbers of $\mathfrak C$ is zero if and only if at least one of the factors is zero.}

Every system $\mathfrak C$ contains the residues $0,1,\ldots,p-1$ of the natural numbers, and for $n=1$ it is identical with that system. If $\alpha$ is fixed and non-zero, then $\alpha\xi$ runs through the whole system together with $\xi$; hence:

\emph{4. Given $\alpha,\beta$ in $\mathfrak C$, with $\alpha\ne0$, there is one and only one $\gamma$ in $\mathfrak C$ satisfying $\alpha\gamma=\beta$.}

Thus division in $\mathfrak C$ is allowed. The system $\mathfrak C$ may therefore be called a finite field. Its degree is $n$, and its modulus is $p$.

If $\alpha\ne0$, multiplication by $\alpha$ permutes all non-zero elements. Multiplying all equations $\alpha\xi=\eta$ gives
\[
 \alpha^{p^n-1}\Pi=\Pi,
\]
where $\Pi$ is the product of all non-zero elements, and hence Fermat's theorem
\begin{equation}
 \alpha^{p^n-1}=1. \tag{3}
\end{equation}
After multiplication by $\alpha$,
\begin{equation}
 \alpha^{p^n}=\alpha, \tag{4}
\end{equation}
which also holds for $\alpha=0$.

If $\alpha\ne0$, there is a least positive exponent $e$ for which
\begin{equation}
 \alpha^e=1 \tag{5}
\end{equation}
holds. Then $\alpha$ is said to belong to the exponent $e$, and $e$ divides $p^n-1$. For the divisors $e$ of $p^n-1$,
\begin{equation}
 \sum \varphi(e)=p^n-1. \tag{6}
\end{equation}
If $\psi(e)$ denotes the number of elements belonging to the exponent $e$, then $\psi(e)=\varphi(e)$ or $0$, and since every non-zero element belongs to one and only one exponent,
\begin{equation}
 \sum \psi(e)=p^n-1. \tag{7}
\end{equation}
Thus $\psi(e)=\varphi(e)$ for every divisor. There are therefore $\varphi(p^n-1)$ elements $\gamma$ belonging to the exponent $p^n-1$, and their powers
\begin{equation}
 1,\gamma,\gamma^2,\ldots,\gamma^{p^n-2} \tag{8}
\end{equation}
exhaust all non-zero elements of $\mathfrak C$. Such elements are primitive roots of $\mathfrak C$.

The elements of $\mathfrak C$ divide into squares and non-squares. If $p=2$, every element is a square. If $p$ is odd, one half of the non-zero elements are squares and the other half are non-squares. Counting zero among the squares:

\emph{6. If $p=2$, all numbers in $\mathfrak C$ are squares. If $p$ is odd, there are $\frac12(p^n+1)$ squares and $\frac12(p^n-1)$ non-squares in $\mathfrak C$.}

For odd $p$, a primitive root is always a non-square; in (8), even exponents give squares and odd exponents give non-squares. Products and quotients of squares are squares; a non-square times a non-zero square is a non-square; and the product of two non-squares is a square.

Finally one will need the theorem:

\emph{7. Every non-square is the sum of two squares.}

For odd $p$, every $\beta$ is a difference of two squares:
\[
 \left(\frac{\beta+1}{2}\right)^2-\left(\frac{\beta-1}{2}\right)^2=\beta.
\]
If $-1$ is a square, this difference is already a sum of two squares. If $-1$ is a non-square, choose a square $\xi^2$ followed by a non-square $\xi^2+1$. For a non-square $\beta$,
\[
 \beta:(\xi^2+1)=\gamma^2
\]
is a square, and therefore
\[
 \beta=(\gamma\xi)^2+\gamma^2.
\]

\section*{\S\,65. Congruence groups in the field $\mathfrak C$.}

In every congruence field $\mathfrak C$ one can form a finite group of linear substitutions by allowing $\alpha,\beta,\gamma,\delta$ to run through all numbers of $\mathfrak C$ and composing two substitutions by
\begin{equation}
 \begin{pmatrix}\alpha&\beta\\\gamma&\delta\end{pmatrix}
 \begin{pmatrix}\alpha'&\beta'\\\gamma'&\delta'\end{pmatrix}
 =
 \begin{pmatrix}
 \alpha\alpha'+\beta\gamma'&\alpha\beta'+\beta\delta'\\
 \gamma\alpha'+\delta\gamma'&\gamma\beta'+\delta\delta'
 \end{pmatrix}. \tag{2}
\end{equation}
Calling $\alpha\delta-\beta\gamma$ the determinant, the determinant of a product is the product of determinants. Imposing
\begin{equation}
 \alpha\delta-\beta\gamma=1 \tag{3}
\end{equation}
therefore gives a group. Its degree is
\begin{equation}
 p^n(p^{2n}-1). \tag{4}
\end{equation}

If $p$ is odd, the two elements
\begin{equation}
 \begin{pmatrix}\alpha&\beta\\\gamma&\delta\end{pmatrix},\qquad
 \begin{pmatrix}-\alpha&-\beta\\-\gamma&-\delta\end{pmatrix} \tag{5}
\end{equation}
are distinct, but not distinct as fractional substitutions. Thus the system decomposes into pairs of this form, which may themselves be regarded as elements of a group of degree
\begin{equation}
 \frac12 p^n(p^{2n}-1). \tag{6}
\end{equation}
This is the congruence group belonging to $\mathfrak C$, denoted by $E$.

Writing the elements of $\mathfrak C$ as
\begin{equation}
 \xi=x_0+x_1\varepsilon+\cdots+x_{n-1}\varepsilon^{n-1}, \tag{7}
\end{equation}
the substitutions
\begin{equation}
 A_0=\begin{pmatrix}0&1\\-1&0\end{pmatrix},\qquad
 B_h=\begin{pmatrix}1&0\\\varepsilon^h&1\end{pmatrix},
 \quad h=0,1,\ldots,n-1, \tag{8}
\end{equation}
generate the group. This follows from
\[
 \begin{pmatrix}1&0\\ \gamma&1\end{pmatrix}
 \begin{pmatrix}1&0\\ \gamma'&1\end{pmatrix}
 =
 \begin{pmatrix}1&0\\ \gamma+\gamma'&1\end{pmatrix},
\]
and hence
\begin{equation}
 \begin{pmatrix}1&0\\ \gamma&1\end{pmatrix}
 =B_0^{c_0}B_1^{c_1}\cdots B_{n-1}^{c_{n-1}}. \tag{9}
\end{equation}
Composition with $A_0$ gives the further elementary forms
\begin{equation}
\begin{gathered}
 A_0\begin{pmatrix}1&0\\ \alpha&1\end{pmatrix}
 =\begin{pmatrix}-\alpha&1\\-1&0\end{pmatrix},\quad
 \begin{pmatrix}1&0\\ \delta&1\end{pmatrix}A_0
 =\begin{pmatrix}0&1\\-1&\delta\end{pmatrix},\\
 A_0\begin{pmatrix}1&0\\ -\beta&1\end{pmatrix}
 \begin{pmatrix}0&-1\\1&0\end{pmatrix}
 =\begin{pmatrix}1&\beta\\0&1\end{pmatrix}.
\end{gathered} \tag{10}
\end{equation}
Furthermore
\begin{equation}
 \begin{pmatrix}\alpha&1\\-1&0\end{pmatrix}
 \begin{pmatrix}1&\beta\\0&1\end{pmatrix}
 \begin{pmatrix}1&0\\\gamma&1\end{pmatrix}
 =
 \begin{pmatrix}\alpha+\gamma+\alpha\beta\gamma&\alpha\beta+1\\
 -\beta\gamma-1&-\beta
 \end{pmatrix}. \tag{11}
\end{equation}
By choosing $\alpha$ and $\gamma$, and then removing the restriction on $\delta$ by multiplying by $A_0$, all substitutions are obtained.

The group $E$ contains a divisor of index $p^n+1$, and it is isomorphic to a permutation group on $p^n+1$ symbols. The fractional expression
\begin{equation}
 \eta=\frac{\alpha\xi+\beta}{\gamma\xi+\delta} \tag{12}
\end{equation}
maps each $\xi$ to a value $\eta$, except where the denominator vanishes. The exception is removed by adjoining an element $\infty$.

\section*{\S\,66. Simplicity of the group $E$.}

The first question in a closer investigation of $E$ concerns possible normal divisors. It can be proved that, apart from two very simple exceptions, no such normal divisor exists.

Assume that $G$ is a normal divisor of $E$ and contains at least one substitution different from the identity,
\begin{equation}
 A=\begin{pmatrix}\alpha&\beta\\\gamma&\delta\end{pmatrix}. \tag{1}
\end{equation}
For every element $T$ of $E$,
\begin{equation}
 TAT^{-1} \tag{2}
\end{equation}
also belongs to $G$. It is therefore enough to show that the generators $A_0,B_h$ occur in $G$.

First, $G$ contains a substitution of the form
\begin{equation}
 \begin{pmatrix}0&\beta\\ \gamma&\delta\end{pmatrix}. \tag{3}
\end{equation}
Indeed,
\begin{equation}
 \begin{pmatrix}\xi&1\\-1&0\end{pmatrix}
 \begin{pmatrix}\alpha&\beta\\\gamma&\delta\end{pmatrix}
 \begin{pmatrix}0&-1\\1&\xi\end{pmatrix}
 =
 \begin{pmatrix}
 \beta\xi+\delta&\beta\xi^2+(\delta-\alpha)\xi-\gamma\\
 -\beta&-\beta\xi+\alpha
 \end{pmatrix}. \tag{4}
\end{equation}
If $\beta\ne0$, choose $\xi$ from $\beta\xi+\delta=0$; if $\beta=0$, one first makes a second application of the same formula.

Second, $G$ contains a substitution of the form
\[
 \begin{pmatrix}0&1\\-1&\xi\end{pmatrix}.
\]
This follows from
\begin{equation}
 \begin{pmatrix}\lambda&\mu\\ \nu&\rho\end{pmatrix}
 \begin{pmatrix}0&\beta\\ \gamma&\delta\end{pmatrix}
 \begin{pmatrix}\rho&-\mu\\-\nu&\lambda\end{pmatrix}
 =
 \begin{pmatrix}
 \mu\rho\gamma-\lambda\nu\beta-\mu\nu\delta&
 -\mu^2\gamma+\lambda^2\beta+\lambda\mu\delta\\
 \rho^2\gamma-\nu^2\beta-\rho\nu\delta&
 -\rho\mu\gamma+\lambda\nu\beta+\lambda\rho\delta
 \end{pmatrix}, \tag{5}
\end{equation}
provided
\begin{equation}
 \lambda\rho-\mu\nu=1,\quad
 \mu\rho\gamma-\lambda\nu\beta-\mu\nu\delta=0,\quad
 -\mu^2\gamma+\lambda^2\beta+\lambda\mu\delta=1 \tag{6}
\end{equation}
can be satisfied. The last two equations imply
\begin{equation}
 \lambda\beta+\mu\delta=\rho,\qquad \mu\gamma=\nu, \tag{7}
\end{equation}
and substitution gives
\begin{equation}
 \lambda^2\beta+\lambda\mu\delta-\mu^2\gamma=1. \tag{8}
\end{equation}
The solvability is obtained by the square and non-square discussion of \S\,64.

Third, except when $n=1,p=2$, $G$ contains
\[
 A_0=\begin{pmatrix}0&1\\-1&0\end{pmatrix}.
\]
First one obtains in $G$
\[
 \begin{pmatrix}0&-1\\1&0\end{pmatrix}
 \begin{pmatrix}0&1\\-1&\xi\end{pmatrix}
 \begin{pmatrix}0&1\\-1&0\end{pmatrix}
 =
 \begin{pmatrix}\xi&1\\-1&0\end{pmatrix},
\]
and therefore
\[
 \begin{pmatrix}0&1\\-1&\xi\end{pmatrix}
 \begin{pmatrix}\xi&1\\-1&0\end{pmatrix}
 =
 \begin{pmatrix}1&0\\2\xi&1\end{pmatrix}.
\]
Consequently,
\[
 \begin{pmatrix}1&0\\2\xi&1\end{pmatrix}
 \begin{pmatrix}0&1\\-1&\xi\end{pmatrix}
 =
 \begin{pmatrix}0&1\\-1&3\xi\end{pmatrix},\qquad
 \begin{pmatrix}1&0\\2\xi&1\end{pmatrix}
 \begin{pmatrix}0&1\\-1&3\xi\end{pmatrix}
 =
 \begin{pmatrix}0&1\\-1&5\xi\end{pmatrix},
\]
and so on, giving
\[
 \begin{pmatrix}0&1\\-1&m\xi\end{pmatrix}
\]
for every odd integer $m$. If $p$ is odd one may take $m=p$ and obtains $A_0$. If $p=2$, let $\rho$ be a non-zero number in $\mathfrak C$; then
\[
 \begin{pmatrix}0&\rho\\-\rho^{-1}&0\end{pmatrix}
 \begin{pmatrix}0&1\\-1&\xi\end{pmatrix}
 \begin{pmatrix}0&-\rho\\\rho^{-1}&0\end{pmatrix}
 =
 \begin{pmatrix}\xi&\rho^2\\-\rho^{-2}&0\end{pmatrix}.
\]
Since here every number is a square, $\rho^2$ may be replaced by $\rho$, and thus the substitution
\[
 \begin{pmatrix}\xi&\rho\\-\rho^{-1}&0\end{pmatrix}
\]
belongs to $G$ for every non-zero $\rho$, and likewise the opposite substitution
\[
 \begin{pmatrix}0&-\rho\\\rho^{-1}&\xi\end{pmatrix}.
\]
For two non-zero numbers $\rho,\sigma$ one has in $G$
\begin{equation}
 \begin{pmatrix}0&\rho^{-1}\sigma^{-1}\\-\rho\sigma&\xi\end{pmatrix}
 \begin{pmatrix}\xi&\sigma^{-1}\\-\sigma&0\end{pmatrix}
 \begin{pmatrix}0&\rho\\-\rho^{-1}&\xi\end{pmatrix}
 =
 \begin{pmatrix}0&1\\-1&\xi\rho(1+\sigma+\sigma\rho)\end{pmatrix}.
 \tag{9}
\end{equation}
If $n>1$, choose $\rho$ and $1+\rho$ non-zero and put $\sigma=-1:(1+\rho)$; then (9) becomes $A_0$. The case $n=1,p=2$ is the first exception.

Fourth, except when $n=1,p=3$, it contains every substitution
\begin{equation}
 \begin{pmatrix}1&0\\ \xi&1\end{pmatrix}. \tag{10}
\end{equation}
The decisive calculation is
\begin{equation}
 \begin{pmatrix}1&0\\ l&1\end{pmatrix}
 \begin{pmatrix}0&\rho\\-\rho^{-1}&0\end{pmatrix}
 \begin{pmatrix}1&0\\ l&1\end{pmatrix}
 \begin{pmatrix}0&1\\-1&0\end{pmatrix}
 =
 \begin{pmatrix}1&0\\ l(\rho^2-1)&1\end{pmatrix}. \tag{11}
\end{equation}
If $\rho^2-1$ can be made non-zero, every value of $\xi$ is obtained.

The remaining exceptional cases are exactly $n=1,p=2$ and $n=1,p=3$. In the latter case the group $E$ has degree $12$ and has the normal divisor of degree $4$
\[
 G=\left(
 \begin{pmatrix}1&0\\0&1\end{pmatrix},
 \begin{pmatrix}0&1\\-1&0\end{pmatrix},
 \begin{pmatrix}-1&-1\\1&0\end{pmatrix},
 \begin{pmatrix}0&-1\\1&1\end{pmatrix}
 \right).
\]
Apart from these two cases the group $E$ is simple. The first simple case listed is
\[
 n=1,\quad p=5,\quad g=60.
\]




\section*{\S\,67. Congruence fields of the second degree.}

In every congruence field $\Cfield_{n,p}$ for the modulus $p$ of degree $n$, the congruence field of the first degree $\Cfield_{1,p}$, consisting of the rational numbers taken modulo $p$, is contained as a divisor. Hence the group of linear substitutions $E_{n,p}$ also contains as a divisor the group consisting of all substitutions of the form
\begin{equation}
 A=\mat{a}{b}{c}{d},\qquad ad-bc=1,
 \tag{1}
\end{equation}
where $a,b,c,d$ are rational integers taken modulo $p$. This group, which occurs in the theory of elliptic functions as the group of modular equations, is the proper object of the present considerations. But for the investigation of this group, and especially for the search for its divisors, it is most advantageous not to consider it independently, but as a divisor of a group $E_{2,p}$.

Let $N$ be a fixed quadratic non-residue of $p$. Then, as noted above, $t^2-N$ is irreducible modulo $p$, and we obtain a congruence field of the second degree by introducing a Galois imaginary number by the equation
\begin{equation}
 \varepsilon^2=N.
 \tag{2}
\end{equation}
For clarity in what follows, lower-case Latin letters denote rational integers modulo $p$, which we here also call real numbers, while Greek letters denote numbers of the field $\Cfield_{2,p}$.

This field $\Cfield_{2,p}$ has the important property that all real numbers are squares in it. If $a$ is a quadratic residue of $p$, then $a=x^2$; if $b$ is a quadratic non-residue, then $bN$ is a quadratic residue, so that $x^2=bN$ can be satisfied, and hence $b=(x\varepsilon^{-1})^2$.

A number of the form $a\varepsilon$ is called purely imaginary, and two numbers $a+b\varepsilon$ and $a-b\varepsilon$ are called conjugate imaginary numbers. The transfer of these names from the ordinary theory of complex numbers to the numbers of $\Cfield_{2,p}$ is justified by the agreement of the rules of calculation and cannot cause confusion, since ordinary complex numbers do not occur in these considerations.

Since for a quadratic non-residue $N$ one has
\[
 N^{\frac12(p-1)}=-1,
\]
it follows from (2) that
\begin{equation}
 \varepsilon^p=-\varepsilon .
 \tag{3}
\end{equation}
Applying the binomial theorem to the $p$th power of
\[
 \alpha=a+b\varepsilon,
\]
and observing that all intermediate binomial coefficients are divisible by $p$, while $a^p=a$ and $b^p=b$, one obtains
\begin{equation}
 \alpha=a+b\varepsilon,
 \qquad \alpha^p=a-b\varepsilon,
 \tag{4}
\end{equation}
and therefore, by multiplication,
\begin{equation}
 \alpha^{p+1}=a^2-Nb^2.
 \tag{5}
\end{equation}
Thus $\alpha^{p+1}$ is always real.

If $\gamma$ is a primitive root of the field $\Cfield_{2,p}$, then $\gamma^r$ is real if and only if $r$ is divisible by $p+1$, and $\gamma^{p+1}$ is a primitive root of the prime number $p$ in the ordinary sense. The necessary and sufficient condition that a number $\alpha$ in $\Cfield_{2,p}$ be real is
\[
 \alpha^p=\alpha .
\]
Every real number can be represented as a $(p+1)$st power of a number in $\Cfield_{2,p}$. The notion of roots of unity also passes over to the numbers of $\Cfield_{2,p}$. If $m$ divides $p^2-1$ and
\[
 \rho=\gamma^{(p^2-1)/m},
\]
then $\rho$ is a primitive $m$th root of unity in $\Cfield_{2,p}$; the other $m$th roots of unity are its powers, and the primitive ones are precisely those powers $\rho^s$ with $s$ prime to $m$. Every non-zero number of $\Cfield_{2,p}$ is a root of unity of degree $p^2-1$.

\section*{\S\,68. The real linear congruence group $L_p$.}

We now pass, equipped with these aids, to the real congruence group $L_p$, consisting of all substitutions
\[
 A=\mat{a}{b}{c}{d}
\]
in which $a,b,c,d$ are real integers modulo $p$ and satisfy $ad-bc=1$. In the earlier notation this group is $E_{1,p}$, and, excluding $p=2$, its order is
\[
 \frac12 p(p^2-1).
\]

If $A,U$ are two elements of a group, the element $U^{-1}AU=A'$ is the transform of $A$ by $U$. The powers of $A'$ are the transforms of the corresponding powers of $A$; transformed elements therefore have the same order.

Let $A$ be any element of $L_p$ and let
\[
 U=\mat{\lambda}{\mu}{\nu}{\rho}
\]
be an element of $E_{2,p}$. We first try to transform $A$ into the normal form
\begin{equation}
 S=\mat{\delta}{0}{0}{\delta^{-1}} .
 \tag{1}
\end{equation}
From
\[
 \mat{a}{b}{c}{d}\mat{\lambda}{\mu}{\nu}{\rho}
 =\mat{\lambda}{\mu}{\nu}{\rho}\mat{\delta}{0}{0}{\delta^{-1}}
\]
one obtains
\begin{equation}
\begin{aligned}
 (a-\delta)\lambda+b\nu&=0,&
 (a-\delta^{-1})\mu+b\rho&=0,\\
 c\lambda+(d-\delta)\nu&=0,&
 c\mu+(d-\delta^{-1})\rho&=0.
\end{aligned}
\tag{2}
\end{equation}
Thus $\delta$ and $\delta^{-1}$ must be the roots of
\begin{equation}
 \delta^2-\delta(a+d)+1=0.
 \tag{3}
\end{equation}
Once they are known, (2) gives the ratios $\lambda:\nu$ and $\mu:\rho$, and the quantities may then be chosen so that $\lambda\rho-\mu\nu=1$.

If $b=c=0$, then $A$ already has the form (1). If, say, $b\ne0$, then (2) gives
\begin{equation}
 \frac{\nu}{\lambda}=-\frac{a-\delta}{b},\qquad
 \frac{\rho}{\mu}=-\frac{a-\delta^{-1}}{b}.
 \tag{4}
\end{equation}
The determinant condition can fail only when $\delta=\delta^{-1}$, that is $\delta=\pm1$, equivalently $a+d=\pm2$. But (3) is always solvable in $\Cfield_{2,p}$, since every real number is a square there:
\begin{equation}
 \delta=\frac{a+d}{2}\pm \sqrt{\left(\frac{a+d}{2}\right)^2-1}.
 \tag{5}
\end{equation}
Consequently $A$ is transformable to $S$ whenever $a+d\ne\pm2$; $\delta$ is real or imaginary according as $(a+d)^2-4$ is a quadratic residue or non-residue of $p$.

In the remaining case $a+d=\pm2$ the form $S$ cannot be produced. Instead $A$ can be brought into
\begin{equation}
 T=\mat{1}{t}{0}{1}
 \tag{6}
\end{equation}
with real $t$. Taking $a+d=2$ first, one has
\begin{equation}
 a-1=-(d-1),\qquad (a-1)(d-1)=bc.
 \tag{7}
\end{equation}
The cases $b=0$ or $c=0$ are immediate, and the case $bc\ne0$ is obtained by the explicit transformation given by (7). Thus every substitution with $a+d=\pm2$ is transformable to $T$.

The orders of the elements follow. There are three cases.

I. For $a+d=\pm2$, the normal form is $T$, and
\[
 T^r=\mat{1}{rt}{0}{1}.
\]
The order is $p$, except for the identity. There are $p^2$ such substitutions, including the identity.

II. If $(a+d)^2-4$ is a quadratic residue, the substitution is transformable to $S$ with real $\delta$. With $\gamma$ a primitive root of $\Cfield_{2,p}$,
\begin{equation}
 \delta=\gamma^{r(p+1)},\qquad
 a+d=\gamma^{r(p+1)}+\gamma^{-r(p+1)},
 \tag{8}
\end{equation}
and $r=1,\ldots,\frac12(p-3)$ gives all values apart from $\pm2$. Such elements have order $\frac12(p-1)$ or a divisor thereof, and their number is
\[
 \frac12p(p-3)(p+1).
\]

III. If $(a+d)^2-4$ is a quadratic non-residue, then $\delta$ is not real but is conjugate to its reciprocal. Thus $\delta^p=\delta^{-1}$ and $\delta^{p+1}=1$, so the order is $\frac12(p+1)$ or a divisor thereof. One writes
\begin{equation}
 \delta=\gamma^{r(p-1)},\qquad
 a+d=\gamma^{r(p-1)}+\gamma^{-r(p-1)},
 \tag{9}
\end{equation}
with $r=1,\ldots,\frac12(p-1)$. This gives
\[
 \frac12p(p-1)^2
\]
elements of the third kind. The three enumerations add up to the order of $L_p$.

\section*{\S\,69. Imaginary form of the group $L_p$.}

The substitutions of the first two kinds have their normal forms $T$ and $S$ inside $L_p$, and they may be transformed into those normal forms by real substitutions. This is not true for substitutions of the third kind.

One can, however, transform the whole group $L_p$ into another group $A_p$ such that $A_p$ is an isomorphic divisor of $E_{2,p}$, the normal forms of transformations of the third kind lie in $A_p$, and every substitution of the third kind in $A_p$ can be transformed into its normal form by substitutions of $A_p$ itself.

For that purpose consider
\[
 S=\mat{\delta}{0}{0}{\delta^{-1}},
\]
where $\delta$ and $\delta^{-1}$ are conjugate imaginary, and choose
\begin{equation}
 R=\frac{1}{\sqrt2}\mat{1}{1}{\varepsilon^{-1}}{-\varepsilon^{-1}}.
 \tag{1}
\end{equation}
With this choice $RSR^{-1}$ is real. If
\[
 A=\mat{a}{b}{c}{d}\in L_p,
\]
then
\begin{equation}
 R^{-1}AR=
 \mat{\alpha}{\beta}{-N\beta'}{\alpha'},
 \tag{2}
\end{equation}
where $\alpha,\alpha'$ and $\beta,\beta'$ are pairs of conjugate imaginary numbers; conversely every substitution of this form satisfying the determinant condition comes from a real substitution. These substitutions form the group $A_p$.

In $A_p$ a substitution of the third kind is transformable to the normal form. If
\begin{equation}
 \alpha\alpha'+N\beta\beta'=1,
 \tag{3}
\end{equation}
then the transformation equations reduce to
\begin{equation}
 (\delta-\alpha)(\delta-\alpha')+N\beta\beta'=0,
 \tag{4}
\end{equation}
together with the corresponding ratio equations; these can be satisfied in conjugate pairs.

Let $\gamma$ be a primitive root of $\Cfield_{2,p}$. For the second and third kinds one may use the normal forms
\begin{equation}
 S_1=\mat{\gamma^{p+1}}{0}{0}{\gamma^{-(p+1)}},
 \qquad
 S'_1=\mat{\gamma^{p-1}}{0}{0}{\gamma^{-(p-1)}}.
 \tag{5}
\end{equation}
Together with the form $T$ for the first kind, this shows that in $L_p$ (or in $A_p$) the substitutions of the first, second and third kinds, including the identity, can be arranged in cycles of
\[
 p,
 \qquad \frac12(p-1),
 \qquad \frac12(p+1)
\]
elements. Two such cycles have no common element except the identity. The numbers of cycles of the three kinds are
\begin{equation}
 p+1,
 \qquad \frac12p(p+1),
 \qquad \frac12p(p-1).
 \tag{6}
\end{equation}

It is also useful to consider a third group in $E_{2,p}$. Choose $v$ with $N=v^{p+1}$. To a substitution of $A_p$ let correspond
\begin{equation}
 B=\mat{\alpha}{\beta}{-\beta'}{\alpha'},
 \qquad \alpha\alpha'+\beta\beta'=1.
 \tag{7}
\end{equation}
The substitutions $B$ form a group $F_p$ isomorphic with $A_p$ and $L_p$. To pass from $B$ to the corresponding real substitution one first reconstructs the corresponding substitution in $A_p$ and then uses $R$. But $F_p$ itself cannot be transformed into $L_p$ or $A_p$ inside the field $\Cfield_{2,p}$.

\section*{\S\,70. Divisors of $L_p$ whose order is divisible by $p$.}

The determination of all divisors of $L_p$ is a problem of the greatest interest. The cyclic subgroups have already been found; their orders are either $p$ or divisors of $\frac12(p-1)$ or $\frac12(p+1)$.

If the order of a divisor $G$ of $L_p$ is divisible by $p$, then $G$ contains a cyclic group of order $p$, and we may transform $G$ so that this cyclic group is
\begin{equation}
 \mat{1}{t}{0}{1},
 \qquad t=0,1,\ldots,p-1.
 \tag{1}
\end{equation}
If $G$ contains in addition a substitution $\mat{a}{b}{c}{d}$ with $c\ne0$, then by conjugating and multiplying with the substitutions (1) one obtains the standard generators of $L_p$; hence $G=L_p$.

Thus a proper divisor with order divisible by $p$ can contain only substitutions of the form
\begin{equation}
 \mat{a}{b}{0}{a^{-1}}.
 \tag{2}
\end{equation}
All such substitutions form a group of order $\frac12p(p-1)$, hence a divisor of index $p+1$. Further divisors are obtained by letting $a$ run only through the powers $a^s$, where $s$ divides $\frac12(p-1)$; their order is
\begin{equation}
 \frac{\frac12p(p-1)}{s}.
 \tag{3}
\end{equation}

As a group of fractional linear substitutions, (2) gives
\[
 \eta=a^2\xi+ab,
\]
so that it fixes the symbol $\infty$ and gives a permutation group on only $p$ letters. Its conjugate subgroups fix the other $p+1$ letters; hence there are $p+1$ such groups.

\section*{\S\,71. Divisors of $L_p$ whose order is not divisible by $p$.}

To find the remaining divisors, whose orders are not divisible by $p$, one may follow step by step the same path used in the determination of the polyhedral groups. We regard $L_p$ as a group of fractional linear substitutions and let the variable $\xi$ range over all $p^2+1$ values of $\Cfield_{2,p}$, including $\infty$. Then
\begin{equation}
 \xi=\frac{a\xi+b}{c\xi+d}
 \tag{1}
\end{equation}
always has two roots; these are the poles of the substitution.

If $G$ has order $n$ not divisible by $p$ and precisely $\nu$ substitutions of $G$ have a given pole $\alpha$, then $\alpha$ is called a $\nu$-fold pole. The substitutions with a fixed $\nu$-fold pole form a cyclic group and can be put in the form
\begin{equation}
 \Theta^{(\mu)}\xi=\gamma^{\mu(p^2-1)/\nu}\xi,
 \qquad \mu=0,1,
 \ldots,\nu-1.
 \tag{2}
\end{equation}
Both poles of a substitution are equally multiple. The poles of the group divide into conjugate systems, and the same counting as for the polyhedral groups gives
\begin{equation}
 2n-2=a(\nu-1)+a'(\nu'-1)+a''(\nu''-1)+\cdots.
 \tag{3}
\end{equation}
Consequently only the same solutions occur as for the cyclic, dihedral, tetrahedral, octahedral and icosahedral groups.

The corresponding divisors of the congruence groups are obtained by substituting the roots of unity of $\Cfield_{2,p}$ into the formulae for the polyhedral groups. Only roots of unity whose degrees divide $p^2-1$ can appear.

I. Cyclic groups:
\begin{equation}
 C_n=\mat{\gamma^{\mu(p^2-1)/(2n)}}{0}{0}{\gamma^{-\mu(p^2-1)/(2n)}},
 \qquad \mu=0,1,\ldots,n-1.
 \tag{4}
\end{equation}
This is real when $n$ divides $\frac12(p-1)$ and lies in $F_p$ when $n$ divides $\frac12(p+1)$.

II. Dihedral groups:
\begin{equation}
 D_m:\quad
 \mat{\gamma^{\mu(p^2-1)/(2m)}}{0}{0}{\gamma^{-\mu(p^2-1)/(2m)}},
 \qquad
 \mat{0}{1}{-1}{0}.
 \tag{5}
\end{equation}
They lie in $L_p$ if $p\equiv1\pmod {2m}$ and in $F_p$ if $p\equiv-1\pmod {2m}$.

III. Tetrahedral groups. Since an eighth root of unity always exists in $\Cfield_{2,p}$, the formulae of the preceding chapter give three generators of a tetrahedral group. Written schematically, they are
\begin{equation}
 \Theta=\mat{\varepsilon}{0}{0}{\varepsilon^{-1}},
 \qquad
 \Psi=\mat{0}{\varepsilon}{-\varepsilon^{-1}}{0},
 \qquad
 X=\frac{1}{\sqrt{-2}}\mat{\varepsilon-\varepsilon^{-1}}{\varepsilon+\varepsilon^{-1}}{\varepsilon+\varepsilon^{-1}}{-\varepsilon+\varepsilon^{-1}}.
 \tag{6}
\end{equation}
According to the residue class of $p$ modulo $8$, the group is either already real or is brought into $L_p$ from $F_p$; in all cases $L_p$ contains a tetrahedral group.

IV. Octahedral groups. Such a group can occur only when $p\equiv\pm1\pmod8$. For $p=7$ one obtains a real octahedral group of index $7$. A convenient set of generators is
\begin{equation}
 \Theta=\mat{2}{0}{0}{4},
 \qquad
 \Psi=\mat{0}{3}{3}{0},
 \qquad
 X=\mat{1}{1}{3}{0}\pmod7.
 \tag{7}
\end{equation}

V. Icosahedral groups. For $p=5$, the whole group $L_p$ is an icosahedral group. For other $p$ an icosahedral group can occur only when $p\equiv\pm1\pmod5$. With a primitive fifth root $q$, the usual formulae give generators of the form
\begin{equation}
 \Theta=\mat{q}{0}{0}{q^{-1}},
 \qquad
 \Psi=\mat{0}{1}{-1}{0},
 \qquad
 X=\mat{q-q^{-1}}{q^2-q^{-2}}{q^2-q^{-2}}{-q+q^{-1}},
 \tag{8}
\end{equation}
up to the normalizing factors already fixed for the icosahedral group. For $p=11$ this gives a real icosahedral subgroup of index $11$.

Thus the possible kinds of divisors of $L_p$ are determined, and all of them have been shown actually to occur. Conjugation inside $L_p$ gives conjugate divisors; imaginary conjugation inside $E_{2,p}$ can give real divisors which are not conjugate inside $L_p$. For $p=7$ and $p=11$ the second class is obtained by replacing the list of quadratic residues in the displayed groups by the corresponding list of non-residues.

\section*{\S\,72. Constitution of the group $L_7$ of order 168.}

Among the simple groups found here, the next after the icosahedral group, which reappears as $L_5$, is the group $L_7$ of order $168$. It contains an octahedral group as a divisor. Since this remarkable group will occur again later, we look more closely at its structure.

In \S\,57 the octahedral group was represented by three elements $\chi,\omega,\Theta$ in the form
\begin{equation}
 \chi^{\lambda}\omega^{\mu}\Theta^{\nu},
 \qquad \lambda=0,1,2,
 \quad \mu=0,1,
 \quad \nu=0,1,2,3.
 \tag{1}
\end{equation}
They satisfy
\begin{equation}
 \omega\chi=\chi^2\omega,
 \quad \Theta\omega=\omega\Theta^3,
 \quad \Theta\chi=\chi^2\omega\Theta^2,
 \quad \Theta^2\chi=\chi\omega\Theta^2.
 \tag{2}
\end{equation}
In particular,
\begin{equation}
 \Theta\chi=\chi^2\omega\Theta^2,
 \quad \Theta^2\chi=\chi\omega\Theta^2,
 \quad \Theta^3\chi=\chi^2\Theta.
 \tag{3}
\end{equation}

For calculation it is convenient to transform the octahedral group modulo $7$ so that the divisor of order $21$ takes the simplest form. The divisor of index $8$ consists of the substitutions
\begin{equation}
 \mat{a}{b}{0}{a^{-1}}\pmod7.
 \tag{4}
\end{equation}
Transforming the earlier substitutions by $\mat{1}{1}{-1}{0}$ gives
\begin{equation}
 \chi=\mat{2}{3}{0}{-3},
 \qquad
 \omega=\mat{1}{-3}{3}{-1},
 \qquad
 \Theta=\mat{1}{3}{-2}{2}.
 \tag{5}
\end{equation}
The octahedral group is then the following table of twenty-four substitutions:
\begin{equation}
\begin{array}{rrrr}
\mat{1}{0}{0}{1},&\mat{1}{3}{-2}{2},&\mat{2}{2}{1}{-2},&\mat{2}{-3}{2}{1},\\[0.7em]
\mat{1}{-3}{3}{-1},&\mat{0}{3}{2}{0},&\mat{-1}{1}{-2}{1},&\mat{3}{1}{-3}{-3},\\[0.7em]
\mat{2}{3}{0}{-3},&\mat{3}{-2}{-1}{1},&\mat{0}{2}{3}{1},&\mat{3}{-3}{1}{-3},\\[0.7em]
\mat{3}{2}{2}{-3},&\mat{1}{1}{-1}{0},&\mat{1}{2}{1}{3},&\mat{-3}{0}{2}{2},\\[0.7em]
\mat{3}{3}{0}{-2},&\mat{-3}{1}{-3}{3},&\mat{2}{0}{-2}{-3},&\mat{-2}{1}{3}{-2},\\[0.7em]
\mat{-2}{2}{1}{2},&\mat{-1}{2}{3}{0},&\mat{2}{1}{3}{2},&\mat{0}{1}{-1}{-1}.
\end{array}
\tag{6}
\end{equation}

To obtain the whole group $L_7$, one adds an element of order seven,
\begin{equation}
 \tau=\mat{1}{1}{0}{1}.
 \tag{7}
\end{equation}
The whole group $L_7$ is then represented by
\begin{equation}
 \tau^{\rho}\chi^{\lambda}\omega^{\mu}\Theta^{\nu},
 \quad \rho=0,1,\ldots,6,
 \quad \lambda=0,1,2,
 \quad \mu=0,1,
 \quad \nu=0,1,2,3.
 \tag{8}
\end{equation}
These elements are all distinct, since no non-trivial power of $\tau$ belongs to the octahedral subgroup.

It remains to give the composition rules. From
\begin{equation}
 \tau^{\rho}=\mat{1}{\rho}{0}{1}
 \tag{9}
\end{equation}
one obtains
\begin{equation}
 \chi\tau^{\rho}=\tau^{4\rho}\chi,
 \qquad
 \chi^2\tau^{\rho}=\tau^{2\rho}\chi^2,
 \tag{10}
\end{equation}
so that all these relations follow from
\begin{equation}
 \chi\tau=\tau^4\chi.
 \tag{11}
\end{equation}
The remaining products, computed from the table, are
\begin{equation}
\begin{array}{ll}
 \omega\tau=\tau^2\chi^2\omega\Theta^2,&
 \Theta\tau=\tau^3\omega\Theta,\\
 \omega\tau^2=\tau\chi^2\Theta^3,&
 \Theta\tau^2=\tau\chi\omega\Theta^3,\\
 \omega\tau^3=\tau^5\chi\Theta^2,&
 \Theta\tau^3=\tau\chi\Theta,\\
 \omega\tau^4=\tau^4\chi^2\Theta,&
 \Theta\tau^4=\tau^6\omega\Theta^2,\\
 \omega\tau^5=\tau^3\chi^2\omega\Theta,&
 \Theta\tau^5=\tau^2\Theta^3,\\
 \omega\tau^6=\tau^6\omega\Theta^3,&
 \Theta\tau^6=\tau^4\chi^2\Theta^2.
\end{array}
\tag{12}
\end{equation}
These twelve relations follow from four fundamental ones, for example
\begin{equation}
 \omega\tau=\tau^2\chi^2\omega\Theta^2,
 \quad
 \omega\tau^3=\tau^5\chi\Theta^2,
 \quad
 \omega\tau^4=\tau^4\chi^2\Theta,
 \quad
 \Theta\tau=\tau^3\omega\Theta.
 \tag{13}
\end{equation}
From them, together with (3) and (10), all others follow. The elements $\omega$ and $\tau$ alone generate the group, since
\begin{equation}
 \Theta^3=\omega\tau\omega\tau^6,
 \qquad
 \chi=\tau^2\Theta\tau^3\omega.
 \tag{14}
\end{equation}

Thus, if four elements $\tau,\chi,\omega,\Theta$ of orders $7,3,2,4$ satisfy the relations
\begin{equation}
\begin{gathered}
 \omega\chi=\chi^2\omega,
 \quad \Theta\omega=\omega\Theta^3,
 \quad \Theta\chi=\chi^2\omega\Theta^2,\\
 \Theta^2\chi=\chi\omega\Theta^2,
 \quad \omega\tau=\tau^2\chi^2\omega\Theta^2,
 \quad \omega\tau^3=\tau^5\chi\Theta^2,\\
 \omega\tau^4=\tau^4\chi^2\Theta,
 \quad \Theta\tau=\tau^3\omega\Theta,
 \quad \chi\tau=\tau^4\chi,
\end{gathered}
\tag{15}
\end{equation}
then the 168 elements
\begin{equation}
 \sigma=\tau^\rho\chi^\lambda\omega^\mu\Theta^\nu,
 \tag{16}
\end{equation}
with $\rho,\lambda,\mu,\nu$ running through complete residue systems modulo $7,3,2,4$, form a simple group of order $168$. Among its divisors are the octahedral subgroup $\chi^\lambda\omega^\mu\Theta^\nu$ of index $7$, the subgroup $\tau^\rho\chi^\lambda$ of order $21$ and index $8$, and the subgroup $\omega^\mu\Theta^\nu$ of order $8$. The whole group can also be written in the orders
\[
 \chi^\lambda\omega^\mu\Theta^\nu\tau^\rho,
 \qquad
 \omega^\mu\Theta^\nu\chi^\lambda\tau^\rho,
 \qquad
 \tau^\rho\omega^\mu\Theta^\nu\chi^\lambda.
\]


\clearpage
\begin{center}
{\Large Third Book.}\\[0.5em]
{\large Applications of Group Theory.}
\end{center}

\begin{center}
{\large Tenth Section.}\\[0.4em]
{\large General theory of metacyclic equations.}
\end{center}

\section*{\S\,73. The resolvents of the composition series.}

From the theorems of general group theory which were treated in the first book of this volume, important algebraic consequences follow when they are applied to the Galois group of an equation.

Let an arbitrary field $\Omega$ now be taken as rationality domain, and let $f(x)=0$ be an equation in this field, without multiple roots, and let $P$ be its Galois group. Suppose that this group $P$ has a normal divisor $Q$. We denote by $n$ the degree of $P$ and by $j$ the index of the divisor $Q$, so that $n:j$ is the degree of $Q$.

We have seen in the first volume (\S\,156) that the group of $f(x)$ is reduced to $Q$ by adjoining a root of an irreducible normal equation of degree $j$, and this auxiliary equation of degree $j$ was called a partial resolvent.

The Galois group of this partial resolvent was obtained as follows. Let $\psi$ denote a function of the roots of $f(x)$ belonging to the group $Q$, and decompose $P$ into the cosets
\[
  P=Q+Qa+Qb+Qc+\cdots ,
\]
where $a,b,c,\ldots$ are certain permutations from $P$. Then the whole coset $Qa$ sends $\psi$ into one and the same function $\psi_a$, and the Galois group of the resolvent consists of the substitutions
\[
  (\psi,\psi),\quad(\psi,\psi_a),\quad(\psi,\psi_b),\quad(\psi,\psi_c),\ldots .
\]
When two of these substitutions are composed, one has to note that
\[
  \psi\longmapsto \psi_a\longmapsto \psi_{ab},
\]
so that
\[
  (\psi,\psi_a)(\psi,\psi_b)=(\psi,\psi_{ab}).
\]
Thus these substitutions are composed in precisely the same way as the cosets in \S\,4 of this volume, and hence:

\begin{enumerate}[label=\arabic*.,leftmargin=2.6em]
\item The Galois group of the partial resolvent belonging to $Q$ is isomorphic to the group $P/Q$ complementary to $Q$.
\end{enumerate}

Now take any composition series of $P$ with its corresponding index series (\S\,6):
\begin{equation}
 P,\;P_1,\;P_2,\ldots,P_{\mu-1},\;1.
 \tag{1}
\end{equation}
By what has just been said, the group of the equation $f=0$ is reduced from $P$ to $P_1$ by adjoining a root of a normal equation of degree $j_1$. It is then reduced to $P_2$ by a root of a normal equation of degree $j_2$, and so on; finally, the equation is completely solved by a root of a normal equation of degree $j_\mu$.

The theorem on the invariance of the index series (\S\,6, I.) now casts a new light on this process of solution.

\begin{enumerate}[label=\arabic*.,leftmargin=2.6em,start=2]
\item To solve the equation $f=0$, one has successively to adjoin one root of a normal equation of each of the degrees $j_1,j_2,\ldots,j_\mu$. The degrees of these resolvents may be changed in their order, but not as a totality.
\end{enumerate}

The numbers $j_1,j_2,\ldots,j_\mu$ are therefore connected far more deeply with the nature of an equation, or more generally with the fields defined by the equation, than is, for example, the degree of the equation. For the degree may be changed in many ways by transformation, when functions of the roots are introduced as new unknowns; the numbers $j_1,j_2,\ldots,j_\mu$, however, remain preserved and are to be regarded as true invariants of the normal field defined by the equation.

Among these invariants is also the degree $n$ of the group $P$ itself, determined by the $j$'s:
\begin{equation}
 n=j_1j_2\cdots j_{\mu-1}j_\mu .
 \tag{2}
\end{equation}
For, by the meaning of the indices, $n:j_1$ is the degree of $P_1$, $n:j_1j_2$ the degree of $P_2$, and so forth; since the last of the degrees is $1$, formula (2) follows.

In general, in a composition series of $P$, each term is a normal divisor only of the immediately preceding one. But individual terms may also occur which are normal divisors of earlier preceding terms. Especially important are those terms which are normal divisors of $P$ itself, and hence also of all the terms of the composition series which precede them.

Let $P_\nu$ be a term of the series (1) which is also a normal divisor of $P$. The index of $P$ with respect to $P_\nu$ is $j_1j_2\cdots j_\nu$, and this product is the degree of the partial resolvent by which the group $P$ is reduced to $P_\nu$; denote this resolvent by $\chi(y)=0$. By theorem 1, the group of this resolvent, which is a normal equation, is obtained in the form $P/P_\nu$.

We can now easily find a composition series for this group with the corresponding index series:
\begin{equation}
 P/P_\nu,\quad P_1/P_\nu,\quad P_2/P_\nu,\ldots,\quad P_{\nu-1}/P_\nu,\quad 1,
 \qquad j_1,j_2,\ldots,j_{\nu-1},j_\nu .
 \tag{3}
\end{equation}
For by theorem 1 of \S\,6, $P_1/P_\nu$ is a normal divisor of $P/P_\nu$ of index $j_1$, and is a greatest normal divisor, since by 2, \S\,6 no normal divisor of $P/P_\nu$ can lie above $P_1/P_\nu$ if no normal divisor of $P$ lies above $P_1$. The same reasoning applies to the following terms of (3).

From this we draw one more important conclusion. We already showed in \S\,7 that, if $Q$ is any normal divisor of $P$, a composition series of $P$ can be found in which $Q$ occurs. If now $R$ is another normal divisor of $P$ and at the same time a divisor of $Q$, then the composition series of $P$ can be arranged so that both $Q$ and $R$ occur in it. Suppose such a composition series is fixed:
\begin{equation}
 P,\;Q',\;Q'',\ldots,\;Q,\ldots,\;R.
 \tag{4}
\end{equation}
From it, by (3), the composition series of the two groups $P/Q$ and $P/R$ are obtained in the forms
\begin{equation}
 \cdots,\quad P/Q,\quad Q'/Q,\quad Q''/Q,\ldots,\quad 1,
 \tag{5}
\end{equation}
\begin{equation}
 \cdots,\quad P/R,\quad Q'/R,\quad Q''/R,\ldots,\quad Q/R,\ldots .
 \tag{6}
\end{equation}
The index of the divisor $Q'/Q$ of $P/Q$ is equal to the index of the divisor $Q'$ of $P$, and hence also to the index of the divisor $Q'/R$ of $P/R$ (\S\,6, 1.); the same holds for the following terms. We therefore state the theorem:

\begin{enumerate}[label=\arabic*.,leftmargin=2.6em,start=3]
\item If $Q$ and $R$ are normal divisors of $P$, and if $R$ is a divisor of $Q$, then the index series of $P/Q$ is a part of the index series of $P/R$.
\end{enumerate}

If the indices $j_1,j_2,\ldots,j_\nu$ are all prime numbers, then the solution of the resolvent $\chi(y)=0$ can be reduced to the solution of a chain of cyclic equations of degrees $j_1,j_2,\ldots,j_\nu$. This resolvent is therefore metacyclic in the sense established in the seventeenth section of the first volume, in which metacyclic equations are identical with those otherwise called algebraically solvable.

An irreducible equation $f(x)=0$ is metacyclic when the index series of its group consists only of primes. At the place just mentioned we investigated the conditions for metacyclic equations of prime degree; we must now carry out these considerations for the general case in which the degree of the equation is arbitrarily composite.

\section*{\S\,74. Metacyclic equations.}

If $f(x)=0$ is an irreducible equation of degree $m$ and $P$ is its Galois group, then $P$ is transitive. Let a composition and index series for $P$ be
\begin{equation}
 P,\;P_1,\;P_2,\ldots,\;P_{\mu-1},\;1 .
 \tag{1}
\end{equation}
Since the last group $1$ is intransitive, an intransitive group must occur at some point in the composition series; let $P_\lambda$ be the first intransitive group in (1). Then all following groups $P_{\lambda+1},P_{\lambda+2},\ldots$, being divisors of $P_\lambda$, are also intransitive.

We recall theorems 1, 2, and 3 of \S\,158 of the first volume. Since $P_\lambda$ is a normal divisor of $P_{\lambda-1}$, it follows from those theorems that $P_{\lambda-1}$ is imprimitive. If, after the necessary adjunctions, the group of $f(x)=0$ has been reduced to $P_{\lambda-1}$, then by the further adjunction of a root of a normal equation of degree $j_\lambda$ the function $f(x)$ splits into several irreducible factors
\begin{equation}
 f(x)=f_1(x)f_2(x)\cdots f_h(x),
 \tag{2}
\end{equation}
all of equal degree, and the number $h$ of them is a divisor of $j_\lambda$. If $m$ denotes the degree of $f(x)$ and $m_1$ the degree of $f_1(x)$, then
\begin{equation}
 m=hm_1 .
 \tag{3}
\end{equation}
By theorem 1 of \S\,158, Vol. I, the equations $f_1=0,f_2=0,\ldots,f_h=0$ all have the same group, which we shall denote by $Q$. If $\alpha,\alpha_1,\alpha_2,\ldots,\alpha_{m_1-1}$ are the roots of $f_1=0$, then we obtain the group $Q$ by collecting the permutations of the $\alpha$ induced by $P_\lambda$. We first have to determine the relation between the degrees $p_\lambda$ and $q$ of the two groups $P_\lambda$ and $Q$.

Several different permutations in $P_\lambda$ may generate the same element in $Q$, that is, they may contain the same permutation of the $\alpha$. If $\pi_1,\pi_2$ are two different permutations from $P_\lambda$ which induce the same permutation among the $\alpha$, then $\pi_2\pi_1^{-1}=\pi_0$ is a permutation which leaves the $\alpha$ at rest, and thus $\pi_2=\pi_0\pi_1$. Conversely, if $\pi_0$ is any permutation from $P_\lambda$ which does not permute the roots $\alpha$, then $\pi_2=\pi_0\pi_1$ produces the same change among the $\alpha$ as $\pi_1$.

It follows that each permutation of the $\alpha$ is produced equally often by the permutations of $P_\lambda$, namely as often as the $\alpha$ remain unchanged by permutations from $P_\lambda$. If this number is denoted by $p_0$, then
\begin{equation}
 p_\lambda=p_0q,
 \tag{4}
\end{equation}
or the degree $p_\lambda$ of $P_\lambda$ is a multiple of the degree $q$ of $Q$.

In \S\,154, 7 of the first volume we proved that the degree of a transitive permutation group is always divisible by the number of letters permuted. Here $f_1(x)$ is irreducible, and therefore the group $Q$ is transitive. Thus $q$ is divisible by the degree of $f_1(x)$, that is, by $m_1$, and hence by (4) also $p_\lambda$ is divisible by $m_1$.

From the meaning of $j_\lambda$, the degree $p_{\lambda-1}$ of $P_{\lambda-1}$ is
\begin{equation}
 p_{\lambda-1}=j_\lambda p_\lambda .
 \tag{5}
\end{equation}
Now $p_\lambda$ is a multiple of $m_1$, $j_\lambda$ is a multiple of $h$, and $hm_1=m$ is the degree of $f(x)$. Hence
\begin{equation}
 p_{\lambda-1}=k m
 \tag{6}
\end{equation}
is a multiple of $m$. But from the meaning of the $j$'s [\S\,73, (2)] one also has
\[
 p_{\lambda-1}=j_\lambda j_{\lambda+1}\cdots j_\mu,
\]
and consequently
\begin{equation}
 j_\lambda j_{\lambda+1}j_{\lambda+2}\cdots j_\mu=km.
 \tag{7}
\end{equation}
This gives a very remarkable consequence.

Suppose that the irreducible function $f(x)$ becomes reducible by successive adjunction of roots of cyclic equations. Then, by \S\,176 of Vol. I, the composition series (1) may be arranged so that the indices $j_1,j_2,\ldots,j_\lambda$ are prime numbers. Since $h$ is a divisor of $j_\lambda$, it follows that $j_\lambda=h$, and by (3) $j_\lambda$ is a prime factor occurring in $m$.

If, besides $j_\lambda$, a second prime $p$ occurs in $m$, then by (7) one of the factors $j_\lambda,j_{\lambda+1},j_{\lambda+2},\ldots,j_\mu$ must be divisible by $p$, so they certainly cannot all be equal to $j_\lambda$. Hence, by theorem III, \S\,7, a composition series of $P_\lambda$,
\begin{equation}
 P_\lambda,\;P_{\lambda+1},\ldots,P_{\mu-1},\;1,
 \tag{8}
\end{equation}
can be found in such a way that among the groups $P_\lambda,P_{\lambda+1},\ldots,P_{\mu-1}$ one, say $P_\nu$, is a normal divisor of $P$. But this $P_\nu$, being a divisor of the intransitive group $P_\lambda$, is itself intransitive, and therefore $P$ itself must be imprimitive (Vol. I, \S\,158, 2). We state this in the following form:

\begin{enumerate}[label=\arabic*.,leftmargin=2.6em]
\item If the degree of an irreducible equation contains several distinct prime factors, then this equation can become reducible by successive adjunction of roots of cyclic equations only if it is imprimitive.
\end{enumerate}

Assuming that such a reduction is possible, we can say something more precise about its manner. Let $P_\nu$ be the first group of the series (8) which is a normal divisor of $P$. Then, by the theorem III, \S\,7 just cited, with a suitable arrangement of the composition series,
\begin{equation}
 j_\lambda=j_{\lambda+1}=\cdots=j_{\nu+1},
 \tag{9}
\end{equation}
so that all are equal to the same prime. The resolvent $\chi=0$ by which $P$ is reduced to $P_\nu$ consequently has a metacyclic group.

Let $A,B,\ldots,S$ be the systems of intransitivity of the group $P_\nu$, and let their number be $s$. Put
\begin{equation}
 m=rs.
 \tag{10}
\end{equation}
Then, by adjoining a root of $\chi=0$, the function $f(x)$ splits into $s$ factors of degree $r$. The $A,B,\ldots,S$ are systems of imprimitivity of $P$ (Vol. I, \S\,158, 2).

Let $Q$ denote the totality of all permutations of $P$ which leave the individual systems $A,B,\ldots,S$ in their places and only interchange the elements within the systems. Then, as we saw in \S\,158 of the first volume, $Q$ is a normal divisor of $P$. By adjoining the roots of an auxiliary equation of degree $s$, $\varphi(y)=0$, the function $f(x)$ splits into $s$ factors of degree $r$, whose roots belong to the individual systems $A,B,\ldots,S$:
\[
 f(x)=f(x,y_\alpha)f(x,y_\beta)f(x,y_\gamma)\cdots .
\]
Since the auxiliary equation $\varphi(y)=0$ reduces the group $P$ to $Q$, the group of $\varphi(y)=0$ is isomorphic to $P/Q$. On the other hand the intransitive group $P_\nu$ is also a divisor of $Q$, since the systems $A,B,\ldots,S$ are not displaced by $P_\nu$. We may therefore apply the theorem \S\,73, 3 to the three groups $P,Q,P_\nu$. Since by hypothesis the index series of $P/P_\nu$ consists only of prime numbers, the same is true of $P/Q$. This group is metacyclic, and hence the auxiliary equation $\varphi(y)=0$ is metacyclic. Thus:

\begin{enumerate}[label=\arabic*.,leftmargin=2.6em,start=2]
\item If an irreducible equation whose degree contains more than one prime factor splits into factors by successive adjunction of radicals, then a splitting into $s$ factors of degree $r$ is brought about by adjoining the roots of a metacyclic equation of degree $s$.
\end{enumerate}

As a special case, this theorem contains Abel's theorem that an irreducible equation whose degree $m$ is not a power of a prime can be solvable by radicals only when it splits into factors after adjoining the roots of a solvable equation of lower degree, whose degree is a divisor of $m$.\footnote{Abel gives the theorem without proof in the memoir \emph{Sur la résolution algébrique des équations}. Oeuvres complètes 1881, Vol. II, p. 217. Compare also the memoir of Galois in Liouville's Journal, Vol. 11.} The question of solving an equation by radicals, or even merely reducing an equation by radicals, is thereby extraordinarily simplified. In fact, further investigation may be restricted to equations whose degree is a prime or a prime power, since all other cases are reduced to these by repeated application of theorem 2.

For example, the question of all irreducible equations of degree $6$ solvable by radicals permits a completely trivial answer: to obtain all metacyclic equations of degree $6$ in any field $\Omega$, adjoin to $\Omega$ a square root and form all cubic equations in the extended field, or adjoin the root of a cubic equation and form all quadratic equations in the extended field.

As an example of such an equation we mention the one treated by Hesse, on which the circular sections of a quadric surface not referred to its principal axes depend. This problem is solved by square roots, after the principal axes have first been determined by solving a cubic equation.\footnote{Hesse, \emph{Ueber die Auflösung derjenigen Gleichungen 6ten Grades etc.} Crelle's Journal, Vol. 41 (1851).}

\section*{\S\,75. Metacyclic equations whose degree is a prime power.}

According to the last theorems, the interest of the further investigations on metacyclic equations is concentrated chiefly on the case in which the degree of the equation is a power $p^k$ of a prime $p$. We assume the equation to be irreducible and restrict ourselves to primitive equations. For the solution of the imprimitive ones is reduced to the successive solution of two, or several, primitive equations whose degrees are likewise powers of $p$ and which, if the original equation is metacyclic, must themselves also be metacyclic.

We now assume, therefore, that $P$ is the group of an irreducible primitive metacyclic equation $f(x)=0$ of degree $p^k$. By Vol. I, \S\,158, 2, not only $P$ itself, but all its normal divisors except the identity group, must still be transitive. We now apply theorem IV of \S\,8, proved there in general: $P$ has a normal divisor $Q$ consisting of mutually commutative elements. If several such divisors are present, we take for $Q$ one whose degree is as small as possible but still greater than $1$. This group $Q$ is therefore still transitive. It can have no proper divisor distinct from the identity which is at the same time a normal divisor of $P$, for otherwise that divisor would take the place of $Q$.

When, by the corresponding adjunctions, the group of our equation is reduced from $P$ to $Q$, the equation $f(x)=0$ has become an Abelian equation in the extended rationality domain. Since it has remained irreducible, by Vol. I, \S\,162 the degree of the equation is equal to the degree of the group; that is, $Q$ is an Abelian group of degree $p^k$. Thus the degrees of all elements of $Q$ are powers of $p$. We shall show that, apart from the identity element, only elements of degree $p$ itself occur in $Q$.

Assume that $p^\lambda$ is the highest degree occurring among the elements of $Q$, and that $\lambda>1$. Then all elements of $Q$ whose degree divides $p^{\lambda-1}$ form a divisor $Q'$ of $Q$ which is different both from $Q$ itself and from the identity group. But this divisor $Q'$ must be a normal divisor of $P$, because if $\pi$ is contained in $Q$ and $\gamma$ in $P$, then $\gamma^{-1}\pi\gamma$, which has the same degree as $\pi$, is likewise contained in $Q$ and therefore belongs to $Q'$ whenever $\pi$ belongs to $Q'$. Since, as observed above, $Q$ can have no proper divisor larger than the identity group and at the same time normal in $P$, we must have $\lambda=1$.

\section*{\S\,76. Representation of the Abelian group $Q$.}

The considerations of the preceding paragraph have shown that in a metacyclic group $P$, when it can occur as the Galois group of an irreducible primitive equation $f(x)=0$ of degree $n=p^k$, an Abelian group $Q$ must be contained as a normal divisor, of degree $p^k$, and having, besides the identity, only elements of degree $p$.

By \S\,9 this group $Q$ can be represented by a basis consisting of $k$ elements $A_1,A_2,\ldots,A_k$ of degree $p$, so that every element of $Q$ receives the form
\begin{equation}
 A_1^{z_1}A_2^{z_2}\cdots A_k^{z_k}.
 \tag{1}
\end{equation}
Here $z_1,z_2,\ldots,z_k$, independently of one another, each run through a complete residue system modulo $p$.

To find the corresponding permutation group, we start from an arbitrary root $x$ of $f(x)$, denote by
\begin{equation}
 [z_1,z_2,\ldots,z_k]
 \tag{2}
\end{equation}
the root into which $x$ is carried by the permutation (1), and stipulate that this symbol shall not change its meaning when the numbers $z_1,z_2,\ldots,z_k$ are changed arbitrarily by multiples of $p$. The root $x$ from which we started then receives the symbol $[0,0,\ldots,0]$, and by (2) every root of $f(x)$ is represented once and only once.

If $x$ is carried by a permutation $A'$ into $x'$, and $x'$ by $A''$ into $x''$, then $x$ is carried by $A'A''$ into $x''$. Hence $[z_1,z_2,\ldots,z_k]$ is carried by the permutation
\begin{equation}
 A=A_1^{\alpha_1}A_2^{\alpha_2}\cdots A_k^{\alpha_k}
 \tag{3}
\end{equation}
into
\[
 [z_1+\alpha_1,\;z_2+\alpha_2,\ldots,z_k+\alpha_k].
\]
Consequently the permutation $A$ may also be denoted by
\begin{equation}
 A=\binom{z}{z+\alpha},
 \tag{4}
\end{equation}
and the whole group $Q$ is obtained by allowing $\alpha_1,\alpha_2,\ldots,\alpha_k$ each to run through a complete residue system modulo $p$.

The group $Q$, whose construction has now been fixed, must be a normal divisor of $P$. From this an all-embracing form can be derived in which all permutations of $P$ must be contained. For this we first need an auxiliary theorem on the analytic representation of permutations.

\section*{\S\,77. Analytic representation of the permutations.}

\textbf{Auxiliary theorem.} If under any permutation of the quantities (2), \S\,76,
\[
 z_1\mapsto z'_1,\quad z_2\mapsto z'_2,\ldots,\quad z_k\mapsto z'_k,
\]
then one can always put
\begin{equation}
\begin{gathered}
 z'_1=\varphi_1(z_1,z_2,\ldots,z_k),\\
 z'_2=\varphi_2(z_1,z_2,\ldots,z_k),\\
 \cdots\\
 z'_k=\varphi_k(z_1,z_2,\ldots,z_k)
\end{gathered}
 \pmod p,
 \tag{1}
\end{equation}
where $\varphi_1,\varphi_2,\ldots,\varphi_k$ are integral rational functions of the variables $z_1,z_2,\ldots,z_k$ with integral coefficients, and with degree not exceeding $p-1$ in any one of the variables.

This theorem is a generalization of the corresponding theorem in \S\,180 of the first volume. There it was proved for $k=1$. Before the general proof, which rests on complete induction, we make the following observation.

Since in (1) the arguments $z_i$ have to be replaced by all combinations of $k$ rows modulo $p$, there are $p^k$ such systems. Each function $\varphi$ has $p^k$ undetermined coefficients, and hence, if the $z'_i$ are given, we obtain $kp^k$ linear congruences for the same number of unknown numbers, namely the coefficients of the $\varphi_i$. It remains only to show that these congruences are independent of one another. In doing this, each function $\varphi$ can be considered separately.

Thus put
\begin{equation}
 z'=\varphi(z_1,z_2,\ldots,z_k)\pmod p.
 \tag{2}
\end{equation}
If for each combination of the $z_i$ the corresponding value of $z'$ is known, then we have $p^k$ linear congruences for determining the coefficients of $\varphi$. It is to be proved that these congruences are always soluble, and we may assume that this possibility has already been established for functions of fewer variables.

Now order $\varphi$ according to powers of the variable $z_1$:
\begin{equation}
 z'=\chi_0+\chi_1z_1+\chi_2z_1^2+\cdots+\chi_{p-1}z_1^{p-1}\pmod p,
 \tag{3}
\end{equation}
where the coefficients $\chi_0,\chi_1,\ldots,\chi_{p-1}$ are functions of the same kind as $\varphi$, except that they depend on one variable fewer. Hold fixed any one combination of $z_2,\ldots,z_k$ and set $z_1=0,1,\ldots,p-1$. Then (3) gives a system of $p$ linear congruences whose determinant is not divisible by $p$. As in \S\,180 of the first volume, one can determine from it the values of $\chi_0,\chi_1,\ldots,\chi_{p-1}$ for this combination of $z_2,\ldots,z_k$. Thus, for each combination of the variables, the values of the function $\chi_i$ are known modulo $p$. By the induction hypothesis the coefficients of the functions $\chi_i$, which depend on only $k-1$ variables, can be determined from these values. The auxiliary theorem is thereby proved.

Consequently every permutation of the roots $[z_1,\ldots,z_k]$ can be represented in the form
\begin{equation}
 \binom{z}{\varphi(z)}
 =
 \binom{z_1,z_2,\ldots,z_k}{\varphi_1(z),\varphi_2(z),\ldots,\varphi_k(z)}
 \tag{4}
\end{equation}
or, more briefly,
\begin{equation}
 \binom{z}{\varphi(z)}.
 \tag{5}
\end{equation}
In the symbol (4) we may replace $z_1,z_2,\ldots$, without changing its meaning, by any other combination $\psi_1,\psi_2,\ldots$. If, by the auxiliary theorem, $z_i=\psi_i(z_1,z_2,\ldots)$, then (4) or (5) can also be written
\begin{equation}
 \binom{\psi_1(z_1,z_2,\ldots),\psi_2(z_1,z_2,\ldots),\ldots}
       {\varphi_1(\psi_1,\psi_2,\ldots),\varphi_2(\psi_1,\psi_2,\ldots),\ldots}
 =
 \binom{\psi(z)}{\varphi[\psi(z)]}.
 \tag{6}
\end{equation}
This leads to the composition of the permutations:
\begin{equation}
 \binom{z}{\psi(z)}\binom{z}{\varphi(z)}
 =
 \binom{z}{\varphi[\psi(z)]}.
 \tag{7}
\end{equation}

\section*{\S\,78. Representation of the metacyclic group $P$.}

Now let the group $Q$ be a normal divisor of the group $P$. That is, if $A$ is a permutation from $Q$ and $B$ a permutation from $P$, there is to be a second permutation $A'$ from $Q$ such that
\begin{equation}
 AB=BA'.
 \tag{1}
\end{equation}
Using the abbreviated notation of the preceding paragraph, put
\[
 A=\binom{z}{z+\alpha},\qquad
 A'=\binom{z}{z+\alpha'},\qquad
 B=\binom{z}{\varphi(z)}.
\]
Then
\[
 AB=\binom{z}{\varphi(z+\alpha)},\qquad
 BA'=\binom{z}{\varphi(z)+\alpha'},
\]
or
\[
 \varphi(z+\alpha)\equiv\varphi(z)+\alpha'.
\]
This congruence is only an abbreviated symbol for the system of congruences modulo $p$:
\begin{equation}
\begin{aligned}
 \varphi_1(z_1+\alpha_1,z_2+\alpha_2,\ldots)&\equiv
        \varphi_1(z_1,z_2,\ldots)+\alpha'_1,\\
 \varphi_2(z_1+\alpha_1,z_2+\alpha_2,\ldots)&\equiv
        \varphi_2(z_1,z_2,\ldots)+\alpha'_2,\\
 &\hspace{4em}\cdots .
\end{aligned}
 \tag{2}
\end{equation}
The system of numbers $\alpha_1,\alpha_2,\ldots$ may here be prescribed arbitrarily modulo $p$; the numbers $\alpha'_1,\alpha'_2,\ldots$ are thereby determined. In the first of the formulae (2), set in turn one of $\alpha_1,\alpha_2,\ldots$ equal to $1$ and all the others equal to $0$. One then obtains
\begin{equation}
\begin{aligned}
 \varphi_1(z_1+1,z_2,\ldots)&\equiv\varphi_1(z_1,z_2,\ldots)+\alpha_{1,1},\\
 \varphi_1(z_1,z_2+1,\ldots)&\equiv\varphi_1(z_1,z_2,\ldots)+\alpha_{1,2},\\
 &\hspace{4em}\cdots .
\end{aligned}
 \tag{3}
\end{equation}
Applying the first of these formulae $\alpha_1$ times, the second $\alpha_2$ times, and so forth, gives
\[
\begin{aligned}
 \varphi_1(z_1+\alpha_1,z_2,\ldots)&\equiv
 \varphi_1(z_1,z_2,\ldots)+\alpha_1\alpha_{1,1},\\
 \varphi_1(z_1,z_2+\alpha_2,\ldots)&\equiv
 \varphi_1(z_1,z_2,\ldots)+\alpha_2\alpha_{1,2}.
\end{aligned}
\]
If now, in the first formula, $z_2$ is replaced by $z_2+\alpha_2$, then the second is applied, and the procedure is continued, the final result is
\begin{equation}
 \varphi_1(z_1+\alpha_1,z_2+\alpha_2,\ldots)
 \equiv
 \varphi_1(z_1,z_2,\ldots)+\alpha_1\alpha_{1,1}
      +\alpha_2\alpha_{1,2}+\cdots .
 \tag{4}
\end{equation}
Set now $z_1,z_2,\ldots=0$, and then write again $z_1,z_2,\ldots$ in place of $\alpha_1,\alpha_2,\ldots$. If $\varphi_1(0,0,\ldots)$ is denoted by $\alpha_1$, (4) gives
\begin{equation}
 \varphi_1(z_1,z_2,\ldots)
 \equiv
 \alpha_1+z_1\alpha_{1,1}+z_2\alpha_{1,2}+\cdots.
 \tag{5}
\end{equation}
Thus $\varphi_1$ must be a linear function. Proceeding in exactly the same way with all the congruences (2), we arrive at the following final result:

\begin{enumerate}[label=\Roman*.,leftmargin=2.8em]
\item All permutations of the Galois group $P$ of a primitive irreducible metacyclic equation of degree $p^k$,
\[
 \binom{z_1,z_2,\ldots,z_k}{z'_1,z'_2,\ldots,z'_k},
\]
are of the form
\begin{equation}
\begin{aligned}
 z'_1&\equiv \alpha_{1,1}z_1+\alpha_{1,2}z_2+\cdots+\alpha_{1,k}z_k+\alpha_1,\\
 z'_2&\equiv \alpha_{2,1}z_1+\alpha_{2,2}z_2+\cdots+\alpha_{2,k}z_k+\alpha_2,\\
 &\hspace{4em}\cdots\\
 z'_k&\equiv \alpha_{k,1}z_1+\alpha_{k,2}z_2+\cdots+\alpha_{k,k}z_k+\alpha_k
\end{aligned}
\pmod p .
 \tag{6}
\end{equation}
\end{enumerate}

The system of congruences (6) represents a permutation if and only if the determinant
\[
 \sum \pm\alpha_{1,1}\alpha_{2,2}\cdots\alpha_{k,k}
\]
is not divisible by $p$, for only then does each system of values $z'$ correspond conversely to a definite system of values $z$ modulo $p$. The totality of the permutations (6) is indeed a group, called the general linear congruence group, and the metacyclic groups $P$ must be contained in it.

Not all linear congruence groups, however, are conversely metacyclic, and to separate those which are is a further problem, solved so far only in special cases.\footnote{C. Jordan, Liouville's Journal de Mathématiques, Series (2), Vol. XIII.} To formulate the problem more precisely, we add the following considerations.

Let the complete linear congruence group (6) be denoted by $R$. From it we separate the group of linear homogeneous substitutions
\begin{equation}
\begin{aligned}
 z'_1&\equiv \alpha_{1,1}z_1+\alpha_{1,2}z_2+\cdots+\alpha_{1,k}z_k,\\
 z'_2&\equiv \alpha_{2,1}z_1+\alpha_{2,2}z_2+\cdots+\alpha_{2,k}z_k,\\
 &\hspace{4em}\cdots\\
 z'_k&\equiv \alpha_{k,1}z_1+\alpha_{k,2}z_2+\cdots+\alpha_{k,k}z_k
\end{aligned}
\pmod p,
 \tag{7}
\end{equation}
with non-vanishing determinant; this group will be called the linear homogeneous congruence group and denoted by $S$. Moreover, (6) contains the $k$-fold cyclic group $Q$ consisting of the substitutions
\begin{equation}
 z'_1=z_1+\alpha_1,\quad
 z'_2=z_2+\alpha_2,\ldots,\quad
 z'_k=z_k+\alpha_k .
 \tag{8}
\end{equation}

If $\sigma$ denotes an element of $S$ and $\gamma$ an element of $Q$, then from the composition rule \S\,77, (7) it follows without difficulty that every substitution $\pi$ of the group $R$ can be represented in one and only one way in each of the two forms
\begin{equation}
 \pi=\sigma\gamma=\gamma'\sigma,
 \tag{9}
\end{equation}
where, if $\sigma$ and $\gamma$ are determined by (7) and (8), $\gamma'$ denotes the substitution
\begin{equation}
 z'_1=z_1+\alpha'_1,\quad
 z'_2=z_2+\alpha'_2,\ldots,\quad
 z'_k=z_k+\alpha'_k,
 \tag{10}
\end{equation}
the $\alpha'_i$ being derived from the $\alpha_i$ by the substitution $\sigma$.

The group $Q$ is a normal divisor of $R$. For, if $\gamma_1$ is any element of $Q$, then repeated application of (9) gives
\[
 \sigma\gamma\gamma_1
 =\gamma'\gamma'_1\sigma
 =\gamma'\gamma'_1\sigma\gamma^{-1}\gamma
 =\gamma'\gamma'_1\gamma'^{-1}\sigma\gamma.
\]
Thus, if $\sigma\gamma=\pi$ and $\gamma'\gamma'_1\gamma'^{-1}=\gamma''_1$, then
\[
 \pi\gamma_1=\gamma''_1\pi,
\]
and consequently
\begin{equation}
 \pi^{-1}Q\pi=Q .
 \tag{11}
\end{equation}

If now $\pi$ runs through any divisor $R_1$ of $R$ which itself contains the group $Q$, then $Q$ is a normal divisor of $R_1$, and the substitution $\sigma$ occurring in (9) must run through a divisor $S_1$ of $S$. We may write
\[
 R_1=S_1Q=QS_1,
\]
and because the substitutions contained in the form (9) are all distinct, the degree of $R_1$ is the product of the degrees of $S_1$ and $Q$.

If $R_2$ is a normal divisor of $R_1$ which still contains $Q$ as normal divisor, then similarly
\[
 R_2=S_2Q,
\]
and $S_2$ is a normal divisor of $S_1$. Conversely, if $S_2$ is a normal divisor of $S_1$, then $S_2Q$ is a normal divisor of $S_1Q$. Hence:

\begin{enumerate}[label=\Roman*.,leftmargin=2.8em,start=2]
\item If $P$ is the Galois group of a primitive irreducible metacyclic equation of degree $p^k$, then $P$ can be represented in the form
\[
 P=TQ,
\]
where $T$ is a metacyclic group contained in the congruence group $S$.
\end{enumerate}

The task of finding all these metacyclic equations is therefore reduced to finding all metacyclic divisors of the homogeneous congruence group $S$.

As an example take the case $p=3,\;k=2$, that is, the metacyclic equations of degree $9$. By the preceding theorem we are concerned with the group $S$ of linear substitutions modulo $3$,
\begin{equation}
 \sigma=\mat{a}{b}{c}{d},
 \tag{12}
\end{equation}
which we considered in \S\,66 from a somewhat different point of view.

There we first considered a divisor $E$ of $S$, consisting of all substitutions $\sigma$ satisfying $ad-bc\equiv1\pmod 3$, and in addition identified the two substitutions
\begin{equation}
 \mat{a}{b}{c}{d},\qquad
 \mat{-a}{-b}{-c}{-d}
 \tag{13}
\end{equation}
as one element. The group so specialized had degree $12$ and possessed a normal divisor of index $3$:
\[
 G=
 \mat{1}{0}{0}{1},\quad
 \mat{0}{1}{-1}{0},\quad
 \mat{-1}{1}{1}{1},\quad
 \mat{1}{1}{1}{-1},
\]
from which $E$ was composed as
\[
 E=G+\mat{1}{0}{1}{1}G+\mat{1}{0}{-1}{1}G .
\]

In the group $S$, however, the two substitutions (13) count as distinct, and in addition the substitutions $\sigma$ for which $ad-bc\equiv -1\pmod 3$ must be added. This raises the degree of the group to $48$. By changing the signs of all elements, $G$ is enlarged to a group of degree $8$, and $E$ to a group of degree $24$, of which the enlarged $G$ is a normal divisor.

Here the whole group $S$ is metacyclic, as is seen from the following composition, in which the enlarged group $G$ is denoted by $S_2$ and the enlarged group $E$ by $S_1$:
\begin{align*}
 S_4&=\mat{1}{0}{0}{1},\ \mat{-1}{0}{0}{-1},\\
 S_3&=S_4+\mat{-1}{1}{1}{1}S_4,\\
 S_2&=S_3+\mat{0}{1}{-1}{0}S_3,\\
 S_1&=S_2+\mat{1}{0}{1}{1}S_2+\mat{1}{0}{-1}{1}S_2,\\
 S&=S_1+\mat{-1}{0}{0}{1}S_1.
\end{align*}
Thus the composition and index series of $S$ is
\[
 \begin{array}{ccccc}
 S,&S_1,&S_2,&S_3,&S_4,\\
 2&3&2&2&
 \end{array}
\]
and we therefore arrive at the result that all equations of degree $9$ with linear congruence group are metacyclic.




\begin{center}
{\large Tenth Section.}\\[0.4em]
{\large General theory of metacyclic equations.}
\end{center}

\section*{\S\,79. The ternary linear congruence group for modulus 2.}

Many interesting relations arise if we consider the ternary linear congruence group $R$ for the modulus $2$. It can be regarded as a transitive permutation group on eight elements, and it must contain the group of the metacyclic equations of the eighth degree. By what has been proved above, the essential point is first to consider the homogeneous congruence group $S$, consisting of the elements
\begin{equation}
 \sigma=\mmat{a}{b}{c}{a_1}{b_1}{c_1}{a_2}{b_2}{c_2},
 \tag{1}
\end{equation}
where the digits $a,b,c,\ldots$ are to be taken modulo $2$, hence as $0$ or $1$, and are compounded according to the rule given in \S\,37. Only those systems of numbers $a,b,c,\ldots$ whose determinant is $1$, that is, odd, are to be considered.

To determine the degree of the group $S$, observe that the first row of $\sigma$ can have seven different forms:
\[
 (1,0,0),\ (0,1,0),\ (0,0,1),\ (0,1,1),\ (1,0,1),\ (1,1,0),\ (1,1,1).
\]
For the first assumption $(a,b,c)=(1,0,0)$ the determinant of $\sigma$ becomes $b_1c_2-c_1b_2$, which can become $1$ in six different ways, namely
\[
 \mat{b_1}{c_1}{b_2}{c_2}=\mat{1}{0}{0}{1},\ \mat{1}{0}{1}{1},\ \mat{0}{1}{1}{0},\ \mat{0}{1}{1}{1},\ \mat{1}{1}{1}{0},\ \mat{1}{1}{0}{1}.
\]
Then $a_1,a_2$ may still be assumed in four different ways. Thus, when the first row is $(1,0,0)$, there are $24$ different forms of $\sigma$. The same number is obtained for the assumptions that the first row is $(0,1,0)$ or $(0,0,1)$. The same result is also obtained for the other assumptions about the first row. For if, for example, it is $(0,1,1)$, then
\[
 a_1(b_2-c_2)-a_2(b_1-c_1)\equiv 1\pmod 2
\]
must hold; this again gives six possibilities for $a_1,b_1-c_1,a_2,b_2-c_2$, and since $c_1,c_2$ can be assumed in four ways, one again obtains $24$ possibilities. Finally, if the first row is $(1,1,1)$, then
\[
 (a_1-c_1)(b_2-c_2)-(a_2-c_2)(b_1-c_1)\equiv 1\pmod 2
\]
must hold, which leads to the same number.

The total number of all distinct substitutions $\sigma$, that is, the degree of $S$, is therefore $168$. We have already encountered the number $168$ once as the degree of a group, namely for the group $L_7$ (\S\,66, 72), and so the conjecture lies close at hand that the groups $S$ and $L_7$ may be isomorphic. If this conjecture is confirmed, the investigation of the group $S$ is substantially simplified, since we already know the divisors of the group $L_7$.

Theorem I of \S\,72 serves to test this conjecture. For this purpose we seek to choose generating elements $\chi,\omega,\Theta,\tau$ of the group $S$ so that they satisfy the characteristic conditions of Theorem I, \S\,72. This can be done in several ways. There is, to be sure, no general procedure for it, but it is easily achieved by a few trials.

First choose arbitrarily an element of seventh degree and form from it the period, for instance
\begin{align}
 \tau&=\mmat{1}{0}{1}{1}{0}{0}{0}{1}{0},&
 \tau^2&=\mmat{1}{1}{1}{1}{0}{1}{1}{0}{0},&
 \tau^3&=\mmat{0}{1}{1}{1}{1}{0}{1}{0}{1},\notag\\
 \tau^4&=\mmat{1}{1}{0}{0}{1}{1}{1}{1}{1},&
 \tau^5&=\mmat{0}{0}{1}{1}{1}{0}{0}{1}{1},&
 \tau^6&=\mmat{0}{1}{0}{0}{0}{1}{1}{1}{0}.
 \tag{2}
\end{align}
Then one seeks, among the elements of second degree, of which there are $21$, one suitable for $\omega$; for example the following one is suitable:
\begin{equation}
 \omega=\mmat{1}{0}{0}{0}{1}{0}{0}{1}{1};
 \tag{3}
\end{equation}
and then, according to the formulae \S\,72, (17), the elements $\Theta$ and $\chi$ are determined:
\begin{align}
 \Theta&=\mmat{1}{1}{1}{0}{1}{0}{0}{1}{1},&
 \Theta^2&=\mmat{1}{1}{0}{0}{1}{0}{0}{0}{1},&
 \Theta^3&=\mmat{1}{0}{1}{0}{1}{0}{0}{1}{1},
 \tag{4}\\
 \chi&=\mmat{1}{0}{0}{0}{0}{1}{0}{1}{1},&
 \chi^2&=\mmat{1}{0}{0}{0}{1}{1}{0}{1}{0}.
 \tag{5}
\end{align}
It is then a matter of a simple calculation to verify the formulae of Theorem I, \S\,72; thereby the agreement of the groups $S$ and $L_7$ is proved. It follows that the group $S$, like $L_7$, is simple.

\section*{\S\,80. Resolvents of the equation of degree 8.}

To make the application to the theory of equations of the eighth degree, we denote, as before, by $Q$ the cyclic group of degree $8$,
\[
 (z_1,z_2,z_3)\longmapsto (z_1+h_1,z_2+h_2,z_3+h_3)\pmod 2,
\]
or, more briefly, $(z_i,z_i+h_i)$, and by $S$ the ternary congruence group considered in \S\,79, and set
\[
 R=SQ,
\]
so that $R$ is the whole ternary linear congruence group for the modulus $2$.

We consider a system of eight independent variables $X_{z_1,z_2,z_3}$, where the indices $z_1,z_2,z_3$ are to be taken modulo $2$. From these we form the seven functions, Lagrange resolvents (Vol. I, \S\,164),
\begin{equation}
 \sum_z (-1)^{\sum z_i\xi_i}X_{z_1,z_2,z_3}=\Psi_{\xi_1,\xi_2,\xi_3},
 \tag{1}
\end{equation}
where, for abbreviation,
\[
 \sum z\xi=z_1\xi_1+z_2\xi_2+z_3\xi_3
\]
is put, and where $\xi_1,\xi_2,\xi_3$ run through a complete residue system modulo $2$, but excluding the combination $0,0,0$.

If we apply to the indices $z_i$ of the quantities $X$ the substitutions of the group $Q$, the $X$ undergo the permutations of a group which we likewise denote by $Q$. The changes of the $\Psi$ corresponding to these permutations follow from (1), if we apply to the indices of the $X$ the substitution $(z_i,z_i+h_i)$:
\begin{equation}
 \sum_z (-1)^{\sum z_i\xi_i}X_{z_1+h_1,z_2+h_2,z_3+h_3}
 =\sum_z (-1)^{\sum (z_i+h_i)\xi_i}X_{z_1,z_2,z_3}
 =(-1)^{\sum h_i\xi_i}\Psi_{\xi_1,\xi_2,\xi_3}.
 \tag{2}
\end{equation}
\begin{enumerate}[label=\arabic*.,leftmargin=2.8em]
\item The functions $\Psi$ are therefore changed by the permutations of $Q$ only in sign, and the squares of the $\Psi$ remain unaltered by $Q$.
\end{enumerate}
It remains to determine the influence of a substitution $\sigma$ from $S$ on the functions $\Psi$. To this end denote by $\rho$ the substitution transposed to $\sigma$, and put
\begin{equation}
 z'=\sigma(z),\qquad \xi=\rho(\xi'),\qquad \xi'=\rho^{-1}(\xi).
 \tag{3}
\end{equation}
Then (\S\,37, 9.)
\begin{equation}
 \sum \xi z=\sum \xi' z'.
 \tag{4}
\end{equation}
If we make the substitutions $\sigma$ of the group $S$ in the indices of $X$, we obtain a permutation group of the $X$ to be denoted again by $S$, and from (1) we get
\[
 \sum_z(-1)^{\sum z_i\xi_i}X_{z'_1,z'_2,z'_3}
 =\sum_z(-1)^{\sum \xi'_i z'_i}X_{z'_1,z'_2,z'_3}.
\]
Since the system of the $z'_i$ runs through the same range of values as the system of the $z_i$, this sum is equal to $\Psi_{\xi'_1,\xi'_2,\xi'_3}$, and we obtain:
\begin{enumerate}[label=\arabic*.,leftmargin=2.8em,start=2]
\item The application of the permutation $\sigma$ to the $X$ in the functions $\Psi$ is equivalent to the application of $\rho^{-1}$ to the indices $\xi_1,\xi_2,\xi_3$.
\end{enumerate}

\section*{\S\,81. The triple systems of the resolvents.}

We now denote the resolvent $\Psi_{\xi_1,\xi_2,\xi_3}$ more briefly by $\Psi_\xi$, and observe that according to formula (2), \S\,80, besides the squares of these functions there are also certain products which admit the permutations of the group $Q$. To these belongs first the product of all the quantities $\Psi$, which we shall denote by $A$:
\begin{equation}
 A=\prod_\xi \Psi_\xi,
 \tag{1}
\end{equation}
and next the products of three $\Psi_\xi\Psi_\eta\Psi_\zeta$, and consequently also their reciprocal values
\begin{equation}
 v={1\over \Psi_\xi\Psi_\eta\Psi_\zeta},
 \tag{2}
\end{equation}
when the indices satisfy the conditions
\begin{equation}
 \xi_1+\eta_1+\zeta_1\equiv0,
 \quad \xi_2+\eta_2+\zeta_2\equiv0,
 \quad \xi_3+\eta_3+\zeta_3\equiv0\pmod2.
 \tag{3}
\end{equation}
Such a system of three functions $\Psi$ we call a triple. Likewise we shall call three index systems $(\xi),(\eta),(\zeta)$ satisfying (3) a triple.

There are altogether seven, and no more, such triples. For among the three systems $(\xi),(\eta),(\zeta)$ no two can be equal, and once two are given the third is uniquely determined. Thus one may take for $(\xi),(\eta)$ each pair among the seven possible systems $(\xi)$, but then each triple is obtained six times. The total number is therefore $7\cdot6:6=7$.

To characterize the separate triple systems, introduce three numbers $\alpha_1,\alpha_2,\alpha_3$ to be taken modulo $2$, not all three vanishing, and observe that the congruence
\begin{equation}
 \alpha_1\xi_1+\alpha_2\xi_2+\alpha_3\xi_3\equiv0\pmod2
 \tag{4}
\end{equation}
is satisfied for three, and only for three, systems of the unknowns $\xi_1,\xi_2,\xi_3$, and that these three form a triple. Thus the triple is determined by the three congruences
\begin{equation}
 \sum\alpha\xi\equiv0,
 \qquad \sum\alpha\eta\equiv0,
 \qquad \sum\alpha\zeta\equiv0\pmod2.
 \tag{5}
\end{equation}
This is seen without lengthy general considerations by considering separately the cases in which only one $\alpha$, or two, or all three $\alpha$ are congruent to $1$ modulo $2$.

Since there are seven different number systems $\alpha_1,\alpha_2,\alpha_3$, the relations (5) determine completely, from each such system of the $\alpha$, one triple; and we may denote the function $v$ unambiguously by $v_{\alpha_1,\alpha_2,\alpha_3}$ or more briefly by $v_\alpha$:
\begin{equation}
 v_\alpha={1\over \Psi_\xi\Psi_\eta\Psi_\zeta}.
 \tag{6}
\end{equation}
\begin{enumerate}[label=\arabic*.,leftmargin=2.8em,start=3]
\item The functions $v_\alpha$, regarded as functions of the $X$, admit the permutations of the group $Q$.
\end{enumerate}
If we apply to the $z$ a substitution $\sigma$ from $S$, then by 2. (\S\,80) the $\xi,\eta,\zeta$ simultaneously undergo the substitution $\rho^{-1}$; but from (5) it follows that the $\alpha$ undergo the substitution $\sigma$, and so we have:
\begin{enumerate}[label=\arabic*.,leftmargin=2.8em,start=4]
\item The application of the substitution $\sigma$ from $S$ to the indices $z$ of $X$ entails the application of the same substitution to the indices $\alpha$ of $v_\alpha$.
\end{enumerate}

If, in the sums (5), where $(\xi),(\eta),(\zeta)$ is the triple belonging to $(\alpha)$, we put in place of $(\alpha)$ another system $(\beta)$, then by (3)
\[
 \sum\beta\xi+\sum\beta\eta+\sum\beta\zeta\equiv0.
\]
If $\beta$ is distinct from $\alpha$, these three sums cannot all be congruent to $0$, because otherwise two different systems $(\alpha),(\beta)$ would lead to the same triple; consequently, of these sums two must be odd and one even. Hence we have the theorem:
\begin{enumerate}[label=\arabic*.,leftmargin=2.8em,start=5]
\item If $(\xi)$ runs through the triple belonging to $(\alpha)$, and if $(\beta)$ is a system different from $(\alpha)$, then among the sums $\sum\beta\xi$ two are odd and one is even.
\end{enumerate}

When in the sum $\sum\alpha\xi$ the system $(\xi)$ runs through all seven possible assumptions, the value $0$ is obtained three times and hence the value $1$ four times. The congruence
\begin{equation}
 \sum\alpha\xi\equiv1\pmod2
 \tag{7}
\end{equation}
therefore has four and only four distinct solutions for $\xi$. The totality of these four solutions we shall call a quadruple. Each triple uniquely determines a quadruple and conversely, and there are therefore also seven different quadruples, completely determined by the number system $(\alpha)$ according to (7). We now prove the theorem:
\begin{enumerate}[label=\arabic*.,leftmargin=2.8em,start=6]
\item If $(\xi)$ runs through the quadruple belonging to $(\alpha)$ and if $(\beta)$ is a system distinct from $(\alpha)$, then among the four sums $\sum\beta\xi$ two are even and two are odd.
\end{enumerate}
For if $(\xi)$ runs through all seven systems, then among the sums $\sum\beta\xi$ there occur three even and four odd; of these the triple belonging to $(\alpha)$ supplies two odd and one even. Thus for the quadruple there remain two even and two odd.

We now change the notation for the quadruples and take, for the system $(\alpha)$, another symbol $(h_1,h_2,h_3)=(h)$. The quadruple belonging to $(h)$ we denote by $(\alpha),(\beta),(\gamma),(\delta)$, so that the four congruences
\begin{equation}
 \sum h\alpha\equiv1,
 \quad \sum h\beta\equiv1,
 \quad \sum h\gamma\equiv1,
 \quad \sum h\delta\equiv1\pmod2
 \tag{8}
\end{equation}
hold. In the denominator of the product of the four functions
\[
 v_\alpha,
 \quad v_\beta,
 \quad v_\gamma,
 \quad v_\delta
\]
there occur, by (6), altogether twelve factors $\Psi_\xi$. Each factor $\Psi_\xi$ occurs as often as the number of even numbers among the sums
\[
 \sum\xi\alpha,
 \quad \sum\xi\beta,
 \quad \sum\xi\gamma,
 \quad \sum\xi\delta
\]
amounts to; that is, $\Psi_h$ does not occur at all, while every other $\Psi_\xi$ occurs twice (by 6). Therefore, if we also use the relation (1), we obtain
\begin{equation}
 \Psi_h^2=A^2v_\alpha v_\beta v_\gamma v_\delta.
 \tag{9}
\end{equation}

\section*{\S\,82. Application to equations of degree 8.}

For the application we assemble the seven triples from which the quadruples can be read as the complements in the full system. The first column contains the symbol $(\alpha)$ to which the triple belongs:
\[
\begin{array}{c|ccc}
100&010&001&011\\
010&100&001&101\\
001&100&010&110\\
011&100&011&111\\
101&010&101&111\\
110&001&110&111\\
111&011&101&110
\end{array}
\]
If we denote the symbols $(\alpha)$ in the order in which they stand in the first column of this table by $1,2,3,4,5,6,7$, then the table for triples and quadruples can be represented more simply as follows:
\[
\begin{array}{c|ccc|cccc}
1&2&3&4&1&5&6&7\\
2&1&3&5&2&4&6&7\\
3&1&2&6&3&4&5&7\\
4&1&4&7&2&3&5&6\\
5&2&5&7&1&3&4&6\\
6&3&6&7&1&2&4&5\\
7&4&5&6&1&2&3&7
\end{array}
\]

In the formulae derived above, we now put for the variables $X_{z_1,z_2,z_3}$ the roots of an equation of degree $8$. We assume that its Galois group $P=TQ$ in the fixed rationality domain $\Omega$ is contained in the linear congruence group $R=SQ$, so that $T$ is a divisor of $S$. All functions admitting the permutations of this group are rational.

To these belongs first the sum of the $X$:
\begin{equation}
 \sum_z X_{z_1,z_2,z_3}=B,
 \tag{1}
\end{equation}
then the quantity $A$ defined by (1) of the preceding paragraph; also the symmetric functions of the seven quantities $\Psi^2$; and, if for the sake of simplicity we add the assumption, to which we shall return, that none of the quantities $\Psi$ vanishes, the symmetric functions of the quantities $v$. Not only the symmetric functions of the $v$ occur, but all functions of the $v$ that admit the permutations of the group $T$.

The seven quantities
\begin{equation}
 v_{1,0,0}=v_1,
 \quad v_{0,1,0}=v_2,
 \quad v_{0,0,1}=v_3,
 \quad v_{0,1,1}=v_4,
 \quad v_{1,0,1}=v_5,
 \quad v_{1,1,0}=v_6,
 \quad v_{1,1,1}=v_7
 \tag{2}
\end{equation}
are then the roots of a rational equation of degree $7$ with group $T$, and by adding equation (1) and the equations [\S\,80, (1)]
\[
 \sum_z(-1)^{\sum z_i\xi_i}X_{z_1,z_2,z_3}=\Psi_{\xi_1,\xi_2,\xi_3},
\]
one obtains, in view of (9) of the preceding paragraph,
\begin{align}
8X_{0,0,0}=A\{&\sqrt{v_1}\sqrt{v_5}\sqrt{v_6}\sqrt{v_7}
+\sqrt{v_2}\sqrt{v_4}\sqrt{v_6}\sqrt{v_7}
+\sqrt{v_3}\sqrt{v_4}\sqrt{v_5}\sqrt{v_7}\notag\\
&+\sqrt{v_2}\sqrt{v_3}\sqrt{v_5}\sqrt{v_6}
+\sqrt{v_1}\sqrt{v_3}\sqrt{v_4}\sqrt{v_6}
+\sqrt{v_1}\sqrt{v_2}\sqrt{v_4}\sqrt{v_5}\notag\\
&+\sqrt{v_1}\sqrt{v_2}\sqrt{v_3}\sqrt{v_7}\}+B.
\tag{3}
\end{align}
There are $15$ changes of signs of the seven square roots $\sqrt v$ under which this expression remains unchanged, namely when all seven signs are changed at once, or when the signs of a triple or of a quadruple are changed. Consequently the expression (3) assumes only eight different values, however the signs of the seven square roots occurring in it may be determined.

\section*{\S\,83. Metacyclic equations of degree 8.}

To find, after these preparations, the primitive metacyclic equations of the eighth degree, we must first determine the metacyclic divisors of the group $S$. The group $S$ itself, as already mentioned, is not metacyclic. Its divisors have already been considered (\S\,70 to 72), and all proper divisors of $S$ are metacyclic. Among them are divisors of degree $21$ and $7$, which we shall denote by $T_{21},T_7$, then divisors of degree $24$, and still other divisors whose degrees divide $24$, which here we shall collect under the symbol $T_{24}$.

These groups $T_{24}$ cannot connect the seven quantities $(\xi)=(\xi_1,\xi_2,\xi_3)$ transitively with one another, since the degree of a transitive permutation group on seven elements must be divisible by $7$ (Vol. I, \S\,154, 7.). We shall show that the permutation groups of the $X$ corresponding to the groups $T_{24}Q$ are imprimitive, and therefore that these groups are to be excluded from our consideration.

Since all permutations $\sigma$ that leave two of the seven elements $(\xi)$ fixed, or only permute them with one another, leave a third element fixed, namely the complement to the triple, two cases are possible:
\begin{enumerate}[label=\arabic*),leftmargin=2.8em]
\item the group $T_{24}$ leaves one element at rest, or
\item the seven elements are decomposed into two parts of three and four elements, each of which is permuted only within itself.
\end{enumerate}
In the latter case the three elements may be regarded as belonging to a triple, and hence the four elements as belonging to a quadruple. For if three elements not belonging to one triple are permuted among themselves, the three complements of every two of these elements to the triple are likewise permuted among themselves, and the one remaining element remains fixed. Thus we return to case 1).

We now add to the seven systems $(\xi)$ the eighth system $(0,0,0)=(0)$, and denote the eight quantities $X_{\xi_1,\xi_2,\xi_3}$ by
\begin{equation}
 0,1,2,3,4,5,6,7.
 \tag{1}
\end{equation}
In case 1), the substitutions of $T_{24}$ leave two of these eight elements, say $0,1$, fixed. There is, moreover, in $Q$ one and only one substitution $\gamma_0$ that interchanges $0$ and $1$.

The two substitutions $1,\gamma_0$ form a group $Q_0$. Any substitution $\gamma_1$ of $Q$ not occurring in $Q_0$ carries $0,1$ into two elements distinct from them, say $2,3$; and through the two substitutions $Q_0\gamma_1$, $0,1$ pass into $2,3$ or into $3,2$. If, further, $\gamma_2$ carries the pair $0,1$ into a third pair $4,5$, then $4,5$ is distinct from both $0,1$ and $2,3$, and so we may decompose $Q$ into the cosets
\[
 Q=Q_0+Q_0\gamma_1+Q_0\gamma_2+Q_0\gamma_3.
\]
Thereby the eight elements (1) are decomposed into four pairs
\begin{equation}
 0,1;\quad 2,3;\quad 4,5;\quad 6,7,
 \tag{2}
\end{equation}
so that through the two substitutions of one of the cosets $Q_0\gamma_i$ the pair $0,1$ goes into one of these four pairs.

If now $h$ is any substitution of the group
\[
 P=T_{24}Q,
\]
then the elements $\tau_0,\tau_1,\tau_2,\tau_3$ of $T_{24}$ and $\gamma'_0,\gamma'_1,\gamma'_2,\gamma'_3$ of $Q$ can be determined so that
\[
 h=\tau_0\gamma'_0,
 \quad \gamma_1h=\tau_1\gamma'_1,
 \quad \gamma_2h=\tau_2\gamma'_2,
 \quad \gamma_3h=\tau_3\gamma'_3,
\]
or
\[
 h=\tau_0\gamma'_0
 =\gamma_1^{-1}\tau_1\gamma'_1
 =\gamma_2^{-1}\tau_2\gamma'_2
 =\gamma_3^{-1}\tau_3\gamma'_3.
\]
This representation shows that $h$ carries any one of the pairs (2) into a pair of the same system. For example, $\gamma_1^{-1}$ carries $2,3$ into $0,1$; $\tau_1$ leaves $0,1$ unaltered; and $\gamma'_1$, which belongs to $Q$, carries $0,1$ into one of the pairs (2).

The group $T_{24}Q$ is therefore imprimitive, in fact in such a way that there are four systems of imprimitivity. The roots $0,1$ satisfy a quadratic equation whose coefficients depend on the roots of a biquadratic equation.

We now pass to case 2), in which the elements of a triple $1,2,3$ and of a quadruple $4,5,6,7$ are connected transitively among themselves by $T_{24}$. The eight quantities (1) then split into two quadruples:
\begin{equation}
 0,1,2,3;\quad 4,5,6,7.
 \tag{3}
\end{equation}
By the substitutions of $Q$, the quantities of one of these quadruples are either merely interchanged among themselves, or they are carried into the quantities of the other quadruple.

To see this, denote, in the earlier notation, $1,2,3$ by $(\xi),(\eta),(\zeta)$. Since these three quantities form a triple, there is [by \S\,81, (5)] a system $(\alpha)$ such that
\begin{equation}
 \sum\alpha\xi\equiv0,
 \quad \sum\alpha\eta\equiv0,
 \quad \sum\alpha\zeta\equiv0\pmod2.
 \tag{4}
\end{equation}
Apply now a substitution from $Q$,
\[
 \begin{pmatrix}
  \xi_1&\xi_2&\xi_3\\
  \xi'_1&\xi'_2&\xi'_3
 \end{pmatrix},
\]
where $\xi'_i=\xi_i+h_i$. It follows from (4) that
\[
 \sum\alpha\xi'\equiv\sum\alpha\eta'\equiv\sum\alpha\zeta'\equiv\sum\alpha h.
\]
Thus $(\xi'),(\eta'),(\zeta')$ form the triple belonging to $(\alpha)$ if $(h)$ itself belongs to this triple, and otherwise $(h),(\xi'),(\eta'),(\zeta')$ form the quadruple belonging to $(\alpha)$.

Let again $h$ denote any substitution from $P$, and let $\gamma$ be a substitution from $Q$ carrying the quadruple $0,1,2,3$ into $4,5,6,7$. Then the elements $\tau',\tau''$ of $T_{24}$ and $\gamma',\gamma''$ of $Q$ can be determined so that
\[
 h=\tau'\gamma'=\gamma^{-1}\tau''\gamma'',
\]
which proves, as above, that the two systems (3) are connected imprimitively by $h$. After adjoining a square root to the rationality domain, the quantities $0,1,2,3$ are then the roots of a biquadratic equation.

These considerations also exclude, for primitive equations of degree $8$, the possibility that one of the seven quantities $\Psi_{\xi_1,\xi_2,\xi_3}$ [\S\,80, (1)] should vanish. These functions are rational functions of the roots $X$; if one of them vanishes, then every permutation of the group of the equation can be applied to the rational equation $\Psi=0$. Hence, by the theorem \S\,80, 2., either all seven functions $\Psi$ must vanish, in which case the roots would be rational, or the substitutions $\rho^{-1}$, applied to the $(\xi)$, must form an intransitive group. The degree of the group could then not be divisible by $7$. But the group of the $\rho^{-1}$ is isomorphic to the group $T$ of the $\sigma$ (\S\,37, 8.), and therefore the degree of $T$ also is not divisible by $7$; consequently $T$ is also intransitive. Hence, as shown above, $TQ$ is imprimitive.

Thus, as groups for primitive metacyclic equations of degree $8$, only the two types
\[
 T_7Q,
 \qquad T_{21}Q
\]
remain. In the first case the quantities $v$ occurring in the formula [\S\,82, (3)] are the roots of a cyclic equation of degree $7$ with group $(z,z+b)$; in the second case they are the roots of a metacyclic equation with group $(z,az+b)$ (Vol. I, \S\,180), where $z,a,b$ are taken modulo $7$ and $a$ runs through the values $1,2,4$.

To obtain an assignment of the $v_i$, as they occur in formula (3), \S\,82, to the $v_z$, where $z$ is the index occurring in the groups $(z,z+b)$ or $(z,az+b)$, we may start from any substitution of seventh order, for example
\[
 \tau=\mmat{1}{0}{1}{1}{0}{0}{0}{1}{0}\quad [\S\,79,(2)].
\]
If we then take $v_{1,0,0}$ for $v_0$, repeated application of $\tau$ to $1,0,0$ gives the following arrangement:
\[
 (1,0,0)=0,
 \quad (1,1,0)=1,
 \quad (1,1,1)=2,
 \quad (0,1,1)=3,
 \quad (1,0,1)=4,
\]
\[
 (0,1,0)=5,
 \quad (0,0,1)=6.
\]
According to \S\,82, (2) one must therefore replace
\[
 v_1,v_2,v_3,v_4,v_5,v_6,v_7
\]
by
\[
 v_0,v_5,v_6,v_3,v_4,v_1,v_2.
\]
Thereby \S\,82, (3) gives
\begin{equation}
 8X=A\sum_{z=0}^{6}\sqrt{v_z}\sqrt{v_{z+1}}\sqrt{v_{z+2}}\sqrt{v_{z+4}}+B,
 \tag{5}
\end{equation}
where the indices of $v$ are to be read modulo $7$. Here $A,B$ are rational quantities, and the $v_z$ are the roots of a cyclic or metacyclic equation of degree $7$ with group $(z,az+b)$.

Conversely, it is not difficult to prove that the eight quantities derived from this expression by changing signs are always the roots of a primitive metacyclic equation of degree $8$; this point will not be pursued more closely here (compare Vol. I, \S\,185).

\section*{\S\,84. Biquadratic equations.}

We have confined ourselves to primitive metacyclic equations of degree $8$, because the imprimitive ones reduce to equations of degree $4$. In order, then, to represent the roots of the latter in the manner of the preceding paragraph, we still have to seek an analogous representation of the roots of biquadratic equations. The same considerations as in \S\,82 apply here, but in a substantially simpler form.

Let us denote the roots of the biquadratic equation by
\begin{equation}
 X_{0,0},\quad X_{1,0},\quad X_{0,1},\quad X_{1,1},
 \tag{1}
\end{equation}
and then the whole permutation group of degree $24$ can be represented as the binary linear congruence group for the modulus $2$:
\begin{equation}
 P=SQ.
 \tag{2}
\end{equation}
The homogeneous group $S$ contained in it has degree $6$ and consists of the substitutions
\begin{equation}
 \sigma=\mat{1}{0}{0}{1},\ \mat{1}{0}{1}{1},\ \mat{0}{1}{1}{0},\ \mat{0}{1}{1}{1},\ \mat{1}{1}{1}{0},\ \mat{1}{1}{0}{1},
 \tag{3}
\end{equation}
and one may take
\begin{equation}
 \sigma_1=\mat{0}{1}{1}{0},
 \qquad
 \sigma_2=\mat{1}{1}{0}{1}
 \tag{4}
\end{equation}
as the generating substitutions of the group $S$.

As in \S\,80, define the three quantities
\begin{equation}
\begin{aligned}
 \Psi_{1,0}&=X_{0,0}-X_{1,0}+X_{0,1}-X_{1,1},\\
 \Psi_{0,1}&=X_{0,0}+X_{1,0}-X_{0,1}-X_{1,1},\\
 \Psi_{1,1}&=X_{0,0}-X_{1,0}-X_{0,1}+X_{1,1}.
\end{aligned}
\tag{5}
\end{equation}
Thus $\Psi_{1,0}^2,\Psi_{0,1}^2,\Psi_{1,1}^2$, and also the product $\Psi_{1,0}\Psi_{0,1}\Psi_{1,1}$, remain unaltered by the substitutions of the cyclic group $Q$.

If the substitution $\sigma$ is applied to the indices $\xi,\eta$ of the $X$, then the indices of the $\Psi$ undergo [as in \S\,80, (2)] the substitution $\rho^{-1}$, where $\rho$ is the substitution transposed to $\sigma$. Corresponding to the substitutions (4), one has
\begin{equation}
 \rho_1^{-1}=\mat{0}{1}{1}{0},
 \qquad
 \rho_2^{-1}=\mat{1}{0}{1}{1}.
 \tag{6}
\end{equation}
Now, for simplicity, suppose that the three quantities (5) are different from zero, and put
\begin{equation}
\begin{aligned}
 v_{1,0}&={\Psi_{0,1}\Psi_{1,1}\over\Psi_{1,0}}={\Psi_{0,1}\Psi_{1,0}\Psi_{1,1}\over\Psi_{1,0}^2},\\[0.25em]
 v_{0,1}&={\Psi_{1,0}\Psi_{1,1}\over\Psi_{0,1}}={\Psi_{0,1}\Psi_{1,0}\Psi_{1,1}\over\Psi_{0,1}^2},\\[0.25em]
 v_{1,1}&={\Psi_{1,0}\Psi_{0,1}\over\Psi_{1,1}}={\Psi_{0,1}\Psi_{1,0}\Psi_{1,1}\over\Psi_{1,1}^2}.
\end{aligned}
\tag{7}
\end{equation}
The quantities $v_{\xi,\eta}$ remain unaltered by the substitutions of $Q$ and undergo, by $\sigma_1,\sigma_2$, the permutations determined by (6):
\[
\begin{array}{c|ccc}
&v_{1,0}&v_{0,1}&v_{1,1}\\
\sigma_1&v_{0,1}&v_{1,0}&v_{1,1}\\
\sigma_2&v_{1,1}&v_{0,1}&v_{1,0}
\end{array}
\]
It follows that the symmetric functions of the three quantities $v_{1,0},v_{0,1},v_{1,1}$ remain unaltered by the permutations $P$, and hence are contained in the rationality domain of the symmetric functions of the $X$. The three quantities $v$ are therefore the roots of a cubic equation, and there would be no difficulty in forming this cubic resolvent of the general biquadratic equation.

From (7) it follows further that
\[
 v_{1,0}v_{0,1}=\Psi_{1,1}^2,
 \quad v_{1,0}v_{1,1}=\Psi_{0,1}^2,
 \quad v_{0,1}v_{1,1}=\Psi_{1,0}^2.
\]
Consequently, using (5), if we also put the rational quantity
\[
 X_{0,0}+X_{1,0}+X_{0,1}+X_{1,1}=a,
\]
we get
\begin{equation}
 4X_{0,0}=a+\sqrt{v_{1,0}}\sqrt{v_{0,1}}
       +\sqrt{v_{1,0}}\sqrt{v_{1,1}}
       +\sqrt{v_{0,1}}\sqrt{v_{1,1}}.
 \tag{8}
\end{equation}
The expression (8) gives only four different values when the signs of the three square roots occurring in it are determined in all possible ways; and conversely it is also easy to show that the four quantities $X$ contained in the form (8), when $v_{1,0},v_{0,1},v_{1,1}$ are the roots of any cubic equation, satisfy a biquadratic equation with rational coefficients.\footnote{This form of the solution of the biquadratic equation is the analogue of Cayley's form of Cardano's formula (Vol. I, \S\,36). I do not know whether this form of the solution of biquadratic equations is otherwise already known in elementary algebra. I first found it in an English problem collection, \emph{Mathematical Tripos, Part I, June 1894}, where the cubic resolvent of the $v$ is also formed.}


\begin{center}
{\large Eleventh Section.}\\[0.4em]
{\large The inflection points of a curve of the third order.}
\end{center}

\section*{\S\,85. Ternary forms and algebraic curves.}

In order, in the further course of the exposition, to present some of the many interesting applications of algebra to geometrical problems, without referring the reader to other aids, the necessary geometrical foundations for these applications are here derived briefly. We restrict ourselves to plane geometry and even there leave aside everything which is not directly required for the applications.\footnote{More extensive material is found in textbooks of analytic geometry, among which George Salmon's \emph{Higher plane curves}, German by Fiedler, should be noted.}

We use homogeneous triangular coordinates. The position of a point $x$ is determined by its distances from the three sides of the coordinate triangle, measured in three arbitrary units or multiplied by three arbitrary constant factors. These quantities are the coordinates of the point $x$, and are denoted by $x_1,x_2,x_3$ or by $y_1,y_2,y_3$. Between them there is at first a linear non-homogeneous relation. By means of this relation the unit can be represented linearly and homogeneously by $x_1,x_2,x_3$, and every non-homogeneous function can be changed into a homogeneous one.

The coordinates $x_1,x_2,x_3$ are now regarded as independent variables, provided only homogeneous forms and equations are used. Not $x_1,x_2,x_3$ themselves, but only the ratios $x_1:x_2:x_3$, determine the point. The value combination $x_1=x_2=x_3=0$ is excluded; $hx_1,hx_2,hx_3$ determine the same point for every $h\ne0$.

Every linear substitution
\begin{equation}
\begin{aligned}
 x_1&=\alpha^{(1)}_1y_1+\alpha^{(2)}_1y_2+\alpha^{(3)}_1y_3,\\
 x_2&=\alpha^{(1)}_2y_1+\alpha^{(2)}_2y_2+\alpha^{(3)}_2y_3,\\
 x_3&=\alpha^{(1)}_3y_1+\alpha^{(2)}_3y_2+\alpha^{(3)}_3y_3,
\end{aligned}
\tag{1}
\end{equation}
whose determinant
\begin{equation}
 r=\sum \pm \alpha^{(1)}_1\alpha^{(2)}_2\alpha^{(3)}_3
 \tag{2}
\end{equation}
is not zero, admits a twofold geometrical interpretation: either as a mapping of two points in the same coordinate system, or as a coordinate transformation of one and the same point. In the next considerations the latter interpretation is retained.

An arbitrary ternary form of degree $n$, $f(x_1,x_2,x_3)$, when set equal to zero, represents a curve of order $n$; a straight line has $n$ intersections with it. By (1) the form $f$ passes into a form of the same degree in $y_1,y_2,y_3$,
\begin{equation}
 f(x_1,x_2,x_3)=\varphi(y_1,y_2,y_3),
 \tag{3}
\end{equation}
and $\varphi=0$ represents the same curve in the new coordinates. Differentiating (3) with respect to $y_1,y_2,y_3$ gives
\begin{equation}
\begin{aligned}
 \varphi'(y_1)&=\alpha^{(1)}_1 f'(x_1)+\alpha^{(1)}_2 f'(x_2)+\alpha^{(1)}_3 f'(x_3),\\
 \varphi'(y_2)&=\alpha^{(2)}_1 f'(x_1)+\alpha^{(2)}_2 f'(x_2)+\alpha^{(2)}_3 f'(x_3),\\
 \varphi'(y_3)&=\alpha^{(3)}_1 f'(x_1)+\alpha^{(3)}_2 f'(x_2)+\alpha^{(3)}_3 f'(x_3).
\end{aligned}
\tag{4}
\end{equation}

\section*{\S\,86. Singular points. Inflection points. Bitangents.}

If for a point $x$ the three equations
\begin{equation}
 f'(x_1)=0,
 \qquad f'(x_2)=0,
 \qquad f'(x_3)=0
 \tag{1}
\end{equation}
are simultaneously satisfied, then by Euler's theorem
\begin{equation}
 nf=x_1f'(x_1)+x_2f'(x_2)+x_3f'(x_3)
 \tag{2}
\end{equation}
the point lies on the curve $f$. The equations \S\,85, (4) also show that this property is preserved under linear transformations. Such points do not occur on every curve of order $n$. Eliminating $x_1,x_2,x_3$ from (1) gives a relation among the coefficients of $f$; this integral rational homogeneous function of the coefficients is called the discriminant of $f$, and is determined only up to a numerical factor. Points of the curve whose coordinates satisfy (1) are called singular points.

If a curve is to have infinitely many singular points, then the partial derivatives must have a common non-constant divisor; by the same considerations this divisor must also have a common divisor with $f$. A curve with infinitely many singular points is therefore reducible, and the points lie on a multiple component of the curve.

For the geometrical interpretation, order $f$ according to powers of $x_3$:
\begin{equation}
 f(x_1,x_2,x_3)=x_3^n f_0+x_3^{n-1}f_1+x_3^{n-2}f_2+\cdots,
 \tag{3}
\end{equation}
where $f_0,f_1,f_2,\ldots$ are binary forms in $x_1,x_2$ and the degree is indicated by the subscript. If a singular point is at the vertex $I_3$, then $f_0$ and $f_1$ must vanish, and the lines through $I_3$ tangent to the curve are obtained from $f_2=0$. If $f_2$ is not identically zero, $I_3$ is a double point; if also $f_2=0$, one obtains a singular point of higher multiplicity.

A point of a curve is called an inflection point if the tangent has three-point contact with the curve there. If $I_3$ is an inflection point and $x_1=0$ is the inflection tangent, then $f$ has the form
\begin{equation}
 f=x_1\varphi+x_3^2\psi,
 \tag{4}
\end{equation}
where $\varphi$ has degree $n-1$ and $\psi$ degree $n-2$. Stronger divisibility expresses four-point contact. A straight line which touches the curve at two distinct points is called a bitangent. If it is chosen as $x_3=0$ and the points of contact are chosen as the two vertices on $x_3=0$, then $f\vert_{x_3=0}$ must have both $x_1^2$ and $x_2^2$ as factors.

\section*{\S\,87. Covariants of ternary forms.}

Let
\[
 f(x_1,x_2,x_3)=\sum a_{i_1i_2i_3}x_1^{i_1}x_2^{i_2}x_3^{i_3}
\]
be a ternary form. A form $C(x,a)$, homogeneous in the variables $x$ and in the coefficients $a$, is called a covariant of $f$ if under the transformation (1) of the variables it transforms by a single factor:
\begin{equation}
 C(y,b)=r^\lambda C(x,a).
 \tag{1}
\end{equation}
If the form is independent of the variables, it is an invariant. The exponent $\lambda$ is the weight.

For every ternary form of degree greater than $1$ there is the Hessian determinant. With the normalization used here we write
\begin{equation}
 H(x)=\det\bigl(f_{ij}\bigr)_{1\le i,j\le3},\qquad f_{ij}=f''(x_i,x_j),
 \tag{2}
\end{equation}
or, up to Weber's chosen numerical factor, $\Delta(x)$. Besides this covariant Weber considers two further general covariants. First, from the second derivatives of $f$ and the first derivatives of the Hessian covariant, one forms
\begin{equation}
 J(x)=\frac1{36}
 \left|\begin{array}{cccc}
 f_{11}&f_{12}&f_{13}&\Delta_1\\
 f_{21}&f_{22}&f_{23}&\Delta_2\\
 f_{31}&f_{32}&f_{33}&\Delta_3\\
 \Delta_1&\Delta_2&\Delta_3&0
 \end{array}\right|,
 \tag{3}
\end{equation}
where $\Delta_i=\Delta'(x_i)$. Then the third covariant is the functional determinant of $f,\Delta,J$:
\begin{equation}
 K(x)=\frac1{9}
 \left|\begin{array}{ccc}
 f_1&\Delta_1&J_1\\
 f_2&\Delta_2&J_2\\
 f_3&\Delta_3&J_3
 \end{array}\right|.
 \tag{4}
\end{equation}
Their transformation laws are
\begin{equation}
 \Delta(y)=r^2\Delta(x),\qquad J(y)=r^6J(x),\qquad K(y)=r^9K(x),
 \tag{5}
\end{equation}
which gives the corresponding weights.

\section*{\S\,88. The Hessian curve.}

The curve represented by $H=0$ is called the Hessian curve of the curve $f=0$. To see its relation with inflection points, write $f$ at the point $I_3$ in the form
\begin{equation}
 f=x_3^n f_0+x_3^{n-1}f_1+x_3^{n-2}f_2+\cdots.
 \tag{1}
\end{equation}
If one inserts the form corresponding to an inflection point $I_3$ with inflection tangent $x_1=0$, then the Hessian covariant, ordered by descending powers of $x_3$, becomes
\begin{equation}
 H=x_3^{3n-6}H_0+x_3^{3n-7}H_1+\cdots,
 \tag{2}
\end{equation}
with
\begin{equation}
 H_0=-2(n-1)^2h^2c.
 \tag{3}
\end{equation}
Thus the Hessian curve passes through $I_3$ precisely when $c=0$, that is, precisely when $I_3$ is an inflection point of the curve $f$. From the expression for $H_1$ one further reads off when the inflection tangent is also tangent to the Hessian curve. For a curve of the third order this four-point contact can occur only if the cubic form becomes reducible.

Thus, for a nonsingular cubic curve, the determination of the inflection points is reduced to the intersections of $f=0$ and $H=0$. Since both curves are cubic, there are nine inflection points, apart from exceptional cases.

\section*{\S\,89. Inflection points of a curve of the third order.}

Let $f=0$ be the equation of a curve of the third order without singular point. Place two vertices of the coordinate triangle, $I_1,I_2$, in two of the nine inflection points, and choose the corresponding inflection tangents as the sides $x_1=0$ and $x_2=0$. Then, when $x_1=0$ or $x_2=0$ is put, $f$ must reduce to one term $hx_3^3$. Hence $f$ has the form
\begin{equation}
 f(x_1,x_2,x_3)=x_1x_2(a_1x_1+a_2x_2+a_3x_3)+hx_3^3.
 \tag{1}
\end{equation}
The straight line
\begin{equation}
 a_1x_1+a_2x_2+a_3x_3=0
 \tag{2}
\end{equation}
is likewise an inflection tangent, and its intersection with $x_3=0$ is the third inflection point on this line.

Let $\varepsilon$ be an imaginary cube root of unity. Introduce two linear functions $z_1,z_2$ by
\begin{equation}
\begin{aligned}
 a_1x_1&=\varepsilon z_1+\varepsilon^2z_2-\tfrac13a_3x_3,\\
 a_2x_2&=\varepsilon^2z_1+\varepsilon z_2-\tfrac13a_3x_3,
\end{aligned}
\tag{3}
\end{equation}
from which also
\begin{equation}
 -(a_1x_1+a_2x_2+a_3x_3)=z_1+z_2-\tfrac13a_3x_3
 \tag{4}
\end{equation}
follows. On substituting in (1), and after multiplying by suitable constant factors, one obtains the symmetric representation
\begin{equation}
 f=\varphi(y_1,y_2,y_3)=h(y_1^3+y_2^3+y_3^3)+6my_1y_2y_3.
 \tag{5}
\end{equation}
Consequently a nonsingular ternary cubic form can be transformed in four different ways into the canonical form (5). Since on each of the lines $y_1=0,y_2=0,y_3=0$ there are three inflection points, the problem of the inflection points is essentially the same as the transformation of the cubic form to canonical form.\footnote{Weber refers here to Newton and Maclaurin on inflection points of cubic curves, Hesse on elimination, and Aronhold on homogeneous functions of the third degree.}

\section*{\S\,90. Transformation of the cubic form to canonical form.}

For the canonical form write
\begin{equation}
 \varphi(y)=h\sum y_i^3+6my_1y_2y_3.
 \tag{1}
\end{equation}
To determine the transformation one first computes the covariants of the canonical form. The Hessian covariant is
\begin{equation}
 \Delta(y)=-hm^2\sum y_i^3+(h^3+2m^3)y_1y_2y_3.
 \tag{2}
\end{equation}
The second covariant is written in the normalization
\begin{equation}
 J(x)=\frac1{36}
 \left|\begin{array}{cccc}
 f_{11}&f_{12}&f_{13}&\Delta_1\\
 f_{21}&f_{22}&f_{23}&\Delta_2\\
 f_{31}&f_{32}&f_{33}&\Delta_3\\
 \Delta_1&\Delta_2&\Delta_3&0
 \end{array}\right|,
 \tag{3}
\end{equation}
and for the canonical form
\begin{align}
 J(y)=&-h^2(h^3+8m^3)^2\sum y_i^3y_j^3
      +3m^2(5h^6+26h^3m^3-4m^6)y_1^2y_2^2y_3^2\notag\\
 &+hm(2h^6+5h^3m^3+20m^6)y_1y_2y_3\sum y_i^3
      +9h^2m^6\left(\sum y_i^3\right)^2 .
 \tag{4}
\end{align}
From the invariant relations
\begin{equation}
 \varphi(y)=f,
 \qquad \Delta(y)=r^2\Delta,
 \qquad J(y)=r^6J
 \tag{5}
\end{equation}
one obtains first
\begin{align}
 (h^3+8m^3)y_1y_2y_3&=m^2f+r^2\Delta,
 \tag{6}\\
 h(h^3+8m^3)\sum y_i^3&=(h^3+2m^3)f-6mr^2\Delta.
 \tag{7}
\end{align}
Then (4) gives
\begin{equation}
\begin{aligned}
 h^2(h^3+8m^3)^2\sum y_i^3y_j^3={}&(2h^3+m^3)m^3f^2
 +m(2h^3-5m^3)r^2f\Delta\\
&+3m^2r^4\Delta^2-r^6J.
\end{aligned}
\tag{8}
\end{equation}

The third covariant is
\begin{equation}
 K(x)=\frac1{9}
 \left|\begin{array}{ccc}
 f_1&\Delta_1&J_1\\
 f_2&\Delta_2&J_2\\
 f_3&\Delta_3&J_3
 \end{array}\right|,
 \tag{9}
\end{equation}
and for the canonical form
\begin{equation}
 K(y)=-h^3(h^3+8m^3)^3
 \left|\begin{array}{ccc}
 1&y_1^3&y_1^6\\
 1&y_2^3&y_2^6\\
 1&y_3^3&y_3^6
 \end{array}\right|.
 \tag{10}
\end{equation}
Since $K(y)=r^9K$, it follows that
\begin{equation}
 r^9K=h^3(h^3+8m^3)^3(y_1^3-y_2^3)(y_1^3-y_3^3)(y_2^3-y_3^3).
 \tag{11}
\end{equation}
If there is no singular point, $h^3+8m^3$ cannot vanish. Put $h=1$. Then (6), (7), (8) give the cubic equation
\begin{equation}
 u^3-Qu^2+Pu-R^3=0,
 \tag{12}
\end{equation}
whose roots are $y_1^3,y_2^3,y_3^3$, with
\begin{align}
 P&={ (2+m^3)m^3f^2+(2-5m^3)mr^2f\Delta+3m^2r^4\Delta^2-r^6J\over (1+8m^3)^2},
 \tag{13}\\
 Q&={ (1+2m^3)f-6mr^2\Delta\over 1+8m^3},
 \tag{14}\\
 R&={m^2f+r^2\Delta\over 1+8m^3}.
 \tag{15}
\end{align}
The discriminant of this cubic equation is
\begin{equation}
 D_1={r^{18}K^2\over(1+8m^3)^6}.
 \tag{16}
\end{equation}

\section*{\S\,91. Invariants of the curve of the third order. Reality.}

For the computation of the invariants Weber considers the determinant of the Hessian covariant of the form
\begin{equation}
 \Delta_{\lambda,\mu}=\lambda f+\mu\Delta.
 \tag{1}
\end{equation}
It can be represented in the form
\begin{equation}
 \Delta_{\lambda,\mu}(x)=Lf+M\Delta,
 \tag{2}
\end{equation}
where $L,M$ are binary forms of the third degree in $\lambda,\mu$. For the canonical form, replacing in \S\,90, (2) the quantities
\[
 h,\qquad m
\]
by
\[
 h(\lambda-6m^2\mu),\qquad m\lambda+(h^3+2m^3)\mu
\]
gives
\begin{equation}
\begin{aligned}
\Delta_{\lambda,\mu}={}&-h(\lambda-6m^2\mu)[m\lambda+(h^3+2m^3)\mu]^2\sum y_i^3\\
&+\{h^3(\lambda-6m^2\mu)^3+2[m\lambda+(h^3+2m^3)\mu]^3\}y_1y_2y_3.
\end{aligned}
\tag{3}
\end{equation}
Consequently
\begin{equation}
\begin{aligned}
 L&=-2\lambda^2\mu(h^3-m^3)m-\lambda\mu^2(h^6-20h^3m^3-8m^6)+8\mu^3m^2(h^3-m^3)^2,\\
 M&=\lambda^3+12\lambda^2\mu m(h^3-m^3)+2\mu^3(h^6-20h^3m^3-8m^6).
\end{aligned}
\tag{4}
\end{equation}
Thus one obtains two invariants of degrees four and six, in canonical form
\begin{equation}
 m(h^3-m^3),
 \qquad h^6-20h^3m^3-8m^6.
 \tag{5}
\end{equation}
Denote the corresponding invariants in the general form by $S,T$, and put $h=1$. Then
\begin{equation}
 m(1-m^3)=r^4S,
 \qquad 1-20m^3-8m^6=r^6T.
 \tag{6}
\end{equation}
For $L,M$ one obtains
\begin{equation}
 L=-2S\lambda^2\mu-T\lambda\mu^2+8S^2\mu^3,
 \qquad M=\lambda^3+12S\lambda^2\mu+2T\mu^3.
 \tag{7}
\end{equation}
Eliminating $m^2$ and $m^5$ gives the equation
\begin{equation}
 27m^8+18Sm^4r^4+Tm^2r^6-S^2r^8=0.
 \tag{8}
\end{equation}
With
\begin{equation}
 z={3m^2\over r^2}
 \tag{9}
\end{equation}
this becomes
\begin{equation}
 z^4+6Sz^2+Tz-3S^2=0.
 \tag{10}
\end{equation}
The discriminant is
\begin{equation}
 D_2=-27(T^2+64S^3)^2=-27D^2,
 \tag{11}
\end{equation}
where $D=T^2+64S^3$. For the canonical form
\begin{equation}
 D=h^3(h^3+8m^3)^3,
 \tag{12}
\end{equation}
and hence $D$ cannot vanish as long as the curve has no singular point. It follows that for a real curve the biquadratic equation has negative discriminant, and therefore two real and two conjugate imaginary roots. For the positive real root $z$, the corresponding values $m$ and $r$ are real, and consequently also $y_1,y_2,y_3$ are real. Thus:

\begin{quote}
Every real ternary cubic form with non-vanishing discriminant can be brought, by a real transformation, into the real canonical form.
\end{quote}

In canonical coordinates the nine inflection points are obtained from
\begin{equation}
 y_1^3+y_2^3+y_3^3+6my_1y_2y_3=0
 \tag{13}
\end{equation}
by putting successively $y_1$, $y_2$, $y_3$ equal to zero. If $\varepsilon$ is an imaginary cube root of unity, their coordinates are
\begin{equation}
\begin{array}{ccc}
(0,1,-1),&(0,1,-\varepsilon),&(0,1,-\varepsilon^2),\\
(-1,0,1),&(-\varepsilon,0,1),&(-\varepsilon^2,0,1),\\
(1,-1,0),&(1,-\varepsilon,0),&(1,-\varepsilon^2,0).
\end{array}
\tag{14}
\end{equation}
The three points in each row lie on the lines $y_1=0$, $y_2=0$, $y_3=0$. The remaining triples of collinear points are arranged as
\begin{equation}
\begin{array}{ccc|ccc|ccc}
(0,1,-1)&(-1,0,1)&(1,-1,0)&(0,1,-1)&(-\varepsilon,0,1)&(1,-\varepsilon^2,0)&(0,1,-1)&(-\varepsilon^2,0,1)&(1,-\varepsilon,0)\\
(0,1,-\varepsilon)&(-\varepsilon,0,1)&(1,-\varepsilon,0)&(0,1,-\varepsilon)&(-\varepsilon^2,0,1)&(1,-1,0)&(0,1,-\varepsilon)&(-1,0,1)&(1,-\varepsilon^2,0)\\
(0,1,-\varepsilon^2)&(-\varepsilon^2,0,1)&(1,-\varepsilon^2,0)&(0,1,-\varepsilon^2)&(-1,0,1)&(1,-\varepsilon,0)&(0,1,-\varepsilon^2)&(-\varepsilon,0,1)&(1,-1,0)
\end{array}
\tag{15}
\end{equation}

For the reality relations, besides (14), the conjugate systems are
\begin{equation}
\begin{array}{ccc}
(0,1,-1),&(-\varepsilon,0,1),&(1,-\varepsilon^2,0)\\
(-1,0,1),&(1,-\varepsilon,0),&(0,1,-\varepsilon^2)\\
(1,-1,0),&(0,1,-\varepsilon),&(-\varepsilon^2,0,1),
\end{array}
\tag{16}
\end{equation}
and
\begin{equation}
\begin{array}{ccc}
(0,1,-1),&(-\varepsilon^2,0,1),&(1,-\varepsilon,0)\\
(-1,0,1),&(1,-\varepsilon^2,0),&(0,1,-\varepsilon)\\
(1,-1,0),&(0,1,-\varepsilon^2),&(-\varepsilon,0,1).
\end{array}
\tag{17}
\end{equation}
A still clearer notation is obtained by letting $\xi$ and $\eta$ each run through a complete residue system modulo $3$. The nine inflection points are denoted by the pairs $(\xi,\eta)$; $(\xi,\eta)$ and $(\eta,\xi)$ are conjugate imaginary points, and $(0,0),(1,1),(2,2)$ are real. The four systems of lines are
\begin{equation}
\begin{array}{ccc@{\qquad}ccc}
(0,0)&(1,2)&(2,1)&(0,0)&(1,1)&(2,2)\\
(1,1)&(2,0)&(0,2)&(1,2)&(2,0)&(0,1)\\
(2,2)&(0,1)&(1,0)&(2,1)&(0,2)&(1,0)\\[0.5em]
(0,0)&(2,0)&(1,0)&(0,0)&(0,2)&(0,1)\\
(1,1)&(0,1)&(2,1)&(1,1)&(1,0)&(1,2)\\
(2,2)&(1,2)&(0,2)&(2,2)&(2,1)&(2,0).
\end{array}
\tag{18}
\end{equation}
Therefore three inflection points $(\xi_1,\eta_1)$, $(\xi_2,\eta_2)$, $(\xi_3,\eta_3)$ lie on a straight line if and only if
\begin{equation}
 \xi_1+\xi_2+\xi_3\equiv0,
 \qquad
 \eta_1+\eta_2+\eta_3\equiv0
 \pmod 3.
 \tag{19}
\end{equation}
The notation of the four systems of straight lines can also be derived from the schema
\begin{equation}
\begin{array}{ccc}
(0,0)&(1,2)&(2,1)\\
(1,1)&(2,0)&(0,2)\\
(2,2)&(0,1)&(1,0),
\end{array}
\tag{20}
\end{equation}
by taking the rows, the columns, and the triples corresponding to the positive and negative terms of the determinant.




\section*{\S\,92. Triple equations.}

The system of the nine inflection points has the property that from any two of them one obtains, by a rational process, a quite definite third one. This appears as a property of the equation of the ninth degree on which the determination of the inflection points depends, for instance of the equation whose roots are the abscissas of these points. We make the following definition.

An equation without equal roots is called a triple equation if every two of its roots determine a third, in such a way that each root of such a triple can be expressed rationally by the two others, and indeed so that in an equation
\begin{equation}
 x_3=\Theta(x_1,x_2)
 \tag{1}
\end{equation}
where \(\Theta\) denotes one fixed rational function, the roots of any triple may be substituted for \(x_1,x_2,x_3\) in arbitrary order. The equation on which the inflection points depend is a triple equation of the ninth degree. There are also triple equations of other degrees, for example of the seventh degree, among which are those considered in the preceding section.\footnote{Triple equations of the ninth degree were first treated by Hesse in the paper ``Algebraische Auflösung derjenigen Gleichung 9ten Grades etc.'', Crelle's Journal, vol. 34 (1847).}\footnote{Cf. Noether, Mathematische Annalen, vol. 15 (1879), ``Ueber die Gleichungen 8ten Grades und ihr Auftreten in der Theorie der Curven vierter Ordnung.''}

We shall here consider only triple equations of the ninth degree, and first show how the notation and arrangement of the roots obtained in the preceding paragraph for the inflection points can be derived from the triple property. Each of the nine roots occurs in four triples, and hence altogether there are \(4\cdot 9:3=12\) different triples.

Take an arbitrary one of these triples and denote its roots by
\[
\trip{(0,0)}{(1,2)}{(2,1)}.
\]
Let \((1,1)\) be an arbitrary fourth root. With the first three it first determines three triples, which we denote by
\[
\begin{array}{ccc}
(1,1)&(0,0)&(2,2),\\
(1,1)&(1,2)&(1,0),\\
(1,1)&(2,1)&(0,1),
\end{array}
\]
and there remain two roots which form the fourth triple with \((1,1)\); we denote it by
\[
\trip{(1,1)}{(2,0)}{(0,2)}.
\]
We next seek the triple determined by \((0,0),(2,0)\). It cannot contain any root which occurs with either of these two roots in one of the triples already fixed; only \((0,1)\) or \((1,0)\) remains. Since \((0,2)\) and \((2,0)\) may still be interchanged, no generality is lost by taking the two further triples
\[
\trip{(0,0)}{(2,0)}{(1,0)},
\qquad
\trip{(0,0)}{(0,2)}{(0,1)}.
\]
Determining the triples which contain \((1,0)\), one obtains
\[
\trip{(2,2)}{(0,1)}{(1,0)},
\qquad
\trip{(1,0)}{(0,2)}{(2,1)},
\qquad
\trip{(0,1)}{(2,0)}{(1,2)}.
\]
Finally the triples containing respectively \((0,2)\) and \((2,0)\) are
\[
\trip{(0,2)}{(1,2)}{(2,2)},
\qquad
\trip{(2,0)}{(2,1)}{(2,2)},
\]
and all twelve triples, hence the arrangement of \S\,91,(20), are thereby determined.

\section*{\S\,93. The group of the triple equations.}

The Galois group \(P\) of the triple equations is now to be determined. It is first clear that this group can contain only those permutations by which every triple again passes into a triple. For if the roots of a triple \(x_1,x_2,x_3\) are carried by a permutation into \(y_1,y_2,y_3\), then, since this permutation is applicable to equation (1),
\[
 y_1=\Theta(y_2,y_3),
\]
and therefore \(y_1\) is the third root of the triple determined by \(y_2,y_3\).

Passing to the notation
\begin{equation}
 x=(\xi,\eta),
 \tag{1}
\end{equation}
it follows from the arrangement of the triples that three roots
\begin{equation}
 x_1=(\xi_1,\eta_1),\qquad x_2=(\xi_2,\eta_2),\qquad x_3=(\xi_3,\eta_3)
 \tag{2}
\end{equation}
form a triple if and only if
\begin{equation}
\begin{aligned}
 \xi_1+\xi_2+\xi_3&\equiv 0,\\
 \eta_1+\eta_2+\eta_3&\equiv 0
\end{aligned}
\qquad (\operatorname{mod}\,3).
\tag{3}
\end{equation}
If a permutation group of the pairs \((\xi,\eta)\), all of whose substitutions belong to the complete linear congruence group modulo 3, is called a linear group, then:

\textbf{1. The group \(P\) of a triple equation is always linear.}

By \S\,77 every permutation of the quantities \((\xi,\eta)\) can in general be represented in the form
\begin{equation}
 \xi'\equiv \varphi(\xi,\eta),\qquad \eta'\equiv \psi(\xi,\eta)
 \qquad(\operatorname{mod}\,3),
 \tag{4}
\end{equation}
where \(\varphi,\psi\) are integral functions with integral coefficients of \(\xi,\eta\), of degree not exceeding 2 in either variable. Ordered according to \(\eta\), they may be written
\begin{equation}
 \xi'=A\eta^2+B\eta+C,
 \qquad
 \eta'=A'\eta^2+B'\eta+C',
 \tag{5}
\end{equation}
where \(A,B,C,A',B',C'\) are integral functions of \(\xi\), of degree at most 2. Holding \(\xi\) fixed and putting \(\eta=0,1,2\), the application to the triple \((\xi,0),(\xi,1),(\xi,2)\) gives that \(A\) and \(A'\) must be congruent to zero for every \(\xi\). The same argument removes the terms in \(\xi^2\). Thus the substitutions have first the form
\[
\begin{aligned}
 \xi'&\equiv m\xi\eta+a\xi+b\eta+c,\\
 \eta'&\equiv m'\xi\eta+a'\xi+b'\eta+c'
\end{aligned}
\qquad(\operatorname{mod}\,3).
\]
Applying these to the triple \((0,0),(1,1),(2,2)\) and putting the sums of the corresponding \(\xi'\)'s and \(\eta'\)'s equal to zero gives \(m=m'=0\). Therefore every permutation of \(P\) is contained in
\begin{equation}
\begin{aligned}
 \xi'&\equiv a\xi+b\eta+c,\\
 \eta'&\equiv a'\xi+b'\eta+c'
\end{aligned}
\qquad(\operatorname{mod}\,3).
\tag{6}
\end{equation}

\textbf{2. Conversely, by every linear substitution of the form (6), three quantities \((\xi,\eta)\) of a triple pass into three quantities \((\xi',\eta')\) which likewise form a triple.}

Indeed for arbitrary three pairs (6) gives
\begin{equation}
\begin{aligned}
 \xi'_1+\xi'_2+\xi'_3&\equiv a(\xi_1+\xi_2+\xi_3)+b(\eta_1+\eta_2+\eta_3),\\
 \eta'_1+\eta'_2+\eta'_3&\equiv a'(\xi_1+\xi_2+\xi_3)+b'(\eta_1+\eta_2+\eta_3),
\end{aligned}
\tag{7}
\end{equation}
so that the congruences (3) imply the same congruences after transformation.

It follows further that every equation of the ninth degree whose Galois group \(P\) is contained in the complete linear congruence group is a triple equation, provided the additional condition is satisfied that \(P\) carries any two assigned roots into any other two assigned roots, that is, provided \(P\) is doubly transitive. For if \(P\) consists only of linear substitutions, then after adjunction of two arbitrary roots \(x_1=(\xi_1,\eta_1)\), \(x_2=(\xi_2,\eta_2)\), the group is reduced to one which also leaves fixed the third root \(x_3=(\xi_3,\eta_3)\) of their triple. Hence \(x_3\) belongs to the new field of rationality and
\begin{equation}
 x_3=\Theta(x_1,x_2).
 \tag{8}
\end{equation}
If \(P\) is doubly transitive, this equation is carried by permutations of \(P\) to every other ordered pair, and hence to every triple. Therefore:

\textbf{3. Every equation of the ninth degree whose group is linear and doubly transitive is a triple equation.}

The permutations of a doubly transitive linear group can carry every triple into every other triple, and in arbitrary order. The hypothesis of double transitivity may also be expressed by saying that the permutations of \(P\) which leave one root fixed still connect the remaining roots transitively. If \(P_0\) denotes the subgroup which leaves a root \(x_0\) fixed, then \(P_0\) must still be transitive on the other roots. If \(P\) is only simply transitive, then \(P_0\) is no longer transitive, and after adjoining \(x_0\) the equation for the remaining eight roots splits into several factors. If \(P\) itself is transitive, \(x_0\) may be any root, since conjugation by a permutation merely changes the notation of the roots.

Among the linear groups, the simply transitive ones have a more special character than the group of the triple equations. The Abelian group \(Q\), consisting of
\begin{equation}
 \xi'\equiv \xi+\alpha,
 \qquad
 \eta'\equiv \eta+\beta
 \qquad(\operatorname{mod}\,3),
 \tag{9}
\end{equation}
is only simply transitive. If \(P_0\) is the subgroup of \(P\) which leaves \((0,0)\) fixed, then \(P_0\) contains only homogeneous substitutions, and is the greatest common divisor of \(P\) and the homogeneous congruence group \(S\); consequently \(P=P_0Q\).

Using the composition series
\[
 S,
 S_1,
 S_2,
 S_3,
 S_4,
\]
one has
\[
 S_4=\left\{
 \mat{1}{0}{0}{1},
 \mat{-1}{0}{0}{-1}
 \right\}.
\]
Under \(S_4\) the roots \((1,1),(2,2)\) only pass into each other. The group \(S_4Q\) may give every ordering of the triple \((0,0),(1,1),(2,2)\), but carries it only into the two further triples \((1,0),(2,1),(0,2)\) and \((0,1),(1,2),(2,0)\). It is therefore only simply transitive and is also imprimitive. By adjoining
\[
 \mat{-1}{1}{1}{1}
\]
one obtains \(S_3\); by adjoining
\[
 \mat{0}{1}{-1}{0}
\]
one obtains \(S_2\). Thus \(S_2Q\), and therefore also \(S_1Q\) and \(SQ\), is doubly transitive. Proper triple equations begin only with \(S_2Q\).

Other linear groups contained in \(SQ\) may also be formed. For instance, with the cyclic group of the third degree
\[
 S'=\left\{
 \mat{1}{0}{0}{1},
 \mat{1}{0}{1}{1},
 \mat{1}{0}{-1}{1}
 \right\}
\]
one obtains the group \(S'Q=QS'\) of order 27. Under \(S'\), \((1,1),(2,2)\) passes into \((1,1),(2,2)\), \((1,2),(2,1)\), and \((1,0),(2,0)\), but not into \((0,1),(0,2)\). This group also is not the group of a triple equation. The special metacyclic equations of the ninth degree and the representation of their roots by Kronecker's principles thus still offer a field for further investigation.

\section*{\S\,94. Reality relations of the triple equations.}

If the field of rationality is real, then certain general conclusions can be drawn concerning the reality of the roots of a triple equation of the ninth degree. If in a triple two real roots occur, then from
\[
 x_3=\Theta(x_1,x_2)=\Theta(x_2,x_1)
\]
the third root must also be real. If therefore any imaginary root is present, then in each of the four triples to which it belongs at least one further imaginary root must occur. Hence the number of imaginary roots is at least five, and, since it must be even, is equal to six.

Two real roots, or two conjugate imaginary roots, always belong to a triple whose third root is real. We shall call a triple real if it contains two real roots or two conjugate imaginary roots. If imaginary roots occur in a triple, the conjugate imaginary roots also form a triple; such two triples are called conjugate imaginary triples. There are then three possibilities:
\begin{enumerate}[label=\arabic*.]
\item All roots real.
\item One real and eight imaginary roots.
\item Three real and six imaginary roots.
\end{enumerate}

In the second case the four pairs of conjugate imaginary roots must determine four triples, all having the same real root as third root. If this real root is denoted by \((0,0)\), the four conjugate pairs must be
\[
(1,2)(2,1),\qquad (1,1)(2,2),\qquad (1,0)(2,0),\qquad (0,1)(0,2).
\]

In the third case the three real roots must form a triple. We take them to be
\begin{equation}
 \trip{(0,0)}{(1,1)}{(2,2)}.
 \tag{1}
\end{equation}
Then in the two rows
\begin{equation}
\begin{array}{ccc}
(1,2)&(0,2)&(1,0)\\
(2,1)&(2,0)&(0,1)
\end{array}
\tag{2}
\end{equation}
conjugate imaginary roots cannot occur in the same row, for otherwise the third root of the corresponding row would be real. It remains to show that the three real triples determined by the conjugate pairs must together contain all three real points. If among the three triples
\[
\trip{(0,0)}{(1,2)}{(2,1)},
\quad
\trip{(0,0)}{(2,0)}{(1,0)},
\quad
\trip{(0,0)}{(0,1)}{(0,2)}
\]
two are real, then the third is real as well; the pairs
\[
(1,2)(2,1),
\qquad
(2,0)(1,0),
\qquad
(0,1)(0,2)
\]
would then be conjugate pairs. But then \((1,1),(2,0),(0,2)\) is a triple, whereas the conjugate roots \((1,1),(1,0),(0,1)\) are not a triple, which is impossible. Therefore each of the three real roots \((0,0),(1,1),(2,2)\) occurs, besides in (1), in one and only one real triple whose other two roots are conjugate imaginary. The third case therefore gives exactly the arrangement of real and imaginary roots already found for the inflection points of the curve of the third order.

That cases 1 and 2 can also occur is shown as follows. From a general equation of the ninth degree one obtains a triple equation by adjoining a function belonging to the linear group \(SQ\). One such function is
\begin{align*}
 V={}&(0,0)(1,1)(2,2)+(0,0)(1,2)(2,1)+(0,0)(1,0)(2,0)\\
 &+(0,0)(0,1)(0,2)+(1,1)(1,2)(1,0)+(1,1)(2,1)(0,1)\\
 &+(1,1)(0,2)(2,0)+(1,2)(0,1)(2,0)+(2,2)(2,1)(2,0)\\
 &+(2,2)(1,2)(0,2)+(2,2)(0,1)(1,0)+(2,1)(0,2)(1,0).
\end{align*}
This function is real in cases 1, 2, and 3. Hence the enlarged field of rationality in which the equation is a triple equation is also real.





\section*{Twelfth Section. The bitangents of a curve of the fourth order.}

\section*{\S\,95. Number of bitangents of a curve of the fourth order.}

A geometrical problem which, because of its many relations with other domains, is of special importance, and which at the same time, in a manner similar to the problem of the inflection points of curves of the third order, offers remarkable algebraic relations, is the problem of the bitangents of curves of the fourth order.

By a \emph{bitangent} of a curve one understands, as was already observed in \S\,86, a straight line which touches the curve at two different points. For curves of order lower than the fourth, bitangents cannot occur. For curves of the fourth order a bitangent has, apart from the points of contact, no point in common with the curve. In special cases the two points of contact may also coincide. Then we have lines with four-point contact.

As an algebraic problem the determination of the bitangents depends on an algebraic equation whose degree is equal to the number of bitangents; the first question is therefore the degree of this equation, that is, the number of bitangents. We restrict ourselves here to curves of the fourth order, although the method which we use is also applicable to curves of higher order. We also assume that the curve of the fourth order is free from singular points.\footnote{Jacobi, Clebsch, Hesse, Steiner, Aronhold, Geiser, Cayley, Salmon, Riemann, Weber, Frobenius, and Noether are the principal references named in the source passage for the theory of bitangents of a plane quartic.}

Let
\begin{equation}
 f(x_1,x_2,x_3)=0
 \tag{1}
\end{equation}
or, more briefly, \(f(x)=0\), be a ternary form of the fourth degree which, when equated to zero, represents a curve of the fourth order without singular point; we shall call it the curve \(f\). If, in this equation, one puts in place of \(x_1,x_2,x_3\) expressions of the form
\[
 x_1+t y_1,\qquad x_2+t y_2,\qquad x_3+t y_3,
\]
there results an equation of the fourth degree with respect to \(t\),
\begin{equation}
 f(x+t y)=0,
 \tag{2}
\end{equation}
whose roots, when \((x)\) and \((y)\) are two fixed points, give, when inserted in \(x+t y\), the coordinates of the intersections of the joining line of the points \((x),(y)\) with the curve \(f\). We shall denote this line by \((x,y)\).

Now order the function \(f(x+t y)\) according to powers of \(t\). This gives, as in Volume I, \S\,60,
\begin{equation}
 f(x+t y)=P_0(x,y)+4tP_1(x,y)+6t^2P_2(x,y)+4t^3P_3(x,y)+t^4P_4(x,y),
 \tag{3}
\end{equation}
where the \(P_\nu(x,y)\) are the polars of \(f(x)\), hence homogeneous functions of degree \(\nu\) in \(y\) and of degree \(4-\nu\) in \(x\). Thus
\begin{align}
 P_0(x,y)&=f(x),\notag\\
 P_1(x,y)&=\frac14\sum_i y_i\frac{\partial f}{\partial x_i}(x),\notag\\
 P_2(x,y)&=\frac1{4\cdot3}\sum_{i,k}y_i y_k\frac{\partial^2 f}{\partial x_i\partial x_k}(x),\notag\\
 P_3(x,y)&=\frac1{4\cdot3\cdot2}\sum_{i,k,l}y_i y_k y_l\frac{\partial^3 f}{\partial x_i\partial x_k\partial x_l}(x),\notag\\
 P_4(x,y)&=f(y).
 \tag{4}
\end{align}

If \(x\) lies on the curve \(f\), then \(P_0=0\), and one root of (2) vanishes. If, in addition, \((y)\) is taken on the tangent at \((x)\), then also \(P_1(x,y)=0\), and two roots of (2) vanish. The other two are determined by
\begin{equation}
 6P_2+4tP_3+t^2P_4=0.
 \tag{5}
\end{equation}
If this quadratic equation has two equal roots, then the line \((x,y)\) will have a second point of contact with the curve \(f\), that is, it will be a bitangent. The condition for this is that the discriminant of (5) vanish, or that
\begin{equation}
 R(x,y)=2P_3(x,y)^2-3P_2(x,y)P_4(x,y)
 \tag{6}
\end{equation}
be equal to zero. The equation \(R=0\) is also satisfied when \((x)\) is an arbitrary point of the curve and \((y)\) coincides with \((x)\); otherwise \(R\) vanishes only when the line \((x,y)\) is a bitangent. Hence:

\emph{The equation \(R(x,y)=0\), when \((x)\) is a point of the curve \(f\) and \((y)\) is a point of the tangent at \((x)\) different from \((x)\), is the necessary and sufficient condition that \((x)\) be a point of contact of a bitangent.}

It remains to derive from this condition another one containing only the \(x\)'s, and satisfied by no points of the curve other than the points of contact of the bitangents. The variables \(y\) satisfy the tangent equation
\begin{equation}
 y_1f'_1(x)+y_2f'_2(x)+y_3f'_3(x)=0.
 \tag{7}
\end{equation}
Let \(b_1,b_2,b_3\) be arbitrary constants and put
\begin{equation}
 b_x=b_1x_1+b_2x_2+b_3x_3.
 \tag{8}
\end{equation}
Together with (7) we add the equation \(b_1y_1+b_2y_2+b_3y_3=0\). Thus \((y)\) is the intersection of the tangent at \((x)\) with the arbitrary line \(b_x=0\), and we obtain
\begin{equation}
 y_1=b_2f'_3-b_3f'_2,
 \quad y_2=b_3f'_1-b_1f'_3,
 \quad y_3=b_1f'_2-b_2f'_1.
 \tag{10}
\end{equation}
Through this substitution \(R(x,y)\) becomes a homogeneous function \(D(x,b)\) of degree twenty in \(x\) and degree six in \(b\). Besides the points of contact of the bitangents, it also vanishes at the four intersections of the line \(b_x\) with the curve \(f\).

If on the tangent we replace \(y\) by \(\mu x+\lambda y\), then, for \(f(x)=0\), one has
\begin{equation}
 R(x,\mu x+\lambda y)=\lambda^6R(x,y).
 \tag{11}
\end{equation}
The point \((\mu x+
\lambda y)\) may be any point of the tangent. Repeating the construction with another arbitrary line
\begin{equation}
 a_x=a_1x_1+a_2x_2+a_3x_3
 \tag{13}
\end{equation}
gives
\begin{equation}
 \frac{D(x,a)}{a_x^6}=\frac{D(x,b)}{b_x^6}.
 \tag{14}
\end{equation}
Since this is valid only on \(f(x)=0\), it yields the identity
\begin{equation}
 b_x^6D(x,a)-a_x^6D(x,b)=f(x)Q(x,a,b).
 \tag{15}
\end{equation}
A comparison in a coordinate system in which \(a_x\) and \(b_x\) are coordinate axes shows that
\[
 Q(x,a,b)=a_x^6Q_1(x)-b_x^6Q_2(x).
\]
Consequently there is a form \(\chi(x,a,b)\) of the fourteenth degree such that
\begin{equation}
 \frac{D(x,a)-f(x)Q_1(x)}{a_x^6}=\chi(x,a,b).
 \tag{16}
\end{equation}
After subtracting the relation obtained by giving \(a,b\) fixed values, one obtains a form \(X(x)\), rationally dependent on the coefficients of \(f\), for which
\begin{equation}
 D(x,a)=a_x^6X(x)+f(x)\Phi(x,a).
 \tag{18}
\end{equation}
Thus the curve
\begin{equation}
 X(x)=0
 \tag{19}
\end{equation}
of the fourteenth degree passes through the contact points of the bitangents and through no other points of \(f=0\). It follows that:

\emph{A curve of the fourth order without singular point has twenty-eight bitangents.}

Indeed, the intersections of \(X=0\) with \(f=0\) amount to \(14\cdot4=56\) points, and the two contact points of each bitangent are counted. The only remaining doubt would be that the curve \(X\) could touch \(f\), or that \(f\) could pass through singular points of \(X\). The following considerations remove this doubt by displaying forms of the equation of the curve for which the existence of twenty-eight distinct bitangents is manifest.

\section*{\S\,96. The Steiner complexes.}

If \(x_1=0\) is the equation of a bitangent of the curve of the fourth order \(f=0\), then, after putting \(x_1=0\), \(f\) must become a square. Hence \(f\) must have the form
\begin{equation}
 f=x_1V-u^2,
 \tag{1}
\end{equation}
where \(V\) is cubic and \(u\) is quadratic. But the function \(f\) can be put into the form (1) in a triply infinite number of ways. If \(p\) is any linear function, then
\[
 f=x_1\bigl(V+2pu+p^2x_1\bigr)-(u+px_1)^2,
\]
again has the form (1).

Let \(y_1=0\) be a second bitangent distinct from \(x_1\). The constants in \(p\) may be chosen so that the conic \(u+px_1=0\) passes through the two points of contact of \(y_1\). Then the cubic form \(V\) must vanish at these same points. Since \(y_1=0\) is a bitangent, the cubic \(V=0\) is touched there by the line \(y_1=0\); hence \(V\) must split and contain \(y_1\) as a factor. Therefore
\begin{equation}
 f=x_1y_1v-u^2,
 \tag{2}
\end{equation}
where \(u,v\) are quadratic forms. Conversely, if \(x_1,y_1\) are any two bitangents of a quartic \(f=0\), then \(f\) can be brought into the form (2).

The form (2), with \(x_1,y_1\) retained, is still obtainable in infinitely many ways. If
\[
 f=x_1y_1v-u^2=x_1y_1v_1-u_1^2,
\]
then
\begin{equation}
 x_1y_1(v-v_1)=(u-u_1)(u+u_1).
 \tag{3}
\end{equation}
Since the intersection of \(x_1\) and \(y_1\) cannot lie on a nonsingular quartic, one of the two factors on the right must be divisible by \(x_1y_1\). After choosing the sign of \(u_1\), put
\[
 u_1=u+\lambda x_1y_1.
\]
Then
\begin{equation}
 f=x_1y_1(v+2\lambda u+\lambda^2x_1y_1)-(u+\lambda x_1y_1)^2.
 \tag{4}
\end{equation}
If \(\lambda\) can be so chosen that the quadratic form \(v+2\lambda u+\lambda^2x_1y_1\) splits into two linear factors, then, with \(u+\lambda x_1y_1=u_1\),
\begin{equation}
 f=x_1y_1x_2y_2-u_1^2.
 \tag{5}
\end{equation}
Thus \(x_2,y_2\) are two further bitangents, and the eight contact points of \(x_1,y_1,x_2,y_2\) lie on the conic \(u_1=0\).

The necessary and sufficient condition that a ternary quadratic form split into two linear factors is the vanishing of its determinant. If \(u,v\) are expressed in the variables \(x_1,y_1\) and a third variable and if their coefficients are denoted by \(b_{ik}\) and \(a_{ik}\), one obtains
\begin{equation}
 \left|\begin{array}{ccc}
 a_{11}+2\lambda b_{11}&a_{12}+2\lambda b_{12}&a_{13}+2\lambda b_{13}\\
 a_{21}+2\lambda b_{21}&a_{22}+2\lambda b_{22}&a_{23}+2\lambda b_{23}+\tfrac12\lambda^2\\
 a_{31}+2\lambda b_{31}&a_{32}+2\lambda b_{32}+\tfrac12\lambda^2&a_{33}+2\lambda b_{33}
 \end{array}\right|=0.
 \tag{6}
\end{equation}
This is an equation of the fifth degree in \(\lambda\), and so it gives five decompositions
\[
 x_2y_2,
 x_3y_3,
 x_4y_4,
 x_5y_5,
 x_6y_6.
\]
Thus:

\emph{To each pair of bitangents \(x_1y_1\) there belong five further pairs \(x_iy_i\) such that the eight points of contact of \(x_1,y_1,x_i,y_i\) lie on a conic.}

The six pairs
\[
 x_1y_1,
 x_2y_2,
 x_3y_3,
 x_4y_4,
 x_5y_5,
 x_6y_6
\]
are called a \emph{Steiner complex}. The word ``complex'' is used here instead of the older term ``Steiner group'', because in algebra the word group has a definite different meaning.

Take three pairs \(x_1y_1,x_2y_2,x_3y_3\) of such a complex. Then
\[
 f=x_1y_1x_2y_2-u_1^2=x_1y_1x_3y_3-u_2^2,
\]
and therefore
\[
 x_1y_1(x_2y_2-x_3y_3)=(u_1-u_2)(u_1+u_2).
\]
As before, it follows that \(u_1-u_2=hx_1y_1\), and hence
\begin{equation}
 4f=4x_1y_1x_2y_2-(x_1y_1+x_2y_2-x_3y_3)^2.
 \tag{7}
\end{equation}
After a harmless change of constant factors this becomes
\begin{align}
 -4f&=(x_1y_1+x_2y_2-x_3y_3)^2-4x_1y_1x_2y_2\notag\\
 &=x_1^2y_1^2+x_2^2y_2^2+x_3^2y_3^2
 -2x_2y_2x_3y_3-2x_1y_1x_2y_2-2x_1y_1x_3y_3.
 \tag{8}
\end{align}
The equation \(f=0\) may also be written in the elegant irrational form
\begin{equation}
 \sqrt{x_1y_1}+\sqrt{x_2y_2}+\sqrt{x_3y_3}=0.
 \tag{9}
\end{equation}
From the symmetry of (9) one obtains:

\emph{The pairs of a Steiner complex have the property that the eight points of contact of any two of the pairs lie on a conic, and that one obtains the same complex regardless of which of the six pairs one starts from.}

Three bitangents whose contact points lie on a conic are said, after Frobenius, to form a \emph{syzygetic triple}; three bitangents whose six contact points do not lie on a conic form an \emph{azygetic triple}. Similarly four bitangents whose eight contact points lie on a conic form a syzygetic quadruple, and a system of bitangents of which every three are azygetic is called an azygetic system.

\emph{Three bitangents occurring in a Steiner complex in such a way that no two of them form one of its pairs, for instance \(x_1,x_2,x_3\), are always azygetic.}

The proof follows from (9). The three lines \(x_1,x_2,x_3\) cannot pass through one point; such a point would be on \(f\) and would therefore be singular. Thus we may use \(x_1,x_2,x_3\) as coordinates and write
\begin{align*}
 y_1&=\alpha_1x_1+\alpha_2x_2+\alpha_3x_3,\\
 y_2&=\beta_1x_1+\beta_2x_2+\beta_3x_3,\\
 y_3&=\gamma_1x_1+\gamma_2x_2+\gamma_3x_3.
\end{align*}
The contact points of \(x_1,x_2,x_3\) are obtained from the three pairs of equations
\begin{align*}
&x_1=0,\quad x_2(\beta_2x_2+\beta_3x_3)-x_3(\gamma_2x_2+\gamma_3x_3)=0,\\
&x_2=0,\quad x_3(\gamma_3x_3+\gamma_1x_1)-x_1(\alpha_3x_3+\alpha_1x_1)=0,\\
&x_3=0,\quad x_1(\alpha_1x_1+\alpha_2x_2)-x_2(\beta_1x_1+\beta_2x_2)=0.
\end{align*}
If these six points lay on a conic, the coefficients of that conic would have to satisfy the relations
\begin{equation}
 a_2=h_1\beta_2,
\quad a_3=-h_1\gamma_3,
\quad a_3=h_2\gamma_3,
\quad a_1=-h_2\alpha_1,
\quad a_1=h_3\alpha_1,
\quad a_2=-h_3\beta_2,
 \tag{10}
\end{equation}
whence \(h_1=-h_2=-h_3=-h_1\), impossible unless all \(h_i\) vanish. Then the conic would pass through the pairwise intersections of \(x_1,x_2,x_3\), whereas it was supposed to pass through their contact points. This proves the theorem.

\end{document}
